\documentclass{ctexbook}
\usepackage{graphicx} % 图片
\usepackage{xcolor} % 颜色
%\usepackage[utf8]{inputenc}
\graphicspath{ {./images/} } % 图片路径
\usepackage{enumitem} % 列表用
%\setlength{\parindent}{0pt}
\usepackage{mathtools} %数学
\usepackage{wasysym}
\usepackage{hyperref}
\usepackage{circuitikz}
\usepackage{multirow}
\usepackage{textcomp}
\usepackage{xstring}
\usepackage{catchfile}
\usepackage{siunitx}
\usepackage[a4paper,left=2cm,top=2.5cm,bottom=2cm,
right=1.5cm]{geometry}
\setmainfont{Noto Serif}
\setsansfont{Noto Sans}
\setmonofont{Noto Mono}
\setCJKmainfont {Noto Serif CJK SC}
% 设置正文罗马族的CJK字体，影响\rmfamily和\textrm 的字体
\setCJKsansfont {Noto Sans SC Regular}
% 设置正文无衬线族的CJK字体，影响\sffamily和\textsf 的字体
\setCJKmonofont {Noto Sans Mono CJK SC}
% 设置正文等宽族的CJK字体，影响\ttfamily 和 \texttt 的字体
\setenumerate[0]{label=(\Alph*)}
\ctexset{
	chapter/name = {第,章},
	chapter/number = \arabic{chapter},
	chapter/numberformat = \color{blue}\zihao{0}\emph,
	section/name = {},
	section/number = \Roman{section},
	%	autoindent=false, %禁用自动调整功能
}
\setlist{nosep}

\CatchFileDef{\headfull}{.git/HEAD.}{}
\StrGobbleRight{\headfull}{1}[\head]
\StrBehind[2]{\head}{/}[\branch]
\IfFileExists{.git/refs/heads/\branch.}{%
	\CatchFileDef{\commit}{.git/refs/heads/\branch.}{}}{%
	\newcommand{\commit}{\dots~(in \emph{packed-refs})}}

\begin{document}%内容开始

\title{\sffamily {\Huge 业余无线电操作技能考核攻略本\\（类别：C）}} %书名
\author{\large \texttt{BG7XTQ（原JJ1DSB）编著}} %作者
\date{\large{编译日期：\texttt \today}\\版本：\texttt{\commit}在分支\texttt{\branch}上} %更新日期
	
\maketitle%标题页
	
%献辞
\thispagestyle{empty}
\vfil
\ \\
\vspace{15em}
\begin{center}
	{\Large 献给Bill。}
\end{center}
%献辞

\newpage
	
\tableofcontents%目录
	
% 前言
%\chapter*{前言}
	
% 编著者的话
\chapter*{编著者的话}

无线电资源是全人类共同的财产。提到无线电，我们再熟悉不过的是日常生活中的手机和Wi-Fi，在军事上，人们利用无线电控制导弹、飞机，%用来杀人
在救险活动上，人们利用无线电辅助实施灾害时的救援，%用来救人
在业余无线电领域，爱好者们互相通信，以提高技能，同时学习新知识。

笔者原本对业余无线电一无所知，因为精通业余无线电的朋友的介绍，才逐渐开始对其有所了解。在日本留学期间，笔者考取了日本的操作证书和电台执照，建立了第一个自己的业余无线电台，开始了业余无线电爱好者的旅途。

在归国后，笔者通过业余无线电操作证考试拿到了A类的操作证。在操作证考试应试学习过程中，笔者深深感到，国内现有的操作证考试应试书籍对于很多小学生读者来说，缺乏细致的解释，题目里的术语艰涩难懂，计算题不知道如何计算，用这些书籍学习的读者，想必难以通过操作证考试。在这样的背景下，笔者萌生了撰写一本老少皆能读懂的操作证考试的应试书籍的想法。

本书在写作过程中，为了让业余无线电知识几乎完全不了解的初学者也能读懂，笔者经过了反复的推敲，尽可能的把复杂的业余无线电知识简单易懂地展现给读者们。本书在解题的过程中，适当地介绍相关的术语，并把重点难点用加粗的字体标出，方便应试者快速记忆概念、理解计算方法。

希望本书能帮助您顺利通过考试。

%咋的，编不出来了？

	
% 第一章 无线电的相关法律 开始
%\include{ham_c_1}
% 第一章 无线电的相关法律 结束
	
% 第二章 无线电通信的方法 开始
%\include{ham_c_2}
% 第二章 无线电通信的方法 结束
	
% 第三章 无线电系统理论 开始
%\include{ham_c_3}
% 第三章 无线电系统理论 结束
	
% 第四章 无线电的安全性 开始
%\include{ham_c_4}
% 第四章 无线电的安全性 结束
	
% 第五章 避免和解决射频干扰 开始
%\include{ham_c_5}
% 第五章 避免和解决射频干扰 结束


\noindent\textbf{试题：}我国现行法律体系中专门针对无线电管理的最高法律文件及其立法机关是：
\begin{enumerate}[leftmargin=3em]
\item 中华人民共和国无线电管理条例，国务院和中央军委
\item 中华人民共和国无线电管理办法，工业和信息化部
\item 中华人民共和国电信条例，国务院
\item 中华人民共和国业余无线电台管理办法，工业和信息化部
\end{enumerate}
\noindent\textbf{解说：}\textbf{}。\\\noindent\textbf{答案：}

\bigskip




\noindent\textbf{试题：}我国现行法律体系中专门针对业余无线电台管理的最高法律文件及其立法机关是：
\begin{enumerate}[leftmargin=3em]
\item 业余无线电台管理办法，工业和信息化部
\item 个人业余无线电台管理暂行办法，国家体委和国家无委
\item 业余无线电台管理暂行规定，国家体委和国家无委
\item 中华人民共和国电信条例，国务院
\end{enumerate}
\noindent\textbf{解说：}\textbf{}。\\\noindent\textbf{答案：}

\bigskip




\noindent\textbf{试题：}我国的无线电主管部门是：
\begin{enumerate}[leftmargin=3em]
\item 各级无线电管理机构
\item 各级体育管理机构
\item 各地业余无线电协会
\item 各地电信管理局
\end{enumerate}
\noindent\textbf{解说：}\textbf{}。\\\noindent\textbf{答案：}

\bigskip




\noindent\textbf{试题：}我国依法负责对业余无线电台实施监督管理的机构是：
\begin{enumerate}[leftmargin=3em]
\item 国家无线电管理机构和地方无线电管理机构
\item 在国家或地方民政部门注册的业余无线电协会
\item 国家体育管理机构和地方体育管理机构
\item 国家和地方公安部门
\end{enumerate}
\noindent\textbf{解说：}\textbf{}。\\\noindent\textbf{答案：}

\bigskip




\noindent\textbf{试题：}《业余无线电台管理办法》所说的“地方无线电管理机构”指的是：
\begin{enumerate}[leftmargin=3em]
\item 省、自治区、直辖市无线电管理机构
\item 地方业余无线电协会或者类似组织机构
\item 地市县（区）及以下各级无线电管理机构
\item 各地方与无线电设备生产销售和无线电应用有关的行政管理机构
\end{enumerate}
\noindent\textbf{解说：}\textbf{}。\\\noindent\textbf{答案：}

\bigskip




\noindent\textbf{试题：}国家鼓励和支持业余无线电台开展下列活动：
\begin{enumerate}[leftmargin=3em]
\item 无线电通信技术研究、普及活动以及突发重大自然灾害等紧急情况下的应急通信活动
\item 休闲娱乐性交谈
\item 机动车辆行车服务性通信活动
\item 作为日常公益活动的通信工具
\end{enumerate}
\noindent\textbf{解说：}\textbf{}。\\\noindent\textbf{答案：}

\bigskip




\noindent\textbf{试题：}关于业余电台管理的正确说法是：
\begin{enumerate}[leftmargin=3em]
\item 依法设置的业余无线电台受国家法律保护
\item 业余无线电爱好者的一切行为都受国家法律保护
\item 通过法律手段限制业余无线电台的设置
\item 在业余电台与其他业务电台遇到干扰纠纷时无条件优先保护其他业务电台
\end{enumerate}
\noindent\textbf{解说：}\textbf{}。\\\noindent\textbf{答案：}

\bigskip




\noindent\textbf{试题：}无线电频率的使用必须得到各级无线电管理机构的批准，基本依据是“无线电频谱资源属于国家所有”，出自于下列法律的第二百五十二条： 
\begin{enumerate}[leftmargin=3em]
\item 中华人民共和国民法典
\item 中华人民共和国宪法
\item 中华人民共和国刑法
\item 中华人民共和国电信法
\end{enumerate}
\noindent\textbf{解说：}\textbf{}。\\\noindent\textbf{答案：}

\bigskip




\noindent\textbf{试题：}我国对无线电管理术语“业余业务”、“卫星业余业务”和“业余无线电台”做出具体定义的法规文件是 
\begin{enumerate}[leftmargin=3em]
\item 中华人民共和国无线电频率划分规定
\item 中华人民共和国无线电管理条例
\item 中华人民共和国电信条例
\item 无线电台执照管理规定
\end{enumerate}
\noindent\textbf{解说：}\textbf{}。\\\noindent\textbf{答案：}

\bigskip




\noindent\textbf{试题：}业余电台的法定用途为：
\begin{enumerate}[leftmargin=3em]
\item 供业余无线电爱好者进行自我训练、相互通信和技术研究
\item 供公民在业余时间进行与个人生活事务有关的通信
\item 供公民在业余时间进行休闲娱乐
\item 供私家车主或者相应组织作为行车安全保障和途中消遣工具
\end{enumerate}
\noindent\textbf{解说：}\textbf{}。\\\noindent\textbf{答案：}

\bigskip




\noindent\textbf{试题：}无线电业余业务是供业余无线电爱好者作下列用途的无线电通信业务：
\begin{enumerate}[leftmargin=3em]
\item 自我训练、相互通信和技术研究
\item 救灾抢险、车队联络和技术学习
\item 娱乐休闲、报告路况和公益服务
\item 技术教学、民兵训练和公益通信
\end{enumerate}
\noindent\textbf{解说：}\textbf{}。\\\noindent\textbf{答案：}

\bigskip




\noindent\textbf{试题：}关于无线电通信的正确说法：
\begin{enumerate}[leftmargin=3em]
\item 无线电通信是指利用无线电波进行的符号、信号、文字、图像、声音或其他信息的传输、发射或接收。
\item 无线电通信包括利用光在内的所有电磁波所进行的各种通信
\item 利用无线电波进行的符号、信号、文字、图像、声音以外的信息传输不属于无线电通信
\item 产生无线电波并用其加热属于无线电通信的一种应用
\end{enumerate}
\noindent\textbf{解说：}\textbf{}。\\\noindent\textbf{答案：}

\bigskip




\noindent\textbf{试题：}无线电波是指：
\begin{enumerate}[leftmargin=3em]
\item 频率为3,000GHz以下的在空间传播的电磁波
\item 频率为3,000GHz以下的所有电磁波
\item 频率为30 Hz至30GHz的在空间传播的电磁波
\item 频率为3,000 Hz至3,000 MHz的电磁波
\end{enumerate}
\noindent\textbf{解说：}\textbf{}。\\\noindent\textbf{答案：}

\bigskip




\noindent\textbf{试题：}个人申请设置具有发信功能的业余无线电台的年龄条件是：
\begin{enumerate}[leftmargin=3em]
\item 年满十八周岁
\item 年满十六周岁
\item 年满十四周岁
\item 具备《业余无线电台操作证书》者申请设置业余无线电台不受年龄限制
\end{enumerate}
\noindent\textbf{解说：}\textbf{}。\\\noindent\textbf{答案：}

\bigskip




\noindent\textbf{试题：}申请设置业余无线电台应当具备的条件有：
\begin{enumerate}[leftmargin=3em]
\item 熟悉无线电管理规定、具备国家规定的操作技术能力、发射设备符合国家技术标准、法律和行政法规规定的其他条件
\item 加入指定协会、具备当地无线电管理机构规定的操作技术能力、发射设备符合国家技术标准、法律和行政法规规定的其他条件
\item 熟悉无线电管理规定、具备国家规定的操作技术能力、发射设备符合国家技术标准、当地无线电管理机构委托的受理机构设置的其他条件
\item 熟悉无线电管理规定、具备当地无线电管理机构委托的考试机构设置的操作技术能力标准、发射设备符合国家技术标准、法律和行政法规规定的其他条件
\end{enumerate}
\noindent\textbf{解说：}\textbf{}。\\\noindent\textbf{答案：}

\bigskip




\noindent\textbf{试题：}使用业余无线电台应当具备的条件有：
\begin{enumerate}[leftmargin=3em]
\item 熟悉无线电管理规定、具备国家规定的操作技术能力并取得相应操作技术能力证明
\item 使用具有发信功能的业余无线电台的，应当年满十八周岁
\item 具备国家或地方无线电管理机构核发的业余无线电台执照
\item 熟悉无线电管理规定、实际上具备国家规定的操作技术能力但不必需取得相应的证明
\end{enumerate}
\noindent\textbf{解说：}\textbf{}。\\\noindent\textbf{答案：}

\bigskip




\noindent\textbf{试题：}按照《业余电台管理办法》规定，申请设置使用配备有多台业余无线电发射设备的业余无线电台，应该：
\begin{enumerate}[leftmargin=3em]
\item 视为一个业余电台，指配一个电台呼号，但所有设备均应经过核定并将参数载入电台执照
\item 视为一个业余电台，指配一个电台呼号，其中只需有一台设备加以核定并将参数载入电台执照
\item 每台设备视为一个业余电台，各指配一个电台呼号，并都应经过核定并将参数载入电台执照
\item 视为一个业余电台，指配一个电台呼号，每个频段选择一台设备加以核定并将参数载入电台执照
\end{enumerate}
\noindent\textbf{解说：}\textbf{}。\\\noindent\textbf{答案：}

\bigskip




\noindent\textbf{试题：}申请设置下列业余无线电台时应在《业余无线电台设置（变更）申请表》 的“台站种类”选择“特殊”类：
\begin{enumerate}[leftmargin=3em]
\item 中继台、信标台、空间台
\item 移动操作的车载台
\item 用于业余卫星通信的地面业余无线电台
\item 需要到外地移动操作的手持台
\end{enumerate}
\noindent\textbf{解说：}\textbf{}。\\\noindent\textbf{答案：}

\bigskip




\noindent\textbf{试题：}申请设置信标台、空间台和技术参数需要超出管理办法规定的特殊业余电台的办法为：
\begin{enumerate}[leftmargin=3em]
\item 在《业余无线电台设置（变更）申请表》 的“台站种类”选择“特殊”类，由地方无线电管理机构受理和初审后交国家无线电管理机构审批
\item 先按设置一般业余电台的办法申请，然后再到本地无线电管理机构办理变更执照核定内容
\item 按照设置一般业余电台的办法申请即可，然后根据需要操作就可以
\item 必须由地方业余无线电协会作为申请单位，经本地无线电管理机构办理批准设台
\end{enumerate}
\noindent\textbf{解说：}\textbf{}。\\\noindent\textbf{答案：}

\bigskip




\noindent\textbf{试题：}设置通信范围涉及两个以上的省、自治区、直辖市或者涉及境外的一般业余无线电台，审批机构是下列中：
\begin{enumerate}[leftmargin=3em]
\item 国家无线电管理机构或其委托的设台地的地方无线电管理机构
\item 设台地地方无线电管理机构
\item 国家无线电管理机构委托的设台地地方无线电民间组织
\item 设台地的地方无线电民间组织
\end{enumerate}
\noindent\textbf{解说：}\textbf{}。\\\noindent\textbf{答案：}

\bigskip




\noindent\textbf{试题：}按照在省、自治区、直辖市范围内通信所申请设置的业余无线电台，如想要将通信范围扩大至涉及两个以上的省、自治区、直辖市或者涉及境外，或者要到设台地以外进行异地发射操作，须办理下列手续：
\begin{enumerate}[leftmargin=3em]
\item 事先向核发执照的无线电管理机构申请办理变更手续，按相关流程经国家无线电管理机构或其委托的设台地的地方无线电管理机构批准后，换发业余无线电台执照
\item 反正已经有了电台执照，可先扩大操作起来，等执照有效期届满时再申请办理变更手续，换发业余无线电台执照
\item 只要不会被发现，可以不申请办理变更手续，悄悄越限操作
\item 反正已经有了电台执照，只需向核发执照的无线电管理机构通报变更情况即可，不必申请办理变更和换发执照
\end{enumerate}
\noindent\textbf{解说：}\textbf{}。\\\noindent\textbf{答案：}

\bigskip




\noindent\textbf{试题：}业余无线电台执照有效期届满后需要继续使用的，应当在下列期限内向核发执照的无线电管理机构申请办理延续手续：
\begin{enumerate}[leftmargin=3em]
\item 有效期届满一个月前
\item 有效期届满二十天前
\item 有效期届满一个月之内
\item 有效期届满三个月之内
\end{enumerate}
\noindent\textbf{解说：}\textbf{}。\\\noindent\textbf{答案：}

\bigskip




\noindent\textbf{试题：}因改进或调整业余发射设备使业余无线电台的技术参数超出其业余无线电台执照所核定的范围时，应当办理下列手续：
\begin{enumerate}[leftmargin=3em]
\item 及时向核发执照的无线电管理机构申请办理变更手续，换发业余无线电台执照
\item 等执照有效期届满时向核发执照的无线电管理机构申请办理变更手续，换发业余无线电台执照
\item 只要设备型号和产品序列号没有改变，不必申请办理变更手续
\item 只需及时向核发执照的无线电管理机构通报变更情况，进行备案即可
\end{enumerate}
\noindent\textbf{解说：}\textbf{}。\\\noindent\textbf{答案：}

\bigskip




\noindent\textbf{试题：}终止使用业余无线电台的，应当向下列机构申请注销执照：
\begin{enumerate}[leftmargin=3em]
\item 核发业余无线电台执照的无线电管理机构
\item 国家无线电管理机构
\item 受国家无线电管理机构委托的地方业余无线电民间组织
\item 受国家无线电管理机构委托的全国性业余无线电民间组织
\end{enumerate}
\noindent\textbf{解说：}\textbf{}。\\\noindent\textbf{答案：}

\bigskip




\noindent\textbf{试题：}业余无线电台专用无线电发射设备的重要特征是：
\begin{enumerate}[leftmargin=3em]
\item 发射频率不得超出业余频段
\item 发射频率必须覆盖所有业余频段
\item 发射方式必须包含调频
\item 必须具有数字对讲方式
\end{enumerate}
\noindent\textbf{解说：}\textbf{}。\\\noindent\textbf{答案：}

\bigskip




\noindent\textbf{试题：}业余无线电发射设备的下列指标必须符合国家的相关规定：
\begin{enumerate}[leftmargin=3em]
\item 频率容限和杂散域发射功率
\item 频率调制频偏和调制度
\item 频率容限和带外发射
\item 指配频带和必要带宽
\end{enumerate}
\noindent\textbf{解说：}\textbf{}。\\\noindent\textbf{答案：}

\bigskip




\noindent\textbf{试题：}业余无线电台使用的发射设备必须符合下列条件：
\begin{enumerate}[leftmargin=3em]
\item 商品设备应当具备《无线电发射设备型号核准证》，自制、改装、拼装设备应通过国家相关技术标准的检测
\item 必须具备《无线电发射设备型号核准证》
\item 商品设备应当具备《无线电发射设备型号核准证》，自制、改装、拼装设备不受限制
\item 国产商品设备应当具备《无线电发射设备型号核准证》，国外商品设备符合国际流行技术标准即可
\end{enumerate}
\noindent\textbf{解说：}\textbf{}。\\\noindent\textbf{答案：}

\bigskip




\noindent\textbf{试题：}对业余无线电台专用无线电发射设备进行型号核准的依据为：
\begin{enumerate}[leftmargin=3em]
\item 国家《无线电频率划分规定》中有关无线电发射设备技术指标的规定
\item 地方无线电管理机构制订的技术标准
\item 经国家认证的检测单位所制订的技术标准
\item 国家关于专业无线电通信发射设备的技术标准
\end{enumerate}
\noindent\textbf{解说：}\textbf{}。\\\noindent\textbf{答案：}

\bigskip




\noindent\textbf{试题：}业余无线电台专用无线电发射设备的发射频率必须满足的条件是：
\begin{enumerate}[leftmargin=3em]
\item 发射频率不能超越业余业务或者卫星业余业务频段
\item 发射频率包含所有业余业务或者卫星业余业务频段
\item 发射频率包含至少一个业余业务或者卫星业余业务频段
\item 发射频率可以在业余频段和非业余频段之间选择
\end{enumerate}
\noindent\textbf{解说：}\textbf{}。\\\noindent\textbf{答案：}

\bigskip




\noindent\textbf{试题：}业余电台的无线电发射设备应符国家规定的下列主要技术指标：
\begin{enumerate}[leftmargin=3em]
\item 符合频率容限、符合杂散发射最大允许功率电平
\item 杂散发射不低于最大允许功率电平、电源电压及频率符合国家电网标准、采用标准天线阻抗
\item 杂散发射不低于最大允许功率电平、频率漂移不低于频率容限、电源利用效率满足节能要求
\item 工作频率范围足够宽、杂散发射不低于最大允许功率电平、带宽大于允许最低值
\end{enumerate}
\noindent\textbf{解说：}\textbf{}。\\\noindent\textbf{答案：}

\bigskip




\noindent\textbf{试题：}频率容限是发射设备的重要指标，通常用下述单位来表示：
\begin{enumerate}[leftmargin=3em]
\item 百万分之几（或者赫兹）
\item dB
\item 瓦
\item 百分之几（或者兆赫）
\end{enumerate}
\noindent\textbf{解说：}\textbf{}。\\\noindent\textbf{答案：}

\bigskip




\noindent\textbf{试题：}杂散域发射功率是发射设备的重要指标，通常用下述单位来表示：
\begin{enumerate}[leftmargin=3em]
\item 绝对功率dBm、低于载波发射功率的分贝值dBc、低于PEP发射功率的相对值dB
\item 绝对功率（瓦）
\item 百分之几
\item 千赫（或者赫芝）
\end{enumerate}
\noindent\textbf{解说：}\textbf{}。\\\noindent\textbf{答案：}

\bigskip




\noindent\textbf{试题：}辐射（radiation）是指任何源的能量流以无线电波的形式向外发出。正确的说法是：
\begin{enumerate}[leftmargin=3em]
\item 闪电产生的电磁波干扰是一种辐射
\item 沿电源线窜入接收机的差模干扰是一种辐射
\item 射频电路中变压器内磁芯里的磁场是一种辐射
\item 射频电路中电容器内极板间的电场是一种辐射
\end{enumerate}
\noindent\textbf{解说：}\textbf{}。\\\noindent\textbf{答案：}

\bigskip




\noindent\textbf{试题：}发射（emission）是指：由无线电发信电台产生的辐射或辐射产物。正确的说法是：
\begin{enumerate}[leftmargin=3em]
\item 业余电台向周围发送的杂散产物是一种发射
\item 无线电接收机本地振荡器辐射的能量是一种发射
\item 医用高频电疗机向周围发送的无线电波能量是一种发射
\item 闪电产生的电磁波干扰是一种发射
\end{enumerate}
\noindent\textbf{解说：}\textbf{}。\\\noindent\textbf{答案：}

\bigskip




\noindent\textbf{试题：}杂散发射是指必要带宽之外的一个或多个频率的发射，其发射电平可降低而不致影响相应信息的传输。一台发射机，工作频率为145.000MHz，但在435.000MHz的频率上也有发射。这种发射属于：
\begin{enumerate}[leftmargin=3em]
\item 杂散发射
\item 带外发射
\item 谐波发射
\item 带内发射
\end{enumerate}
\noindent\textbf{解说：}\textbf{}。\\\noindent\textbf{答案：}

\bigskip




\noindent\textbf{试题：}业余无线电专用发射设备必须满足的主要技术指标要求包括：
\begin{enumerate}[leftmargin=3em]
\item 频率容限和杂散辐射不超过限值，发射频率不超出国家规定的业余频率
\item 频率容限不低于限值，杂散辐射不超过限值，发射频率不超出国家规定的业余频率
\item 频率容限和杂散辐射不超过限值，发射频率包括国家规定的业余频率
\item 发射功率不低于功率限额，输出阻抗符合工业标准
\end{enumerate}
\noindent\textbf{解说：}\textbf{}。\\\noindent\textbf{答案：}

\bigskip




\noindent\textbf{试题：}业余无线电台使用的频率应当符合下述规定：
\begin{enumerate}[leftmargin=3em]
\item 《中华人民共和国无线电频率划分规定》
\item ITU《无线电规则》第IV节“频率划分表”
\item IARU三区“频率规划”
\item 一般业余无线电书籍所叙述的频率
\end{enumerate}
\noindent\textbf{解说：}\textbf{}。\\\noindent\textbf{答案：}

\bigskip




\noindent\textbf{试题：}业余无线电台在业余业务、卫星业余业务作为次要业务使用的频率或者与其他主要业务共同使用的频率上发射操作时，应当注意：
\begin{enumerate}[leftmargin=3em]
\item 遵守无线电管理机构对该频率的使用规定
\item 首先守听频率是否已由其他业务电台占用，如听不到，即可按照先来先用的原则放心使用
\item 只要遵守了《中华人民共和国无线电频率划分规定》的有关规定即可放心使用
\item 可以任意使用，但在遇到其他业务电台使用时要主动避让
\end{enumerate}
\noindent\textbf{解说：}\textbf{}。\\\noindent\textbf{答案：}

\bigskip




\noindent\textbf{试题：}关于业余频率的使用，正确的叙述是：
\begin{enumerate}[leftmargin=3em]
\item 业余无线电台在无线电管理机构核准其使用的频段内，享有平等的频率使用权
\item 任何业余无线电台在任何频段都享有平等的频率使用权
\item 业余无线电台在无线电管理机构核准其使用的频段内，不同类别的业余电台享有不同优先程度的频率使用权
\item 依法成立的地方业余无线电民间组织的业余电台，在其常用的台网频率上享有比其他个人设置的业余电台优先的使用权
\end{enumerate}
\noindent\textbf{解说：}\textbf{}。\\\noindent\textbf{答案：}

\bigskip




\noindent\textbf{试题：}在无线电管理中，由国家将某个特定的频带列入频率划分表，规定该频带可在指定的条件下供业余业余业务或者卫星业余业务使用，这个过程称为：
\begin{enumerate}[leftmargin=3em]
\item 划分
\item 分配
\item 指配
\item 授权
\end{enumerate}
\noindent\textbf{解说：}\textbf{}。\\\noindent\textbf{答案：}

\bigskip




\noindent\textbf{试题：}在无线电管理中，将无线电频率或频道规定由一个或多个部门，在指定的区域内供地面或空间无线电通信业务在指定条件下使用，这个过程称为：
\begin{enumerate}[leftmargin=3em]
\item 分配
\item 划分
\item 指配
\item 授权
\end{enumerate}
\noindent\textbf{解说：}\textbf{}。\\\noindent\textbf{答案：}

\bigskip




\noindent\textbf{试题：}在无线电管理中，将无线电频率或频道批准给具体的业余无线电台在规定条件下使用，这个过程称为：
\begin{enumerate}[leftmargin=3em]
\item 指配
\item 划分
\item 分配
\item 授权
\end{enumerate}
\noindent\textbf{解说：}\textbf{}。\\\noindent\textbf{答案：}

\bigskip




\noindent\textbf{试题：}必要带宽（necessary bandwidth）是指：对给定的发射类别而言，其恰好足以保证在相应速率及在指定条件下具有所要求质量的信息传输的所需带宽。业余电台单边带话音通信SSB、低速莫尔斯电码通信CW、调频话音通信FM和残余边带业余电视VSB ATV的必要带宽分别是： 
\begin{enumerate}[leftmargin=3em]
\item 3000Hz、400Hz、12.5kHz、5MH以上
\item 3000Hz、400Hz、5MHz以上、12.5kHz
\item 5MHz、3000Hz、400Hz、12.5kHz
\item 12.5kHz、5MHz以上、400Hz、2700Hz
\end{enumerate}
\noindent\textbf{解说：}\textbf{}。\\\noindent\textbf{答案：}

\bigskip




\noindent\textbf{试题：}带外发射是指由于调制过程而产生的、刚超出必要带宽的一个或多个频率的发射，但杂散发射除外。如果单边带发射机在发送由两个音频信号组成的双音信号时，产生的三阶互调分量刚超出了必要带宽3kHz。这种分量的发射属于：
\begin{enumerate}[leftmargin=3em]
\item 带外发射
\item 杂散发射
\item 谐波发射
\item 带内发射
\end{enumerate}
\noindent\textbf{解说：}\textbf{}。\\\noindent\textbf{答案：}

\bigskip




\noindent\textbf{试题：}在频率划分表中，一个频带被标明划分给多种业务时，这些业务被分为下述类别：
\begin{enumerate}[leftmargin=3em]
\item 主要业务和次要业务
\item 业余业务和非业余业务
\item 民用业务和军用业务
\item 安全业务和一般业务
\end{enumerate}
\noindent\textbf{解说：}\textbf{}。\\\noindent\textbf{答案：}

\bigskip




\noindent\textbf{试题：}在频率划分表中，当一个频段划分给业余业务或卫星业余业务和多个其他业务，并且业余业务和卫星业余业务作为次要业务时，业余无线电台应该遵循的规则是：
\begin{enumerate}[leftmargin=3em]
\item 不得对主要业务电台产生有害干扰
\item 可要求保护不受来自主要业务电台的有害干扰
\item 不得对来自同一业务或其他次要业务电台的有害干扰提出保护要求
\item 容许因设备技术问题对主要业务电台产生短时间有害干扰
\end{enumerate}
\noindent\textbf{解说：}\textbf{}。\\\noindent\textbf{答案：}

\bigskip




\noindent\textbf{试题：}在频率划分表中，当一个频段划分给业余业务或卫星业余业务和多个其他业务，并且业余业务和卫星业余业务作为次要业务时，业余无线电台遵循的规则是：
\begin{enumerate}[leftmargin=3em]
\item 不得对来自主要业务电台的有害干扰提出保护要求
\item 可要求保护不受来自主要业务电台的有害干扰
\item 不得对来自同一业务或其他次要业务电台的有害干扰提出保护要求
\item 容许因设备技术问题对主要业务电台产生短时间有害干扰
\end{enumerate}
\noindent\textbf{解说：}\textbf{}。\\\noindent\textbf{答案：}

\bigskip




\noindent\textbf{试题：}在频率划分表中，当一个频段划分给业余业务或卫星业余业务和多个其他业务，并且业余业务和卫星业余业务作为次要业务时，业余无线电台遵循的规则是：
\begin{enumerate}[leftmargin=3em]
\item 可要求保护不受来自同一业务或其他次要业务电台的有害干扰
\item 可要求保护不受来自主要业务电台的有害干扰
\item 不得对来自同一业务或其他次要业务电台的有害干扰提出保护要求
\item 容许因设备技术问题对主要业务电台产生短时间有害干扰
\end{enumerate}
\noindent\textbf{解说：}\textbf{}。\\\noindent\textbf{答案：}

\bigskip




\noindent\textbf{试题：}分配给业余业务的某频段的频率下限为F1，业余电台实际可以工作的发信频率应为：
\begin{enumerate}[leftmargin=3em]
\item F1+信号下边带的频率宽度
\item F1
\item F1-信号下边带的频率宽度
\item F1-2×信号下边带的频率宽度
\end{enumerate}
\noindent\textbf{解说：}\textbf{}。\\\noindent\textbf{答案：}

\bigskip




\noindent\textbf{试题：}分配给业余业务的某频段的频率上限为F2，业余电台实际可以工作的发信频率应为：
\begin{enumerate}[leftmargin=3em]
\item F2-信号上边带的频率宽度
\item F2
\item F2+信号上边带的频率宽度
\item F2-2×信号上边带的频率宽度
\end{enumerate}
\noindent\textbf{解说：}\textbf{}。\\\noindent\textbf{答案：}

\bigskip




\noindent\textbf{试题：}为了满足我国《无线电频率划分规定》“电台的技术特性”关于无线电通信“把带宽保持在技术状态和该项业务的性质所允许的最低值上”的要求，业余电台操作者应了解各种通信方式的必要带宽。决定必要带宽的因素是：
\begin{enumerate}[leftmargin=3em]
\item 所要传输的信息速率越高、整个通信系统的噪声干扰越大，必要带宽越宽
\item 发射设备的功率越大，必要带宽越宽
\item 接收设备的灵敏度越高，必要带宽越宽
\item 通信距离越近，必要带宽越宽
\end{enumerate}
\noindent\textbf{解说：}\textbf{}。\\\noindent\textbf{答案：}

\bigskip




\noindent\textbf{试题：}我国分配给业余业务和卫星业余业务专用的频段有：
\begin{enumerate}[leftmargin=3em]
\item 7MHz、14MHz、21MHz、28MHz、47GHz频段
\item 7MHz、14MHz、21MHz、28MHz、144MHz频段
\item 3.5MHz、14MHz、21MHz、28MHz、10GHz频段
\item 7MHz、14MHz、28MHz、144MHz、430MHz频段
\end{enumerate}
\noindent\textbf{解说：}\textbf{}。\\\noindent\textbf{答案：}

\bigskip




\noindent\textbf{试题：}我国分配给业余业务和卫星业余业务与其他业务共用、并且业余业务和卫星业余业务作为主要业务之一的30MHz以下频段有：
\begin{enumerate}[leftmargin=3em]
\item 1.8MHz、3.5MHz、14.25MHz、18.068MHz、24.89MHz频段
\item 3.5MHz、7MHz、14.25MHz、21MHz、24.89MHz频段
\item 3.5MHz、10.1MHz、14.25MHz、18.068MHz、29.7MHz频段
\item 1.8MHz、10.1MHz、14.25MHz、18.068MHz、21.45MHz频段
\end{enumerate}
\noindent\textbf{解说：}\textbf{}。\\\noindent\textbf{答案：}

\bigskip




\noindent\textbf{试题：}[Q]我国在VHF和UHF范围内分配给业余业务和卫星业余业务与其他业务共用并设业务类别为主要业务与次要业务，以下那些频率分配给业余业务和卫星业余业务与其他业务共用并且业余业务和卫星业余业务作为主要业务： 
\begin{enumerate}[leftmargin=3em]
\item 50MHz、144MHz
\item 144MHz、430MHz
\item 50MHz、430MHz
\item 220MHz、430MHz
\end{enumerate}
\noindent\textbf{解说：}\textbf{}。\\\noindent\textbf{答案：}

\bigskip




\noindent\textbf{试题：}我国分配给业余业务和卫星业余业务与其他业务共用、并且业余业务和卫星业余业务作为唯一主要业务的频段的个数以及在3GHz以下的该类频段分别为：
\begin{enumerate}[leftmargin=3em]
\item 3个，144-146MHz
\item 4个，7.0-7.2MHz
\item 5个，50-54MHz
\item 5个，28-29.7MHz
\end{enumerate}
\noindent\textbf{解说：}\textbf{}。\\\noindent\textbf{答案：}

\bigskip




\noindent\textbf{试题：}我国分配给业余业务和卫星业余业务与其他业务共用、并且业余业务和卫星业余业务作为次要业务的1200MHz以下频段有：
\begin{enumerate}[leftmargin=3em]
\item 135.7kHz、10.1MHz、430MHz
\item 3.5MHz、7MHz、50MHz
\item 3.5MHz、18.068MHz、144MHz
\item 10.1MHz、24.89MHz、430MHz
\end{enumerate}
\noindent\textbf{解说：}\textbf{}。\\\noindent\textbf{答案：}

\bigskip




\noindent\textbf{试题：}俗称的6米业余波段的频率范围以及业余业务和卫星业余业务的使用状态分别为：
\begin{enumerate}[leftmargin=3em]
\item 50-54MHz，主要业务
\item 50-52MHz，次要业务
\item 51-54MHz，专用
\item 52-56MHz，次要业务
\end{enumerate}
\noindent\textbf{解说：}\textbf{}。\\\noindent\textbf{答案：}

\bigskip




\noindent\textbf{试题：}俗称的2米业余波段的频率范围以及我国业余业务和卫星业余业务的使用状态分别为：
\begin{enumerate}[leftmargin=3em]
\item 144-148MHz；其中144-146MHz为唯一主要业务，146-148MHz为与其他业务共同作为主要业务
\item 144-146MHz；专用
\item 144-148MHz；其中144-146MHz为专用，146-148MHz为次要业务
\item 144-148MHz；次要业务
\end{enumerate}
\noindent\textbf{解说：}\textbf{}。\\\noindent\textbf{答案：}

\bigskip




\noindent\textbf{试题：}俗称的0.7米业余波段的频率范围以及业余业务和卫星业余业务的使用状态分别为：
\begin{enumerate}[leftmargin=3em]
\item 430-440MHz，次要业务
\item 430-440MHz，主要业务
\item 430-440MHz，专用
\item 420-470MHz，次要业务
\end{enumerate}
\noindent\textbf{解说：}\textbf{}。\\\noindent\textbf{答案：}

\bigskip




\noindent\textbf{试题：}在我国和多数其他国家的频率分配中，业余业务在430-440MHz频段中作为次要业务与其他业务共用。这个频段中我国分配的主要业务是：
\begin{enumerate}[leftmargin=3em]
\item 无线电定位和航空无线电导航
\item 固定业务
\item 移动业务
\item 水上移动和航空移动
\end{enumerate}
\noindent\textbf{解说：}\textbf{}。\\\noindent\textbf{答案：}

\bigskip




\noindent\textbf{试题：}VHF业余无线电台在144MHz频段进行本地联络时应避免占用的频率为：
\begin{enumerate}[leftmargin=3em]
\item 144-144.035MHz和145.8-146MHz
\item 144.035-145.8MHz
\item 144.050-144.053MHz和145.100-145.750MHz
\item 144.035-144.053MHz和145.550-145.750MHz
\end{enumerate}
\noindent\textbf{解说：}\textbf{}。\\\noindent\textbf{答案：}

\bigskip




\noindent\textbf{试题：}UHF业余无线电台在430MHz频段进行本地联络时应避免占用的频率为：
\begin{enumerate}[leftmargin=3em]
\item 431.9-432.240MHz和435-438MHz
\item 430-431.9MHz和432.240-435MHz
\item 431-432MHz和438-440MHz
\item 430-431.2MHz和435-436MHz
\end{enumerate}
\noindent\textbf{解说：}\textbf{}。\\\noindent\textbf{答案：}

\bigskip




\noindent\textbf{试题：}430MHz业余频段中留给业余卫星通信使用，话音及其他通信方式不应占用的频率段为：
\begin{enumerate}[leftmargin=3em]
\item 435MHz至438MHz
\item 432MHz至434MHz
\item 438MHz至439MHz
\item 433MHz至435MHz
\end{enumerate}
\noindent\textbf{解说：}\textbf{}。\\\noindent\textbf{答案：}

\bigskip




\noindent\textbf{试题：}144MHz业余频段中留给业余卫星通信使用，话音及其他通信方式不应占用的频率段为：
\begin{enumerate}[leftmargin=3em]
\item 145.8MHz至146MHz
\item 144.8MHz至145MHz
\item 144.2MHz至144.5MHz
\item 145.4MHz至144.6MHz
\end{enumerate}
\noindent\textbf{解说：}\textbf{}。\\\noindent\textbf{答案：}

\bigskip




\noindent\textbf{试题：}我国无线电频率划分表划分给业余业务使用的最接近无线宽带WiFi频率的频带为2,300-2,450MHz，属于无线电频谱的下列频带（波段）：
\begin{enumerate}[leftmargin=3em]
\item 特高频（分米波）
\item 甚高频（米波）
\item 高频（短波）
\item 超高频（厘米波）
\end{enumerate}
\noindent\textbf{解说：}\textbf{}。\\\noindent\textbf{答案：}

\bigskip




\noindent\textbf{试题：}我国无线电频率划分表划分给业余业务使用的最接近无线宽带WiFi频率的频带为2,300-2,450MHz，属于无线电频谱的下列频带（波段）：
\begin{enumerate}[leftmargin=3em]
\item UHF
\item VHF
\item HF
\item SHF
\end{enumerate}
\noindent\textbf{解说：}\textbf{}。\\\noindent\textbf{答案：}

\bigskip




\noindent\textbf{试题：}我国无线电频率划分表划分给业余业务使用的最接近C波段卫星电视广播频率的频带是5.650-5.850GHz，属于无线电频谱的下列频带（波段）：
\begin{enumerate}[leftmargin=3em]
\item 超高频（厘米波）
\item 特高频（分米波）
\item 极高频（毫米波）
\item 甚高频（米波）
\end{enumerate}
\noindent\textbf{解说：}\textbf{}。\\\noindent\textbf{答案：}

\bigskip




\noindent\textbf{试题：}我国无线电频率划分表划分给业余业务使用的最接近C波段卫星电视广播频率的频带是5.650-5.850GHz，属于无线电频谱的下列频带（波段）：
\begin{enumerate}[leftmargin=3em]
\item SHF
\item UHF
\item EHF
\item VHF
\end{enumerate}
\noindent\textbf{解说：}\textbf{}。\\\noindent\textbf{答案：}

\bigskip




\noindent\textbf{试题：}我国无线电频率划分表划分给业余业务使用的最接近Ku波段卫星电视广播频率的频带10-10.5GHz，属于无线电频谱的下列频带（波段）：
\begin{enumerate}[leftmargin=3em]
\item 超高频（厘米波）
\item 特高频（分米波）
\item 极高频（毫米波）
\item 甚高频（米波）
\end{enumerate}
\noindent\textbf{解说：}\textbf{}。\\\noindent\textbf{答案：}

\bigskip




\noindent\textbf{试题：}我国无线电频率划分表划分给业余业务使用的最接近Ku波段卫星电视广播频率的频带10-10.5GHz，属于无线电频谱的下列频带（波段）：
\begin{enumerate}[leftmargin=3em]
\item SHF
\item UHF
\item EHF
\item VHF
\end{enumerate}
\noindent\textbf{解说：}\textbf{}。\\\noindent\textbf{答案：}

\bigskip




\noindent\textbf{试题：}我国无线电频率划分表划分给业余业务使用的最高频带为241GHz-250GHz，属于无线电频谱的下列频带（波段）：
\begin{enumerate}[leftmargin=3em]
\item 极高频（毫米波）
\item 超高频（厘米波）
\item 至高频（丝米波或亚毫米波）
\item 特高频（分米波）
\end{enumerate}
\noindent\textbf{解说：}\textbf{}。\\\noindent\textbf{答案：}

\bigskip




\noindent\textbf{试题：}我国无线电频率划分表划分给业余业务使用的最高频带为241GHz-250GHz，属于无线电频谱的下列频带（波段）：
\begin{enumerate}[leftmargin=3em]
\item EHF
\item SHF
\item THF
\item UHF
\end{enumerate}
\noindent\textbf{解说：}\textbf{}。\\\noindent\textbf{答案：}

\bigskip




\noindent\textbf{试题：}我国无线电频率划分表划分给业余业务使用的最低频带为135.7-137.8 kHz，属于无线电频谱的下列频带（波段）：
\begin{enumerate}[leftmargin=3em]
\item 低频（长波）
\item 甚低频（甚长波）
\item 特低频（特长波）
\item 超低频（超长波）
\end{enumerate}
\noindent\textbf{解说：}\textbf{}。\\\noindent\textbf{答案：}

\bigskip




\noindent\textbf{试题：}我国无线电频率划分表划分给业余业务使用的最低频带为135.7-137.8 kHz，属于无线电频谱的下列频带（波段）：
\begin{enumerate}[leftmargin=3em]
\item LF
\item VLF
\item ULF
\item SLF
\end{enumerate}
\noindent\textbf{解说：}\textbf{}。\\\noindent\textbf{答案：}

\bigskip




\noindent\textbf{试题：}我国无线电频率划分表划分给业余业务使用的1,800kHz-2,000kHz属于无线电频谱的下列频带（波段）：
\begin{enumerate}[leftmargin=3em]
\item 中频（中波）
\item 甚高频（米波）
\item 低频（长波）
\item 高频（短波）
\end{enumerate}
\noindent\textbf{解说：}\textbf{}。\\\noindent\textbf{答案：}

\bigskip




\noindent\textbf{试题：}我国无线电频率划分表划分给业余业务使用的1,800kHz-2,000kHz属于无线电频谱的下列频带（波段）：
\begin{enumerate}[leftmargin=3em]
\item MF
\item HF
\item VHF
\item LF
\end{enumerate}
\noindent\textbf{解说：}\textbf{}。\\\noindent\textbf{答案：}

\bigskip




\noindent\textbf{试题：}我国无线电频率划分表划分给业余业务使用的28MHz-29.7MHz属于无线电频谱的下列频带（波段）：
\begin{enumerate}[leftmargin=3em]
\item 高频（短波）
\item 中频（中波）
\item 甚高频（米波）
\item 低频（长波）
\end{enumerate}
\noindent\textbf{解说：}\textbf{}。\\\noindent\textbf{答案：}

\bigskip




\noindent\textbf{试题：}我国无线电频率划分表划分给业余业务使用的28MHz-29.7MHz属于无线电频谱的下列频带（波段）：
\begin{enumerate}[leftmargin=3em]
\item HF
\item MF
\item VHF
\item LF
\end{enumerate}
\noindent\textbf{解说：}\textbf{}。\\\noindent\textbf{答案：}

\bigskip




\noindent\textbf{试题：}我国无线电频率划分表划分给业余业务使用的50MHz-54MHz属于无线电频谱的下列频带（波段）：
\begin{enumerate}[leftmargin=3em]
\item 甚高频（米波）
\item 高频（短波）
\item 超高频（厘米波）
\item 特高频（分米波）
\end{enumerate}
\noindent\textbf{解说：}\textbf{}。\\\noindent\textbf{答案：}

\bigskip




\noindent\textbf{试题：}我国无线电频率划分表划分给业余业务使用的50MHz-54MHz属于无线电频谱的下列频带（波段）：
\begin{enumerate}[leftmargin=3em]
\item VHF
\item HF
\item SHF
\item UHF
\end{enumerate}
\noindent\textbf{解说：}\textbf{}。\\\noindent\textbf{答案：}

\bigskip




\noindent\textbf{试题：}我国分配给业余业务和卫星业余业务的频率最高的频段及其共用类别是：
\begin{enumerate}[leftmargin=3em]
\item 248-250GHz，唯一主要业务
\item 134-136GHz，唯一主要业务
\item 47-47.2GHz，专用业务
\item 77.5-78GHz，主要业务
\end{enumerate}
\noindent\textbf{解说：}\textbf{}。\\\noindent\textbf{答案：}

\bigskip




\noindent\textbf{试题：}国际业余无线电界把WARC-76增加分配给业余业务和卫星业余业务的三个HF频段俗称为WARC频段，它们的频率范围是：
\begin{enumerate}[leftmargin=3em]
\item 10.1-10.15MHz、18.068-18.168MHz、24.89-24.99MHz
\item 10.068-10.168MHz、18.1-18.15MHz、24.89-24.99MHz
\item 10.1-10.15MHz、18.89-18.99MHz、24.068-24.168MHz
\item 10.89-10.88MHz、18.1-18.15MHz、24.068-24.168MHz
\end{enumerate}
\noindent\textbf{解说：}\textbf{}。\\\noindent\textbf{答案：}

\bigskip




\noindent\textbf{试题：}我国和ITU分配给业余业务和卫星业余业务的属于低频（LF）和中频（MF）范围的频段分别为：
\begin{enumerate}[leftmargin=3em]
\item 我国 135.7-137.8kHz；ITU 137.8-138.6kHz和472-479kHz
\item 我国和ITU均为 135.7-137.8kHz
\item 我国和ITU均为472-479kHz
\item 我国 135.7-137.8kHz；ITU为 137.8-139.6kHz
\end{enumerate}
\noindent\textbf{解说：}\textbf{}。\\\noindent\textbf{答案：}

\bigskip




\noindent\textbf{试题：}俗称的40米业余波段，其在ITU1、2、3区的频率范围以及业余业务和卫星业余业务的使用状态分别为：
\begin{enumerate}[leftmargin=3em]
\item 7.0-7.2MHz、7.0-7.3MHz、7.0-7.2MHz，专用
\item 7.0-7.3MHz、7.0-7.3MHz、7.0-7.3MHz，专用
\item 7.0-7.3MHz、7.0-7.3MHz、7.0-7.2MHz，专用
\item 7.0-7.1MHz、7.0-7.2MHz、7.0-7.3MHz，专用
\end{enumerate}
\noindent\textbf{解说：}\textbf{}。\\\noindent\textbf{答案：}

\bigskip




\noindent\textbf{试题：}俗称的160米业余波段的频率范围以及业余业务和卫星业余业务的使用状态分别为：
\begin{enumerate}[leftmargin=3em]
\item 1800-2000kHz，主要业务
\item 1800-1900kHz，次要业务
\item 1900-2000kHz，主要业务
\item 1700-1900kHz，专用业务
\end{enumerate}
\noindent\textbf{解说：}\textbf{}。\\\noindent\textbf{答案：}

\bigskip




\noindent\textbf{试题：}俗称的80米业余波段的频率范围以及业余业务和卫星业余业务的使用状态分别为：
\begin{enumerate}[leftmargin=3em]
\item 3.5-3.9MHz，主要业务
\item 3.5-3.9MHz，次要业务
\item 3.5-4.0MHz，主要业务
\item 3.5-3.6MHz，专用业务
\end{enumerate}
\noindent\textbf{解说：}\textbf{}。\\\noindent\textbf{答案：}

\bigskip




\noindent\textbf{试题：}俗称的20米业余波段的频率范围以及业余业务和卫星业余业务的使用状态分别为：
\begin{enumerate}[leftmargin=3em]
\item 14-14.25MHz为专用，14.25-14.35为主要业务
\item 14-14.15MHz为专用，14.15-14.25为主要业务
\item 14-14.35MHz为专用，14.35-14.45为主要业务
\item 14-14.35MHz，专用
\end{enumerate}
\noindent\textbf{解说：}\textbf{}。\\\noindent\textbf{答案：}

\bigskip




\noindent\textbf{试题：}俗称的15米业余波段的频率范围以及业余业务和卫星业余业务的使用状态分别为：
\begin{enumerate}[leftmargin=3em]
\item 21-21.45MHz，专用
\item 21-21.45MHz，主要业务
\item 21-21.35MHz，专用
\item 21-21.45MHz，次要业务
\end{enumerate}
\noindent\textbf{解说：}\textbf{}。\\\noindent\textbf{答案：}

\bigskip




\noindent\textbf{试题：}俗称的10米业余波段的频率范围以及业余业务和卫星业余业务的使用状态分别为：
\begin{enumerate}[leftmargin=3em]
\item 28-29.7MHz，专用
\item 28-29.7MHz，主要业务
\item 28-29.6MHz，专用
\item 28-30MHz，次要业务
\end{enumerate}
\noindent\textbf{解说：}\textbf{}。\\\noindent\textbf{答案：}

\bigskip




\noindent\textbf{试题：}商用业余单边带电台的语音信号通道为300-3000Hz，高音频部分衰减很大。如果希望对设备进行改装，以将信号带宽扩展到30-16000Hz，大大提升信号质量。可行的方案是：
\begin{enumerate}[leftmargin=3em]
\item 放弃该目标，不应超越业务性质所允许的必要带宽最低值
\item 配用广播级超高音质话筒
\item 将原单边带信道的晶体滤波器更换为带宽更宽的滤波器
\item 在话筒电路中增加分段可控音频均衡电路以提升话筒信号的高音分量
\end{enumerate}
\noindent\textbf{解说：}\textbf{}。\\\noindent\textbf{答案：}

\bigskip




\noindent\textbf{试题：}短波业余电台应避免在IARU信标工作频率±500Hz的范围内发射电波。这些频率是：
\begin{enumerate}[leftmargin=3em]
\item 14.100MHz、18.110MHz、21.150MHz、24.930MHz、28.200MHz
\item 7.100MHz、10.070MHz、14.100MHz、21.100MHz、28.200MHz
\item 7.150MHz、14.110MHz、18.150MHz、21.150MHz、28.150MHz
\item 14.150MHz、18.100MHz、21.200MHz、24.930MHz、28.200MHz
\end{enumerate}
\noindent\textbf{解说：}\textbf{}。\\\noindent\textbf{答案：}

\bigskip




\noindent\textbf{试题：}不能用于通话的HF业余频段为：
\begin{enumerate}[leftmargin=3em]
\item 10MHz业余频段
\item 18MHz业余频段
\item 14MHz业余频段
\item 1.8MHz业余频段
\end{enumerate}
\noindent\textbf{解说：}\textbf{}。\\\noindent\textbf{答案：}

\bigskip




\noindent\textbf{试题：}我国短波业余电台在7MHz频段进行LSB通话时可以实际占用的频率为：
\begin{enumerate}[leftmargin=3em]
\item 7.030-7.200MHz
\item 7.000-7.100MHz
\item 7.023-7.200MHz
\item 7.000-7.200MHz
\end{enumerate}
\noindent\textbf{解说：}\textbf{}。\\\noindent\textbf{答案：}

\bigskip




\noindent\textbf{试题：}短波业余电台在14MHz频段进行USB通话时可以实际占用的频率为：
\begin{enumerate}[leftmargin=3em]
\item 14.100-14.350MHz
\item 14.030-14.350MHz
\item 14.000-14.250MHz
\item 14.070-14.250MHz
\end{enumerate}
\noindent\textbf{解说：}\textbf{}。\\\noindent\textbf{答案：}

\bigskip




\noindent\textbf{试题：}短波业余电台在18MHz频段进行USB通话时可以实际占用的频率为：
\begin{enumerate}[leftmargin=3em]
\item 18.1105-18.168MHz
\item 18.110-18.170MHz
\item 18.068-18.186MHz
\item 18.1005-18.180MHz
\end{enumerate}
\noindent\textbf{解说：}\textbf{}。\\\noindent\textbf{答案：}

\bigskip




\noindent\textbf{试题：}短波业余电台在21MHz频段进行USB通话时可以实际占用的频率为：
\begin{enumerate}[leftmargin=3em]
\item 21.125-21.45MHz，除去21.1495-21.1505
\item 21.125-21.45MHz
\item 21-21.45MHz
\item 21-21.35MHz
\end{enumerate}
\noindent\textbf{解说：}\textbf{}。\\\noindent\textbf{答案：}

\bigskip




\noindent\textbf{试题：}短波业余电台在24MHz频段进行USB通话时可以实际占用的频率为：
\begin{enumerate}[leftmargin=3em]
\item 24.9305-24.99MHz
\item 24.928-24.988MHz
\item 24.890-24.98MHz
\item 24.9205-24.99MHz
\end{enumerate}
\noindent\textbf{解说：}\textbf{}。\\\noindent\textbf{答案：}

\bigskip




\noindent\textbf{试题：}短波业余电台在29MHz频段进行USB通话时可以实际占用的频率为：
\begin{enumerate}[leftmargin=3em]
\item 28.3-29.3MHz
\item 28-29.7MHz
\item 28.250-29.7MHz
\item 28.2-29.6MHz
\end{enumerate}
\noindent\textbf{解说：}\textbf{}。\\\noindent\textbf{答案：}

\bigskip




\noindent\textbf{试题：}短波业余电台在29MHz频段进行FM通话时可以实际占用的频率为：
\begin{enumerate}[leftmargin=3em]
\item 29.51-29.7MHz
\item 29.3-29.7MHz
\item 28.3-29.510MHz z
\item 28-29.7MHz
\end{enumerate}
\noindent\textbf{解说：}\textbf{}。\\\noindent\textbf{答案：}

\bigskip




\noindent\textbf{试题：}为什么不能在低于1.2GHz的业余频段进行常规的ATV通信？
\begin{enumerate}[leftmargin=3em]
\item ATV通信需占5MHz以上带宽，较低业余频段不足以容纳
\item 多数话音通信集中在较低频段，易对ATV通信产生严重干扰
\item 多数话音通信集中在较低频段，而ATV画面发射时间通常较长，所以要主动避让
\item 频率越高，ATV传输画面的质量越稳定
\end{enumerate}
\noindent\textbf{解说：}\textbf{}。\\\noindent\textbf{答案：}

\bigskip




\noindent\textbf{试题：}按照有关规定，144MHz和430MHz频段业余中继台的上下行频差应分别为：
\begin{enumerate}[leftmargin=3em]
\item 0.6MHz、5MHz。
\item 2MHz、10MHz
\item 5MHz、5MHz
\item 7.5MHz、10MHz
\end{enumerate}
\noindent\textbf{解说：}\textbf{}。\\\noindent\textbf{答案：}

\bigskip




\noindent\textbf{试题：}28MHz业余频段中留给业余卫星通信、话音及其他通信方式不应占用的频率段为：
\begin{enumerate}[leftmargin=3em]
\item 29.3MMz至29.51MHz
\item 28.3MHz至28.61MHz
\item 28.7MHz至28.95MHz
\item 29.15MHz至29.35MHz
\end{enumerate}
\noindent\textbf{解说：}\textbf{}。\\\noindent\textbf{答案：}

\bigskip




\noindent\textbf{试题：}《业余无线电台管理办法》规定业余无线电台设置、正确使用业余无线电台呼号的办法规是：
\begin{enumerate}[leftmargin=3em]
\item 业余无线电台应当在每次通信建立及结束时，主动报出本台呼号，在发射过程中至少每十分钟报出本台呼号一次；对于通信对方，也应使用对方电台的呼号加以标识
\item 业余电台在和熟悉的通信对象联络、已经从信号特征确认双方业余电台身份时，可以省略呼号的发送
\item 业余电台在通信中可以用姓名、代号、适当的别名或者法规定呼号的部分数字和字母代替完整的业余电台呼号作为电台的标识
\item 业余电台在通信中可以用自造的呼号作为无线电管理机构指配的业余电台呼号的补充，一起作为电台的标识
\end{enumerate}
\noindent\textbf{解说：}\textbf{}。\\\noindent\textbf{答案：}

\bigskip




\noindent\textbf{试题：}业余无线电台应当在每次通信建立及结束时，主动报出本台呼号，在发射过程中至少每十分钟报出本台呼号一次。这里的“呼号”是指：
\begin{enumerate}[leftmargin=3em]
\item 完整的电台呼号，如在设台地以外的地点进行异地发射操作，还应在前面加上字母B、操作地分区号和符号“/”
\item 可以是完整的电台呼号，也可以是完整电台呼号的任何一部分
\item 可以是完整的电台呼号，也可以是电台呼号的分区号加后缀
\item 一般指无线电管理机构指配的电台呼号，但也可以是对方能够理解的民间自创呼号、代号、代码等
\end{enumerate}
\noindent\textbf{解说：}\textbf{}。\\\noindent\textbf{答案：}

\bigskip




\noindent\textbf{试题：}业余无线电台呼号的指配流程是：
\begin{enumerate}[leftmargin=3em]
\item 无线电管理机构核发业余无线电台执照时，同时指配业余无线电台呼号
\item 在向无线电管理机构委托的受理服务机构提交设台申请窗口后，由服务机构指配呼号
\item 无线电管理机构核发业余无线电台执照后，由申请人再向其申请指配呼号
\item 业余无线电台设台人在提交设台申请的同时提出所要求指配的呼号，经服务机构同意后，报无线电管理机构正式指配
\end{enumerate}
\noindent\textbf{解说：}\textbf{}。\\\noindent\textbf{答案：}

\bigskip




\noindent\textbf{试题：}业余无线电爱好者对业已指配给自己的电台呼号不满意，是否可以申请另行指配业余无线电台呼号？
\begin{enumerate}[leftmargin=3em]
\item 不可以。核发业余无线电台执照的无线电管理机构已经为申请人指配业余无线电台呼号的，不另行指配其他业余无线电台呼号
\item 更新所设置的业余无线电台类别时可以申请另行指配业余无线电台呼号
\item 可以申请另行指配业余无线电台呼号，但须缴纳额外的费用
\item 业余无线电台执照有效期届满、设台人向核发执照的无线电管理机构申请办理延续手续时可以申请另行指配业余无线电台呼号
\end{enumerate}
\noindent\textbf{解说：}\textbf{}。\\\noindent\textbf{答案：}

\bigskip




\noindent\textbf{试题：}各地业余无线电台呼号前缀字母和后缀字符的可用范围的确定方法是：
\begin{enumerate}[leftmargin=3em]
\item 由国家无线电管理机构编制和分配
\item 地方无线电管理机构根据当地呼号资源的使用情况自行分配
\item 地方无线电民间组织提出建议，当地无线电管理机构批准
\item 由业余无线电爱好者根据需求提出建议，当地无线电管理机构批准
\end{enumerate}
\noindent\textbf{解说：}\textbf{}。\\\noindent\textbf{答案：}

\bigskip




\noindent\textbf{试题：}业余无线电爱好者是否可要求设台地所在地方无线电管理机构给予指配超出业已分配给该地方的前缀字母和后缀字符可用范围的业余无线电台呼号？
\begin{enumerate}[leftmargin=3em]
\item 不能，特殊业余无线电台呼号只能由国家无线电管理机构指配
\item 可以，但只限于与在当地所举办的大型国际或国家级活动有关的特殊电台
\item 可以，但只限于与当地政府组织的大型科技活动有关的特殊电台
\item 可以，但只限于当地业余无线电台参加国际重要业余无线电活动的特殊情况
\end{enumerate}
\noindent\textbf{解说：}\textbf{}。\\\noindent\textbf{答案：}

\bigskip




\noindent\textbf{试题：}经地方无线电管理机构批准设置的业余无线电台，设台地迁入其他省、自治区或者直辖市时，应办理的手续为：
\begin{enumerate}[leftmargin=3em]
\item 先到原核发执照的无线电管理机构办理申请注销原业余无线电台，再到迁入地的地方无线电管理机构办理申请设置业余无线电台的手续
\item 持原电台执照直接到迁入地的地方无线电管理机构申请办理变更手续
\item 持原电台执照直接到原核发执照的无线电管理机构申请办理变更手续
\item 不需要办理任何手续
\end{enumerate}
\noindent\textbf{解说：}\textbf{}。\\\noindent\textbf{答案：}

\bigskip




\noindent\textbf{试题：}经国家无线电管理机构批准设置的业余无线电台，设台地迁入其他省、自治区或者直辖市时，应办理的手续为：
\begin{enumerate}[leftmargin=3em]
\item 先到原核发执照的无线电管理机构申请办理注销手续，缴回原电台执照，领取国家无线电管理机构已批准设台的证明，凭证明到迁入地的地方无线电管理机构完成申请变更手续，领取新电台执照
\item 持原电台执照直接到迁入地的地方无线电管理机构申请办理变更手续
\item 持原电台执照直接到原核发执照的无线电管理机构申请办理变更手续
\item 不需要办理任何手续
\end{enumerate}
\noindent\textbf{解说：}\textbf{}。\\\noindent\textbf{答案：}

\bigskip




\noindent\textbf{试题：}设台地迁入其他省、自治区或者直辖市时，业余电台呼号的指配方法为：
\begin{enumerate}[leftmargin=3em]
\item 由设台人选择：方法一，注销原电台呼号，指配迁入地的新电台呼号；方法二，申请在迁入地继续指配原来的电台呼号
\item 必须继续指配原来的电台呼号
\item 必须指配迁入地的新电台呼号
\item 可以在保留原电台呼号的同时申请指配迁入地的新电台呼号
\end{enumerate}
\noindent\textbf{解说：}\textbf{}。\\\noindent\textbf{答案：}

\bigskip




\noindent\textbf{试题：}设台地迁入其他省、自治区或者直辖市时，申请在迁入地继续指配原来的电台呼号的手续为：
\begin{enumerate}[leftmargin=3em]
\item 先到原核发执照的无线电管理机构申请办理注销手续，缴回原电台执照，取得由迁入地指配原业余无线电台呼号的书面同意，再到迁入地的地方无线电管理机构办理相应的手续、重新指配原电台呼号，领取新的电台执照
\item 不需办理任何手续即可把原电台呼号带到迁入地继续使用
\item 只需到原核发执照的无线电管理机构申请申请办理为迁移后的电台继续使用原电台呼号的全部手续
\item 只需到迁入地的地方无线电管理机构申请申请办理为迁移后的电台继续使用原电台呼号的全部手续
\end{enumerate}
\noindent\textbf{解说：}\textbf{}。\\\noindent\textbf{答案：}

\bigskip




\noindent\textbf{试题：}在实际通信中，是否可以把本台呼号的地区号码加后缀视作《业余电台管理办法》所说的“本台呼号”？
\begin{enumerate}[leftmargin=3em]
\item 不可以。不完整呼号不具有呼号的属性，不能视作呼号
\item 在熟悉的友台之间呼叫和联络中可以把不完整呼号视作“呼号”
\item 在VHF/UHF频段进行本地呼叫和联络时可以把不完整呼号视作“呼号”
\item 在HF频段进行国内呼叫和联络时可以把不完整呼号视作“呼号”
\end{enumerate}
\noindent\textbf{解说：}\textbf{}。\\\noindent\textbf{答案：}

\bigskip




\noindent\textbf{试题：}由国家无线电管理机构批准设台的北京火腿的电台呼号为BH1AAA，把电台带到西安去使用，则本台呼号应该为：
\begin{enumerate}[leftmargin=3em]
\item B9/BH1AAA
\item BH1AAA/9
\item B1/BH9AAA
\item BH1AAA/B9
\end{enumerate}
\noindent\textbf{解说：}\textbf{}。\\\noindent\textbf{答案：}

\bigskip




\noindent\textbf{试题：}某业余无线电爱好者，自己所设置的业余无线电台呼号为BH1ZZZ。现该爱好者到业余无线电台BH9YYY做客并在该台进行发射操作。应当使用的呼号为：
\begin{enumerate}[leftmargin=3em]
\item BH9YYY或者B9/BH1ZZZ
\item BH1ZZZ或者B9/BH1ZZZ
\item BH1ZZZ/9或者BH1ZZZ/BH9
\item BH9/BH1ZZZ或者BH1ZZZ
\end{enumerate}
\noindent\textbf{解说：}\textbf{}。\\\noindent\textbf{答案：}

\bigskip




\noindent\textbf{试题：}某业余无线电爱好者，自己所设置的业余无线电台呼号为BH1ZZZ。现该爱好者将自己的业余无线电台带到湖南进行异地发射操作。应当使用的呼号为：
\begin{enumerate}[leftmargin=3em]
\item B7/BH1ZZZ
\item BH1ZZZ
\item BH1ZZZ/B7
\item BH7/BH1ZZZ
\end{enumerate}
\noindent\textbf{解说：}\textbf{}。\\\noindent\textbf{答案：}

\bigskip




\noindent\textbf{试题：}某业余无线电爱好者，自己所设置的业余无线电台呼号为BH1ZZZ。现该爱好者到业余无线电台BH3YYY做客并并在该台进行发射操作。这种发射操作在业余无线电台管理中称为：
\begin{enumerate}[leftmargin=3em]
\item 客席发射操作
\item 异地发射操作
\item 违章发射操作
\item 移动发射操作
\end{enumerate}
\noindent\textbf{解说：}\textbf{}。\\\noindent\textbf{答案：}

\bigskip




\noindent\textbf{试题：}某业余无线电爱好者，自己所设置的业余无线电台呼号为BH1ZZZ。现该爱好者将自己的业余无线电台带到广东进行发射操作。这种发射操作在业余无线电台管理中称为：
\begin{enumerate}[leftmargin=3em]
\item 异地发射操作
\item 客席发射操作
\item 违章发射操作
\item 临时发射操作
\end{enumerate}
\noindent\textbf{解说：}\textbf{}。\\\noindent\textbf{答案：}

\bigskip




\noindent\textbf{试题：}BH1ZZZ由北京迁入河北省，并办妥了由河北无线电管理机构指配使用原电台呼号的全部手续，领取了新的业余无线电台执照。该台在日常通信时应使用呼号：
\begin{enumerate}[leftmargin=3em]
\item B3/BH1ZZZ
\item 固定台址发射操作用BH1ZZZ，移动发射操用B3/BH1ZZZ
\item 可任选使用呼号B3/BH1ZZZ或者BH1ZZZ
\item BH1ZZZ
\end{enumerate}
\noindent\textbf{解说：}\textbf{}。\\\noindent\textbf{答案：}

\bigskip




\noindent\textbf{试题：}发射类别（class of emission）是指用标准符号标示的某发射的一组特性，例如主载波调制方式，调制信号，被发送信息的类型以及其他适用的信号特性。表示CW报的发射类别是：
\begin{enumerate}[leftmargin=3em]
\item A1A
\item J3E
\item F2B
\item G2B
\end{enumerate}
\noindent\textbf{解说：}\textbf{}。\\\noindent\textbf{答案：}

\bigskip




\noindent\textbf{试题：}发射类别（class of emission）是指用标准符号标示的某发射的一组特性，例如主载波调制方式，调制信号，被发送信息的类型以及其他适用的信号特性。表示单边带话的发射类别是：
\begin{enumerate}[leftmargin=3em]
\item J3E
\item A1A
\item F2B
\item G2B
\end{enumerate}
\noindent\textbf{解说：}\textbf{}。\\\noindent\textbf{答案：}

\bigskip




\noindent\textbf{试题：}发射类别（class of emission）是指用标准符号标示的某发射的一组特性，例如主载波调制方式，调制信号，被发送信息的类型以及其他适用的信号特性。表示用单边带话传输的RTTY信号的发射类别是：
\begin{enumerate}[leftmargin=3em]
\item F2B
\item A1A
\item J3E
\item G2B
\end{enumerate}
\noindent\textbf{解说：}\textbf{}。\\\noindent\textbf{答案：}

\bigskip




\noindent\textbf{试题：}发射类别（class of emission）是指用标准符号标示的某发射的一组特性，例如主载波调制方式，调制信号，被发送信息的类型以及其他适用的信号特性。表示用单边带话传输的PSK31信号的发射类别是：
\begin{enumerate}[leftmargin=3em]
\item G2B
\item A1A
\item J3E
\item F2B
\end{enumerate}
\noindent\textbf{解说：}\textbf{}。\\\noindent\textbf{答案：}

\bigskip




\noindent\textbf{试题：}发射类别（class of emission）是指用标准符号标示的某发射的一组特性，例如主载波调制方式，调制信号，被发送信息的类型以及其他适用的信号特性。表示用单边带话传输的SSTV信号的发射类别是：
\begin{enumerate}[leftmargin=3em]
\item F3F
\item J3E
\item F2B
\item F3E
\end{enumerate}
\noindent\textbf{解说：}\textbf{}。\\\noindent\textbf{答案：}

\bigskip




\noindent\textbf{试题：}发射类别（class of emission）是指用标准符号标示的某发射的一组特性，例如主载波调制方式，调制信号，被发送信息的类型以及其他适用的信号特性。表示调频话的发射类别是：
\begin{enumerate}[leftmargin=3em]
\item F3E
\item F3F
\item J3E
\item F2B
\end{enumerate}
\noindent\textbf{解说：}\textbf{}。\\\noindent\textbf{答案：}

\bigskip




\noindent\textbf{试题：}发射类别（class of emission）是指用标准符号标示的某发射的一组特性，例如主载波调制方式，调制信号，被发送信息的类型以及其他适用的信号特性。表示用调频话传输的RTTY信号的发射类别是：
\begin{enumerate}[leftmargin=3em]
\item F2B
\item F3E
\item J3E
\item F3F
\end{enumerate}
\noindent\textbf{解说：}\textbf{}。\\\noindent\textbf{答案：}

\bigskip




\noindent\textbf{试题：}发射类别是指用标准符号标示的某发射的一组特性，例如主载波调制方式，调制信号，被发送信息的类型以及其他适用的信号特性。表示用调频话传输的SSTV信号的发射类别是：
\begin{enumerate}[leftmargin=3em]
\item F3F
\item F2B
\item J3E
\item F3E
\end{enumerate}
\noindent\textbf{解说：}\textbf{}。\\\noindent\textbf{答案：}

\bigskip




\noindent\textbf{试题：}发射类别是指用标准符号标示的某发射的一组特性，例如主载波调制方式，调制信号，被发送信息的类型以及其他适用的信号特性。表示调频ATV信号的发射类别是： 
\begin{enumerate}[leftmargin=3em]
\item F3F
\item F2B
\item F3E
\item J3E
\end{enumerate}
\noindent\textbf{解说：}\textbf{}。\\\noindent\textbf{答案：}

\bigskip




\noindent\textbf{试题：}某业余电台操作者听到业余专用频率上出现某种显然出自非业余电台的人为干扰发射，于是按下话筒向该发射者宣传无线电管理法规知识。对这种做法的评论应该是：
\begin{enumerate}[leftmargin=3em]
\item 错误；违反“业余无线电台的通信对象应当限于业余无线电台”规定。
\item 正确；但有点条乱，不予提倡
\item 正确；抓机遇宣传法规，应该提倡
\item 正确；但需注意态度耐心、用语文明
\end{enumerate}
\noindent\textbf{解说：}\textbf{}。\\\noindent\textbf{答案：}

\bigskip




\noindent\textbf{试题：}在业余无线电台中转发广播电台、互联网聊天、电话通话、其他电台的联络信号，这类行为的性质是：
\begin{enumerate}[leftmargin=3em]
\item 错误行为；违反“业余无线电台的通信对象应当限于业余无线电台”规定，因为通信中产生信息的一方不是通信业余无线电台本身
\item 正确行为；既然可以联络，不必要限制向话筒送什么内容
\item 如果转发的目的是进行技术调试、用转发信号作为测试信号的话，就是正常行为
\item 不算错误但也不值得提倡
\end{enumerate}
\noindent\textbf{解说：}\textbf{}。\\\noindent\textbf{答案：}

\bigskip




\noindent\textbf{试题：}业余电台在通信中为其他人或者单位、组织转达信息。对这种做法的评论应该是：
\begin{enumerate}[leftmargin=3em]
\item 违法行为；违反“业余无线电台的通信对象应当限于业余无线电台”的规定
\item 只要所转达的信息在内容上不违反《业余电台管理办法规》的禁止规定就是合法行为
\item 只要转达信息是无偿的，就是合法行为
\item 只要所转达的信息是有利于社会的公益信息，就是合法行为
\end{enumerate}
\noindent\textbf{解说：}\textbf{}。\\\noindent\textbf{答案：}

\bigskip




\noindent\textbf{试题：}某业余无线电协会在发射操作中向其会员播发公益性通知和技术训练讲座，但未得到相应无线电管理机构的批准。对这种做法的评论应该是：
\begin{enumerate}[leftmargin=3em]
\item 违法行为；违反“未经核发业余无线电台执照的无线电管理机构批准，业余无线电台不得以任何方式进行广播或者发射通播性质的信号”的规定
\item 只要所播发的通知或讲座有利于当地业余无线电爱好者技术水平的提高，不能算违法行为
\item 只要所播发的通知或讲座有利于当地业余无线电应急通信训练，不能算违法行为
\item 只要所播发的通知或讲座是涉及宣传业余电台管理知识的，不能算违法行为
\end{enumerate}
\noindent\textbf{解说：}\textbf{}。\\\noindent\textbf{答案：}

\bigskip




\noindent\textbf{试题：}关于业余无线电台在通信过程中使用的语言，正确的做法为：
\begin{enumerate}[leftmargin=3em]
\item 任何时候都应当使用明语及业余无线电领域公认的缩略语和简语
\item 可以使用虽然不是所有火腿通用、但在某些火腿圈子内部有一定可懂度的新编缩略语或暗语
\item 语言要创新，可以使用自创的特殊缩略语，虽开始时象是暗语，用多了就会变明语
\item 可提倡使用稀有语言或方言，尽量使特定通信对象以外的业余无线电台听不懂，以减少他台呼叫和插入的机会
\end{enumerate}
\noindent\textbf{解说：}\textbf{}。\\\noindent\textbf{答案：}

\bigskip




\noindent\textbf{试题：}业余无线电台实验新的编码、调制方式、数字通信协议或者交换尚未公开格式的数据文件，正确做法是：
\begin{enumerate}[leftmargin=3em]
\item 事先尽可能采取各种办法向信号可能覆盖范围内的业余无线电爱好者公开有关技术细节，并提交给核发其业余无线电台执照的地方无线电管理机构
\item 事先尽可能采取各种办法向信号可能覆盖范围内的业余无线电爱好者公开有关技术细节，但不必提交给核发其业余无线电台执照的地方无线电管理机构
\item 应事先提交给核发其业余无线电台执照的地方无线电管理机构，但不必向其他业余无线电爱好者公开有关技术细节
\item 不必事先公开或者提交给核发其业余无线电台执照的地方无线电管理机构，以后再说
\end{enumerate}
\noindent\textbf{解说：}\textbf{}。\\\noindent\textbf{答案：}

\bigskip




\noindent\textbf{试题：}由国家无线电管理机构审批的业余无线电台在设台地以外的地点进行异地发射操作时，应该注意：
\begin{enumerate}[leftmargin=3em]
\item 既要符合业余电台执照所核定的各项参数约束，又要遵守操作所在地的地方无线电管理机构的相关规定
\item 遵守的限制以业余电台执照所核定的各项参数和核发其业余电台执照的地方无线电管理机构的规定为准，与操作所在地的规定无关
\item 遵守的限制以操作所在地的地方无线电管理机构的相关规定为准，与核发电台执照的地方无线电管理机构的规定无关
\item 有了电台执照就是万事大吉，不必认真了解和遵守什么具体规定
\end{enumerate}
\noindent\textbf{解说：}\textbf{}。\\\noindent\textbf{答案：}

\bigskip




\noindent\textbf{试题：}具备国家无线电管理机构规定的操作技术能力并具有法律规定有效证明文件、但还没有获准设置自己的业余电台的人是否可以到业余电台进行发射操作？答案是：
\begin{enumerate}[leftmargin=3em]
\item 可以。使用所操作业余电台的呼号，由该业余电台的设台人对操作不妥而造成的有害干负责
\item 可以。因为自己没有呼号，只能在通信中使用临时自编的呼号，或用姓名代替呼号
\item 不可以
\item 青少年可以，成人不可以
\end{enumerate}
\noindent\textbf{解说：}\textbf{}。\\\noindent\textbf{答案：}

\bigskip




\noindent\textbf{试题：}尚未考得《业余电台操作证书》的人在接受业余电台培训中实习发射操作应遵守的条件是什么？
\begin{enumerate}[leftmargin=3em]
\item 必须已接受法规等基础培训、必须由电台负责人现场辅导、必须在执照核定范围以及国家规定的操作权限内、进行短时间体验性发射操作实习
\item 只要业余电台设置人或者其技术负责人能确认实际上已经具备操作技术能力，可以独立进行发射操作，并能为其操作不善造成的后果负责，可以独立发射操作
\item 尚未取得关于具备操作技术能力有效证明文件的人任何情况下都不可以进行发射操作
\item 尚未取得关于具备操作技术能力有效证明文件者如为青少年，可以在集体业余电台独立操作，如为成人则任何情况下都不可以进行发射操作
\end{enumerate}
\noindent\textbf{解说：}\textbf{}。\\\noindent\textbf{答案：}

\bigskip




\noindent\textbf{试题：}业余无线电台设置人应对其无线电发射设备担负的法定责任为：
\begin{enumerate}[leftmargin=3em]
\item 应当确保其无线电发射设备处于正常工作状态，避免对其他无线电业务造成有害干扰
\item 应当确保其无线电发射设备随最先进型号更新，为其他业余电台树立求新的榜样
\item 应当确保其无线电发射设备达到最大发射功率，以克服其他无线电业务的干扰
\item 应当确保其无线电发射设备经常处于工作状态，以提高业余频率的实际占用度
\end{enumerate}
\noindent\textbf{解说：}\textbf{}。\\\noindent\textbf{答案：}

\bigskip




\noindent\textbf{试题：}业余无线电爱好者使用业余无线电收信设备应遵守的规定为：
\begin{enumerate}[leftmargin=3em]
\item 不得接收与业余业务和卫星业余业务无关的信号
\item 只要不造成对其他业务的无线电干扰，接收无线电信号没有限制
\item 只要不被查出来，可以接收任何无线电信号
\item 只要出于个人对信息的兴趣而不涉及赢利，可以接收任何无线电信号
\end{enumerate}
\noindent\textbf{解说：}\textbf{}。\\\noindent\textbf{答案：}

\bigskip




\noindent\textbf{试题：}业余无线电爱好者无意接收到非业余业务和卫星业余业务的信息时，应遵守的规则为：
\begin{enumerate}[leftmargin=3em]
\item 不得传播、公布
\item 只可以在业余无线电台间共享，不得在其他场合公开
\item 只可以用非无线电方式在业余无线电爱好者之间交流，不得以无线电方式转发
\item 既然自己可以收到，别人也一定可以收到，当然可以传播、公布或者利用
\end{enumerate}
\noindent\textbf{解说：}\textbf{}。\\\noindent\textbf{答案：}

\bigskip




\noindent\textbf{试题：}业余无线电台是否可以发射从广播电台收到的信号、音像节目的录音，或者故意转送电台周围的声音？
\begin{enumerate}[leftmargin=3em]
\item 不可以，不得发送与业余业务和卫星业余业务无关的信号
\item 可以，因为该类信息没有保密性
\item 可以，用于显示自己发射设备的信号质量
\item 可以，用于提起其他有业余无线电台操作员精神，防止乏困
\end{enumerate}
\noindent\textbf{解说：}\textbf{}。\\\noindent\textbf{答案：}

\bigskip




\noindent\textbf{试题：}国家对于利用业余无线电台从事发布、传播违反法律或者公共道德的信息的行为的态度是；
\begin{enumerate}[leftmargin=3em]
\item 禁止
\item 不提倡
\item 容忍
\item 不可以发布传、播违法信息。但违反公共道德的信息属于水平问题，不鼓励就是了
\end{enumerate}
\noindent\textbf{解说：}\textbf{}。\\\noindent\textbf{答案：}

\bigskip




\noindent\textbf{试题：}出租车安装业余电台并用来传递有关载客的信息，这种行为的性质是：
\begin{enumerate}[leftmargin=3em]
\item 违法行为，违反了严禁利用业余无线电台从事从事商业或者其他营利活动的规定
\item 不太好，因为占用了其他业余电台通信的频率
\item 只要不影响其他业余电台的正常通信就可以
\item 只要管理部门不来查处就可以
\end{enumerate}
\noindent\textbf{解说：}\textbf{}。\\\noindent\textbf{答案：}

\bigskip




\noindent\textbf{试题：}利用业余无线电台通信来促销业余无线电产品或者推动与业余无线电活动有关的其他商业性活动，对这类行为的态度应该是：
\begin{enumerate}[leftmargin=3em]
\item 禁止
\item 不提倡但也不禁止，毕竟有利于业余无线电活动发展
\item 只要是业余无线电民间组织是获利方，即使从事商业或其他营利活动，应支持
\item 如果设台人或者设台单位本身是以这类经营为生的，应适当理解和容忍
\end{enumerate}
\noindent\textbf{解说：}\textbf{}。\\\noindent\textbf{答案：}

\bigskip




\noindent\textbf{试题：}利用自己的业余电台强信号故意压制其他业余电台的正常通信，或者在业余无线电频率上转播音乐或广播节目，这些行为的性质属于：
\begin{enumerate}[leftmargin=3em]
\item 违法行为，违反了严禁阻碍其他无线电台通信的规定
\item 不妥行为，没有考虑到他人的乐趣
\item 正常现象，社会上一些人素质就是如此，应该谅解
\item 不文明行为，对其他业余电台不够礼貌
\end{enumerate}
\noindent\textbf{解说：}\textbf{}。\\\noindent\textbf{答案：}

\bigskip




\noindent\textbf{试题：}业余无线电活动是否有序开展，会影响整个社会的无线电通信的安全和有效，使用不当甚至会导致生命财产损失。业余无线电爱好者在这方的法定责任是：
\begin{enumerate}[leftmargin=3em]
\item 业余无线电台设置、使用人应当加强自律
\item 个人没有责任，只能依靠管理部门的监督检查和违法查处
\item 个人没有责任，只能依靠业余无线电民间组织充当“协管”
\item 有了电台执照，日常一切言行当然可以带到电台通信中，无责任可言
\end{enumerate}
\noindent\textbf{解说：}\textbf{}。\\\noindent\textbf{答案：}

\bigskip




\noindent\textbf{试题：}国际电联规定的确定发射电台辐射功率的原则为：
\begin{enumerate}[leftmargin=3em]
\item 发射电台只应辐射为保证满意服务所必要的功率
\item 发射电台应辐射尽量大的功率以提供尽量好的信号质量
\item HF频段发射电台应辐射尽量大的功率，VHF频段发射电台应辐射尽量小的功率
\item VHF/UHF频段发射电台应辐射尽量大的功率，HF频段发射电台应辐射尽量小的功率
\end{enumerate}
\noindent\textbf{解说：}\textbf{}。\\\noindent\textbf{答案：}

\bigskip




\noindent\textbf{试题：}业余电台通信受到违法电台或者不明电台的有害干扰。正确的做法是：
\begin{enumerate}[leftmargin=3em]
\item 不予理睬，收集有关信息并向无线电管理机构举报
\item 在频率上向其宣传无线电管理法，要求其停止干扰
\item 立即报告无线电管理机构进行干涉
\item 用大功率信号对其进行压制
\end{enumerate}
\noindent\textbf{解说：}\textbf{}。\\\noindent\textbf{答案：}

\bigskip




\noindent\textbf{试题：}按照我国规定，购置使用公众对讲机不需取得批准。业余无线电爱好者需要与公众对讲机用户通信时应该：
\begin{enumerate}[leftmargin=3em]
\item 业余无线电台不能用于与公众对讲机通信
\item 将业余无线电台设置到公众对讲机的频率，以不大于业余无线电台执照核定的发射功率与之通信
\item 将业余无线电台设置到公众对讲机的频率，以不大于0.5W的发射功率与之通信
\item 将业余无线电台设置到公众对讲机的频率，但只能进行由业余无线电台到公众对讲机的单向发信
\end{enumerate}
\noindent\textbf{解说：}\textbf{}。\\\noindent\textbf{答案：}

\bigskip




\noindent\textbf{试题：}关于业余无线电台的应急通信，正确的叙述是：
\begin{enumerate}[leftmargin=3em]
\item 在突发重大自然灾害等紧急情况下，业余无线电台才可以和非业余无线电台进行规定内容的通信
\item 在日常应急通信训练中，业余无线电台可以和各种非业余无线电台进行通信
\item 在日常应急通信训练中，业余无线电台可以和地方公益性救援团体的非业余无线电台进行通信
\item 在日常应急通信训练中，业余无线电台可以和地方公益性救援团体的非业余无线电台进行通信，但须经当地业余无线电协会同意
\end{enumerate}
\noindent\textbf{解说：}\textbf{}。\\\noindent\textbf{答案：}

\bigskip




\noindent\textbf{试题：}业余无线电台允许与非业余无线电台通信的条件是：
\begin{enumerate}[leftmargin=3em]
\item 在突发重大自然灾害等紧急情况下，内容限于与抢险救灾直接相关的紧急事务或者应急救援相关部门交办的任务
\item 在当地政府或非盈利机构组织的公益活动中，内容限于与公益事务或者相关的活动组织机构交办的任务
\item 在青少年科技教育活动中，仅可与青少年非业余无线电台通信，内容限于与青少年科技教育直接有关的事务
\item 在无线电技术研究中，仅可与具备其他业务电台执照的对象通信，内容限于技术实验所需的信号
\end{enumerate}
\noindent\textbf{解说：}\textbf{}。\\\noindent\textbf{答案：}

\bigskip




\noindent\textbf{试题：}关于业余无线电台的应急通信，正确的叙述是：
\begin{enumerate}[leftmargin=3em]
\item 在突发重大自然灾害等紧急情况下，业余无线电台的通信内容可以涉及应急救援相关部门交办的任务
\item 在平时的任何时侯，业余无线电台的通信内容可以涉及任何政府组织和非盈利机构交办的任务
\item 平时在专门的应急通信训练活动中，业余无线电台的通信内容可以涉及应急救援相关部门和组织机构交办的任务
\item 在日常公益性社会活动中，业余无线电台的通信内容可以涉及各种公益机构交办的任务
\end{enumerate}
\noindent\textbf{解说：}\textbf{}。\\\noindent\textbf{答案：}

\bigskip




\noindent\textbf{试题：}法规和国际业余无线电惯例要求业余电台日志记载的必要基本内容是：
\begin{enumerate}[leftmargin=3em]
\item 通信时间、通信频率、通信模式、对方呼号、双方信号报告
\item 通信对方姓名、对方所在国家或城市、通信模式、双方信号报告
\item 通信时间、通信频率、双方收发信设备和天线、对方台址
\item 通信时间、通信模式、对方信号报告、对方台址、对方天气
\end{enumerate}
\noindent\textbf{解说：}\textbf{}。\\\noindent\textbf{答案：}

\bigskip




\noindent\textbf{试题：}法规和国际业余无线电惯例要求业余电台日志记载的必要基本内容是：
\begin{enumerate}[leftmargin=3em]
\item DATE、TIME、FREQ、MODE、CALL（对方）、RST（双方）
\item DATE、FREQ、QTH（对方）、RIG（对方）、RST（双方）、WX（对方）
\item DATE、TIME、MODE、CALL（对方）、QTH（对方）、RST（双方）
\item CALL（通信对方）、T IME、FREQ、RIG（对方）、RST（双方）、PWR（双方）
\end{enumerate}
\noindent\textbf{解说：}\textbf{}。\\\noindent\textbf{答案：}

\bigskip




\noindent\textbf{试题：}填写和邮寄QSL卡片时的正确做法有：
\begin{enumerate}[leftmargin=3em]
\item 迫切需要方卡回寄卡片时，应直接向对方地址邮寄卡片并附加SASE
\item 填写错误时应划去或使用涂改液覆盖错误内容并加以改正
\item 自己的邮寄地址与电台的发射地点不同时，应在QTH栏目内填明详细邮寄地址
\item 通过卡片管理局寄出卡片并希望对方回卡时，应在卡片上注明PSE QSL DIRECT
\end{enumerate}
\noindent\textbf{解说：}\textbf{}。\\\noindent\textbf{答案：}

\bigskip




\noindent\textbf{试题：}关于QSL卡片的正确用法是：
\begin{enumerate}[leftmargin=3em]
\item 不是作为联络或收听证明而交换QSL卡片时，应填上“Eye ball QSO”等有关说明，不应赠送空白卡片
\item 空白QSL卡片可以当做照片或者名片，任意赠送、交换、散发
\item 出于火腿互相帮助的目的，虽然对方没有联络到自己，也可以发去确认联络的QSL卡片
\item 如果在联络中没有听清对方呼号，可以在寄发QSL卡片的对方台名栏中填写对方操作员姓名
\end{enumerate}
\noindent\textbf{解说：}\textbf{}。\\\noindent\textbf{答案：}

\bigskip




\noindent\textbf{试题：}负责协调国际无线电管理的政府间组织是：
\begin{enumerate}[leftmargin=3em]
\item 国际电信联盟
\item 国际业余无线电联盟
\item 联合国大会
\item 联合国科教文组织
\end{enumerate}
\noindent\textbf{解说：}\textbf{}。\\\noindent\textbf{答案：}

\bigskip




\noindent\textbf{试题：}国际业余无线电联盟支持会员协会在发展业余无线电时把它看成一种：
\begin{enumerate}[leftmargin=3em]
\item 国家资源
\item 拉动经济的措施
\item 休闲娱乐手段
\item 国防后备队伍
\end{enumerate}
\noindent\textbf{解说：}\textbf{}。\\\noindent\textbf{答案：}

\bigskip




\noindent\textbf{试题：}国际业余无线电联盟由下列元素组成：
\begin{enumerate}[leftmargin=3em]
\item 成员协会、理事会、地区组织
\item 理事会、地区组织、个人会员
\item 理事会、地区分会、专业委员会
\item 个人会员、成员协会、理事会、地区组织
\end{enumerate}
\noindent\textbf{解说：}\textbf{}。\\\noindent\textbf{答案：}

\bigskip




\noindent\textbf{试题：}ITU对于5至30MHz短波频率的态度是：（RR 4.11、 4.11）
\begin{enumerate}[leftmargin=3em]
\item 该段频率特别适用于远距离通信，因此不鼓励本地通信使用，从而防止对远距离通信的有害干扰
\item 在保证优先满足本地通信的前提下，允许远距离通信适当使用这一段频率
\item 同样鼓励本地和远距离通信尽量使用该段频率，提高这一段频谱的使用率
\item 鼓励在这些频段建立大功率中继站，以提高本地通信的覆盖效果
\end{enumerate}
\noindent\textbf{解说：}\textbf{}。\\\noindent\textbf{答案：}

\bigskip




\noindent\textbf{试题：}ITU《无线电规则》禁止所有电台发射下列种类的电波：
\begin{enumerate}[leftmargin=3em]
\item 减幅波（阻尼波）
\item 调幅波
\item 调频波
\item 调相波
\end{enumerate}
\noindent\textbf{解说：}\textbf{}。\\\noindent\textbf{答案：}

\bigskip




\noindent\textbf{试题：}如果收到国外寄来的QSL卡片中夹带有一张或多张IRC，应该：
\begin{enumerate}[leftmargin=3em]
\item 尽快检查电台日志确认联络的真实性，并通过邮局直接向对方地址寄出自己的QSL卡片
\item 不必确认联络的真实性，尽快通过邮局直接向对方地址寄出自己的QSL卡片
\item 尽快检查电台日志确认联络的真实性，并通过国内的QSL卡片管理局寄出自己的QSL卡片
\item 不必理会
\end{enumerate}
\noindent\textbf{解说：}\textbf{}。\\\noindent\textbf{答案：}

\bigskip




\noindent\textbf{试题：}业余中继台的设置和技术参数等应满足下列关键条件：
\begin{enumerate}[leftmargin=3em]
\item 符合国家以及设台地的地方无线电管理机构的规定
\item 符合设台地的地方业余无线电民间组织的规划
\item 仅需符合设台地的地方无线电管理机构的规划及相关规定
\item 符合申请人关于设置中继台的客观需求和技术考虑
\end{enumerate}
\noindent\textbf{解说：}\textbf{}。\\\noindent\textbf{答案：}

\bigskip




\noindent\textbf{试题：}业余中继台必备的技术措施为：
\begin{enumerate}[leftmargin=3em]
\item 设专人负责监控和管理工作，配备有效的遥控手段，保证造成有害干扰时及时停止发射
\item 技术加密措施，防止未经设台人允许的业余无线电台启用中继
\item 尽量提高发射功率，以便压制覆盖区内的其他强信号干扰
\item 设热备份系统，保证不间断工作
\end{enumerate}
\noindent\textbf{解说：}\textbf{}。\\\noindent\textbf{答案：}

\bigskip




\noindent\textbf{试题：}某团体依法设置了一部业余中继台。其正确做法是：
\begin{enumerate}[leftmargin=3em]
\item 向其覆盖区域内的所有业余无线电台提供平等的服务，并将使用业余中继台所需的各项技术参数公开
\item 中继台是设台者出资建设和维护的，因此仅供经设置者允许的业余电台使用
\item 中继台是设台者出资建设和维护的，因此仅供本团体成员优先使用，空闲时方供其他业余电台使用
\item 为保证中继台正常运行，要求覆盖区内所有业余电台缴纳维护成本，否则不准使用
\end{enumerate}
\noindent\textbf{解说：}\textbf{}。\\\noindent\textbf{答案：}

\bigskip




\noindent\textbf{试题：}选择144MHz或430MHz业余模拟调频中继台同频段收发频差的原则是：
\begin{enumerate}[leftmargin=3em]
\item 采用业余无线电标准频差，即144MHz频段600kHz，430MHz频段5MHz
\item 尽量采用非标准频差以阻止一般业余无线电台占用
\item 采用经常变换频差的办法减少占用度
\item 可以在国家《无线电频率划分规定》所规定业余频率范围内任意选择
\end{enumerate}
\noindent\textbf{解说：}\textbf{}。\\\noindent\textbf{答案：}

\bigskip




\noindent\textbf{试题：}业余中继台的使用原则是：
\begin{enumerate}[leftmargin=3em]
\item 除必要的短暂通信外，应保持业余中继台具有足够的空闲时间，以便随时响应突发灾害应急呼叫
\item 应使中继台尽量处于接近饱和的忙碌状态，提高使用效率
\item 鼓励业余无线电民间组织（协会）通过中继台向当地会员发布通知
\item 鼓励青少年学生通过中继台交流解题方法和学习心得
\end{enumerate}
\noindent\textbf{解说：}\textbf{}。\\\noindent\textbf{答案：}

\bigskip




\noindent\textbf{试题：}如果你知道另一个电台的呼号，想要在中继上呼叫他，你应该怎么做？
\begin{enumerate}[leftmargin=3em]
\item 呼叫对方的呼号，并报出自己的呼号
\item 呼叫“break break”，然后说出对方的呼号
\item 呼叫“CQ”三次，然后说出对方的呼号
\item 等待，直到你要呼叫的电台呼叫CQ后，立刻回答他
\end{enumerate}
\noindent\textbf{解说：}\textbf{}。\\\noindent\textbf{答案：}

\bigskip




\noindent\textbf{试题：}业余无线电台设置、使用人应当接受下列机构对业余无线电台及其使用情况的监督检查：
\begin{enumerate}[leftmargin=3em]
\item 无线电管理机构或者其委托单位的监督检查
\item 业余无线电民间组织的独立监督检查
\item 单位或所在居委会、村民委员会、物主委员会的监督检查
\item 国家计量监督部门的监督检查
\end{enumerate}
\noindent\textbf{解说：}\textbf{}。\\\noindent\textbf{答案：}

\bigskip




\noindent\textbf{试题：}对擅自设置、使用业余无线电台的单位或个人，国家无线电管理机构或者地方无线电管理机构可以根据其具体情况给予下列处罚：
\begin{enumerate}[leftmargin=3em]
\item 警告、查封或者没收设备、没收非法所得；情节严重的，可以并处一千元以上，五千元以下的罚款
\item 劝告拆除非法设置的电台；情节严重的，可以并处警告、查封或者没收设备
\item 责令停止使用非法设置的电台；情节严重的，可以并处警告、查封或者没收设备
\item 责令停止使用非法设置的电台并作出书面检查；情节严重的，可以并处一千元以下的罚款
\end{enumerate}
\noindent\textbf{解说：}\textbf{}。\\\noindent\textbf{答案：}

\bigskip




\noindent\textbf{试题：}业余电台干扰无线电业务的，国家无线电管理机构或者地方无线电管理机构可以根据其具体情况给予设置业余无线电台的单位或个人下列处罚：
\begin{enumerate}[leftmargin=3em]
\item 警告、查封或者没收设备、没收非法所得；情节严重的，可以并处一千元以上，五千元以下的罚款
\item 劝告拆除非法设置的电台；情节严重的，可以并处警告、查封或者没收设备
\item 责令停止使用非法设置的电台；情节严重的，可以并处警告、查封或者没收设备
\item 责令停止使用非法设置的电台并作出书面检查；情节严重的，可以并处一千元以下的罚款
\end{enumerate}
\noindent\textbf{解说：}\textbf{}。\\\noindent\textbf{答案：}

\bigskip




\noindent\textbf{试题：}业余电台随意变更核定项目、发送和接收与业余无线电无关的信号的，国家无线电管理机构或者地方无线电管理机构可以根据其具体情况给予设置业余无线电台的单位或个人下列处罚：
\begin{enumerate}[leftmargin=3em]
\item 警告、查封或者没收设备、没收非法所得；情节严重的，可以并处一千元以上，五千元以下的罚款
\item 劝告拆除非法设置的电台；情节严重的，可以并处警告、查封或者没收设备
\item 责令停止使用非法设置的电台；情节严重的，可以并处警告、查封或者没收设备
\item 责令停止使用非法设置的电台并作出书面检查；情节严重的，可以并处一千元以下的罚款
\end{enumerate}
\noindent\textbf{解说：}\textbf{}。\\\noindent\textbf{答案：}

\bigskip




\noindent\textbf{试题：}超出核定范围使用频率或者有其他违反频率管理有关规定的行为的，无线电管理机构可以根据其具体情况给予设置业余无线电台的单位或个人下列处罚：
\begin{enumerate}[leftmargin=3em]
\item 责令限期改正，可以处警告或者三万元以下的罚款
\item 责令限期改正，可以处警告或者一千元以上，五千元以下的罚款
\item 责令限期改正，可以处警告或者一千元以下的罚款
\item 责令限期改正，情节严重的，可以并处警告、查封或者没收设备
\end{enumerate}
\noindent\textbf{解说：}\textbf{}。\\\noindent\textbf{答案：}

\bigskip




\noindent\textbf{试题：}对涂改、仿制、伪造、倒卖、出租、出借业余无线电台执照，或者以其他形式非法转让业余无线电台执照的，无线电管理机构可以给予下列处罚：
\begin{enumerate}[leftmargin=3em]
\item 应当责令限期改正，可以处警告或者三万元以下的罚款
\item 应当责令限期改正，可以处警告或者一千元以上，五千元以下的罚款
\item 应当责令限期改正，可以处警告或者一千元以下的罚款
\item 应当责令限期改正，情节严重的，可以并处警告、查封或者没收设备
\end{enumerate}
\noindent\textbf{解说：}\textbf{}。\\\noindent\textbf{答案：}

\bigskip




\noindent\textbf{试题：}对盗用、出租、出借、转让、私自编制或者违法使用业余无线电台呼号的，无线电管理机构可以给予下列处罚：
\begin{enumerate}[leftmargin=3em]
\item 应当责令限期改正，可以处警告或者三万元以下的罚款
\item 应当责令限期改正，可以处警告或者一千元以上，五千元以下的罚款
\item 应当责令限期改正，可以处警告或者一千元以下的罚款
\item 应当责令限期改正，情节严重的，可以并处警告、查封或者没收设备
\end{enumerate}
\noindent\textbf{解说：}\textbf{}。\\\noindent\textbf{答案：}

\bigskip




\noindent\textbf{试题：}对以不正当手段取得业余无线电台执照的，无线电管理机构可以给予下列处罚：
\begin{enumerate}[leftmargin=3em]
\item 责令限期改正，可以处警告或者三万元以下的罚款
\item 责令限期改正，可以处警告或者一千元以上，五千元以下的罚款
\item 责令限期改正，可以处警告或者一千元以下的罚款
\item 责令限期改正，情节严重的，可以并处警告、查封或者没收设备
\end{enumerate}
\noindent\textbf{解说：}\textbf{}。\\\noindent\textbf{答案：}

\bigskip




\noindent\textbf{试题：}对向负责监督检查的无线电管理机构隐瞒有关情况、提供虚假材料或者拒绝提供反映其活动情况的真实材料的，无线电管理机构可以给予下列处罚：
\begin{enumerate}[leftmargin=3em]
\item 责令限期改正，可以处警告或者三万元以下的罚款
\item 责令限期改正，可以处警告或者一千元以上，五千元以下的罚款
\item 责令限期改正，可以处警告或者一千元以下的罚款
\item 责令限期改正，情节严重的，可以并处警告、查封或者没收设备
\end{enumerate}
\noindent\textbf{解说：}\textbf{}。\\\noindent\textbf{答案：}

\bigskip




\noindent\textbf{试题：}对违法使用业余无线电台造成严重后果的，无线电管理机构可以给予下列处罚：
\begin{enumerate}[leftmargin=3em]
\item 应当责令限期改正，可以处警告或者三万元以下的罚款
\item 应当责令限期改正，可以处警告或者一千元以上，五千元以下的罚款
\item 应当责令限期改正，可以处警告或者一千元以下的罚款
\item 应当责令限期改正，情节严重的，可以并处警告、查封或者没收设备
\end{enumerate}
\noindent\textbf{解说：}\textbf{}。\\\noindent\textbf{答案：}

\bigskip




\noindent\textbf{试题：}违反国家规定，擅自设置、使用无线电台（站），或者擅自占用频率，经责令停止使用后拒不停止使用，干扰无线电通信正常进行，造成严重后果的的，可被判犯扰乱无线电通信管理秩序罪，处三年以下有期徒刑、拘役或者管制，并处或者单处罚金。这个规定出自于下列法规律：
\begin{enumerate}[leftmargin=3em]
\item 中华人民共和国刑法
\item 中华人民共和国民法通则
\item 中华人民共和国无线电管理条例
\item 中华人民共和国电信法
\end{enumerate}
\noindent\textbf{解说：}\textbf{}。\\\noindent\textbf{答案：}

\bigskip




\noindent\textbf{试题：}无线电管制是指在下列范围内依法采取的对无线电波的发射、辐射和传播实施的强制性管理：
\begin{enumerate}[leftmargin=3em]
\item 在特定时间和特定区域内
\item 在全国范围、所有时间内
\item 在特定范围、所有时间内
\item 在例行范围和例行时间内
\end{enumerate}
\noindent\textbf{解说：}\textbf{}。\\\noindent\textbf{答案：}

\bigskip




\noindent\textbf{试题：}无线电管制是指在特定时间和特定区域内，依法采取的下列性质的管理：
\begin{enumerate}[leftmargin=3em]
\item 对无线电波的发射、辐射和传播实施的强制性管理
\item 对无线电波的发射、辐射实施的指导和行业自律性管理
\item 对无线电发射设备的生产、销售实施的强制性管理
\item 对无线电发射设备的生产、销售实施的指导和行业自律性管理
\end{enumerate}
\noindent\textbf{解说：}\textbf{}。\\\noindent\textbf{答案：}

\bigskip




\noindent\textbf{试题：}在特定时间和特定区域内实施无线电管制时，与业余无线电有关的管理措施包括：
\begin{enumerate}[leftmargin=3em]
\item 限制或者禁止业余无线电台（站）的使用，以及对特定的无线电频率实施技术阻断等
\item 限制或者禁止业余无线电台设备的生产和销售
\item 限制、但不会禁止业余无线电台（站）的使用
\item 依法设置的业余电台不在管制范围之内
\end{enumerate}
\noindent\textbf{解说：}\textbf{}。\\\noindent\textbf{答案：}

\bigskip




\noindent\textbf{试题：}决定实施无线电管制的机构为：
\begin{enumerate}[leftmargin=3em]
\item 在全国范围内或者跨省、自治区、直辖市实施，由国务院和中央军事委员会决定。在省、自治区、直辖市范围内实施，由省、自治区、直辖市人民政府和相关军区决定
\item 在全国范围内或者跨省、自治区、直辖市实施，由国家无线电管理机构决定。在省、自治区、直辖市范围内实施，由相关地方无线电管理机构决定
\item 在地、市、县实施，由地、市、县人民政府决定
\item 在单位、居民区实施，由单位上级业务主管机构和区人民政府共同决定
\end{enumerate}
\noindent\textbf{解说：}\textbf{}。\\\noindent\textbf{答案：}

\bigskip




\noindent\textbf{试题：}违反无线电管制命令和无线电管制指令的，由下列机构依法进行处罚：
\begin{enumerate}[leftmargin=3em]
\item 国家无线电管理机构或者省、自治区、直辖市无线电管理机构；违反治安管理规定者由公安机关处罚
\item 城管、工商、交通联合执法
\item 当地业余无线电协会
\item 所在军区派出的专门机构
\end{enumerate}
\noindent\textbf{解说：}\textbf{}。\\\noindent\textbf{答案：}

\bigskip




\noindent\textbf{试题：}业余电台违反无线电管制命令和无线电管制指令的，可以依法规受到下列处罚：
\begin{enumerate}[leftmargin=3em]
\item  责令改正；拒不改正的，关闭、查封、暂扣或者拆除相关设备；情节严重的，吊销电台执照；违反治安管理规定的，由公安机关处罚
\item  处警告或者三万元以下的罚款
\item  处警告或者一千元以上，五千元以下的罚款 
\item  责令改正；并开除业余无线电协会会籍、罚没无线电通信设备
\end{enumerate}
\noindent\textbf{解说：}\textbf{}。\\\noindent\textbf{答案：}

\bigskip




\noindent\textbf{试题：}C类业余无线电台允许发射的频率为：
\begin{enumerate}[leftmargin=3em]
\item 各业余业务和卫星业余业务频段
\item 各VHF和UHF频段
\item 限于HF、VHF、UHF范围内的各业余业务和卫星业余业务频段
\item 30-3000MHz范围外的各业余业务和卫星业余业务频段
\end{enumerate}
\noindent\textbf{解说：}\textbf{}。\\\noindent\textbf{答案：}

\bigskip




\noindent\textbf{试题：}C类业余无线电台允许发射的最大发射功率为不大于：
\begin{enumerate}[leftmargin=3em]
\item 30MHz以下业余频段不大于1000瓦，30MHz以上业余频段不大于25瓦
\item 30MHz以上业余频段不大于1000瓦，30MHz以下业余频段不大于25瓦
\item 25瓦
\item 100瓦
\end{enumerate}
\noindent\textbf{解说：}\textbf{}。\\\noindent\textbf{答案：}

\bigskip




\noindent\textbf{试题：}负责组织C类业余无线电台所需操作技术能力的验证的机构是：
\begin{enumerate}[leftmargin=3em]
\item 国家无线电管理机构
\item 地方无线电管理机构
\item 国家无线电管理机构和地方无线电管理机构
\item 地方业余无线电协会
\end{enumerate}
\noindent\textbf{解说：}\textbf{}。\\\noindent\textbf{答案：}

\bigskip




\noindent\textbf{试题：}业余无线电测向活动所用的隐蔽台属于下列范畴的设备：
\begin{enumerate}[leftmargin=3em]
\item 属于业余无线电台设备，设置使用须向无线电管理机构提出申请并得到批准
\item 属于特殊免执照发射设备，不需取得任何行政许可即可任意设置使用
\item 不属于无线电管理对象，由体育主管部门或者体育民间组织的默认即可设置使用
\item 属于体育或游戏器材，其性能指标不需经过无线电检测部门检测合格
\end{enumerate}
\noindent\textbf{解说：}\textbf{}。\\\noindent\textbf{答案：}

\bigskip




\noindent\textbf{试题：}我国规定允许C类业余电台在30MHz以下使用的最大发射功率为1000瓦。正确理解为：
\begin{enumerate}[leftmargin=3em]
\item 仅在参加国际竞赛、进行特殊技术试验等必需使用最大功率时才可使用
\item 可以在日常联络中使用最大功率以求对方给出更好的信号报告
\item 可以在日常联络遇频率上电台太多的情况下用来取得较强的话语权
\item 可以在遇到初学者操作不规范时用最大功率插入、加以纠正和指导
\end{enumerate}
\noindent\textbf{解说：}\textbf{}。\\\noindent\textbf{答案：}

\bigskip




\noindent\textbf{试题：}自制业余无线电发射设备，在经无线电检测机构检测合格并取得电台执照之前，调试时天线输出端应连接（或串联必要的仪表后连接）：
\begin{enumerate}[leftmargin=3em]
\item 假负载
\item VSWR严格等于1:1的驻波天线
\item VSWR严格等于1:1的行波天线
\item 测试专用的标准环形天线
\end{enumerate}
\noindent\textbf{解说：}\textbf{}。\\\noindent\textbf{答案：}

\bigskip




\noindent\textbf{试题：}某俱乐部约定了一个成员业余电台之间交流技术的网络频率，当遇有其他业余电台按通信惯例要求参加通信时，处理原则应为：
\begin{enumerate}[leftmargin=3em]
\item 无条件欢迎加入，因为任何核准的业余电台对频率享有平等的频率使用权
\item 要求其他业余电台在任何时间都不得使用俱乐部自己约定的专用通信频率
\item 要求其他业余电台在俱乐部成员结束网络通信后再使用该频率
\item 由俱乐部网络控制台决定是其他业余电台是否可以加入
\end{enumerate}
\noindent\textbf{解说：}\textbf{}。\\\noindent\textbf{答案：}

\bigskip




\noindent\textbf{试题：}业余电台在发起呼叫前不可缺少的操作步骤是：
\begin{enumerate}[leftmargin=3em]
\item 先守听一段时间，确保没有其他电台正在使用频率
\item 检查发射功率是否达到设备的额定输出功率
\item 先用礼貌的语言请其他电台让出频率
\item 先用吹话筒、吹口哨等方法发出连续信号检查天线驻波比
\end{enumerate}
\noindent\textbf{解说：}\textbf{}。\\\noindent\textbf{答案：}

\bigskip




\noindent\textbf{试题：}业余电台在发射调试信号进行发射功率和天线驻波比等检查时必须注意做到的是：
\begin{enumerate}[leftmargin=3em]
\item 先将频率设置到无人使用的空闲频率、偏离常用的热点频率
\item 先将天线的发射方向指向正北
\item 先将收发信机的语音压缩功能打开
\item 话筒离嘴距离在2公分以上，电键按键时间不短于5秒钟
\end{enumerate}
\noindent\textbf{解说：}\textbf{}。\\\noindent\textbf{答案：}

\bigskip




\noindent\textbf{试题：}单边带业余电台在测试检查天线驻波比需要发射平稳的连续信号。文明的作法是：
\begin{enumerate}[leftmargin=3em]
\item 先将电台设为CW方式按电键，或者设为AM或FM方式按PTT键（不对话筒说话），产生连续载波，测试结束后设回SSB方式
\item 将电台设为SSB方式，用平稳的气流对话筒吹口哨
\item 将电台设为SSB方式，深呼吸后用平稳的气流对话筒发长音“啊”
\item 将电台设为SSB方式，深呼吸后用平稳的气流对话筒发长音“嘻”
\end{enumerate}
\noindent\textbf{解说：}\textbf{}。\\\noindent\textbf{答案：}

\bigskip




\noindent\textbf{试题：}业余电台发起呼叫前应先守听一段时间，如没有听到信号，应再询问“有人使用频率吗”？确认没有应答方能发起呼叫。下列英语短句中不能正确表达这一询问的是：
\begin{enumerate}[leftmargin=3em]
\item Calling you, Roger？
\item Is the frequency in use？
\item Is any body in the frequency？
\item Any body here？
\end{enumerate}
\noindent\textbf{解说：}\textbf{}。\\\noindent\textbf{答案：}

\bigskip




\noindent\textbf{试题：}业余电台发起呼叫前应先守听一段时间，如没有听到信号，应再询问“有人使用频率吗”？确认没有应答方能发起呼叫。用CW表达这一询问的方法是：
\begin{enumerate}[leftmargin=3em]
\item QRL？
\item QRX？
\item QRZ？
\item QRV？
\end{enumerate}
\noindent\textbf{解说：}\textbf{}。\\\noindent\textbf{答案：}

\bigskip




\noindent\textbf{试题：}业余电台BH1ZZZ用话音发起CQ呼叫的正确格式为：
\begin{enumerate}[leftmargin=3em]
\item CQ、CQ、CQ。BH1ZZZ呼叫。Bravo Hotel One Zulu Zulu Zulu呼叫，BH1ZZZ呼叫。听到请回答。
\item CQ、CQ、CQ。听到请回答。
\item CQ、CQ、CQ。我是1ZZZ。听到请回答
\item CQ、CQ、CQ，CQ、CQ、CQ，CQ、CQ、CQ。BH1ZZZ呼叫。请过来。
\end{enumerate}
\noindent\textbf{解说：}\textbf{}。\\\noindent\textbf{答案：}

\bigskip




\noindent\textbf{试题：}业余电台BH1ZZZ用话音发起CQ呼叫的正确格式为：
\begin{enumerate}[leftmargin=3em]
\item CQ CQ CQ.This is BH1ZZZ. Bravo Hotel One Zulu Zulu Zulu, BH1ZZZ is calling. I’m standing by.
\item CQ CQ CQ. Go ahead please.
\item CQ CQ CQ. This One Zulu Zulu Zulu. Over.
\item CQ CQ CQ, CQ CQ CQ, CQ CQ CQ. This is BH1ZZZ. Back to you.
\end{enumerate}
\noindent\textbf{解说：}\textbf{}。\\\noindent\textbf{答案：}

\bigskip




\noindent\textbf{试题：}业余电台BH1ZZZ用话音呼叫BH8YYY的正确格式为：
\begin{enumerate}[leftmargin=3em]
\item BH8YYY、BH8YYY、BH8YYY。BH1ZZZ呼叫。Bravo Hotel One Zulu Zulu Zulu，BH1ZZZ呼叫。听到请回答。
\item BH8YYY。我是BH1ZZZ，我是BH1ZZZ，我是BH1ZZZ。听到请回答。
\item BH8YYY、BH8YYY、BH8YYY。我是1ZZZ。听到请回答
\item 8YYY、8YYY、8YYY。BH1ZZZ呼叫。请过来。
\end{enumerate}
\noindent\textbf{解说：}\textbf{}。\\\noindent\textbf{答案：}

\bigskip




\noindent\textbf{试题：}业余电台BH1ZZZ用话音呼叫BH8YYY的正确格式为：
\begin{enumerate}[leftmargin=3em]
\item Bravo Hotel Eight Yankee Yankee Yankee, Bravo Hotel Eight Yankee Yankee Yankee, Bravo Hotel Eight Yankee Yankee Yankee.This is Bravo Hotel One Zulu Zulu Zulu. Bravo Hotel One Zulu Zulu Zulu, Bravo Hotel One Zulu Zulu Zulu is calling. I’m standing by.
\item Bravo Hotel Eight Yankee Yankee Yankee, Bravo Hotel Eight Yankee Yankee Yankee, Bravo Hotel Eight Yankee Yankee Yankee. Go ahead please.
\item BH8YYY, BH8YYY, BH8YYY. This One Zulu Zulu Zulu. Come in please.
\item 8YYY, 8YYY, YYY. This is BH1ZZZ. Over.
\end{enumerate}
\noindent\textbf{解说：}\textbf{}。\\\noindent\textbf{答案：}

\bigskip




\noindent\textbf{试题：}BH1ZZZ希望加入两个电台正在通信中的谈话，正确的方法为：
\begin{enumerate}[leftmargin=3em]
\item 在双方对话的间隙，短暂发射一次“Break in！”或“插入！”，如得到响应，再说明本台呼号 “BH1ZZZ请求插入”，等对方正式表示邀请后，方能加入
\item 在一方正在发射期间，短暂插入一次“Break in”，向正在收听的一方发出插入请求
\item 短暂发射一次“Break in！”或“插入！”，如对方无反应，应加大功率反复作此发射
\item 只要双方都是自己熟悉的业余电台操作员，可直接插入谈话，不必拘泥礼节
\end{enumerate}
\noindent\textbf{解说：}\textbf{}。\\\noindent\textbf{答案：}

\bigskip




\noindent\textbf{试题：}以请求插入的方式加入两个电台正在通信中的谈话，应满足的起码条件是：
\begin{enumerate}[leftmargin=3em]
\item 确认自己的加入不会影响原通信双方的乐趣
\item 自己的信号质量不亚于原通信双方
\item 自己的操作技巧不亚于原通信双方
\item 自己拥有比原通信双方更有吸引力的谈话内容
\end{enumerate}
\noindent\textbf{解说：}\textbf{}。\\\noindent\textbf{答案：}

\bigskip




\noindent\textbf{试题：}业余电台之间进行通信，必须相互正确发送和接收的信息为：
\begin{enumerate}[leftmargin=3em]
\item 本台呼号、对方呼号、信号报告
\item 本台呼号、对方呼号、QTH
\item 本台呼号、信号报告、QTH
\item 对方呼号、信号报告、设备情况
\end{enumerate}
\noindent\textbf{解说：}\textbf{}。\\\noindent\textbf{答案：}

\bigskip




\noindent\textbf{试题：}如何回答一个CQ呼叫？
\begin{enumerate}[leftmargin=3em]
\item 先报出对方的呼号，再报出自己的呼号
\item 先报出自己的呼号，再报出对方的呼号
\item 说：“CQ”，并报出对方的呼号
\item 先给出信号报告，再报出自己的呼号
\end{enumerate}
\noindent\textbf{解说：}\textbf{}。\\\noindent\textbf{答案：}

\bigskip




\noindent\textbf{试题：}当一部电台在呼叫CQ时，他的意思是？
\begin{enumerate}[leftmargin=3em]
\item 非特指地呼叫任何一部电台
\item 此电台正在测试天线，不需要任何电台回答这个呼叫
\item 只有被呼叫的电台可以回答，其他人不能回答
\item 呼叫重庆的电台
\end{enumerate}
\noindent\textbf{解说：}\textbf{}。\\\noindent\textbf{答案：}

\bigskip




\noindent\textbf{试题：}如果其他电台报告你在2米波段的信号刚才非常强，但是突然变弱或不可辨，这时你应当怎么做？
\begin{enumerate}[leftmargin=3em]
\item 稍稍移动一下自己的位置，有时信号无规律反射造成的多径效应可能导致失真
\item 打开哑音发射功能
\item 请对方电台调整自己的静噪设置
\item 将你电台中的镍氢电池换成锂电池
\end{enumerate}
\noindent\textbf{解说：}\textbf{}。\\\noindent\textbf{答案：}

\bigskip




\noindent\textbf{试题：}下列哪种方式可以让你快速切换到一个你经常使用的频率？
\begin{enumerate}[leftmargin=3em]
\item 将这个频率作为一个频道存储在电台中
\item 打开哑音输出
\item 关闭哑音输出
\item 使用快速扫描模式来切换到那个频率
\end{enumerate}
\noindent\textbf{解说：}\textbf{}。\\\noindent\textbf{答案：}

\bigskip




\noindent\textbf{试题：}参加DX网络通信有助于与一些稀有电台建立通信。正确做法是：
\begin{enumerate}[leftmargin=3em]
\item 事前了解网络规则，未经主控台允许不能随意发起呼叫，根据主控台要求进行登录，然后需随时注意主控台的安排，在主控台安排DX电台呼叫自己时及时回答联络
\item 听到DX网络通信后，应抓住机会立即对听到的电台发起呼叫
\item 当两个电台在网络主控台安排下互联联络时，自己可以通过Break in插入通信
\item DX网络时间内肯定有很多DX电台在守听，利用该频点呼叫CQ定有收获
\end{enumerate}
\noindent\textbf{解说：}\textbf{}。\\\noindent\textbf{答案：}

\bigskip




\noindent\textbf{试题：}某业余电台CW信号的片段为“by1cra test up 5”，意思是：
\begin{enumerate}[leftmargin=3em]
\item 该台的呼号为BY1CRA，正在参加竞赛，接收频率从发射频率向上偏移5kHz
\item 该台的呼号为BY1CRA，正在试验信号为UP5的发射设备
\item 该台正在试验设备并且呼叫BY1CRA，请对方在高5kHz的频率回答
\item 该台的呼号为BY1CRA，正在试验信号，并希望呼号前缀为UP5的业余电台回答
\end{enumerate}
\noindent\textbf{解说：}\textbf{}。\\\noindent\textbf{答案：}

\bigskip




\noindent\textbf{试题：}业余电台异频收发操作方式（split）的主要适用场合为：
\begin{enumerate}[leftmargin=3em]
\item 本台呼叫受到大量电台回答（pile-up）时将收发信号错开以避免干扰，或者通信双方处于不同的国际频率分配或频率规划区域时
\item 两个业余电台不想使其他业余电台听到完整的对话内容时
\item 一个业余电台使用独立的收信和发信设备时
\item 通信双方之一使用频率固定的发信设备时
\end{enumerate}
\noindent\textbf{解说：}\textbf{}。\\\noindent\textbf{答案：}

\bigskip




\noindent\textbf{试题：}在HF频段DX竞赛中回答正在呼叫CQ的参赛电台的方法是，先通过守听搞清其守听频率，然后在其发送“QRZ？”或者在其CQ呼叫停止后进行如下操作：
\begin{enumerate}[leftmargin=3em]
\item 在其守听频率上发送本台的完整呼号
\item 在其守听频率上发送对方完整呼号和本台的完整呼号
\item 在其呼叫频率上发送对方完整呼号和本台的完整呼号
\item 在其呼叫频率上发送本台的完整呼号
\end{enumerate}
\noindent\textbf{解说：}\textbf{}。\\\noindent\textbf{答案：}

\bigskip




\noindent\textbf{试题：}14022KHz有很多电台争相报出自己的呼号，原来是想呼叫发射频率为14020KHz的某稀有台。如要加入对该稀有台的呼叫，应该：
\begin{enumerate}[leftmargin=3em]
\item 守听14020KHz，在稀有台结束和其他电台联络或者呼叫CQ和QRZ时，在14022KHz快速准确地发送自己的呼号
\item 守听14020KHz，在稀有台结束和其他电台联络或者呼叫CQ和QRZ时，在14020KHz快速准确地发送自己的呼号
\item 在14022KHz不断发送自己的呼号
\item 在14020KHz呼叫该稀有台
\end{enumerate}
\noindent\textbf{解说：}\textbf{}。\\\noindent\textbf{答案：}

\bigskip




\noindent\textbf{试题：}法规要求业余电台在通信建立及结束时主动发送本台呼号。允许用发送呼号的一部分来代替发送完整呼号的情况是：
\begin{enumerate}[leftmargin=3em]
\item 在任何情况下都必须用完整呼号作为电台标识
\item 在熟悉的友台之间呼叫可以仅使用呼号后缀作为电台标识
\item 在VHF/UHF频段进行本地呼叫时可以仅使用呼号后缀作为电台标识
\item 在HF频段进行国内呼叫时可以仅使用呼号后缀作为电台标识
\end{enumerate}
\noindent\textbf{解说：}\textbf{}。\\\noindent\textbf{答案：}

\bigskip




\noindent\textbf{试题：}下列哪一项，是在竞赛中与其他电台通联时的良好习惯？
\begin{enumerate}[leftmargin=3em]
\item 尽可能不发送除了呼号和必要的信号报告以外的其他信息
\item 为了确保自己在对方的日志中，要和他通联两次
\item 如果频率上有电台堆叠，要确保每次只报自己呼号的后两位
\item 三项全部正确
\end{enumerate}
\noindent\textbf{解说：}\textbf{}。\\\noindent\textbf{答案：}

\bigskip




\noindent\textbf{试题：}“你和我还有事吗”的业余无线电通信Q简语为：
\begin{enumerate}[leftmargin=3em]
\item QRU？
\item QRL？
\item QRB？
\item QRU
\end{enumerate}
\noindent\textbf{解说：}\textbf{}。\\\noindent\textbf{答案：}

\bigskip




\noindent\textbf{试题：}“我和你无事了”的业余无线电通信Q简语为：
\begin{enumerate}[leftmargin=3em]
\item QRU
\item QRL
\item QRS
\item QRU？
\end{enumerate}
\noindent\textbf{解说：}\textbf{}。\\\noindent\textbf{答案：}

\bigskip




\noindent\textbf{试题：}“谁在呼叫我”的业余无线电通信Q简语为：
\begin{enumerate}[leftmargin=3em]
\item QRZ？
\item QRZ
\item QRA？
\item QSL？
\end{enumerate}
\noindent\textbf{解说：}\textbf{}。\\\noindent\textbf{答案：}

\bigskip




\noindent\textbf{试题：}“要我增加功率吗”的业余无线电通信Q简语为：
\begin{enumerate}[leftmargin=3em]
\item QRO？
\item QSO？
\item QRO
\item QRS？
\end{enumerate}
\noindent\textbf{解说：}\textbf{}。\\\noindent\textbf{答案：}

\bigskip




\noindent\textbf{试题：}“要我减小功率吗”的业余无线电通信Q简语为：
\begin{enumerate}[leftmargin=3em]
\item QRP？
\item QSP？
\item QRP
\item QRS？
\end{enumerate}
\noindent\textbf{解说：}\textbf{}。\\\noindent\textbf{答案：}

\bigskip




\noindent\textbf{试题：}“我能直接和×××电台通信”的业余无线电通信Q简语为：
\begin{enumerate}[leftmargin=3em]
\item QSO　×××
\item QRV　×××
\item QSP　×××
\item QRU　×××
\end{enumerate}
\noindent\textbf{解说：}\textbf{}。\\\noindent\textbf{答案：}

\bigskip




\noindent\textbf{试题：}“你能直接和×××电台通信吗”的业余无线电通信Q简语为：
\begin{enumerate}[leftmargin=3em]
\item QSO　××× ？
\item QRV　××× ？
\item QRL　××× ？
\item QRT　××× ？
\end{enumerate}
\noindent\textbf{解说：}\textbf{}。\\\noindent\textbf{答案：}

\bigskip




\noindent\textbf{试题：}“我遇到他台干扰”的业余无线电通信Q简语为：
\begin{enumerate}[leftmargin=3em]
\item QRM
\item QSM
\item QRN
\item QSB？
\end{enumerate}
\noindent\textbf{解说：}\textbf{}。\\\noindent\textbf{答案：}

\bigskip




\noindent\textbf{试题：}“你遇到他台干扰吗”的业余无线电通信Q简语为：
\begin{enumerate}[leftmargin=3em]
\item QRM？
\item QSM？
\item QSN？
\item QSD？
\end{enumerate}
\noindent\textbf{解说：}\textbf{}。\\\noindent\textbf{答案：}

\bigskip




\noindent\textbf{试题：}“你遇到天电干扰吗”的业余无线电通信Q简语为：
\begin{enumerate}[leftmargin=3em]
\item QRN？
\item QSM？
\item QSN？
\item QRV？
\end{enumerate}
\noindent\textbf{解说：}\textbf{}。\\\noindent\textbf{答案：}

\bigskip




\noindent\textbf{试题：}“要我加快发送速度吗”的业余无线电通信Q简语为：
\begin{enumerate}[leftmargin=3em]
\item QRQ？
\item QSQ？
\item QRS？
\item QRT？
\end{enumerate}
\noindent\textbf{解说：}\textbf{}。\\\noindent\textbf{答案：}

\bigskip




\noindent\textbf{试题：}“请加快发送速度”的业余无线电通信Q简语为：
\begin{enumerate}[leftmargin=3em]
\item QRQ
\item QSM
\item QSV
\item QSQ
\end{enumerate}
\noindent\textbf{解说：}\textbf{}。\\\noindent\textbf{答案：}

\bigskip




\noindent\textbf{试题：}“要我减慢发送速度吗”的业余无线电通信Q简语为：
\begin{enumerate}[leftmargin=3em]
\item QRS？
\item QSQ？
\item QRQ？
\item QRT？
\end{enumerate}
\noindent\textbf{解说：}\textbf{}。\\\noindent\textbf{答案：}

\bigskip




\noindent\textbf{试题：}“请减慢发送速度”的业余无线电通信Q简语为：
\begin{enumerate}[leftmargin=3em]
\item QRS
\item QSM
\item QSV
\item QRO
\end{enumerate}
\noindent\textbf{解说：}\textbf{}。\\\noindent\textbf{答案：}

\bigskip




\noindent\textbf{试题：}“你是否已准备好”的业余无线电通信Q简语为：
\begin{enumerate}[leftmargin=3em]
\item QRV？
\item QSV？
\item QRL？
\item QRU？
\end{enumerate}
\noindent\textbf{解说：}\textbf{}。\\\noindent\textbf{答案：}

\bigskip




\noindent\textbf{试题：}“我已准备好”的业余无线电通信Q简语为：
\begin{enumerate}[leftmargin=3em]
\item QRV
\item QSV
\item QRL
\item QRU
\end{enumerate}
\noindent\textbf{解说：}\textbf{}。\\\noindent\textbf{答案：}

\bigskip




\noindent\textbf{试题：}“要我停止发送吗”的业余无线电通信Q简语为：
\begin{enumerate}[leftmargin=3em]
\item QRT？
\item QST？
\item QRT
\item QSX
\end{enumerate}
\noindent\textbf{解说：}\textbf{}。\\\noindent\textbf{答案：}

\bigskip




\noindent\textbf{试题：}“请停止发送”的业余无线电通信Q简语为：
\begin{enumerate}[leftmargin=3em]
\item QRT
\item QST
\item QRV
\item QSB
\end{enumerate}
\noindent\textbf{解说：}\textbf{}。\\\noindent\textbf{答案：}

\bigskip




\noindent\textbf{试题：}“我的信号有衰落吗”的业余无线电通信Q简语为：
\begin{enumerate}[leftmargin=3em]
\item QSB？
\item QSD？
\item QRB？
\item QSP？
\end{enumerate}
\noindent\textbf{解说：}\textbf{}。\\\noindent\textbf{答案：}

\bigskip




\noindent\textbf{试题：}“你的信号有衰落”的业余无线电通信Q简语为：
\begin{enumerate}[leftmargin=3em]
\item QSB
\item QSX
\item QRE
\item QSP
\end{enumerate}
\noindent\textbf{解说：}\textbf{}。\\\noindent\textbf{答案：}

\bigskip




\noindent\textbf{试题：}“我发报的手法有毛病吗”的业余无线电通信Q简语为：
\begin{enumerate}[leftmargin=3em]
\item QSD？
\item QSB？
\item QRT？
\item QSV？
\end{enumerate}
\noindent\textbf{解说：}\textbf{}。\\\noindent\textbf{答案：}

\bigskip




\noindent\textbf{试题：}“你发报的手法有毛病”的业余无线电通信Q简语为：
\begin{enumerate}[leftmargin=3em]
\item QSD
\item QRC
\item QSU
\item QSK
\end{enumerate}
\noindent\textbf{解说：}\textbf{}。\\\noindent\textbf{答案：}

\bigskip




\noindent\textbf{试题：}“你正忙着吗”的业余无线电通信Q简语为：
\begin{enumerate}[leftmargin=3em]
\item QRL？
\item QRX？
\item QSU？
\item QRU？
\end{enumerate}
\noindent\textbf{解说：}\textbf{}。\\\noindent\textbf{答案：}

\bigskip




\noindent\textbf{试题：}“我正忙着”的业余无线电通信Q简语为：
\begin{enumerate}[leftmargin=3em]
\item QRL
\item QRX
\item QSV
\item QRB
\end{enumerate}
\noindent\textbf{解说：}\textbf{}。\\\noindent\textbf{答案：}

\bigskip




\noindent\textbf{试题：}“能在你的信号间隙中接收吗（即QSK插入方式）”的业余无线电通信Q简语为：
\begin{enumerate}[leftmargin=3em]
\item QSK？
\item QRF？
\item QSG？
\item QRJ？
\end{enumerate}
\noindent\textbf{解说：}\textbf{}。\\\noindent\textbf{答案：}

\bigskip




\noindent\textbf{试题：}“我给你收据（QSL卡片）、我已收妥”的业余无线电通信Q简语为：
\begin{enumerate}[leftmargin=3em]
\item QSL
\item QRG
\item QSX
\item QRV
\end{enumerate}
\noindent\textbf{解说：}\textbf{}。\\\noindent\textbf{答案：}

\bigskip




\noindent\textbf{试题：}“我的电台位置是××××”的业余无线电通信Q简语为：
\begin{enumerate}[leftmargin=3em]
\item QTH　××××
\item QRD　××××
\item QSL　××××
\item QSP　××××
\end{enumerate}
\noindent\textbf{解说：}\textbf{}。\\\noindent\textbf{答案：}

\bigskip




\noindent\textbf{试题：}“我遇到天电干扰”的业余无线电通信Q简语为：
\begin{enumerate}[leftmargin=3em]
\item QRN
\item QST
\item QSN
\item QRM
\end{enumerate}
\noindent\textbf{解说：}\textbf{}。\\\noindent\textbf{答案：}

\bigskip




\noindent\textbf{试题：}“我在发射的信号间隙中接收（即QSK插入方式）”的业余无线电通信Q简语为：
\begin{enumerate}[leftmargin=3em]
\item QSK
\item QRK
\item QSM
\item QRE
\end{enumerate}
\noindent\textbf{解说：}\textbf{}。\\\noindent\textbf{答案：}

\bigskip




\noindent\textbf{试题：}“你能给我收据（或QSL卡片）吗”的业余无线电通信Q简语为：
\begin{enumerate}[leftmargin=3em]
\item QSL？
\item QRL？
\item QSA？
\item QSD？
\end{enumerate}
\noindent\textbf{解说：}\textbf{}。\\\noindent\textbf{答案：}

\bigskip




\noindent\textbf{试题：}“你能传信到×××电台吗”的业余无线电通信Q简语为：
\begin{enumerate}[leftmargin=3em]
\item QSP ×××？
\item QRD ×××？
\item QSX ×××？
\item QRV ×××？
\end{enumerate}
\noindent\textbf{解说：}\textbf{}。\\\noindent\textbf{答案：}

\bigskip




\noindent\textbf{试题：}“我能传信到×××电台”的业余无线电通信Q简语为：
\begin{enumerate}[leftmargin=3em]
\item QSP ×××
\item QSH ×××
\item QRL ×××
\item QRP ×××
\end{enumerate}
\noindent\textbf{解说：}\textbf{}。\\\noindent\textbf{答案：}

\bigskip




\noindent\textbf{试题：}“你将在nnnn KHz（或MHz）频率守听×××电台吗”的业余无线电通信Q简语为：
\begin{enumerate}[leftmargin=3em]
\item QSX ××× ON nnnn KHz（或MHz）？
\item QRU ××× ON nnnn KHz（或MHz）？
\item QSL ××× ON nnnn KHz（或MHz）？
\item QRV ××× ON nnnn KHz（或MHz）？
\end{enumerate}
\noindent\textbf{解说：}\textbf{}。\\\noindent\textbf{答案：}

\bigskip




\noindent\textbf{试题：}“我将在nnnn KHz（或MHz）频率守听×××电台”的业余无线电通信Q简语为：
\begin{enumerate}[leftmargin=3em]
\item QSX ××× ON nnnn KHz（或MHz）
\item QRX ××× ON nnnn KHz（或MHz）
\item QRZ ××× ON nnnn KHz（或MHz）
\item QSP ××× ON nnnn KHz（或MHz）
\end{enumerate}
\noindent\textbf{解说：}\textbf{}。\\\noindent\textbf{答案：}

\bigskip




\noindent\textbf{试题：}“要我将频率改到nnnn频率吗”的业余无线电通信Q简语为：
\begin{enumerate}[leftmargin=3em]
\item QSY nnnn KHz（或MHz）？
\item QRY nnnn KHz（或MHz）？
\item QSV nnnn KHz（或MHz）？
\item QRV nnnn KHz（或MHz）？
\end{enumerate}
\noindent\textbf{解说：}\textbf{}。\\\noindent\textbf{答案：}

\bigskip




\noindent\textbf{试题：}“请将频率改到nnnn频率”的业余无线电通信Q简语为：
\begin{enumerate}[leftmargin=3em]
\item QSY nnnn KHz（或MHz）
\item QRY nnnn KHz（或MHz）
\item QSU nnnn KHz（或MHz）
\item QRO nnnn KHz（或MHz）
\end{enumerate}
\noindent\textbf{解说：}\textbf{}。\\\noindent\textbf{答案：}

\bigskip




\noindent\textbf{试题：}“你的电台位置在哪里”的业余无线电通信Q简语为：
\begin{enumerate}[leftmargin=3em]
\item QTH？
\item QRA？
\item QSA？
\item QSZ？
\end{enumerate}
\noindent\textbf{解说：}\textbf{}。\\\noindent\textbf{答案：}

\bigskip




\noindent\textbf{试题：}“我的信号强度如何”的业余无线电通信Q简语为：
\begin{enumerate}[leftmargin=3em]
\item QSA？
\item QSB？
\item QSD？
\item QTU？
\end{enumerate}
\noindent\textbf{解说：}\textbf{}。\\\noindent\textbf{答案：}

\bigskip




\noindent\textbf{试题：}“你的信号强度为×级（1-5级）”的业余无线电通信Q简语为：
\begin{enumerate}[leftmargin=3em]
\item QSA ×
\item QSL ×
\item QRM ×
\item QRY
\end{enumerate}
\noindent\textbf{解说：}\textbf{}。\\\noindent\textbf{答案：}

\bigskip




\noindent\textbf{试题：}业余无线电通信常用缩语“ABT”的意思是：
\begin{enumerate}[leftmargin=3em]
\item 关于、大约
\item 在…之上
\item 电池
\item 衰减
\end{enumerate}
\noindent\textbf{解说：}\textbf{}。\\\noindent\textbf{答案：}

\bigskip




\noindent\textbf{试题：}“地址”的业余无线电通信常用缩语是：
\begin{enumerate}[leftmargin=3em]
\item ADR或ADDR
\item ABT
\item ABV
\item ATT
\end{enumerate}
\noindent\textbf{解说：}\textbf{}。\\\noindent\textbf{答案：}

\bigskip




\noindent\textbf{试题：}业余无线电常用缩语“ATT”的意思是：
\begin{enumerate}[leftmargin=3em]
\item 衰减
\item 关于、大约
\item 衰落
\item 地址
\end{enumerate}
\noindent\textbf{解说：}\textbf{}。\\\noindent\textbf{答案：}

\bigskip




\noindent\textbf{试题：}业余无线电常用缩语“PWR”的意思是：
\begin{enumerate}[leftmargin=3em]
\item 功率
\item 中继台
\item 报告
\item 读写
\end{enumerate}
\noindent\textbf{解说：}\textbf{}。\\\noindent\textbf{答案：}

\bigskip




\noindent\textbf{试题：}“再”、“再来一次”的业余无线电通信常用缩语是：
\begin{enumerate}[leftmargin=3em]
\item AGN
\item ABT
\item ABV
\item GA
\end{enumerate}
\noindent\textbf{解说：}\textbf{}。\\\noindent\textbf{答案：}

\bigskip




\noindent\textbf{试题：}业余无线电通信常用缩语“GA”的意思是：
\begin{enumerate}[leftmargin=3em]
\item 继续、请过来
\item 公安
\item 垂直地网天线
\item 姑娘
\end{enumerate}
\noindent\textbf{解说：}\textbf{}。\\\noindent\textbf{答案：}

\bigskip




\noindent\textbf{试题：}业余无线电通信常用缩语“AHR”的意思是：
\begin{enumerate}[leftmargin=3em]
\item 另一个
\item 天线
\item 这里
\item 地址
\end{enumerate}
\noindent\textbf{解说：}\textbf{}。\\\noindent\textbf{答案：}

\bigskip




\noindent\textbf{试题：}“天线”的业余无线电通信常用缩语是：
\begin{enumerate}[leftmargin=3em]
\item ANT
\item ATT
\item ATN
\item ATR
\end{enumerate}
\noindent\textbf{解说：}\textbf{}。\\\noindent\textbf{答案：}

\bigskip




\noindent\textbf{试题：}业余无线电常用缩语“ARDF”的意思是：
\begin{enumerate}[leftmargin=3em]
\item 业余无线电测向
\item 天线调谐器、天调
\item 地址
\item 天线测试仪
\end{enumerate}
\noindent\textbf{解说：}\textbf{}。\\\noindent\textbf{答案：}

\bigskip




\noindent\textbf{试题：}“收听”的业余无线电常用缩语是：
\begin{enumerate}[leftmargin=3em]
\item KP
\item RCV
\item GA
\item HR
\end{enumerate}
\noindent\textbf{解说：}\textbf{}。\\\noindent\textbf{答案：}

\bigskip




\noindent\textbf{试题：}业余无线电常用缩语“HST”的意思是：
\begin{enumerate}[leftmargin=3em]
\item 快速收发报
\item 这里、听到
\item 信号报告
\item 通播
\end{enumerate}
\noindent\textbf{解说：}\textbf{}。\\\noindent\textbf{答案：}

\bigskip




\noindent\textbf{试题：}业余无线电CW通信常用缩语“AS”（经常连发在一起）的意思是：
\begin{enumerate}[leftmargin=3em]
\item 请稍等
\item 天线
\item 关于
\item 全部
\end{enumerate}
\noindent\textbf{解说：}\textbf{}。\\\noindent\textbf{答案：}

\bigskip




\noindent\textbf{试题：}业余无线电通信常用缩语“AS”的意思有：
\begin{enumerate}[leftmargin=3em]
\item 请稍等、亚洲、如同
\item 天线开关
\item 关于
\item 回答
\end{enumerate}
\noindent\textbf{解说：}\textbf{}。\\\noindent\textbf{答案：}

\bigskip




\noindent\textbf{试题：}业余无线电通信常用词语“BEST”的意思是：
\begin{enumerate}[leftmargin=3em]
\item 最好的
\item 电池组
\item 信号报告
\item 通播电报
\end{enumerate}
\noindent\textbf{解说：}\textbf{}。\\\noindent\textbf{答案：}

\bigskip




\noindent\textbf{试题：}业余无线电通信常用缩语“BJT”的意思是：
\begin{enumerate}[leftmargin=3em]
\item 北京时间
\item 双极型半导体管
\item 结型场效应半导体管
\item 双基极二极管
\end{enumerate}
\noindent\textbf{解说：}\textbf{}。\\\noindent\textbf{答案：}

\bigskip




\noindent\textbf{试题：}业余无线电通信常用缩语“BK”的意思是：
\begin{enumerate}[leftmargin=3em]
\item 插入、打断
\item 结束工作
\item 请马上回答
\item 千字节（单位）
\end{enumerate}
\noindent\textbf{解说：}\textbf{}。\\\noindent\textbf{答案：}

\bigskip




\noindent\textbf{试题：}“QSL卡片管理局”的业余无线电通信常用缩语是：
\begin{enumerate}[leftmargin=3em]
\item BURO
\item BOUR
\item BRUO
\item BROU
\end{enumerate}
\noindent\textbf{解说：}\textbf{}。\\\noindent\textbf{答案：}

\bigskip




\noindent\textbf{试题：}“遇到”、“见面”的业余无线电通信常用缩语是：
\begin{enumerate}[leftmargin=3em]
\item C
\item MTRS
\item SE
\item MET
\end{enumerate}
\noindent\textbf{解说：}\textbf{}。\\\noindent\textbf{答案：}

\bigskip




\noindent\textbf{试题：}业余无线电通信常用缩语“CFM”的意思是：
\begin{enumerate}[leftmargin=3em]
\item 确认
\item 法拉
\item 调频
\item 呼叫
\end{enumerate}
\noindent\textbf{解说：}\textbf{}。\\\noindent\textbf{答案：}

\bigskip




\noindent\textbf{试题：}业余无线电通信常用词语“CHEERIO”的意思是：
\begin{enumerate}[leftmargin=3em]
\item 再会、祝贺
\item 圣诞节
\item 英文字符
\item 再次见面
\end{enumerate}
\noindent\textbf{解说：}\textbf{}。\\\noindent\textbf{答案：}

\bigskip




\noindent\textbf{试题：}业余无线电通信常用缩语“CL”、“CLS”、“CLG”的意思分别是：
\begin{enumerate}[leftmargin=3em]
\item 关闭（或呼叫）、呼号、呼叫
\item 呼号、清除、关闭
\item 计数、云层、呼号
\item 确认、呼叫、清除
\end{enumerate}
\noindent\textbf{解说：}\textbf{}。\\\noindent\textbf{答案：}

\bigskip




\noindent\textbf{试题：}业余无线电通信常用词语“DATE”的意思是：
\begin{enumerate}[leftmargin=3em]
\item 日期
\item 时间
\item 频率
\item 地址
\end{enumerate}
\noindent\textbf{解说：}\textbf{}。\\\noindent\textbf{答案：}

\bigskip




\noindent\textbf{试题：}业余无线电通信常用缩语“DR”的意思是：
\begin{enumerate}[leftmargin=3em]
\item 亲爱的
\item 远距离
\item 从
\item 二极管
\end{enumerate}
\noindent\textbf{解说：}\textbf{}。\\\noindent\textbf{答案：}

\bigskip




\noindent\textbf{试题：}单元（常用于天线振子）的业余无线电通信常用缩语是：
\begin{enumerate}[leftmargin=3em]
\item EL、ELE、ELS
\item YAGI
\item ANT
\item UNIT
\end{enumerate}
\noindent\textbf{解说：}\textbf{}。\\\noindent\textbf{答案：}

\bigskip




\noindent\textbf{试题：}业余无线电CW通信常用缩语“ES”的意思是：
\begin{enumerate}[leftmargin=3em]
\item 和
\item 是
\item 从
\item 请等待
\end{enumerate}
\noindent\textbf{解说：}\textbf{}。\\\noindent\textbf{答案：}

\bigskip




\noindent\textbf{试题：}业余无线电通信常用缩语“FB”的意思是：
\begin{enumerate}[leftmargin=3em]
\item 很好的
\item 腐败
\item 美好的祝愿
\item 再见
\end{enumerate}
\noindent\textbf{解说：}\textbf{}。\\\noindent\textbf{答案：}

\bigskip




\noindent\textbf{试题：}“频率”的业余无线电通信常用缩语是：
\begin{enumerate}[leftmargin=3em]
\item FREQ
\item FER
\item TUNE
\item FIND
\end{enumerate}
\noindent\textbf{解说：}\textbf{}。\\\noindent\textbf{答案：}

\bigskip




\noindent\textbf{试题：}业余无线电通信常用缩语“GND”的意思是：
\begin{enumerate}[leftmargin=3em]
\item 地线，地面
\item 格林威治时间
\item 好运气
\item 高兴
\end{enumerate}
\noindent\textbf{解说：}\textbf{}。\\\noindent\textbf{答案：}

\bigskip




\noindent\textbf{试题：}业余无线电通信常用缩语“OM”的意思是：
\begin{enumerate}[leftmargin=3em]
\item 老朋友
\item 欧姆
\item 或者
\item 老人
\end{enumerate}
\noindent\textbf{解说：}\textbf{}。\\\noindent\textbf{答案：}

\bigskip




\noindent\textbf{试题：}“电台设备”的业余无线电通信常用缩语是：
\begin{enumerate}[leftmargin=3em]
\item RIG
\item REG
\item EQP
\item SB
\end{enumerate}
\noindent\textbf{解说：}\textbf{}。\\\noindent\textbf{答案：}

\bigskip




\noindent\textbf{试题：}业余无线电通信常用词语“FINE”的意思是：
\begin{enumerate}[leftmargin=3em]
\item 好的，精细的
\item 调谐
\item 发现
\item 确认
\end{enumerate}
\noindent\textbf{解说：}\textbf{}。\\\noindent\textbf{答案：}

\bigskip




\noindent\textbf{试题：}业余无线电通信常用缩语“FR”、“FER”的意思是：
\begin{enumerate}[leftmargin=3em]
\item 为了，对于
\item 好的，精细的
\item 希望
\item 频率
\end{enumerate}
\noindent\textbf{解说：}\textbf{}。\\\noindent\textbf{答案：}

\bigskip




\noindent\textbf{试题：}“下午好”的业余无线电通信常用缩语是：
\begin{enumerate}[leftmargin=3em]
\item GA
\item GE
\item GN
\item GM
\end{enumerate}
\noindent\textbf{解说：}\textbf{}。\\\noindent\textbf{答案：}

\bigskip




\noindent\textbf{试题：}“早晨好”的业余无线电通信常用缩语是：
\begin{enumerate}[leftmargin=3em]
\item GM
\item GA
\item GL
\item GB
\end{enumerate}
\noindent\textbf{解说：}\textbf{}。\\\noindent\textbf{答案：}

\bigskip




\noindent\textbf{试题：}“晚上好”的业余无线电通信常用缩语是：
\begin{enumerate}[leftmargin=3em]
\item GE
\item GM
\item GN
\item GA
\end{enumerate}
\noindent\textbf{解说：}\textbf{}。\\\noindent\textbf{答案：}

\bigskip




\noindent\textbf{试题：}业余无线电通信常用缩语“GN”的意思是：
\begin{enumerate}[leftmargin=3em]
\item 晚安
\item 早晨好
\item 好运气
\item 高兴
\end{enumerate}
\noindent\textbf{解说：}\textbf{}。\\\noindent\textbf{答案：}

\bigskip




\noindent\textbf{试题：}“再见”的业余无线电通信常用缩语是：
\begin{enumerate}[leftmargin=3em]
\item GB
\item GE
\item GL
\item GA
\end{enumerate}
\noindent\textbf{解说：}\textbf{}。\\\noindent\textbf{答案：}

\bigskip




\noindent\textbf{试题：}业余无线电通信常用缩语“GL”的意思是：
\begin{enumerate}[leftmargin=3em]
\item 好运气
\item 早安
\item 再见
\item 晚安
\end{enumerate}
\noindent\textbf{解说：}\textbf{}。\\\noindent\textbf{答案：}

\bigskip




\noindent\textbf{试题：}业余无线电通信常用缩语“GLD”的意思是：
\begin{enumerate}[leftmargin=3em]
\item 高兴
\item 好运气
\item 再见
\item 地线，地面
\end{enumerate}
\noindent\textbf{解说：}\textbf{}。\\\noindent\textbf{答案：}

\bigskip




\noindent\textbf{试题：}业余无线电通信常用缩语“GMT”的意思是：
\begin{enumerate}[leftmargin=3em]
\item 格林威治时间
\item 地线，地面
\item 好运气
\item 高兴
\end{enumerate}
\noindent\textbf{解说：}\textbf{}。\\\noindent\textbf{答案：}

\bigskip




\noindent\textbf{试题：}“抄收”的业余无线电通信常用缩语是：
\begin{enumerate}[leftmargin=3em]
\item CPI
\item HPY
\item HPI
\item CFM
\end{enumerate}
\noindent\textbf{解说：}\textbf{}。\\\noindent\textbf{答案：}

\bigskip




\noindent\textbf{试题：}“希望”的业余无线电通信常用缩语是：
\begin{enumerate}[leftmargin=3em]
\item HPE
\item HPY
\item HPI
\item CPI
\end{enumerate}
\noindent\textbf{解说：}\textbf{}。\\\noindent\textbf{答案：}

\bigskip




\noindent\textbf{试题：}业余无线电通信常用缩语“HPY”、“HPI”的意思是：
\begin{enumerate}[leftmargin=3em]
\item 幸福
\item 希望
\item 抄收
\item 这里
\end{enumerate}
\noindent\textbf{解说：}\textbf{}。\\\noindent\textbf{答案：}

\bigskip




\noindent\textbf{试题：}业余无线电通信常用缩语“HR”的意思是：
\begin{enumerate}[leftmargin=3em]
\item 这里、听到
\item 幸福
\item 希望
\item 号码
\end{enumerate}
\noindent\textbf{解说：}\textbf{}。\\\noindent\textbf{答案：}

\bigskip




\noindent\textbf{试题：}“怎样”、“如何”的业余无线电通信常用缩语是：
\begin{enumerate}[leftmargin=3em]
\item HW
\item CW
\item HR
\item HPI
\end{enumerate}
\noindent\textbf{解说：}\textbf{}。\\\noindent\textbf{答案：}

\bigskip




\noindent\textbf{试题：}“很多”的业余无线电通信常用缩语是：
\begin{enumerate}[leftmargin=3em]
\item MNY、MNI
\item ALL
\item VY
\item NAME
\end{enumerate}
\noindent\textbf{解说：}\textbf{}。\\\noindent\textbf{答案：}

\bigskip




\noindent\textbf{试题：}业余无线电通信常用缩语“MTRS”的意思是：
\begin{enumerate}[leftmargin=3em]
\item 米
\item 先生
\item 太太
\item 小姐
\end{enumerate}
\noindent\textbf{解说：}\textbf{}。\\\noindent\textbf{答案：}

\bigskip




\noindent\textbf{试题：}“方式”的业余无线电通信常用英语是：
\begin{enumerate}[leftmargin=3em]
\item MODE
\item NAME
\item MTRS
\item NICE
\end{enumerate}
\noindent\textbf{解说：}\textbf{}。\\\noindent\textbf{答案：}

\bigskip




\noindent\textbf{试题：}“名字”的业余无线电通信常用英语是：
\begin{enumerate}[leftmargin=3em]
\item NAME
\item MODE
\item NICE
\item MNI
\end{enumerate}
\noindent\textbf{解说：}\textbf{}。\\\noindent\textbf{答案：}

\bigskip




\noindent\textbf{试题：}业余无线电通信常用词语“NICE”的意思是：
\begin{enumerate}[leftmargin=3em]
\item 良好的
\item 方式
\item 名字
\item 鼠标
\end{enumerate}
\noindent\textbf{解说：}\textbf{}。\\\noindent\textbf{答案：}

\bigskip




\noindent\textbf{试题：}业余无线电通信常用缩语“NW”的意思是：
\begin{enumerate}[leftmargin=3em]
\item 现在
\item 不
\item 新的
\item 怎样
\end{enumerate}
\noindent\textbf{解说：}\textbf{}。\\\noindent\textbf{答案：}

\bigskip




\noindent\textbf{试题：}“操作员”的业余无线电通信常用缩语是：
\begin{enumerate}[leftmargin=3em]
\item OP、OPR
\item RPT
\item OM
\item CPI
\end{enumerate}
\noindent\textbf{解说：}\textbf{}。\\\noindent\textbf{答案：}

\bigskip




\noindent\textbf{试题：}“邮政信箱”的业余无线电通信常用缩语是：
\begin{enumerate}[leftmargin=3em]
\item P O BOX
\item MAIL
\item BURO
\item QTH
\end{enumerate}
\noindent\textbf{解说：}\textbf{}。\\\noindent\textbf{答案：}

\bigskip




\noindent\textbf{试题：}业余无线电通信常用缩语“RMKS”的意思是：
\begin{enumerate}[leftmargin=3em]
\item 备注、注释
\item 报告
\item 中继台
\item 业余无线电测向
\end{enumerate}
\noindent\textbf{解说：}\textbf{}。\\\noindent\textbf{答案：}

\bigskip




\noindent\textbf{试题：}“收信机”的业余无线电通信常用缩语是：
\begin{enumerate}[leftmargin=3em]
\item RCVR，RX
\item XCVR
\item XMTR
\item RMKS
\end{enumerate}
\noindent\textbf{解说：}\textbf{}。\\\noindent\textbf{答案：}

\bigskip




\noindent\textbf{试题：}“发信机”的业余无线电通信常用缩语是：
\begin{enumerate}[leftmargin=3em]
\item TX、XMTR
\item XTL
\item VXO
\item VXCO
\end{enumerate}
\noindent\textbf{解说：}\textbf{}。\\\noindent\textbf{答案：}

\bigskip




\noindent\textbf{试题：}“收发信机”的业余无线电通信常用缩语是：
\begin{enumerate}[leftmargin=3em]
\item XCVR
\item XVTR
\item XMTR
\item XTL
\end{enumerate}
\noindent\textbf{解说：}\textbf{}。\\\noindent\textbf{答案：}

\bigskip




\noindent\textbf{试题：}业余无线电通信常用缩语“WX”的意思是：
\begin{enumerate}[leftmargin=3em]
\item 天气
\item 瓦特
\item 联络、工作
\item 星期
\end{enumerate}
\noindent\textbf{解说：}\textbf{}。\\\noindent\textbf{答案：}

\bigskip




\noindent\textbf{试题：}业余无线电通信常用缩语“73”的意思是：
\begin{enumerate}[leftmargin=3em]
\item 向对方的致意、美好的祝愿
\item 再见
\item 希望下次再见
\item 谢谢你
\end{enumerate}
\noindent\textbf{解说：}\textbf{}。\\\noindent\textbf{答案：}

\bigskip




\noindent\textbf{试题：}业余无线电通话常用语“Roger”的用法是：
\begin{enumerate}[leftmargin=3em]
\item 回答起始语，相当于“明白”，仅在已完全抄收对方刚才发送的信息时使用
\item 回答起始语，相当于“听到”，用于能听到对方信号、但不一定能全部抄收的情况
\item 回答起始语，表示开始发话了，任何情况都可使用
\item 惯用口头语，相当于电话的“喂”，仅提起注意，不包含任何意义
\end{enumerate}
\noindent\textbf{解说：}\textbf{}。\\\noindent\textbf{答案：}

\bigskip




\noindent\textbf{试题：}业余电台通信中常用到缩写“SASE”，其意义是：
\begin{enumerate}[leftmargin=3em]
\item 写好收信人地址的信封
\item 国际邮资券
\item 请尽快寄出QSL卡片
\item 请勿通过卡片管理局交换QSL卡
\end{enumerate}
\noindent\textbf{解说：}\textbf{}。\\\noindent\textbf{答案：}

\bigskip




\noindent\textbf{试题：}“报告”的业余无线电通信常用缩语是：
\begin{enumerate}[leftmargin=3em]
\item RPRT
\item PRT
\item RMKS
\item MSG
\end{enumerate}
\noindent\textbf{解说：}\textbf{}。\\\noindent\textbf{答案：}

\bigskip




\noindent\textbf{试题：}业余无线电通信常用缩语“SK”（通常在CW中连在一起拍发）的意思是：
\begin{enumerate}[leftmargin=3em]
\item 结束通信
\item 开关
\item 谢谢
\item 下次再见
\end{enumerate}
\noindent\textbf{解说：}\textbf{}。\\\noindent\textbf{答案：}

\bigskip




\noindent\textbf{试题：}“对不起”的业余无线电通信常用缩语是：
\begin{enumerate}[leftmargin=3em]
\item SRI，SRY
\item TNX
\item SK
\item AS
\end{enumerate}
\noindent\textbf{解说：}\textbf{}。\\\noindent\textbf{答案：}

\bigskip




\noindent\textbf{试题：}“电台”的业余无线电通信常用缩语是：
\begin{enumerate}[leftmargin=3em]
\item STN
\item QTH
\item SRY
\item ANT
\end{enumerate}
\noindent\textbf{解说：}\textbf{}。\\\noindent\textbf{答案：}

\bigskip




\noindent\textbf{试题：}业余无线电通信常用缩语“SURE”的意思是：
\begin{enumerate}[leftmargin=3em]
\item 确实
\item 电台
\item 对不起
\item 短波收听者
\end{enumerate}
\noindent\textbf{解说：}\textbf{}。\\\noindent\textbf{答案：}

\bigskip




\noindent\textbf{试题：}业余无线电常用缩语“SWL”的意思是：
\begin{enumerate}[leftmargin=3em]
\item 短波收听者
\item 确认
\item 高兴
\item 对不起
\end{enumerate}
\noindent\textbf{解说：}\textbf{}。\\\noindent\textbf{答案：}

\bigskip




\noindent\textbf{试题：}“温度”的业余无线电通信常用缩语是：
\begin{enumerate}[leftmargin=3em]
\item TEMP
\item TMPO
\item TUNE
\item WX
\end{enumerate}
\noindent\textbf{解说：}\textbf{}。\\\noindent\textbf{答案：}

\bigskip




\noindent\textbf{试题：}“谢谢”的业余无线电通信常用缩语是：
\begin{enumerate}[leftmargin=3em]
\item TNX，TKS
\item SRI，SRY
\item 73
\item TRY
\end{enumerate}
\noindent\textbf{解说：}\textbf{}。\\\noindent\textbf{答案：}

\bigskip




\noindent\textbf{试题：}业余无线电通信常用缩语“TU”的意思是：
\begin{enumerate}[leftmargin=3em]
\item 谢谢你
\item 发信机
\item 天调、天线调谐器
\item 电子管
\end{enumerate}
\noindent\textbf{解说：}\textbf{}。\\\noindent\textbf{答案：}

\bigskip




\noindent\textbf{试题：}“世界协调时”的业余无线电通信常用缩语是：
\begin{enumerate}[leftmargin=3em]
\item UTC
\item TUC
\item UCT
\item CUT
\end{enumerate}
\noindent\textbf{解说：}\textbf{}。\\\noindent\textbf{答案：}

\bigskip




\noindent\textbf{试题：}业余无线电通信常用缩语“VIA”的意思是：
\begin{enumerate}[leftmargin=3em]
\item 经、由
\item 美国之音
\item 声控
\item 邮寄
\end{enumerate}
\noindent\textbf{解说：}\textbf{}。\\\noindent\textbf{答案：}

\bigskip




\noindent\textbf{试题：}“很”、“非常”的业余无线电通信常用缩语是：
\begin{enumerate}[leftmargin=3em]
\item VY
\item MNI
\item VIA
\item ALL
\end{enumerate}
\noindent\textbf{解说：}\textbf{}。\\\noindent\textbf{答案：}

\bigskip




\noindent\textbf{试题：}业余无线电通信常用缩语“WK”的意思是：
\begin{enumerate}[leftmargin=3em]
\item 星期、工作
\item 圣诞节
\item 瓦特
\item 插入
\end{enumerate}
\noindent\textbf{解说：}\textbf{}。\\\noindent\textbf{答案：}

\bigskip




\noindent\textbf{试题：}业余无线电通信常用缩语“WKD”的意思是：
\begin{enumerate}[leftmargin=3em]
\item 联络过、工作过
\item 星期
\item 天气
\item 圣诞节
\end{enumerate}
\noindent\textbf{解说：}\textbf{}。\\\noindent\textbf{答案：}

\bigskip




\noindent\textbf{试题：}业余无线电通信常用缩语“WTS”的意思是：
\begin{enumerate}[leftmargin=3em]
\item 瓦特
\item 工作、联络
\item 星期
\item 天气
\end{enumerate}
\noindent\textbf{解说：}\textbf{}。\\\noindent\textbf{答案：}

\bigskip




\noindent\textbf{试题：}业余无线电通信常用缩语“XMAS”的意思是：
\begin{enumerate}[leftmargin=3em]
\item 圣诞节
\item 发信机
\item 收发信机
\item 晶体
\end{enumerate}
\noindent\textbf{解说：}\textbf{}。\\\noindent\textbf{答案：}

\bigskip




\noindent\textbf{试题：}业余无线电通信常用缩语“XYL”的意思是：
\begin{enumerate}[leftmargin=3em]
\item 妻子、已婚女子
\item 晶体
\item 姑娘
\item 发信机
\end{enumerate}
\noindent\textbf{解说：}\textbf{}。\\\noindent\textbf{答案：}

\bigskip




\noindent\textbf{试题：}业余无线电通信常用缩语“YL”的意思是：
\begin{enumerate}[leftmargin=3em]
\item 小姐、女士
\item 好运气
\item 你的
\item 呼叫
\end{enumerate}
\noindent\textbf{解说：}\textbf{}。\\\noindent\textbf{答案：}

\bigskip




\noindent\textbf{试题：}业余无线电通信常用缩语“TU”的意思是：
\begin{enumerate}[leftmargin=3em]
\item 谢谢你
\item 结束联络
\item 再见
\item 美好的祝愿
\end{enumerate}
\noindent\textbf{解说：}\textbf{}。\\\noindent\textbf{答案：}

\bigskip




\noindent\textbf{试题：}“你的”或者“你是”的业余无线电通信常用缩语是：
\begin{enumerate}[leftmargin=3em]
\item UR
\item TU
\item FB
\item US
\end{enumerate}
\noindent\textbf{解说：}\textbf{}。\\\noindent\textbf{答案：}

\bigskip




\noindent\textbf{试题：}业余无线电通信常用缩语“88”的意思是：
\begin{enumerate}[leftmargin=3em]
\item 向对方异性操作员的致意、美好的祝愿
\item 再见
\item 祝对方发达、发财
\item 谢谢你
\end{enumerate}
\noindent\textbf{解说：}\textbf{}。\\\noindent\textbf{答案：}

\bigskip




\noindent\textbf{试题：}业余无线电通信中常用的天线种类的缩写DP代表：
\begin{enumerate}[leftmargin=3em]
\item 偶极天线
\item 长线天线
\item 定向天线
\item 垂直天线
\end{enumerate}
\noindent\textbf{解说：}\textbf{}。\\\noindent\textbf{答案：}

\bigskip




\noindent\textbf{试题：}业余无线电通信中常用的天线种类的缩写LW代表：
\begin{enumerate}[leftmargin=3em]
\item 长线天线
\item 偶极天线
\item 定向天线
\item 垂直天线
\end{enumerate}
\noindent\textbf{解说：}\textbf{}。\\\noindent\textbf{答案：}

\bigskip




\noindent\textbf{试题：}业余无线电通信中常用的天线种类的缩写GP代表：
\begin{enumerate}[leftmargin=3em]
\item 垂直接地天线
\item 对数周期天线
\item 偶极天线
\item 定向天线
\end{enumerate}
\noindent\textbf{解说：}\textbf{}。\\\noindent\textbf{答案：}

\bigskip




\noindent\textbf{试题：}业余无线电通信中常用的天线种类的缩写BEAM代表：
\begin{enumerate}[leftmargin=3em]
\item 定向天线
\item 专指八木天线
\item 偶极天线
\item 垂直天线
\end{enumerate}
\noindent\textbf{解说：}\textbf{}。\\\noindent\textbf{答案：}

\bigskip




\noindent\textbf{试题：}业余无线电通信中常用的天线种类的缩写YAGI代表：
\begin{enumerate}[leftmargin=3em]
\item 八木天线
\item 定向天线
\item 偶极天线
\item 垂直天线
\end{enumerate}
\noindent\textbf{解说：}\textbf{}。\\\noindent\textbf{答案：}

\bigskip




\noindent\textbf{试题：}业余无线电通信中常用的天线种类的缩写VER代表：
\begin{enumerate}[leftmargin=3em]
\item 垂直天线
\item 垂直接地天线
\item 定向天线
\item 偶极天线
\end{enumerate}
\noindent\textbf{解说：}\textbf{}。\\\noindent\textbf{答案：}

\bigskip




\noindent\textbf{试题：}业余无线电技术常提到的天线种类的缩写LP代表：
\begin{enumerate}[leftmargin=3em]
\item 对数周期天线
\item 垂直天线
\item 垂直接地天线
\item 定向天线
\end{enumerate}
\noindent\textbf{解说：}\textbf{}。\\\noindent\textbf{答案：}

\bigskip




\noindent\textbf{试题：}业余无线电通信方式缩写CW的英文原词意义是：
\begin{enumerate}[leftmargin=3em]
\item 等幅电报
\item 莫尔斯编码
\item 幅度键控
\item 移频键控
\end{enumerate}
\noindent\textbf{解说：}\textbf{}。\\\noindent\textbf{答案：}

\bigskip




\noindent\textbf{试题：}为了便于计算时间，将地球划分为若干个时区，各理论时区的划分方法是：
\begin{enumerate}[leftmargin=3em]
\item 全球划分为24个时区，每个理论时区宽度为经度15度，本初子午线通过0区的中心
\item 全球划分为12个时区，每个理论时区宽度为经度30度，本初子午线通过0区的中心
\item 全球划分为24个时区，每个理论时区宽度为经度15度，其边界为东西经度为15的整倍数的子午线
\item 全球划分为12个时区，每个理论时区宽度为经度30度，其边界为东西经度为30的整倍数的子午线
\end{enumerate}
\noindent\textbf{解说：}\textbf{}。\\\noindent\textbf{答案：}

\bigskip




\noindent\textbf{试题：}为了便于计算时间，将地球划分为若干个时区，各理论时区的命名规则是：
\begin{enumerate}[leftmargin=3em]
\item 本初子午线通过其中心的为0区，向东依次为东1区、东2区…东12区，向西依次为西1区、西2区…西12区
\item 本初子午线通过中心的为0区，向东依次为1区、2区…24区
\item 本初子午线通过中心的为0区，向西依次为1区、2区…24区
\item 本初子午线通过中心的为0区，如向东数则依次称为东1区、东2区…东24区，如向西数则依次称为西1区、西2区…西24区
\end{enumerate}
\noindent\textbf{解说：}\textbf{}。\\\noindent\textbf{答案：}

\bigskip




\noindent\textbf{试题：}为了便于计算时间，将地球划分为若干个时区，北京的情况是：
\begin{enumerate}[leftmargin=3em]
\item 北京处于东8区，地方时间比0时区的时间早8小时
\item 北京处于东8区，地方时间比0时区的时间晚8小时
\item 北京处于西8区，地方时间比0时区的时间早8小时
\item 北京处于西8区，地方时间比0时区的时间晚8小时
\end{enumerate}
\noindent\textbf{解说：}\textbf{}。\\\noindent\textbf{答案：}

\bigskip




\noindent\textbf{试题：}为了便于计算时间，将地球划分为若干个时区，其理论分区为每区宽经度15度。北京、西安和乌鲁木齐实际所属的时区应为：
\begin{enumerate}[leftmargin=3em]
\item 世界上实际使用法定分区，北京、西安、乌鲁木齐都属于东8区
\item 根据所在经度推算，北京、西安、乌鲁木齐分别处于东8区、东7区和东6区
\item 根据所在经度推算，北京、西安、乌鲁木齐分别处于西8区、西7区和西6区
\item 根据所在经度推算，北京、西安、乌鲁木齐分别处于东6区、东7区和东8区
\end{enumerate}
\noindent\textbf{解说：}\textbf{}。\\\noindent\textbf{答案：}

\bigskip




\noindent\textbf{试题：}已知北京时间，相应的UTC时间应为：
\begin{enumerate}[leftmargin=3em]
\item 北京时间的小时数减8，如小时数小于0，则小时数加24，日期改为前一天。
\item 北京时间的小时数减8，如小时数小于0，则小时数加24，日期改为后一天。
\item 北京时间的小时数加8，如小时数大于24，则小时数减24，日期改为前一天。
\item 北京时间的小时数加8，如小时数大于24，则小时数减24，日期改为后一天。
\end{enumerate}
\noindent\textbf{解说：}\textbf{}。\\\noindent\textbf{答案：}

\bigskip




\noindent\textbf{试题：}已知UTC时间，相应的北京时间应为：
\begin{enumerate}[leftmargin=3em]
\item UTC时间的小时数加8，如小时数大于24，则小时数减24，日期改为后一天。
\item 北京时间的小时数减8，如小时数小于0，则小时数加24，日期改为后一天。
\item 北京时间的小时数加8，如小时数大于24，则小时数减24，日期改为前一天。
\item 北京时间的小时数减8，如小时数小于0，则小时数加24，日期改为前一天。
\end{enumerate}
\noindent\textbf{解说：}\textbf{}。\\\noindent\textbf{答案：}

\bigskip




\noindent\textbf{试题：}已知某业余电台处于西N时区（N为0-12间的整数），该台的当地时间应比北京时间：
\begin{enumerate}[leftmargin=3em]
\item 晚8+N小时
\item 晚8-N小时
\item 早8+N小时
\item 早8-N小时
\end{enumerate}
\noindent\textbf{解说：}\textbf{}。\\\noindent\textbf{答案：}

\bigskip




\noindent\textbf{试题：}已知某业余电台处于东N时区（N为0-12间的整数），该台的当地时间应比北京时间：
\begin{enumerate}[leftmargin=3em]
\item 晚8-N小时
\item 晚8+N小时
\item 早8-N小时
\item 早8+N小时
\end{enumerate}
\noindent\textbf{解说：}\textbf{}。\\\noindent\textbf{答案：}

\bigskip




\noindent\textbf{试题：}为划分无线电频率，国际电信联盟《无线电规则》进行了如下的区域划分：
\begin{enumerate}[leftmargin=3em]
\item 将世界划分为3个区域，中国位于第3区
\item 将世界划分为40个区域，中国位于第24、25区
\item 将世界划分为89个区域，中国位于第33、42、43、44、45、50区
\item 将世界划分为17个区域，中国位于第8区
\end{enumerate}
\noindent\textbf{解说：}\textbf{}。\\\noindent\textbf{答案：}

\bigskip




\noindent\textbf{试题：}在业余无线电通信中，经常用到把全球分为三个区域的分区办法。制定该分区的国际机构及其公布的文件分别为：
\begin{enumerate}[leftmargin=3em]
\item 国际电信联盟ITU，《无线电规则》
\item 美国业余无线电协会ARRL，《业余无线电手册》
\item 国际业余无线电协会IARU，《IARU新闻》
\item 美国《CQ》杂志，《WAZ奖状规则》
\end{enumerate}
\noindent\textbf{解说：}\textbf{}。\\\noindent\textbf{答案：}

\bigskip




\noindent\textbf{试题：}ITU的区域划分有一套详细的规则，粗略地描述大体是：
\begin{enumerate}[leftmargin=3em]
\item 欧洲、俄罗斯亚洲部分、蒙古及部分西北亚国家为一区，南北美洲为二区，亚洲（除俄罗斯、蒙古和部分西北亚洲国家）和大洋洲为三区
\item 欧洲、俄罗斯亚洲部分、蒙古及部分西北亚国家为一区，亚洲（除俄罗斯、蒙古和部分西北亚洲国家）和大洋洲为二区，南北美洲为三区，
\item 南北美洲为一区，欧洲、俄罗斯亚洲部分、蒙古及部分西北亚国家为二区，亚洲（除俄罗斯、蒙古和部分西北亚洲国家）和大洋洲为三区
\item 南北美洲为一区，亚洲（除俄罗斯、蒙古和部分西北亚洲国家）和大洋洲为二区，欧洲、俄罗斯亚洲部分、蒙古及部分西北亚国家为三区
\end{enumerate}
\noindent\textbf{解说：}\textbf{}。\\\noindent\textbf{答案：}

\bigskip




\noindent\textbf{试题：}业余无线电通信计算成绩时，经常用到“CQ分区”。制定该分区的民间机构及其公布的文件分别为：
\begin{enumerate}[leftmargin=3em]
\item 美国《CQ》杂志，《WAZ奖状规则》
\item 美国业余无线电协会ARRL，《业余无线电手册》
\item 国际业余无线电协会ITU，《IARU新闻》
\item 英国业余无线电协会RSGB，《无线电通信》杂志
\end{enumerate}
\noindent\textbf{解说：}\textbf{}。\\\noindent\textbf{答案：}

\bigskip




\noindent\textbf{试题：}我国所属的“CQ分区”有：
\begin{enumerate}[leftmargin=3em]
\item 23、24、27
\item 42、43、44
\item 23、24
\item 42、43、44、50
\end{enumerate}
\noindent\textbf{解说：}\textbf{}。\\\noindent\textbf{答案：}

\bigskip




\noindent\textbf{试题：}我国黄岩岛、东沙岛、钓鱼岛分别属于“CQ分区”的：
\begin{enumerate}[leftmargin=3em]
\item 27、24、24
\item 24、24、25
\item 27、27、24
\item 44、44、50
\end{enumerate}
\noindent\textbf{解说：}\textbf{}。\\\noindent\textbf{答案：}

\bigskip




\noindent\textbf{试题：}“ITU分区”是IARU的活动计算通信成绩的基础。我国所属的“ITU分区”有：
\begin{enumerate}[leftmargin=3em]
\item 33、42、43、44、50
\item 33、42、43、44
\item 23、24
\item 23、24、27
\end{enumerate}
\noindent\textbf{解说：}\textbf{}。\\\noindent\textbf{答案：}

\bigskip




\noindent\textbf{试题：}“ITU分区”是IARU的活动计算通信成绩的基础。我国黄岩岛、东沙岛、钓鱼岛分别属于“ITU分区”的：
\begin{enumerate}[leftmargin=3em]
\item 50、44、44
\item 24、24、25
\item 44、44、44
\item 27、24、24
\end{enumerate}
\noindent\textbf{解说：}\textbf{}。\\\noindent\textbf{答案：}

\bigskip




\noindent\textbf{试题：}业余无线电通信梅登海德网格定位系统（Maidenhead Grid Square Locator）是一种：
\begin{enumerate}[leftmargin=3em]
\item 根据经纬度坐标对地球表面进行网格划分和命名，用以标示地理位置的系统
\item 卫星定位系统
\item 根据国际呼号系列对地球表面进行网格划分和命名，用以标示地理位置的系统
\item 根据国际政治行政区划对地球表面进行网格划分和命名，用以标示地理位置的系统
\end{enumerate}
\noindent\textbf{解说：}\textbf{}。\\\noindent\textbf{答案：}

\bigskip




\noindent\textbf{试题：}业余无线电通信常用的梅登海德网格定位系统网格名称的格式为：
\begin{enumerate}[leftmargin=3em]
\item 2个字母和2位数字、2个字母和2位数字再加2个字母
\item 4位数字或者6位数字
\item 4个字母或者6个字母
\item 呼号前缀字母加2位数字和2个字母
\end{enumerate}
\noindent\textbf{解说：}\textbf{}。\\\noindent\textbf{答案：}

\bigskip




\noindent\textbf{试题：}业余无线电通信常用的梅登海德网格定位系统网格名称的长度是4字符或6字符，两者定位精度不同，差别为：
\begin{enumerate}[leftmargin=3em]
\item 两者网格大小不同，4字符网格为经度2度和纬度1度，6字符网格为经度5分和纬度2.5分
\item 4字符网格精确到国家分区，6字符网格精确到国家的城市或县乡
\item 4字符网格根据国际呼号系列区分，6字符网格在4字符基础上加以经纬度细分
\item 4字符网格名称用于HF频段通信，6字符网格名称用于VHF/UHF通信
\end{enumerate}
\noindent\textbf{解说：}\textbf{}。\\\noindent\textbf{答案：}

\bigskip




\noindent\textbf{试题：}以下呼号前缀中，所属CQ分区与埃及相同的是：
\begin{enumerate}[leftmargin=3em]
\item 5A
\item 	5W
\item 	5T
\item 	5X
\end{enumerate}
\noindent\textbf{解说：}\textbf{}。\\\noindent\textbf{答案：}

\bigskip




\noindent\textbf{试题：}对于中国HAM来说，属于既稀有又困难的是：
\begin{enumerate}[leftmargin=3em]
\item KP5A
\item VK0HR
\item P5/4L4FN
\item KP1A
\end{enumerate}
\noindent\textbf{解说：}\textbf{}。\\\noindent\textbf{答案：}

\bigskip




\noindent\textbf{试题：}3V、4X、5A、6Y字头所代表的国家是：
\begin{enumerate}[leftmargin=3em]
\item Tunis、Israel、Libya、Jamaica
\item Guinea、Israel、Fiji Islands、Senegal
\item Libya、Israel、Jamaica、Guinea
\item Jamaica、Israel、Libya、Senegal
\end{enumerate}
\noindent\textbf{解说：}\textbf{}。\\\noindent\textbf{答案：}

\bigskip




\noindent\textbf{试题：}业余无线电通信所说的“网格定位”是什么意思？
\begin{enumerate}[leftmargin=3em]
\item 一个由一串字母和数字确定的地理位置
\item 一个由一串字母和数字确定的方位角和仰角
\item 用来调谐末级功放的设备
\item 用于无线电测向运动的设备
\end{enumerate}
\noindent\textbf{解说：}\textbf{}。\\\noindent\textbf{答案：}

\bigskip




\noindent\textbf{试题：}收发信机面板上或设置菜单中的符号VOX代表什么功能？
\begin{enumerate}[leftmargin=3em]
\item 发信机声控，接入后将根据对话筒有无语音输入的判别自动控制收发转换
\item 自动天线调谐，对天线电路的电压驻波比进行检测并进行自动补偿，以维持最小驻波比
\item 发信自动电平控制，对射频输出电平进行检测并反馈控制，以维持其在适当限度之内
\item 发信自动音量控制，对音频输入电平进行检测并反馈控制，以维持其在适当限度之内
\end{enumerate}
\noindent\textbf{解说：}\textbf{}。\\\noindent\textbf{答案：}

\bigskip




\noindent\textbf{试题：}收发信机中的PTT是指什么信号？
\begin{enumerate}[leftmargin=3em]
\item 按键发射，有信号（一般为对地接通）时发射机由等待转为发射
\item 发信语音压缩，对音频输入电平进行检测并反馈控制，以提升语音包络幅度较小的部分
\item 收信机前置放大器，在接收微弱信号时接入（此时某些技术指标可能低于额定值）
\item 自动天线调谐，对天线电路的电压驻波比进行检测并进行自动补偿，以维持最小驻波比
\end{enumerate}
\noindent\textbf{解说：}\textbf{}。\\\noindent\textbf{答案：}

\bigskip




\noindent\textbf{试题：}收发信机面板上或设置菜单中的符号SQL代表什么功能？
\begin{enumerate}[leftmargin=3em]
\item 静噪控制，检测到接收信号低于一定电平时关断音频输出
\item 发信语音压缩，对音频输入电平进行检测并反馈控制，以提升语音包络幅度较小的部分
\item 收信机前置放大器，在接收微弱信号时接入（此时某些技术指标可能低于额定值）
\item 自动天线调谐，对天线电路的电压驻波比进行检测并进行自动补偿，以维持最小驻波比
\end{enumerate}
\noindent\textbf{解说：}\textbf{}。\\\noindent\textbf{答案：}

\bigskip




\noindent\textbf{试题：}有些调频接收机的参数设置菜单有NFM和WFM两种选择。它们的含义是：
\begin{enumerate}[leftmargin=3em]
\item NFM为窄带调频方式，适用于信道带宽25kHz/12.5kHz的通信信号；WFM为宽带调频方式，适用于接收信道带宽180kHz左右的广播信号
\item NFM代表数字化语音方式，WFM代表模拟语音方式
\item NFM为调频通信本地方式（较低灵敏度），WFM为调频通信远程方式（最高灵敏度）
\item NFM为单频率守候方式，WFM为双频率守候方式
\end{enumerate}
\noindent\textbf{解说：}\textbf{}。\\\noindent\textbf{答案：}

\bigskip




\noindent\textbf{试题：}某些对讲机具有发送DTMF码的功能。缩写DTMF指的是：
\begin{enumerate}[leftmargin=3em]
\item 双音多频编码，由8个音调频率中的两个频率组合成的控制信号，代表16种状态之一，用于遥控和传输数字等简单字符
\item 亚音调静噪，即从67-250.3Hz的38个亚音调频率中选取一个作为选通信号，代表38种状态之一，接收机没有收到特定的选通信号时自动关闭音频输出
\item 数字设备识别码，即在松开PTT按键时自动发送一串代表设备代号的二进制数据
\item 自动静噪，即在接收机没有收到信号时自动关闭音频输出
\end{enumerate}
\noindent\textbf{解说：}\textbf{}。\\\noindent\textbf{答案：}

\bigskip




\noindent\textbf{试题：}某些对讲机具有发送CTCSS码的功能。缩写CTCSS指的是：
\begin{enumerate}[leftmargin=3em]
\item 亚音调静噪，即从67-250.3Hz的38个亚音调频率中选取一个作为选通信号，代表38种状态之一，接收机没有收到特定的选通信号时自动关闭音频输出
\item 双音多频编码，由8个音调频率中的两个频率组合成的控制信号，代表16种状态之一，用于遥控和传输数字等简单字符
\item 数字设备识别码，即在松开PTT按键时自动发送一串代表设备代号的二进制数据
\item 自动静噪，即在接收机没有收到信号时自动关闭音频输出
\end{enumerate}
\noindent\textbf{解说：}\textbf{}。\\\noindent\textbf{答案：}

\bigskip




\noindent\textbf{试题：}关于是否可以在FM话音通信时单凭接收机听到对方语音的音量大小来准确判断对方信号的强弱，正确答案及其理由是：
\begin{enumerate}[leftmargin=3em]
\item 不能。因为鉴频输出大小只取决于射频信号的频偏，而且正常信号的幅度会被限幅电路切齐到同样大小
\item 不能。因为信号越强，自动增益控制作用也越强，增益的急剧减小使声音反而被压低
\item 能。最后的信号是接收到的射频信号经过放大处理得到的，当然信号强声音越大
\item 能。调频信号越强，频偏也必然越大，解调后的声音也越大
\end{enumerate}
\noindent\textbf{解说：}\textbf{}。\\\noindent\textbf{答案：}

\bigskip




\noindent\textbf{试题：}用设置在NFM方式的对讲机接收WFM信号，其效果为：
\begin{enumerate}[leftmargin=3em]
\item 可以听到信号，但当调制信号幅度较大、音调较高时会发生明显非线性失真
\item 听不到信号，但接收到信号时调频噪声会变得寂静
\item 可以正常听到信号，但声音的高音频部分衰减较大，缺乏高音
\item 可以正常听到信号，但声音比较小
\end{enumerate}
\noindent\textbf{解说：}\textbf{}。\\\noindent\textbf{答案：}

\bigskip




\noindent\textbf{试题：}用设置在WFM方式的对讲机接收NFM信号，其效果为：
\begin{enumerate}[leftmargin=3em]
\item 可以正常听到信号，但声音比较小
\item 可以听到信号，但当调制信号幅度较大、音调较高时会发生明显非线性失真
\item 听不到信号，但接收到信号时调频噪声会变得寂静
\item 可以正常听到信号，但声音的高音频部分衰减较大，缺乏高音
\end{enumerate}
\noindent\textbf{解说：}\textbf{}。\\\noindent\textbf{答案：}

\bigskip




\noindent\textbf{试题：}调频接收机没有接收到信号时，会输出强烈的噪声。关于这种噪声的描述是：
\begin{enumerate}[leftmargin=3em]
\item 由天线背景噪声和机内电路噪声的随机频率变化经鉴频形成，其大小与天线接收到的背景噪声幅度无关
\item 由天线接收到的背景噪声的随机幅度变化经放大形成，其大小与天线背景噪声电压成正比
\item 由天线接收到的背景噪声的随机幅度变化经放大形成，其大小与天线背景噪声电压的平方成正比
\item 由天线接收到的背景噪声的随机幅度变化经放大形成，其大小与天线背景噪声电压的平方根成正比
\end{enumerate}
\noindent\textbf{解说：}\textbf{}。\\\noindent\textbf{答案：}

\bigskip




\noindent\textbf{试题：}如果业余中继台发射机被断断续续的干扰信号所启动，夹杂着不清楚的语音，根据覆盖区内其他业余电台的监听，确定中继台上行频率并没有电台工作。则：
\begin{enumerate}[leftmargin=3em]
\item 可能是中继台附近的两个其他发射机的强信号在中继台上行频率造成了互调干扰
\item 肯定是中继台接收机受到了人为恶意干扰
\item 可能是中继台接收机发生了寄生振荡
\item 可能是中继台发射机发生了寄生振荡
\end{enumerate}
\noindent\textbf{解说：}\textbf{}。\\\noindent\textbf{答案：}

\bigskip




\noindent\textbf{试题：}业余电台在进行业余卫星通信时使用超过常规要求的发射功率，造成的结果以及对这种做法的态度是：
\begin{enumerate}[leftmargin=3em]
\item 过强的上行信号会使卫星转发器压低对其他信道的转发功率，严重影响别人通信；必须反对
\item 上行功率越大，转发的效果越好，通信范围越大；可提倡
\item 上行功率超过一定值对通信效果改善不大，但并无明显坏处；无所谓
\item 上行功率太大造成浪费和电磁污染；不提倡
\end{enumerate}
\noindent\textbf{解说：}\textbf{}。\\\noindent\textbf{答案：}

\bigskip




\noindent\textbf{试题：}即使在空旷平地，接收到的本地VHF/UHF信号强度也可能会随着接收位置的移动而发生变化，最主要的可能原因是：
\begin{enumerate}[leftmargin=3em]
\item 直射和经地面反射等多条路径到达的电波相位不同，互相叠加或抵消造成衰落（多径效应）
\item 收发位置之间的空气流动造成电波折射不均匀而飘移
\item 不同接收位置大地导电率有差异
\item 移动过程中设备与大地之间的分布电容发生微小的变动
\end{enumerate}
\noindent\textbf{解说：}\textbf{}。\\\noindent\textbf{答案：}

\bigskip




\noindent\textbf{试题：}在相距不远的两点接收同一个远方信号，信号强度发生很大差别，且差别随两点间距离的增大呈周期性变化。这是因为：
\begin{enumerate}[leftmargin=3em]
\item 多径传播，各路径到达的信号相位延迟不同而互相干涉
\item 发射机到两接收点传播距离不同造成传播路径衰耗不同
\item 大气扰动影响
\item 地磁影响影响
\end{enumerate}
\noindent\textbf{解说：}\textbf{}。\\\noindent\textbf{答案：}

\bigskip




\noindent\textbf{试题：}下列哪种设备可以用来代替普通的扬声器，可在嘈杂的环境中更好地抄收语音信号？
\begin{enumerate}[leftmargin=3em]
\item 耳机
\item 低通滤波器
\item 视频显示器
\item 吊杆话筒
\end{enumerate}
\noindent\textbf{解说：}\textbf{}。\\\noindent\textbf{答案：}

\bigskip




\noindent\textbf{试题：}收发信机中的静噪控制的目的是什么？
\begin{enumerate}[leftmargin=3em]
\item 在没有信号的情况下，关闭音频输出，使其不会输出噪音。
\item 控制发射机的输出功率
\item 可以进行自动增益控制
\item 使接收机的输出音量调到最大
\end{enumerate}
\noindent\textbf{解说：}\textbf{}。\\\noindent\textbf{答案：}

\bigskip




\noindent\textbf{试题：}如果对方报告你的调频电台发射的信号听起来失真严重、可辨度差，可能的原因是：
\begin{enumerate}[leftmargin=3em]
\item 三项都可能
\item 电台的电源电压不足
\item 电台所处的位置不好
\item 电台的发射频率不准确
\end{enumerate}
\noindent\textbf{解说：}\textbf{}。\\\noindent\textbf{答案：}

\bigskip




\noindent\textbf{试题：}一些VHF/UHF业余无线电调频手持对讲机或车载台的设置菜单中有“全频偏”和“半频偏”的选择，其表示的意义是：
\begin{enumerate}[leftmargin=3em]
\item 分别表示信道间隔为25kHz或者12.5kHz
\item 全频偏方式下发射频率的误差比半频偏方式大一倍
\item 全频偏方式适用于收发异频的中继台通信，半频偏方式适用于同频对讲通信
\item 全频偏方式语音信号经过压缩，半频偏方式语音信号只压缩低音频分量
\end{enumerate}
\noindent\textbf{解说：}\textbf{}。\\\noindent\textbf{答案：}

\bigskip




\noindent\textbf{试题：}发现有业余电台的发射操作技巧不够规范，但还不至于造成严重的干扰和影响，正确的做法是：
\begin{enumerate}[leftmargin=3em]
\item 通过电话、邮件等方式提出善意的改进建议
\item 立即在频率上当面加以指出和纠正
\item 立即报告无线电管理机构进行干涉
\item 立即报告当地业余无线电协会，由其总部电台到频率上进行纠察
\end{enumerate}
\noindent\textbf{解说：}\textbf{}。\\\noindent\textbf{答案：}

\bigskip




\noindent\textbf{试题：}接收机设置项目中缩写“NB”和“SQL”的中文简称和作用是：
\begin{enumerate}[leftmargin=3em]
\item NB为“抑噪”，切除高于平均信号的大幅度突发脉冲噪声；SQL为“静噪”，信噪比达不到一定水平时自动关闭音频输出
\item NB和SQL都是指“抑噪”，收不到有用信号时自动关断背景噪声
\item NB和SQL都是指“静噪”，切除高于平均信号的大幅度突发脉冲噪声
\item NB和SQL都是指“静噪”，收不到带有预期的特定控制信号时自动关断音频输出
\end{enumerate}
\noindent\textbf{解说：}\textbf{}。\\\noindent\textbf{答案：}

\bigskip




\noindent\textbf{试题：}单边带发信机语音压缩的作用是：
\begin{enumerate}[leftmargin=3em]
\item 压低较强语音信号的幅度、提升较弱信号的幅度，以改善小幅度语音在接收端的信噪比
\item 压低较弱语音信号的幅度、提升较强信号的幅度，以增加语音的动态范围和抑扬顿挫感
\item 压低语音信号的低频分量，提升高频分量，以增加信号的带宽，使高音更加细腻
\item 压缩信号所占用的频谱宽度，提高无线电频谱的利用率
\end{enumerate}
\noindent\textbf{解说：}\textbf{}。\\\noindent\textbf{答案：}

\bigskip




\noindent\textbf{试题：}收发信机面板上的符号ALC代表什么功能？
\begin{enumerate}[leftmargin=3em]
\item 发信自动电平控制，对射频输出电平进行检测并反馈控制，以维持其在适当限度之内
\item 发信自动音量控制，对音频输入电平进行检测并反馈控制，以维持其在适当限度之内
\item 自动天线调谐，对天线电路的电压驻波比进行检测并进行自动补偿，以维持最小驻波比
\item 自动频率控制，对发射频率的漂移进行检测并反馈控制，以维持准确的工作频率
\end{enumerate}
\noindent\textbf{解说：}\textbf{}。\\\noindent\textbf{答案：}

\bigskip




\noindent\textbf{试题：}收发信机面板上的符号AT代表什么功能？
\begin{enumerate}[leftmargin=3em]
\item 自动天线调谐，对天线电路的电压驻波比进行检测并进行自动补偿，以维持最小驻波比
\item 发信自动音量控制，对音频输入电平进行检测并反馈控制，以维持其在适当限度之内
\item 发信自动电平控制，对射频输出电平进行检测并反馈控制，以维持其在适当限度之内
\item 自动频率控制，对发射频率的漂移进行检测并反馈控制，以维持准确的工作频率
\end{enumerate}
\noindent\textbf{解说：}\textbf{}。\\\noindent\textbf{答案：}

\bigskip




\noindent\textbf{试题：}收发信机面板上的符号ATT代表什么功能？
\begin{enumerate}[leftmargin=3em]
\item 收信机输入衰减器，在接收大信号时接入，使信号不致过大而使前级电路过载
\item 自动天线调谐，对天线电路的电压驻波比进行检测并进行自动补偿，以维持最小驻波比
\item 发信自动电平控制，对射频输出电平进行检测并反馈控制，以维持其在适当限度之内
\item 发信自动音量控制，对音频输入电平进行检测并反馈控制，以维持其在适当限度之内
\end{enumerate}
\noindent\textbf{解说：}\textbf{}。\\\noindent\textbf{答案：}

\bigskip




\noindent\textbf{试题：}收发信机面板上的符号AGC代表什么功能？
\begin{enumerate}[leftmargin=3em]
\item 收信机自动增益控制，对中频级信号电平进行检测并反馈控制，防止电路过载
\item 收信自动音量控制，对音频输出电平进行检测并反馈控制，以维持其在适当限度之内
\item 自动天线调谐，对天线电路的电压驻波比进行检测并进行自动补偿，以维持最小驻波比
\item 发信自动电平控制，对射频输出电平进行检测并反馈控制，以维持其在适当限度之内
\end{enumerate}
\noindent\textbf{解说：}\textbf{}。\\\noindent\textbf{答案：}

\bigskip




\noindent\textbf{试题：}收发信机面板上的符号PRE代表什么功能？
\begin{enumerate}[leftmargin=3em]
\item 收信机前置放大器，在接收微弱信号时接入（此时某些技术指标可能低于额定值）
\item 自动天线调谐，对天线电路的电压驻波比进行检测并进行自动补偿，以维持最小驻波比
\item 发信自动电平控制，对射频输出电平进行检测并反馈控制，以维持其在适当限度之内
\item 发信语音压缩，对音频输入电平进行检测并反馈控制，以提升语音包络幅度较小的部分
\end{enumerate}
\noindent\textbf{解说：}\textbf{}。\\\noindent\textbf{答案：}

\bigskip




\noindent\textbf{试题：}收发信机面板上的符号PROC代表什么功能？
\begin{enumerate}[leftmargin=3em]
\item 发信语音压缩，对音频输入电平进行检测并反馈控制，以提升语音包络幅度较小的部分
\item 收信机前置放大器，在接收微弱信号时接入（此时某些技术指标可能低于额定值）
\item 自动天线调谐，对天线电路的电压驻波比进行检测并进行自动补偿，以维持最小驻波比
\item 发信自动电平控制，对射频输出电平进行检测并反馈控制，以维持其在适当限度之内
\end{enumerate}
\noindent\textbf{解说：}\textbf{}。\\\noindent\textbf{答案：}

\bigskip




\noindent\textbf{试题：}业余收发信机面饭上RIT的中文名称和代表的意义是：
\begin{enumerate}[leftmargin=3em]
\item 接收增量调谐，在接收频率的主调谐不变的基础上，对接收频率进行附加微调
\item 发射增量调谐，在发射频率的主调谐不变的基础上，对发射频率进行附加微调
\item 异频收发，接收和发射使用互相独立的频率
\item 清除信道频率存贮器的所有数据
\end{enumerate}
\noindent\textbf{解说：}\textbf{}。\\\noindent\textbf{答案：}

\bigskip




\noindent\textbf{试题：}业余收发信机面饭上XIT的中文名称和代表的意义是：
\begin{enumerate}[leftmargin=3em]
\item 发射增量调谐，在发射频率的主调谐不变的基础上，对发射频率进行附加微调
\item 接收增量调谐，在接收频率的主调谐不变的基础上，对接收频率进行附加微调
\item 异频收发，接收和发射使用互相独立的频率
\item 清除信道频率存贮器的所有数据
\end{enumerate}
\noindent\textbf{解说：}\textbf{}。\\\noindent\textbf{答案：}

\bigskip




\noindent\textbf{试题：}业余收发信机面饭上SPLIT的中文名称和代表的意义是：
\begin{enumerate}[leftmargin=3em]
\item 异频收发，接收和发射使用互相独立的频率
\item 发射增量调谐，在发射频率的主调谐不变的基础上，对发射频率进行附加微调
\item 接收增量调谐，在接收频率的主调谐不变的基础上，对接收频率进行附加微调
\item 清除信道频率存贮器的所有数据
\end{enumerate}
\noindent\textbf{解说：}\textbf{}。\\\noindent\textbf{答案：}

\bigskip




\noindent\textbf{试题：}应该关闭接收机AGC功能的情况是：
\begin{enumerate}[leftmargin=3em]
\item 有用微弱信号和强干扰同时出现时
\item 接收微弱信号时
\item 接收特强信号时
\item 有用强信号中夹杂着微弱干扰
\end{enumerate}
\noindent\textbf{解说：}\textbf{}。\\\noindent\textbf{答案：}

\bigskip




\noindent\textbf{试题：}应选择较短的接收机AGC时间常数的情况是：
\begin{enumerate}[leftmargin=3em]
\item 接收FM/FSK/PSK等包络幅度恒定的信号
\item 接收SSB和AM等正常信号的包络幅度不断变化的信号
\item 将接收机应用于“比幅度法规”测向时
\item 有用微弱信号和强干扰同时出现时
\end{enumerate}
\noindent\textbf{解说：}\textbf{}。\\\noindent\textbf{答案：}

\bigskip




\noindent\textbf{试题：}应选择较长的接收机AGC时间常数的情况是：
\begin{enumerate}[leftmargin=3em]
\item 接收SSB和AM等正常信号的包络幅度不断变化的信号
\item 接收FM/FSK/PSK等包络幅度恒定的信号遇到变化很快的传播衰落时
\item 将接收机应用于“比幅度法规”测向时
\item 有用微弱信号和强干扰同时出现时
\end{enumerate}
\noindent\textbf{解说：}\textbf{}。\\\noindent\textbf{答案：}

\bigskip




\noindent\textbf{试题：}收听射频/中频增益和音频增益分开控制的通信接收机时，较好的设置方法是：
\begin{enumerate}[leftmargin=3em]
\item 信号特弱时尽量把射频/中频增益开到最大，信号特强时尽量把音频增益开到最大，然后从低到高调整另一个增益以得到适当的音量
\item 信号特弱时尽量把音频增益开到最大，信号特强时尽量把射频/中频增益开到最大，然后从低到高调整另一个增益以得到适当的音量
\item 任何情况下都应将射频/中频增益放在中间位置，然后从低到高调整音频增益以得到适当的音量
\item 任何情况下都应将音频增益放在中间位置，然后从低到高调整射频/中频增益以得到适当的音量
\end{enumerate}
\noindent\textbf{解说：}\textbf{}。\\\noindent\textbf{答案：}

\bigskip




\noindent\textbf{试题：}如果短波业余电台附近环境有不可避免的强烈噪声源影响接收微弱信号，合理的做法是：
\begin{enumerate}[leftmargin=3em]
\item 尽量只呼叫和回答能听到的电台，必须发起CQ呼叫时应降低功率
\item 发起CQ呼叫时应增大功率，以便压倒环境噪声
\item 发起正常满功率CQ呼叫，并应尽量调小接收机射频增益
\item 发起正常满功率CQ呼叫，并设置好接收机的NB、AGC等功能
\end{enumerate}
\noindent\textbf{解说：}\textbf{}。\\\noindent\textbf{答案：}

\bigskip




\noindent\textbf{试题：}某业余通信接收机的中频滤波器带宽有100Hz、400Hz、2.7kHz和6kHz几挡选择。如果要为接收SSB、AM、PSK31和CW方式的信号分别从中选择合适的挡位，应该依次为：
\begin{enumerate}[leftmargin=3em]
\item 2.7kHz、6kHz、100Hz、400Hz
\item 6kHz、2.7kHz、400Hz、100Hz
\item 2.7kHz、400Hz、6kHz、100Hz
\item 2.7kHz、100Hz、6kHz、400Hz
\end{enumerate}
\noindent\textbf{解说：}\textbf{}。\\\noindent\textbf{答案：}

\bigskip




\noindent\textbf{试题：}向无线电发射机的话筒输入端施加电压幅度10倍于规定值的音频信号，可能会得到什么效果？
\begin{enumerate}[leftmargin=3em]
\item 电路失真产生谐波和互调产物，加大占用带宽，而通信效果很差
\item 有效延长通信距离，使信号听起来更加饱满
\item 有效延长通信距离，但信号听起来带有失真
\item 使发射功率超过额定值而烧坏发射机末级射频功率器件
\end{enumerate}
\noindent\textbf{解说：}\textbf{}。\\\noindent\textbf{答案：}

\bigskip




\noindent\textbf{试题：}在接收通信电台的信号时，通常用射频增益和音频增益旋钮就可以获得适当的灵敏度和音量。为什么很多业余无线电接收机还要另设一个衰减（ATT）开关？
\begin{enumerate}[leftmargin=3em]
\item 特强带外干扰可使前级器件进入非线性区而产生互调干扰，需在接收机最前端加入衰减
\item 通常增益控制旋钮的控制范围不够宽，ATT开关用来加宽增益控制范围
\item 遇特强的带内干扰时，打开ATT开关可防止音量过大而损坏扬声器或耳机	
\item 避免本台发射机的强信号损坏接收电路
\end{enumerate}
\noindent\textbf{解说：}\textbf{}。\\\noindent\textbf{答案：}

\bigskip




\noindent\textbf{试题：}单边带发信机中发信自动电平控制ALC的主要作用是：
\begin{enumerate}[leftmargin=3em]
\item 防止过驱动带来的调制失真
\item 改善发信频率的稳定度
\item 实现天线电路阻抗的自动匹配
\item 防止话筒过于灵敏造成的背景噪音
\end{enumerate}
\noindent\textbf{解说：}\textbf{}。\\\noindent\textbf{答案：}

\bigskip




\noindent\textbf{试题：}业余电台发射单边带语音信号中，语音虽然基本正常，但操作员周围噪杂的声音很响，应该：
\begin{enumerate}[leftmargin=3em]
\item 调低发射机的话筒增益
\item 调低发射机的射频输出功率
\item 重新调整发射机天线电路的匹配
\item 重新调整发信机的自动电平控制（ALC）
\end{enumerate}
\noindent\textbf{解说：}\textbf{}。\\\noindent\textbf{答案：}

\bigskip




\noindent\textbf{试题：}通常的成品收发信机上都有AGC功能，下面那种用法是正确的？
\begin{enumerate}[leftmargin=3em]
\item 通常在SSB模式时，使用慢速或中速 AGC
\item FM 模式使用快速 AGC
\item 在接收小信号且有临频大信号干扰时关闭
\item 将AGC 设为AUTO ，可以适应各种调制模式的接收
\end{enumerate}
\noindent\textbf{解说：}\textbf{}。\\\noindent\textbf{答案：}

\bigskip




\noindent\textbf{试题：}如果将发射机的话筒增益设置得过高会导致什么问题？
\begin{enumerate}[leftmargin=3em]
\item 发射机发射的信号可能会失真
\item 发射机的输出功率将会特别高
\item 发射机的频率会变得不稳定
\item 驻波比会增加
\end{enumerate}
\noindent\textbf{解说：}\textbf{}。\\\noindent\textbf{答案：}

\bigskip




\noindent\textbf{试题：}电台的下列哪一项控制功能可以使听起来音调偏高或偏低的SSB语音信号变得正常？
\begin{enumerate}[leftmargin=3em]
\item RIT功能
\item 带宽选择
\item 哑音静噪
\item 自动增益控制或限幅器
\end{enumerate}
\noindent\textbf{解说：}\textbf{}。\\\noindent\textbf{答案：}

\bigskip




\noindent\textbf{试题：}1.9MHz和3.5MHz等短波低频段业余DX通信经常使用到名词“灰线”。这是指：
\begin{enumerate}[leftmargin=3em]
\item 地球自转时白昼和黑夜交界的区域
\item 地球有极光活动的地区
\item 连接地球具有相同最高可用频率的地点所形成的线
\item 地球上通信两端之间的大圆连线
\end{enumerate}
\noindent\textbf{解说：}\textbf{}。\\\noindent\textbf{答案：}

\bigskip




\noindent\textbf{试题：}业余无线电爱好者经常用“灰线”预测1.9MHz和3.5MHz等短波低频段的DX通信效果，基本原理是：
\begin{enumerate}[leftmargin=3em]
\item 通信两端同时处于灰线时最有利于建立天波的多跳反射途径
\item 通信两端分别处于灰线两侧4000公里以外的对称点时传播最好
\item 通信两端分别处于灰线同一侧4000公里以外的两点时传播最好
\item 通信两端应避免同时处于灰线才能取得更稳定的传播
\end{enumerate}
\noindent\textbf{解说：}\textbf{}。\\\noindent\textbf{答案：}

\bigskip




\noindent\textbf{试题：}DX SUMMIT上出现这样一行：“PY2SS 28005.3 KH6/N6GQ Simplex, fading signal in SA 2040 15 Feb 2012”，以下描述正确的是：
\begin{enumerate}[leftmargin=3em]
\item 巴西电台收听夏威夷电台效果不好
\item 巴西电台收听关岛电台效果很好
\item 巴西电台收听夏威夷电台效果很好
\item 巴西电台收听关岛电台效果不好
\end{enumerate}
\noindent\textbf{解说：}\textbf{}。\\\noindent\textbf{答案：}

\bigskip




\noindent\textbf{试题：}DX SUMMIT上出现这样一行：JA2WS 14009.6 E51WWA Calling CQ, no takers, Cook Is0052 14 Feb 2008。其正确的意思是：
\begin{enumerate}[leftmargin=3em]
\item E51WWA在库克岛，呼叫CQ，没人应答
\item E51WWA在库克岛，呼叫CQ，很多人应答
\item JA2WS 在库克岛，呼叫CQ，没人应答
\item JA2WS在库克岛，呼叫CQ，很多人应答
\end{enumerate}
\noindent\textbf{解说：}\textbf{}。\\\noindent\textbf{答案：}

\bigskip




\noindent\textbf{试题：}当最高可用频率（MUF）为20MHz时哪个波段具有较大的成功通联机会？
\begin{enumerate}[leftmargin=3em]
\item 17m
\item 20m
\item 15m
\item 12m
\end{enumerate}
\noindent\textbf{解说：}\textbf{}。\\\noindent\textbf{答案：}

\bigskip




\noindent\textbf{试题：}IOTA活动在20米波段SSB模式和CW模式常用的频率是多少？
\begin{enumerate}[leftmargin=3em]
\item 14.260 MHz / 14.040 MHz
\item 14.190 MHz / 14. 040 MHz
\item 14.190 MHz / 14.060 MHz
\item 14.260 MHz / 14.060 MHz
\end{enumerate}
\noindent\textbf{解说：}\textbf{}。\\\noindent\textbf{答案：}

\bigskip




\noindent\textbf{试题：}NCDXF信标台在20米、17米、15米和10米波段的的频率是？
\begin{enumerate}[leftmargin=3em]
\item 14.100 / 18.110 / 21.150 / 28.200 MHz
\item 14.100 / 18.110 /21.150 / 28.100 MHz
\item 14.100 / 18.150 /21.150 / 28.200 MHz
\item 14.100 / 18.100 / 21.150 /28.200 MHz
\end{enumerate}
\noindent\textbf{解说：}\textbf{}。\\\noindent\textbf{答案：}

\bigskip




\noindent\textbf{试题：}在业余无线电中，莫尔斯电码的信号速度经常用WPM作为单位，其计量方法为：
\begin{enumerate}[leftmargin=3em]
\item 每分钟包含的无间隙“点”的个数除以50即为WPM数
\item 每分钟实际拍发的英文字符数即为WPM数
\item 每分钟实际拍发的中文字符数即为WPM数
\item 每分钟包含的无间隙“点”的个数除以60即为WPM数
\end{enumerate}
\noindent\textbf{解说：}\textbf{}。\\\noindent\textbf{答案：}

\bigskip




\noindent\textbf{试题：}在业余无线电中，莫尔斯电码的信号速度经常用WPM作为单位，其计量方法为：
\begin{enumerate}[leftmargin=3em]
\item 每分钟拍发连续英文单词“PARIS”的个数即为WPM数
\item 每分钟实际拍发的英文字符数即为WPM数
\item 每分钟实际拍发的中文字符数即为WPM数
\item 每分钟拍发连续“点”的个数除以60即为WPM数
\end{enumerate}
\noindent\textbf{解说：}\textbf{}。\\\noindent\textbf{答案：}

\bigskip




\noindent\textbf{试题：}在业余无线电中，计算莫尔斯电码的WPM速度时采用的信号时值标准（以一个“点”信号的时间为比较基准）为，点信号、划信号、字符内点划信号的间隔、字符之间的间隔、单词（组）之间的间隔分别为：
\begin{enumerate}[leftmargin=3em]
\item 1、3、1、3、7
\item 1、3、1、3、5
\item 1、5、1、5、7
\item 1、5、2、3、3
\end{enumerate}
\noindent\textbf{解说：}\textbf{}。\\\noindent\textbf{答案：}

\bigskip




\noindent\textbf{试题：}SSTV或ATV接收到的图像有时发生左右倾斜为平行四边形或斜粗条纹的现象，是因为：
\begin{enumerate}[leftmargin=3em]
\item 双方的扫描行频不一致。
\item 双方的行频、帧频一致，但是帧扫描的初始时刻不同步
\item 双方设备没有准确调谐在同一频率上
\item 双方设备没有正确设置调制极性
\end{enumerate}
\noindent\textbf{解说：}\textbf{}。\\\noindent\textbf{答案：}

\bigskip




\noindent\textbf{试题：}SSTV 或ATV接收到的图像有时发生中间出现水平粗黑条纹、上下半帧图像错位的现象，是因为：
\begin{enumerate}[leftmargin=3em]
\item 双方的行频、帧频一致，但是帧扫描的初始时刻不同步
\item 双方的扫描行频不一致。
\item 双方设备没有准确调谐在同一频率上
\item 双方设备没有正确设置调制极性
\end{enumerate}
\noindent\textbf{解说：}\textbf{}。\\\noindent\textbf{答案：}

\bigskip




\noindent\textbf{试题：}业余无线电慢扫描电视传送的是：
\begin{enumerate}[leftmargin=3em]
\item 逐行扫描的静止图像
\item 交叉扫描的静止图像
\item 逐行扫描的活动图像
\item 交叉扫描的活动图像
\end{enumerate}
\noindent\textbf{解说：}\textbf{}。\\\noindent\textbf{答案：}

\bigskip




\noindent\textbf{试题：}业余无线电模拟慢扫描电视图像的常用行分辨率为：
\begin{enumerate}[leftmargin=3em]
\item 每帧128或256行
\item 每帧320或640行
\item 每帧640或720行
\item 每帧768或1024行
\end{enumerate}
\noindent\textbf{解说：}\textbf{}。\\\noindent\textbf{答案：}

\bigskip




\noindent\textbf{试题：}业余无线电模拟慢扫描电视图像的音频基带调制方式以及图像信息和调制信号的对应关系为：
\begin{enumerate}[leftmargin=3em]
\item 调频，白电平为高频率（2300Hz），黑电平为低频率（1500Hz），同步电平为最低频率（1200Hz）
\item 调幅，白电平为高电平（87.5%），黑电平为低电平（30%），同步电平为最低电平（25%）
\item 调幅，白电平为低电平（75%），黑电平为高电平（75%），同步电平为最高电平（100%）
\item 调幅，白电平为低频率（1200Hz），黑电平为高频率（2000Hz），同步电平为最高频率（2300Hz）
\end{enumerate}
\noindent\textbf{解说：}\textbf{}。\\\noindent\textbf{答案：}

\bigskip




\noindent\textbf{试题：}某业余电台要通过短波话音信道传输分辨率为1024x768像素甚至更为精细的电路图，适当的通信方式为： 
\begin{enumerate}[leftmargin=3em]
\item 数字慢扫描电视（DSSTV）或无线电传真（FAX）
\item 业余电视（ATV）
\item 模拟慢扫描电视（SSTV）
\item 无线电传（RTTY）
\end{enumerate}
\noindent\textbf{解说：}\textbf{}。\\\noindent\textbf{答案：}

\bigskip




\noindent\textbf{试题：}在监听Robot、Scotte、Martin、SC2等方式的业余SSTV通信时，可以听到信号中伴有一种非常规则的每秒钟重复数次的短促的“咯、咯”声，它们是：
\begin{enumerate}[leftmargin=3em]
\item 行同步信号
\item 一种超低频寄生振荡
\item 标准频率信标
\item 图像信号和载频之间的差拍
\end{enumerate}
\noindent\textbf{解说：}\textbf{}。\\\noindent\textbf{答案：}

\bigskip




\noindent\textbf{试题：}在业余模拟SSTV通信中，Robot、Scottie、AVT、Martin等不同传送方式有什么差别？
\begin{enumerate}[leftmargin=3em]
\item 像素分辨率、像素调制方法、行频、帧频等参数不同，适用于不同画质和环境条件
\item 占用频带宽度不同
\item 要求的接收机灵敏度不同
\item 操作难度不同
\end{enumerate}
\noindent\textbf{解说：}\textbf{}。\\\noindent\textbf{答案：}

\bigskip




\noindent\textbf{试题：}深空（deep space)是指离地球的距离约等于或大于下列值的空间：
\begin{enumerate}[leftmargin=3em]
\item 200万千米
\item 20万千米
\item 10万米
\item 15000米
\end{enumerate}
\noindent\textbf{解说：}\textbf{}。\\\noindent\textbf{答案：}

\bigskip




\noindent\textbf{试题：}业余地球卫星的轨道倾角（inclination）是指包含：
\begin{enumerate}[leftmargin=3em]
\item 轨道平面与地球赤道平面的夹角，由地球赤道平面在轨道升交点按逆时针方向计决定
\item 轨道平面与地球赤道平面的夹角，由地球赤道平面在轨道升交点按顺时针方向计决定
\item 轨道平面与地球赤道平面的夹角，由地球赤道平面在轨道降交点按逆时针方向计决定
\item 轨道平面与地球赤道平面的夹角，由地球赤道平面在轨道降交点按顺时针方向计决定
\end{enumerate}
\noindent\textbf{解说：}\textbf{}。\\\noindent\textbf{答案：}

\bigskip




\noindent\textbf{试题：}业余卫星的周期（period）是指：
\begin{enumerate}[leftmargin=3em]
\item 连续两次经过其轨道上的某特定点的间隔时间
\item 连续两次经过地球上某观测点的间隔时间
\item 连续两次经过地球赤道的间隔时间
\item 连续两次经过地球北极的间隔时间
\end{enumerate}
\noindent\textbf{解说：}\textbf{}。\\\noindent\textbf{答案：}

\bigskip




\noindent\textbf{试题：}地球同步（geosynchronous）卫星是指：
\begin{enumerate}[leftmargin=3em]
\item 运行周期等于地球自转周期的地球卫星
\item 所经过地点的地方时基本相同的卫星
\item 其轨道平面通过地球北极和南极地区的卫星
\item 瞬时轨道平面与太阳始终保持固定取向的卫星
\end{enumerate}
\noindent\textbf{解说：}\textbf{}。\\\noindent\textbf{答案：}

\bigskip




\noindent\textbf{试题：}太阳同步（轨道）（Sun-synchronousorbit）卫星是指：
\begin{enumerate}[leftmargin=3em]
\item 瞬时轨道平面与太阳始终保持固定取向的卫星
\item 对地球保持大致相对静止的卫星
\item 圆形及顺行轨道位于地球赤道平面上，并对地球保持相对静止的卫星
\item 运行周期等于地球自转周期的地球卫星
\end{enumerate}
\noindent\textbf{解说：}\textbf{}。\\\noindent\textbf{答案：}

\bigskip




\noindent\textbf{试题：}地球(geostationary)静止卫星是指：
\begin{enumerate}[leftmargin=3em]
\item 对地球保持大致相对静止的卫星
\item 所经过地点的地方时基本相同的卫星
\item 其轨道平面通过地球北极和南极地区的卫星
\item 所有的地球同步卫星
\end{enumerate}
\noindent\textbf{解说：}\textbf{}。\\\noindent\textbf{答案：}

\bigskip




\noindent\textbf{试题：}卫星升轨道（ascending pass）的方向是指轨道的下列区段：
\begin{enumerate}[leftmargin=3em]
\item 由南向北
\item 由北向南
\item 由东北向西南
\item 由东南向西北
\end{enumerate}
\noindent\textbf{解说：}\textbf{}。\\\noindent\textbf{答案：}

\bigskip




\noindent\textbf{试题：}卫星降轨道（descending pass）的方向是指轨道的下列区段：
\begin{enumerate}[leftmargin=3em]
\item 由北向南
\item 由南向北
\item 由东北向西南
\item 由东南向西北
\end{enumerate}
\noindent\textbf{解说：}\textbf{}。\\\noindent\textbf{答案：}

\bigskip




\noindent\textbf{试题：}卫星的周期是指：
\begin{enumerate}[leftmargin=3em]
\item 卫星沿轨道绕地球一周所需的时间
\item 卫星绕质心自旋一周所需的时间
\item 卫星两次正好从地面某一点的正上方通过的间隔时间
\item 卫星随地球绕太阳一周所需的时间
\end{enumerate}
\noindent\textbf{解说：}\textbf{}。\\\noindent\textbf{答案：}

\bigskip




\noindent\textbf{试题：}在业余无线电领域中，缩写AMSAT是指：
\begin{enumerate}[leftmargin=3em]
\item 美国的业余无线电卫星公司，一个专注于实验卫星的设计、制造、运行和推进空间教育的非营利性志愿者组织
\item 所有业余无线电卫星的总称
\item 某一系列业余无线电卫星的总称
\item 某一颗业余无线电卫星的名称
\end{enumerate}
\noindent\textbf{解说：}\textbf{}。\\\noindent\textbf{答案：}

\bigskip




\noindent\textbf{试题：}在描述业余无线电卫星时经常用到缩写OSCAR（奥斯卡），这个名称的由来是：
\begin{enumerate}[leftmargin=3em]
\item “搭载有业余无线电装置的地球轨道卫星”的英文缩写
\item 按照发明者奥斯卡的方案设计制造的业余卫星
\item 纪念业余无线电技术先驱者奥斯卡
\item 电影奥斯卡金奖的基金所赞助的业余卫星活动
\end{enumerate}
\noindent\textbf{解说：}\textbf{}。\\\noindent\textbf{答案：}

\bigskip




\noindent\textbf{试题：}我国发射的第一颗业余卫星的发射年份、名称、国际OSCAR系列号和转发器模式分别为：
\begin{enumerate}[leftmargin=3em]
\item 2009年，希望一号（XW-1），HO68，V/U（J）
\item 2008年，希望一号（XW-1），HO68，V/U（J）
\item 2010年，希望一号（XW-1），HO68，V/U（J）
\item 2008年，希望一号（XW-1），无OSCAR编号，B
\end{enumerate}
\noindent\textbf{解说：}\textbf{}。\\\noindent\textbf{答案：}

\bigskip




\noindent\textbf{试题：}业余卫星一般采用低轨道或者椭圆轨道设计，主要原因是：
\begin{enumerate}[leftmargin=3em]
\item 发射成本低
\item 卫星转发器功率可以增加
\item 增强地球吸引力以保持运转稳定
\item 增加卫星转发器覆盖面积
\end{enumerate}
\noindent\textbf{解说：}\textbf{}。\\\noindent\textbf{答案：}

\bigskip




\noindent\textbf{试题：}卫星的线性转发器是指：
\begin{enumerate}[leftmargin=3em]
\item 卫星上将接收到的一定带宽内的所有上行信号在频谱上平移（或平移加反转）后加以转发的中继设备
\item 接收到的上行信号经线性放大后解调处理，再重新调制并转发的卫星中继设备
\item 接收到的上行信号解调处理后经线性放大，再重新调制并转发的卫星中继设备
\item 接收到的上行信号解调处理和重新调制后，经线性放大再转发的卫星中继设备
\end{enumerate}
\noindent\textbf{解说：}\textbf{}。\\\noindent\textbf{答案：}

\bigskip




\noindent\textbf{试题：}卫星线性转发器的特点是：
\begin{enumerate}[leftmargin=3em]
\item 转发工作频带范围内的所有信号而不改变原有的调制特征，可为各种常用业余无线电通信方式进行中继
\item 转发器采用线性放大电路，转发的信号失真小
\item 先对接收信号解调，经整形、线性放大后重新调制转发，提高下行信号高信噪比
\item 转发信号的幅度与上行信号幅度成线性正比，只能为CW、SSB和AM等带有振幅调制特征的通信方式进行中继
\end{enumerate}
\noindent\textbf{解说：}\textbf{}。\\\noindent\textbf{答案：}

\bigskip




\noindent\textbf{试题：}有时接收到的卫星下行信号强度会出现周期性的快衰落，这是因为：
\begin{enumerate}[leftmargin=3em]
\item 卫星翻滚造成收发信天线极化的失配
\item 电波受电离层的影响
\item 电波受高空气流的影响
\item 卫星电源不够稳定
\end{enumerate}
\noindent\textbf{解说：}\textbf{}。\\\noindent\textbf{答案：}

\bigskip




\noindent\textbf{试题：}在安装卫星轨道预测软件后，在预测卫星位置之前必须进行的准备工作是：
\begin{enumerate}[leftmargin=3em]
\item 下载待预测卫星的开普勒轨道根数（轨道要素）
\item 输入当地的磁偏角和磁倾角
\item 下载最近的太阳黑子平均数
\item 下载最近的F107射电辐射通量
\end{enumerate}
\noindent\textbf{解说：}\textbf{}。\\\noindent\textbf{答案：}

\bigskip




\noindent\textbf{试题：}业余低轨卫星的转发器覆盖范围有限。利用这类卫星进行全球性业余无线电通信的解决方法是：
\begin{enumerate}[leftmargin=3em]
\item 由卫星对上行数据进行存贮和转发
\item 由地面站进行地面中继
\item 地面业余电台换用更大功率的发射机
\item 增加地面业余电台的天线高度
\end{enumerate}
\noindent\textbf{解说：}\textbf{}。\\\noindent\textbf{答案：}

\bigskip




\noindent\textbf{试题：}业余卫星通信中所说的A、B、K、JA、JD、S、T等工作方式，表示的是：
\begin{enumerate}[leftmargin=3em]
\item 业余无线电转发器的上下行通信所用的业余频段以及所支持的调制方式
\item 发射业余卫星的国家的呼号前缀
\item 业余无线电转发器的发射功率等级
\item 业余卫星的某些轨道参数
\end{enumerate}
\noindent\textbf{解说：}\textbf{}。\\\noindent\textbf{答案：}

\bigskip




\noindent\textbf{试题：}申请WAS奖状时，不能用以下哪个呼号的卡片作为通联证明卡片：
\begin{enumerate}[leftmargin=3em]
\item VA3EJ
\item AC6HY
\item N8DX
\item KC1XX
\end{enumerate}
\noindent\textbf{解说：}\textbf{}。\\\noindent\textbf{答案：}

\bigskip




\noindent\textbf{试题：}申请5BWAZ奖状，至少要提交多少个不同分区的卡片：
\begin{enumerate}[leftmargin=3em]
\item 198
\item 199
\item 200
\item 40
\end{enumerate}
\noindent\textbf{解说：}\textbf{}。\\\noindent\textbf{答案：}

\bigskip




\noindent\textbf{试题：}在CQWW系列比赛中，哪些组别不能使用CW SKIMMER作为辅助工具使用：
\begin{enumerate}[leftmargin=3em]
\item SOAB
\item SOAB-A
\item SOSB-20
\item MS
\end{enumerate}
\noindent\textbf{解说：}\textbf{}。\\\noindent\textbf{答案：}

\bigskip




\noindent\textbf{试题：}卫星和6米的WAZ奖状需要申请者提交至少多少个区的QSL卡片：
\begin{enumerate}[leftmargin=3em]
\item 25
\item 30
\item 35
\item 40
\end{enumerate}
\noindent\textbf{解说：}\textbf{}。\\\noindent\textbf{答案：}

\bigskip




\noindent\textbf{试题：}160m的WAZ奖状需要申请者提交至少多少个区的QSL卡片：
\begin{enumerate}[leftmargin=3em]
\item 25
\item 30
\item 35
\item 40
\end{enumerate}
\noindent\textbf{解说：}\textbf{}。\\\noindent\textbf{答案：}

\bigskip




\noindent\textbf{试题：}BG1DOB在IARU HF World Championship竞赛中通联到某个电台，信号报告为599JARL，那么这个电台的呼号有可能是：
\begin{enumerate}[leftmargin=3em]
\item 8N1HQ
\item BA4RF
\item B4HQ
\item DA1HQ
\end{enumerate}
\noindent\textbf{解说：}\textbf{}。\\\noindent\textbf{答案：}

\bigskip




\noindent\textbf{试题：}北京20岁女大学生参加IARU HF World Championship单人组比赛，她发给别人的信号报告有可能是：
\begin{enumerate}[leftmargin=3em]
\item 5944
\item 5900
\item 5924
\item 59001
\end{enumerate}
\noindent\textbf{解说：}\textbf{}。\\\noindent\textbf{答案：}

\bigskip




\noindent\textbf{试题：}北京20岁女大学生参加ALL ASIAN SSB单人组比赛，她发给别人的信号报告有可能是：
\begin{enumerate}[leftmargin=3em]
\item 5900
\item 5944
\item 5924
\item 59001
\end{enumerate}
\noindent\textbf{解说：}\textbf{}。\\\noindent\textbf{答案：}

\bigskip




\noindent\textbf{试题：}北京20岁女大学生参加CQ WW CW竞赛，她发给别人的信号报告有可能是：
\begin{enumerate}[leftmargin=3em]
\item 59924
\item 59900
\item 59944
\item 599001
\end{enumerate}
\noindent\textbf{解说：}\textbf{}。\\\noindent\textbf{答案：}

\bigskip




\noindent\textbf{试题：}以北京为中心与DXCC实体3X、HZ、G、OE之间的大圆距离由远及近的排序为：
\begin{enumerate}[leftmargin=3em]
\item 3X、G、OE、HZ
\item 3X、HZ、OE、G
\item G、OE、3X、HZ
\item HZ、OE、3X、G
\end{enumerate}
\noindent\textbf{解说：}\textbf{}。\\\noindent\textbf{答案：}

\bigskip




\noindent\textbf{试题：}某业余电台在参加比赛，对方报来的交换信息是599CN，这个业余电台参加的比赛是：
\begin{enumerate}[leftmargin=3em]
\item IARU ARRL 160M CONTEST
\item CQWW CW
\item ALL ASIAN DX CONTEST
\item CQWW WPX
\end{enumerate}
\noindent\textbf{解说：}\textbf{}。\\\noindent\textbf{答案：}

\bigskip




\noindent\textbf{试题：}按照竞赛惯例，以下频率中不能用于IOTA竞赛的有？
\begin{enumerate}[leftmargin=3em]
\item 3560-3600 kHz
\item 7040-7060 kHz
\item 14060-14125 kHz
\item 14300-14350 kHz
\end{enumerate}
\noindent\textbf{解说：}\textbf{}。\\\noindent\textbf{答案：}

\bigskip




\noindent\textbf{试题：}下述通信不属于电信（telecommunication）范畴：
\begin{enumerate}[leftmargin=3em]
\item 邮政通信
\item 有线通信
\item 无线电通信
\item 光通信
\end{enumerate}
\noindent\textbf{解说：}\textbf{}。\\\noindent\textbf{答案：}

\bigskip




\noindent\textbf{试题：}关于地面无线电通信（terrestrial radio communication）的正确说法是：
\begin{enumerate}[leftmargin=3em]
\item 航空通信属于地面无线电通信
\item 水上通信不属于地面无线电通信
\item 地面业余电台之间所进行的任何无线电通信都属于地面无线电通信
\item 依靠电离层反射的天波所进行的通信不属于地面无线电通信
\end{enumerate}
\noindent\textbf{解说：}\textbf{}。\\\noindent\textbf{答案：}

\bigskip




\noindent\textbf{试题：}下列情况会产生减幅波辐射：
\begin{enumerate}[leftmargin=3em]
\item 电路接触点打火
\item 对讲机按键发射
\item 电视机本振泄漏
\item 医用高频加热器泄漏
\end{enumerate}
\noindent\textbf{解说：}\textbf{}。\\\noindent\textbf{答案：}

\bigskip




\noindent\textbf{试题：}“电磁感应现象”是指：
\begin{enumerate}[leftmargin=3em]
\item 变化的磁场可以产生感应电流
\item 足够高的电压可以产生感应磁场
\item 足够大的静止磁场可以产生感应电流
\item 摩擦产生的静电可以感应出直流磁场
\end{enumerate}
\noindent\textbf{解说：}\textbf{}。\\\noindent\textbf{答案：}

\bigskip




\noindent\textbf{试题：}无线电技术中常遇到的术语“互感现象”是指：
\begin{enumerate}[leftmargin=3em]
\item 两个回路电流变化，通过磁场耦合而相互产生感应电动势的现象
\item 两个元器件连接在一起时各自的阻抗对流过对方的电流产生的相互影响
\item 两个部件的电压形成的电场互相叠加的现象
\item 两种不同频率的信号在空间叠加造成的干扰现象
\end{enumerate}
\noindent\textbf{解说：}\textbf{}。\\\noindent\textbf{答案：}

\bigskip




\noindent\textbf{试题：}高档无线电通信设备的说明书经常会出现缩写TCXO或OCXO。它们是分别指：
\begin{enumerate}[leftmargin=3em]
\item 温度补偿式晶体振荡器、恒温控制式晶体振荡器
\item 中频数字信号处理部件、音频数字信号处理部件
\item 超低温前置射频放大器、超线性射频功率放大器
\item 通信模式自动识别功能、干扰信号自动抵消功能
\end{enumerate}
\noindent\textbf{解说：}\textbf{}。\\\noindent\textbf{答案：}

\bigskip




\noindent\textbf{试题：}无线电原理经常用到缩写VFO，它代表：
\begin{enumerate}[leftmargin=3em]
\item 可变频率振荡器
\item 压控振荡器
\item 石英晶体元件
\item 可变频率石英振荡器
\end{enumerate}
\noindent\textbf{解说：}\textbf{}。\\\noindent\textbf{答案：}

\bigskip




\noindent\textbf{试题：}无线电原理经常用到缩写VXO，它代表：
\begin{enumerate}[leftmargin=3em]
\item 可变石英振荡器
\item 可变频率振荡器
\item 石英晶体元件
\item 压控振荡器
\end{enumerate}
\noindent\textbf{解说：}\textbf{}。\\\noindent\textbf{答案：}

\bigskip




\noindent\textbf{试题：}无线电原理经常用到缩写VCO，它代表：
\begin{enumerate}[leftmargin=3em]
\item 压控振荡器
\item 可变石英振荡器
\item 可变频率振荡器
\item 石英晶体元件
\end{enumerate}
\noindent\textbf{解说：}\textbf{}。\\\noindent\textbf{答案：}

\bigskip




\noindent\textbf{试题：}无线电原理经常用到缩写XTAL，它代表：
\begin{enumerate}[leftmargin=3em]
\item 石英晶体元件
\item 可变频率振荡器
\item 晶体振荡器
\item 压控振荡器
\end{enumerate}
\noindent\textbf{解说：}\textbf{}。\\\noindent\textbf{答案：}

\bigskip




\noindent\textbf{试题：}能够确定直流电路中任何一个两端元件工作状况的基本参数包括：
\begin{enumerate}[leftmargin=3em]
\item 通过电流、两端电压、电阻
\item 频率、电荷量、电场强度
\item 波长、电压、电容量
\item 驻波比、功率、阻抗
\end{enumerate}
\noindent\textbf{解说：}\textbf{}。\\\noindent\textbf{答案：}

\bigskip




\noindent\textbf{试题：}物理量“电流”描述的是：
\begin{enumerate}[leftmargin=3em]
\item 单位时间内流过电路的电荷数量
\item 流过电路的累计电荷数量
\item 电源所能供出的最大的电荷数量
\item 电荷在电路导体内的运动速度
\end{enumerate}
\noindent\textbf{解说：}\textbf{}。\\\noindent\textbf{答案：}

\bigskip




\noindent\textbf{试题：}物理量“电压”描述的是：
\begin{enumerate}[leftmargin=3em]
\item 电源加在电路两端驱动电子流动的力量大小
\item 电源把其它形式的能量转化为电能的能力
\item 单位时间内流过电路的电荷数量
\item 电源所能供出的最大的电荷数量
\end{enumerate}
\noindent\textbf{解说：}\textbf{}。\\\noindent\textbf{答案：}

\bigskip




\noindent\textbf{试题：}物理量“电动势”描述的是：
\begin{enumerate}[leftmargin=3em]
\item 电源把其它形式的能量转化为电能的能力
\item 电源加在电路两端驱动电子流动的力量大小
\item 单位时间内流过电路的电荷数量
\item 电源所能供出的最大的电荷数量
\end{enumerate}
\noindent\textbf{解说：}\textbf{}。\\\noindent\textbf{答案：}

\bigskip




\noindent\textbf{试题：}物理量“电阻”描述的是：
\begin{enumerate}[leftmargin=3em]
\item 电路对电流的阻碍力大小
\item 电路阻力所消耗的能量多少
\item 电路阻断电流所需要的过度时间
\item 电流克服电路阻力的能力大小
\end{enumerate}
\noindent\textbf{解说：}\textbf{}。\\\noindent\textbf{答案：}

\bigskip




\noindent\textbf{试题：}物理量“功率”描述的是：
\begin{enumerate}[leftmargin=3em]
\item 单位时间里电所能够做的功
\item 电源总共能够做的功
\item 负载总共消耗的能量
\item 电源所能供出的最大的电荷数量
\end{enumerate}
\noindent\textbf{解说：}\textbf{}。\\\noindent\textbf{答案：}

\bigskip




\noindent\textbf{试题：}电动势的单位是：
\begin{enumerate}[leftmargin=3em]
\item 伏（特）
\item 安（培）
\item 瓦（特）
\item 欧（姆）
\end{enumerate}
\noindent\textbf{解说：}\textbf{}。\\\noindent\textbf{答案：}

\bigskip




\noindent\textbf{试题：}电源两端电动势的方向为：
\begin{enumerate}[leftmargin=3em]
\item 从电源的负极到正极
\item 从电源的正极到负极
\item 取决于负载电阻和电源内阻的相对大小
\item 与电源的电动势方向相同
\end{enumerate}
\noindent\textbf{解说：}\textbf{}。\\\noindent\textbf{答案：}

\bigskip




\noindent\textbf{试题：}将一个电阻为R的负载接到电压为U的电源上。负载中的电流Ｉ和负载消耗的功率Ｐ分别为：（U、I、R、P的单位分别为伏特、安培、欧姆、瓦特，”x＾m”表示“x的m次方”）
\begin{enumerate}[leftmargin=3em]
\item I = U/R，P = U＾2/R
\item I = U×R，P =U/R
\item I = R/U，P =U×R
\item I = U×R，P =U＾2×R
\end{enumerate}
\noindent\textbf{解说：}\textbf{}。\\\noindent\textbf{答案：}

\bigskip




\noindent\textbf{试题：}一个电阻为R的负载中流过的电流为Ｉ。其两端的电压U所消耗的功率Ｐ分别为：（U、I、R、P的单位分别为伏特、安培、欧姆、瓦特，”x＾m”表示“x的m次方”）
\begin{enumerate}[leftmargin=3em]
\item U = I×R，P = I＾2×R
\item U = I / R，P = I＾2/R
\item U = R / I，P = R / I＾2
\item U = I + R，P =I ×R
\end{enumerate}
\noindent\textbf{解说：}\textbf{}。\\\noindent\textbf{答案：}

\bigskip




\noindent\textbf{试题：}一个电阻负载两端电压为U，流过的电流为I。它的电阻R和所消耗的功率P分别为：（U、I、R、P的单位分别为伏特、安培、欧姆、瓦特，”x＾m”表示“x的m次方”）
\begin{enumerate}[leftmargin=3em]
\item R = U/I，P = U×I
\item R = I / U，P = U＾2/I
\item R = U×I，P = U/ I＾2
\item R = U＾2/I，P = U＾2×I
\end{enumerate}
\noindent\textbf{解说：}\textbf{}。\\\noindent\textbf{答案：}

\bigskip




\noindent\textbf{试题：}一个电阻负载两端电压为U，所消耗的功率为P。流过负载的电流I和负载的电阻R分别为：（U、I、R、P的单位分别为伏特、安培、欧姆、瓦特，”x＾m”表示“x的m次方”）
\begin{enumerate}[leftmargin=3em]
\item I = P/U，R = U＾2/P
\item I = P×U，R = P/U
\item I = U/P，R = P/U
\item I = P/ U＾2，R = P×U
\end{enumerate}
\noindent\textbf{解说：}\textbf{}。\\\noindent\textbf{答案：}

\bigskip




\noindent\textbf{试题：}有阻值分别为R1和R2的两个负载，其中R1的电阻值是R2的N倍，把它们并联后接到电源上，则：（”x＾m”表示“x的m次方”）
\begin{enumerate}[leftmargin=3em]
\item 流过R1的电流是R2的1/N，R1消耗的功率是R2的1/N
\item 流过R1的电流是R2的N倍，R1消耗的功率是R2的N＾2倍
\item 流过R1的电流与R2的相同，R1消耗的功率是R2的1/N＾2
\item 流过R1的电流与R2的相同，R1消耗的功率是R2的N倍
\end{enumerate}
\noindent\textbf{解说：}\textbf{}。\\\noindent\textbf{答案：}

\bigskip




\noindent\textbf{试题：}有阻值分别为R1和R2的两个负载，其中R1的电阻值是R2的N倍，把它们并联后接到电源上，则：（”x＾m”表示“x的m次方”）
\begin{enumerate}[leftmargin=3em]
\item R1两端的电压与R2的相同，R1消耗的功率是R2的1/N
\item R1两端的电压与R2的相同，R1消耗的功率是R2的N＾2倍
\item R1两端的电压是R2的1/N，R1消耗的功率是R2的1/N＾2
\item R1两端的电压是R2的N倍，R1消耗的功率是R2的N＾2倍
\end{enumerate}
\noindent\textbf{解说：}\textbf{}。\\\noindent\textbf{答案：}

\bigskip




\noindent\textbf{试题：}有阻值分别为R1和R2的两个负载，其中R1的电阻值是R2的N倍，把它们串联后接到电源上，则：（”x＾m”表示“x的m次方”）
\begin{enumerate}[leftmargin=3em]
\item 流过R1的电流与R2的相同，R1消耗的功率是R2的N倍
\item 流过R1的电流与R2的相同，R1消耗的功率是R2的1/N
\item 流过R1的电流是R2的1/N，R1消耗的功率是R2的1/N＾2
\item 流过R1的电流是R2的N倍，R1消耗的功率是R2的N＾2倍
\end{enumerate}
\noindent\textbf{解说：}\textbf{}。\\\noindent\textbf{答案：}

\bigskip




\noindent\textbf{试题：}有阻值分别为R1和R2的两个负载，其中R1的电阻值是R2的N倍，把它们串联后接到电源上，则：（”x＾m”表示“x的m次方”）
\begin{enumerate}[leftmargin=3em]
\item R1两端的电压是R2的N倍，R1消耗的功率是R2的N倍
\item R1两端的电压是R2的1/N，R1消耗的功率是R2的1/N＾2
\item R1两端的电压与R2的相同，R1消耗的功率是R2的1/N
\item R1两端的电压与R2的相同，R1消耗的功率是R2的N＾2倍
\end{enumerate}
\noindent\textbf{解说：}\textbf{}。\\\noindent\textbf{答案：}

\bigskip




\noindent\textbf{试题：}已知A、B两个设备的工作电压相同，A的耗电功率是B的N倍。则：（”x＾m”表示“x的m次方”）
\begin{enumerate}[leftmargin=3em]
\item A的工作电流是B的N倍
\item A的工作电流是B的N＾(1/2)倍
\item A的工作电流是B的N＾2倍
\item A的工作电流是B的1/N倍
\end{enumerate}
\noindent\textbf{解说：}\textbf{}。\\\noindent\textbf{答案：}

\bigskip




\noindent\textbf{试题：}已知A、B两个设备的工作电压相同，A的额定电流是B的N倍。则：（”x＾m”表示“x的m次方”）
\begin{enumerate}[leftmargin=3em]
\item A的耗电功率是B的N倍
\item A的耗电功率是B的N＾(1/2)倍
\item A的耗电功率是B的N＾2倍
\item A的耗电功率是B的1/N倍
\end{enumerate}
\noindent\textbf{解说：}\textbf{}。\\\noindent\textbf{答案：}

\bigskip




\noindent\textbf{试题：}将N个相同的电阻负载串联后接到电源上。与每个负载单独接到电源相比：（”x＾m”表示“x的m次方”）
\begin{enumerate}[leftmargin=3em]
\item 串联后流过每个电阻的电流减少到1/N，每个电阻的耗电功率减少到1/N＾2
\item 串联后流过每个电阻的电流减少到1/N，每个电阻的耗电功率减少到1/N
\item 串联后流过每个电阻的电流不变，每个电阻的耗电功率减少到1/N
\item 串联后流过每个电阻的电流增加到N倍，每个电阻的耗电功率增加到N＾2倍
\end{enumerate}
\noindent\textbf{解说：}\textbf{}。\\\noindent\textbf{答案：}

\bigskip




\noindent\textbf{试题：}将N个相同的电阻负载串联后接到电源上。与每个负载单独接到电源相比：（”x＾m”表示“x的m次方”）
\begin{enumerate}[leftmargin=3em]
\item 串联后每个电阻两端的电压减少到1/N，每个电阻的耗电功率减少到1/ N＾2
\item 串联后每个电阻两端的电压减少到1/N，每个电阻的耗电功率减少到1/N
\item 串联后每个电阻两端的电压不变，每个电阻的耗电功率减少到1/N
\item 串联后每个电阻两端的电压增加到N倍，每个电阻的耗电功率增加到N＾2倍
\end{enumerate}
\noindent\textbf{解说：}\textbf{}。\\\noindent\textbf{答案：}

\bigskip




\noindent\textbf{试题：}将N个相同的电阻负载并联后接到电源上。与每个负载单独接到电源相比：（”x＾m”表示“x的m次方”）
\begin{enumerate}[leftmargin=3em]
\item 并联后流过每个电阻的电流不变，所有电阻的总耗电功率增加到N倍
\item 并联后流过每个电阻的电流不变，所有电阻的总耗电功率增加到N＾2倍
\item 并联后流过每个电阻的电流增加到N倍，两个电阻的总耗电功率增加到N＾2倍
\item 并联后流过每个电阻的电流减少到1/N，两个电阻的总耗电功率减少到1/N＾2
\end{enumerate}
\noindent\textbf{解说：}\textbf{}。\\\noindent\textbf{答案：}

\bigskip




\noindent\textbf{试题：}将N个相同的电阻负载并联后接到电源上。与每个负载单独接到电源相比：（”x＾m”表示“x的m次方”）
\begin{enumerate}[leftmargin=3em]
\item 并联后每个电阻两端的电压不变，所有电阻的总耗电功率增加到N倍
\item 并联后每个电阻两端的电压不变，所有电阻的总耗电功率增加到N＾2倍
\item 并联后每个电阻两端的电压增加到N倍，两个电阻的总耗电功率增加到N＾2倍
\item 并联后每个电阻两端的电压减少到1/N，两个电阻的总耗电功率减少到1/N＾2
\end{enumerate}
\noindent\textbf{解说：}\textbf{}。\\\noindent\textbf{答案：}

\bigskip




\noindent\textbf{试题：}一个电阻负载，如果将其两端的工作电压提高百分之N，则：（”x＾m”表示“x的m次方”）
\begin{enumerate}[leftmargin=3em]
\item 耗电量增加到原来的[百分之（100+N）]＾2
\item 耗电量增加到原来的百分之（100+N）
\item 耗电量比原来的增加百分之N
\item 耗电量比原来的增加[百分之N]＾2
\end{enumerate}
\noindent\textbf{解说：}\textbf{}。\\\noindent\textbf{答案：}

\bigskip




\noindent\textbf{试题：}一个电阻负载，如果将其两端的工作电压降低百分之N，则：（”x＾m”表示“x的m次方”）
\begin{enumerate}[leftmargin=3em]
\item 耗电量减少到原来的[百分之（100-N）]＾2
\item 耗电量减少到原来的百分之（100-N）
\item 耗电量比原来的减少百分之N
\item 耗电量比原来的减少[百分之N]＾2
\end{enumerate}
\noindent\textbf{解说：}\textbf{}。\\\noindent\textbf{答案：}

\bigskip




\noindent\textbf{试题：}分别用电压为120V的蓄电池组和电压最大值为120V的交流变压器驱动同样的电阻负载，哪一个电阻每分钟发出的热量多？
\begin{enumerate}[leftmargin=3em]
\item 蓄电池驱动的电阻所发的热量是交流变压器上的电阻的2倍左右
\item 蓄电池驱动的电阻所发的热量是交流变压器上的电阻的0.7倍左右
\item 蓄电池驱动的电阻所发的热量是交流变压器上的电阻的1.4倍左右
\item 两个电源驱动的电阻发热相同
\end{enumerate}
\noindent\textbf{解说：}\textbf{}。\\\noindent\textbf{答案：}

\bigskip




\noindent\textbf{试题：}分别用电压为120V的蓄电池组和电压有效值为120V的交流变压器驱动同样的电阻负载，哪一个电阻每分钟发出的热量多？
\begin{enumerate}[leftmargin=3em]
\item 两个电源驱动的电阻发热相同
\item 蓄电池驱动的电阻所发的热量是交流变压器上的电阻的1.4倍左右
\item 蓄电池驱动的电阻所发的热量是交流变压器上的电阻的0.7倍左右
\item 蓄电池驱动的电阻所发的热量是交流变压器上的电阻的2倍左右
\end{enumerate}
\noindent\textbf{解说：}\textbf{}。\\\noindent\textbf{答案：}

\bigskip




\noindent\textbf{试题：}分别用电压为120V的蓄电池组和电压有效值为120V的交流变压器串联二极管后驱动同样的电阻负载，哪一个电阻每分钟发出的热量多？（忽略二极管的电压降）
\begin{enumerate}[leftmargin=3em]
\item 蓄电池驱动的电阻所发的热量是交流变压器上的电阻的2倍左右
\item 蓄电池驱动的电阻所发的热量是交流变压器上的电阻的1.4倍左右
\item 蓄电池驱动的电阻所发的热量是交流变压器上的电阻的0.7倍左右
\item 两个电源驱动的电阻发热相同
\end{enumerate}
\noindent\textbf{解说：}\textbf{}。\\\noindent\textbf{答案：}

\bigskip




\noindent\textbf{试题：}分别用电压为120V的蓄电池组和电压最大值为120V的交流变压器经过带电容滤波的全波整流电路驱动同样的电阻负载，哪一个电阻每分钟发出的热量多？（忽略整流器的电压降）
\begin{enumerate}[leftmargin=3em]
\item 两个电源驱动的电阻发热相同
\item 蓄电池驱动的电阻所发的热量是交流变压器上的电阻的2倍左右
\item 蓄电池驱动的电阻所发的热量是交流变压器上的电阻的1.4倍左右
\item 蓄电池驱动的电阻所发的热量是交流变压器上的电阻的0.7倍左右
\end{enumerate}
\noindent\textbf{解说：}\textbf{}。\\\noindent\textbf{答案：}

\bigskip




\noindent\textbf{试题：}分别用电压有效值为120V、频率为50Hz的交流电源和电压有效值为120V、频率为10kHz的方波电源驱动同样的电阻负载，哪一个电阻每分钟发出的热量多？
\begin{enumerate}[leftmargin=3em]
\item 两个电源驱动的电阻发热相同
\item 10kHz电路电阻所发的热量是50Hz电路电阻的5倍左右
\item 10kHz电路电阻所发的热量是50Hz电路电阻的1/5倍左右
\item 10kHz电路电阻所发的热量是50Hz电路电阻的200倍左右
\end{enumerate}
\noindent\textbf{解说：}\textbf{}。\\\noindent\textbf{答案：}

\bigskip




\noindent\textbf{试题：}正弦交流电压或电流的峰值（peak value）是指：（”x＾m”表示“x的m次方”）
\begin{enumerate}[leftmargin=3em]
\item 从零点算起的最大值
\item 一个周期内瞬时值的平均乘以值2＾(1/2)
\item 负半周最大幅度与正半周最大幅度的平均值
\item 负半周最大幅度与正半周最大幅度的差值的二次方
\end{enumerate}
\noindent\textbf{解说：}\textbf{}。\\\noindent\textbf{答案：}

\bigskip




\noindent\textbf{试题：}正弦交流电压或电流的峰值峰-峰值（peak to peak）是指：
\begin{enumerate}[leftmargin=3em]
\item 从负半周峰值到正半周峰值之间的差
\item 从零点算起的最大值
\item 负半周最大幅度与正半周最大幅度的差值的二次方
\item 负半周最大幅度与正半周最大幅度的差值的平方根
\end{enumerate}
\noindent\textbf{解说：}\textbf{}。\\\noindent\textbf{答案：}

\bigskip




\noindent\textbf{试题：}任意交流电压的有效值是指：（”x＾m”表示“x的m次方”）
\begin{enumerate}[leftmargin=3em]
\item 在同一电阻上可以转换出与该交流电压效果相同的热量的直流电压
\item 最终转换成应用所需的有用能量的那部分电压值
\item 电压的平均值乘以2＾(1/2)
\item 电压的峰值除以2＾(1/2)
\end{enumerate}
\noindent\textbf{解说：}\textbf{}。\\\noindent\textbf{答案：}

\bigskip




\noindent\textbf{试题：}最小值为-50伏、峰-峰值为100伏的方波电压，其有效值电压为：
\begin{enumerate}[leftmargin=3em]
\item 50伏
\item 约70.7伏
\item 约35.4伏
\item 100伏
\end{enumerate}
\noindent\textbf{解说：}\textbf{}。\\\noindent\textbf{答案：}

\bigskip




\noindent\textbf{试题：}最小值为0、峰-峰值为100伏的方波电压，其有效值电压为：
\begin{enumerate}[leftmargin=3em]
\item 50伏
\item 约70.7伏
\item 约35.4伏
\item 100伏
\end{enumerate}
\noindent\textbf{解说：}\textbf{}。\\\noindent\textbf{答案：}

\bigskip




\noindent\textbf{试题：}最小值为-50伏、峰-峰值为100伏的三角波电压，其有效值电压为：
\begin{enumerate}[leftmargin=3em]
\item 25伏
\item 50伏
\item 约70.7伏
\item 100伏
\end{enumerate}
\noindent\textbf{解说：}\textbf{}。\\\noindent\textbf{答案：}

\bigskip




\noindent\textbf{试题：}最小值为0、峰-峰值为100伏的三角波电压，其有效值电压为：
\begin{enumerate}[leftmargin=3em]
\item 50伏
\item 约70.7伏
\item 25伏
\item 100伏
\end{enumerate}
\noindent\textbf{解说：}\textbf{}。\\\noindent\textbf{答案：}

\bigskip




\noindent\textbf{试题：}最小值为-50伏、峰-峰值为100伏的方波电压，其平均值电压为：
\begin{enumerate}[leftmargin=3em]
\item 0
\item 50伏
\item 约70.7伏
\item 100伏
\end{enumerate}
\noindent\textbf{解说：}\textbf{}。\\\noindent\textbf{答案：}

\bigskip




\noindent\textbf{试题：}最小值为0、峰-峰值为100伏的方波电压，其平均值电压为：
\begin{enumerate}[leftmargin=3em]
\item 50伏
\item 约70.7伏
\item 约35.4伏
\item 100伏
\end{enumerate}
\noindent\textbf{解说：}\textbf{}。\\\noindent\textbf{答案：}

\bigskip




\noindent\textbf{试题：}最小值为-50伏、峰-峰值为100伏的三角波电压，其平均值电压为：
\begin{enumerate}[leftmargin=3em]
\item 0
\item 25伏
\item 50伏
\item 约70.7伏
\end{enumerate}
\noindent\textbf{解说：}\textbf{}。\\\noindent\textbf{答案：}

\bigskip




\noindent\textbf{试题：}最小值为0、峰-峰值为100伏的三角波电压，其平均值电压为：
\begin{enumerate}[leftmargin=3em]
\item 50伏
\item 约70.7伏
\item 25伏
\item 100伏
\end{enumerate}
\noindent\textbf{解说：}\textbf{}。\\\noindent\textbf{答案：}

\bigskip




\noindent\textbf{试题：}把两个幅度相等、相位相差360°的正弦电压信号源相串联，得到的是：
\begin{enumerate}[leftmargin=3em]
\item 幅度为单个信号源的2倍、相位与原信号源相同的正弦电压
\item 电压为0
\item 幅度与单个信号源的相同、相位与原信号源相差180°的正弦电压
\item 幅度与单个信号源的相同、频率比原信号高一倍的正弦电压
\end{enumerate}
\noindent\textbf{解说：}\textbf{}。\\\noindent\textbf{答案：}

\bigskip




\noindent\textbf{试题：}把两个幅度相等、相位相差180°的正弦电压信号源相串联，得到的是：
\begin{enumerate}[leftmargin=3em]
\item 电压为0
\item 幅度为单个信号源的2倍、相位与原信号源相同的正弦电压
\item 幅度与单个信号源的相同、相位与原信号源相差90°的正弦电压
\item 幅度与单个信号源的相同、频率比原信号高一倍的正弦电压
\end{enumerate}
\noindent\textbf{解说：}\textbf{}。\\\noindent\textbf{答案：}

\bigskip




\noindent\textbf{试题：}把两个幅度相等、相位相差90°的正弦电压信号源相串联，得到的是：
\begin{enumerate}[leftmargin=3em]
\item 幅度为单个信号源的1.41倍、相位与原信号源相差45°的正弦电压
\item 幅度与单个信号源的相同、相位与原信号源相差45°的正弦电压
\item 幅度为单个信号源的2倍、相位与原信号源相同的正弦电压
\item 幅度与单个信号源的相同、频率比原信号高一倍的正弦电压
\end{enumerate}
\noindent\textbf{解说：}\textbf{}。\\\noindent\textbf{答案：}

\bigskip




\noindent\textbf{试题：}二极管在无线电技术里广泛用于整流或检波，主要是利用了它的下列性能：
\begin{enumerate}[leftmargin=3em]
\item 单向导电
\item 电流放大
\item 高频滤波
\item 过流保护
\end{enumerate}
\noindent\textbf{解说：}\textbf{}。\\\noindent\textbf{答案：}

\bigskip




\noindent\textbf{试题：}三极电子管的三个基本电极是：
\begin{enumerate}[leftmargin=3em]
\item 阴极（灯丝）、栅极、屏极（板极）
\item 发射极、基极、集电极
\item 源极、栅极、漏极
\item 正极、负极、接地极
\end{enumerate}
\noindent\textbf{解说：}\textbf{}。\\\noindent\textbf{答案：}

\bigskip




\noindent\textbf{试题：}双极型半导体三极管的三个基本电极是：
\begin{enumerate}[leftmargin=3em]
\item 发射极、基极、集电极
\item 阴极（灯丝）、栅极、屏极（板极）
\item 源极、栅极、漏极
\item 正极、负极、接地极
\end{enumerate}
\noindent\textbf{解说：}\textbf{}。\\\noindent\textbf{答案：}

\bigskip




\noindent\textbf{试题：}场效应型半导体三极管的三个基本电极是：
\begin{enumerate}[leftmargin=3em]
\item 源极、栅极、漏极
\item 阴极（灯丝）、栅极、屏极（板极）
\item 发射极、基极、集电极
\item 正极、负极、接地极
\end{enumerate}
\noindent\textbf{解说：}\textbf{}。\\\noindent\textbf{答案：}

\bigskip




\noindent\textbf{试题：}现代无线电电路大量应用表面装焊元器件（Surface Mounted Devices）。其规格参数0402、0603、0805、1206等是指：
\begin{enumerate}[leftmargin=3em]
\item 元器件的长和宽（以单位：毫米或0.1英寸）
\item 容量值（单位：微法）
\item 电阻值（单位：千欧）
\item 电感值（单位：微亨）
\end{enumerate}
\noindent\textbf{解说：}\textbf{}。\\\noindent\textbf{答案：}

\bigskip




\noindent\textbf{试题：}业余无线电收发信机所用的半导体器件经常会在型号或者使用说明中用一组字母加以描述，例如 SIP、DIP、SOT、SOP、SSOP、QFP等，它们用来说明：
\begin{enumerate}[leftmargin=3em]
\item 器件的封装形式
\item 器件的工作温度范围
\item 器件的输入输出电平范围
\item 器件的工作频率范围
\end{enumerate}
\noindent\textbf{解说：}\textbf{}。\\\noindent\textbf{答案：}

\bigskip




\noindent\textbf{试题：}有些无线电收发信机在外接13.8V直流电源插座后正向串联一个大电流二极管作为保护。所利用的二极管的特性和动作原理为：
\begin{enumerate}[leftmargin=3em]
\item 二极管的单向导电特性；在接反电源极性时自动阻断电源
\item 大电流二极管的散热性能；帮助电路发散热量，防止设备过热
\item 二极管导电区的非线性；在电源过压时加大压降，起到稳压作用
\item 二极管导电区的非线性；在设备电流超限时加大压降，起到限流作用
\end{enumerate}
\noindent\textbf{解说：}\textbf{}。\\\noindent\textbf{答案：}

\bigskip




\noindent\textbf{试题：}有些无线电收发信机在外接13.8V直流电源插座后串联一个熔断器，并在电源线的设备端反向跨接一个大电流二极管作为保护。所利用的二极管的特性和动作原理为：
\begin{enumerate}[leftmargin=3em]
\item 二极管的单向导电特性；接反电源极性时造成很大短路电路，烧断熔丝，造成电源永久切断
\item 大电流二极管的散热性能；帮助电路发散热量，防止设备过热
\item 二极管导电区的非线性；在电源过压时加大压降，起到稳压作用
\item 二极管导电区的非线性；在设备电流超限时加大压降，起到限流作用
\end{enumerate}
\noindent\textbf{解说：}\textbf{}。\\\noindent\textbf{答案：}

\bigskip




\noindent\textbf{试题：}用一副臂长约10cm的小偶极天线并联一个半导体二极管和直流微安表，做一个简单的天线场强表。应选用硅二极管还是锗二极管，正确的考虑是：
\begin{enumerate}[leftmargin=3em]
\item 锗、硅二极管的起始导通电压分别为0.3和0.7V，应选用锗二极管使场强表比较灵敏
\item 所有二极管都具有同样的单向导电特性，采用锗、硅二极管效果完全相同
\item 硅二极管反向击穿电压比较高，应选用硅管使场强表延长使用寿命
\item 电路过于简单，无论使用锗、硅二极管，都无法工作
\end{enumerate}
\noindent\textbf{解说：}\textbf{}。\\\noindent\textbf{答案：}

\bigskip




\noindent\textbf{试题：}将发光二极管接到直流电源上时的注意点是：
\begin{enumerate}[leftmargin=3em]
\item 正向连接，需串联限流电阻
\item 反向连接，需串联限流电阻
\item 正向无所谓，但需并联限流电阻
\item 正向直接连接到电源上
\end{enumerate}
\noindent\textbf{解说：}\textbf{}。\\\noindent\textbf{答案：}

\bigskip




\noindent\textbf{试题：}三极电子管至少应有几个引出电极？
\begin{enumerate}[leftmargin=3em]
\item 4个
\item 2个
\item 3个
\item 5个
\end{enumerate}
\noindent\textbf{解说：}\textbf{}。\\\noindent\textbf{答案：}

\bigskip




\noindent\textbf{试题：}三极电子管的放大作用表现在：
\begin{enumerate}[leftmargin=3em]
\item 较小的栅极信号电压变化可以引起较大的屏极电流的成比例变化
\item 较小的阴极电流变化就可以引起较大的屏极电压的成比例变化
\item 较小的栅极信号电流变化可以引起较大的屏极电流的成比例变化
\item 较小的栅极信号电流变化可以引起较大的屏极电压的成比例变化
\end{enumerate}
\noindent\textbf{解说：}\textbf{}。\\\noindent\textbf{答案：}

\bigskip




\noindent\textbf{试题：}场效应半导体管的放大作用表现在：
\begin{enumerate}[leftmargin=3em]
\item 较小的栅极信号电压变化可以引起较大的漏极电流的成比例变化
\item 较小的源极信号电流变化可以引起较大的漏极电压的成比例变化
\item 较小的栅极信号电流变化可以引起较大的漏极电流的成比例变化
\item 较小的栅极信号电流变化可以引起较小的漏极电流的成比例变化
\end{enumerate}
\noindent\textbf{解说：}\textbf{}。\\\noindent\textbf{答案：}

\bigskip




\noindent\textbf{试题：}双极型半导体管的放大作用表现在：
\begin{enumerate}[leftmargin=3em]
\item 较小的基极信号电流变化可以引起较大的集电极电流的成比例变化
\item 较小的发射极信号电流变化可以引起较大的集电极电流的成比例变化
\item 较小的基极信号电压变化可以引起较大的集电极电压的成比例变化
\item 较小的基极信号电压变化可以引起较大的集电极电流的成比例变化
\end{enumerate}
\noindent\textbf{解说：}\textbf{}。\\\noindent\textbf{答案：}

\bigskip




\noindent\textbf{试题：}描述三极电子管放大能力的最常用参数、单位及符号是：
\begin{enumerate}[leftmargin=3em]
\item 互导（跨导）； 单位为毫安/伏；符号为“gm”
\item 电流放大倍数；没有单位量纲；符号为“β”
\item 互导（跨导）； 单位为“倍”；符号为“A”
\item 电压放大倍数；没有单位量纲；符号为“k”
\end{enumerate}
\noindent\textbf{解说：}\textbf{}。\\\noindent\textbf{答案：}

\bigskip




\noindent\textbf{试题：}描述场效应半导体三极管放大能力的最常用参数、单位及符号是：
\begin{enumerate}[leftmargin=3em]
\item 互导（跨导）；； 单位为毫安/伏；符号为“gm”
\item 电流放大倍数；没有单位量纲；符号为“β”
\item 互导（跨导）； 单位为“倍”；符号为“A”
\item 电压放大倍数；没有单位量纲；符号为“k”
\end{enumerate}
\noindent\textbf{解说：}\textbf{}。\\\noindent\textbf{答案：}

\bigskip




\noindent\textbf{试题：}描述双极型半导体三极管放大能力的最常用参数、单位及符号是：
\begin{enumerate}[leftmargin=3em]
\item 电流放大倍数；没有单位量纲；符号为“β”
\item 互导（跨导）； 单位为毫安/伏；符号为“gm”
\item 互导（跨导）； 单位为“倍”；符号为“A”
\item 电压放大倍数；没有单位量纲；符号为“k”
\end{enumerate}
\noindent\textbf{解说：}\textbf{}。\\\noindent\textbf{答案：}

\bigskip




\noindent\textbf{试题：}按输入阻抗由低到高排列为：
\begin{enumerate}[leftmargin=3em]
\item 双极型半导体三极管、结型场效应三极管、金属绝缘栅场效应三极管
\item 双极型半导体三极管、金属绝缘栅场效应三极管、结型场效应三极管
\item 金属绝缘栅场效应三极管、双极型半导体三极管、结型场效应三极管
\item 结型场效应三极管、双极型半导体三极管、金属绝缘栅场效应三极管
\end{enumerate}
\noindent\textbf{解说：}\textbf{}。\\\noindent\textbf{答案：}

\bigskip




\noindent\textbf{试题：}下列哪一种元器件可以用一个较小的电流来控制较大的电流？
\begin{enumerate}[leftmargin=3em]
\item 晶体管
\item 电感
\item 电阻
\item 电容
\end{enumerate}
\noindent\textbf{解说：}\textbf{}。\\\noindent\textbf{答案：}

\bigskip




\noindent\textbf{试题：}下列哪一种元器件只允许单方向的电流流动？
\begin{enumerate}[leftmargin=3em]
\item 二极管
\item 熔断器
\item 电阻
\item 稳压元件
\end{enumerate}
\noindent\textbf{解说：}\textbf{}。\\\noindent\textbf{答案：}

\bigskip




\noindent\textbf{试题：}下列哪一种元器件既可以当作电子开关又可以当作放大器使用？
\begin{enumerate}[leftmargin=3em]
\item 晶体管
\item 电位器
\item 单刀双掷开关
\item 电压表
\end{enumerate}
\noindent\textbf{解说：}\textbf{}。\\\noindent\textbf{答案：}

\bigskip




\noindent\textbf{试题：}下列哪一种元器件由三层半导体材料制作而成？
\begin{enumerate}[leftmargin=3em]
\item 双极型晶体管
\item 发电机
\item 三极电子管
\item 五栅变频管
\end{enumerate}
\noindent\textbf{解说：}\textbf{}。\\\noindent\textbf{答案：}

\bigskip




\noindent\textbf{试题：}下列哪一种元器件可以放大信号？
\begin{enumerate}[leftmargin=3em]
\item 晶体管
\item 可变电阻
\item 电解电容
\item 多芯电池
\end{enumerate}
\noindent\textbf{解说：}\textbf{}。\\\noindent\textbf{答案：}

\bigskip




\noindent\textbf{试题：}通常如何从外表识别半导体二极管的负极？
\begin{enumerate}[leftmargin=3em]
\item 靠近负极的元件表面上通常印有环线标记
\item 负极上通常标有“负极”字样
\item 负极上通常标有字母“C”
\item 负极上通常标有星号“*”
\end{enumerate}
\noindent\textbf{解说：}\textbf{}。\\\noindent\textbf{答案：}

\bigskip




\noindent\textbf{试题：}二极管的两个电极分别叫什么？
\begin{enumerate}[leftmargin=3em]
\item 正极和负极
\item 源极和漏极
\item 加极和减极
\item 门极和基极
\end{enumerate}
\noindent\textbf{解说：}\textbf{}。\\\noindent\textbf{答案：}

\bigskip




\noindent\textbf{试题：}哪一种半导体器件拥有“发射极”？
\begin{enumerate}[leftmargin=3em]
\item 双极型晶体管
\item 场效应管
\item 硅二极管
\item 桥式整流器
\end{enumerate}
\noindent\textbf{解说：}\textbf{}。\\\noindent\textbf{答案：}

\bigskip




\noindent\textbf{试题：}哪一种半导体器件拥有“栅极”？
\begin{enumerate}[leftmargin=3em]
\item 场效应管
\item 双极型晶体管
\item 硅二极管
\item 桥式整流器
\end{enumerate}
\noindent\textbf{解说：}\textbf{}。\\\noindent\textbf{答案：}

\bigskip




\noindent\textbf{试题：}下列哪一个名字用来描述晶体管放大信号的能力？
\begin{enumerate}[leftmargin=3em]
\item 增益
\item 前向电阻
\item 前向压降
\item 导通电阻
\end{enumerate}
\noindent\textbf{解说：}\textbf{}。\\\noindent\textbf{答案：}

\bigskip




\noindent\textbf{试题：}集成电路是指：
\begin{enumerate}[leftmargin=3em]
\item 将一个电路的大量元器件集合于一个单晶片上所制成的器件
\item 多极继电器
\item 把多个电阻或电容元件堆积在一起
\item 变压器
\end{enumerate}
\noindent\textbf{解说：}\textbf{}。\\\noindent\textbf{答案：}

\bigskip




\noindent\textbf{试题：}电阻的色环从靠近边缘的一个数起，分别为棕、橙、橙、黑，其阻值为：
\begin{enumerate}[leftmargin=3em]
\item 13千欧
\item 3700欧
\item 4.7兆欧
\item 250欧
\end{enumerate}
\noindent\textbf{解说：}\textbf{}。\\\noindent\textbf{答案：}

\bigskip




\noindent\textbf{试题：}电阻的色环从靠近边缘的一个数起，分别为黄、白、红、银，其阻值为：
\begin{enumerate}[leftmargin=3em]
\item 4900欧
\item 49千欧
\item 4.7兆欧
\item 390千欧
\end{enumerate}
\noindent\textbf{解说：}\textbf{}。\\\noindent\textbf{答案：}

\bigskip




\noindent\textbf{试题：}电阻的色环从靠近边缘的一个数起，分别为棕、橙、橙、黑、银，其阻值为：
\begin{enumerate}[leftmargin=3em]
\item 133欧
\item 13千欧
\item 4.7兆欧
\item 130千欧
\end{enumerate}
\noindent\textbf{解说：}\textbf{}。\\\noindent\textbf{答案：}

\bigskip




\noindent\textbf{试题：}电阻的色环分别依次为金、红、黑、灰、橙，其阻值为：
\begin{enumerate}[leftmargin=3em]
\item 38000欧
\item 208千欧
\item 28千欧
\item 2070欧
\end{enumerate}
\noindent\textbf{解说：}\textbf{}。\\\noindent\textbf{答案：}

\bigskip




\noindent\textbf{试题：}电阻的色环分别依次为金、棕、黑、白、黄，其阻值为：
\begin{enumerate}[leftmargin=3em]
\item 4900欧
\item 1090千欧
\item 109千欧
\item 190千欧
\end{enumerate}
\noindent\textbf{解说：}\textbf{}。\\\noindent\textbf{答案：}

\bigskip




\noindent\textbf{试题：}把一个电感值为100微亨的线圈跨接在电路上，不致严重妨碍原电路工作或造成安全危险的是：
\begin{enumerate}[leftmargin=3em]
\item 跨接在HF收发信机天线插座两端
\item 跨接在220伏交流电源插座两端
\item 跨接在12伏蓄电池两端
\item 跨接在输出阻抗为600欧的音响信号线两端
\end{enumerate}
\noindent\textbf{解说：}\textbf{}。\\\noindent\textbf{答案：}

\bigskip




\noindent\textbf{试题：}把一个电容量为1000微法的电容器跨接在电路上，不致严重妨碍原电路工作或造成安全危险的是：
\begin{enumerate}[leftmargin=3em]
\item 跨接在12伏蓄电池两端
\item 跨接在HF收发信机天线插座两端
\item 跨接在输出阻抗为600欧的音响信号线两端
\item 跨接在220伏交流电源插座两端
\end{enumerate}
\noindent\textbf{解说：}\textbf{}。\\\noindent\textbf{答案：}

\bigskip




\noindent\textbf{试题：}在业余收发信机的常见元器件中，标有耐压指标的是：
\begin{enumerate}[leftmargin=3em]
\item 电容
\item 电阻
\item 熔丝
\item 电感
\end{enumerate}
\noindent\textbf{解说：}\textbf{}。\\\noindent\textbf{答案：}

\bigskip




\noindent\textbf{试题：}在业余收发信机的常见元器件中，以额定耗散功率指标分类的是：
\begin{enumerate}[leftmargin=3em]
\item 电阻
\item 电容
\item 熔丝
\item 电感
\end{enumerate}
\noindent\textbf{解说：}\textbf{}。\\\noindent\textbf{答案：}

\bigskip




\noindent\textbf{试题：}在业余收发信机的常见元器件中，标有额定电流指标的是：
\begin{enumerate}[leftmargin=3em]
\item 熔丝
\item 电阻
\item 电容
\item 电感
\end{enumerate}
\noindent\textbf{解说：}\textbf{}。\\\noindent\textbf{答案：}

\bigskip




\noindent\textbf{试题：}电路图中，电解电容器所标注的电容量值的缺省单位是：
\begin{enumerate}[leftmargin=3em]
\item 法拉
\item 毫法拉
\item 皮法拉
\item 微法拉
\end{enumerate}
\noindent\textbf{解说：}\textbf{}。\\\noindent\textbf{答案：}

\bigskip




\noindent\textbf{试题：}电路图中，非电解电容器所标注的大于1的整数或小数电容量值的缺省单位是：
\begin{enumerate}[leftmargin=3em]
\item 皮法拉
\item 微法拉
\item 法拉
\item 毫法拉
\end{enumerate}
\noindent\textbf{解说：}\textbf{}。\\\noindent\textbf{答案：}

\bigskip




\noindent\textbf{试题：}在电路图中或元件上，非电解电容器的容量标注为3位数字时，读取容量的方法为：
\begin{enumerate}[leftmargin=3em]
\item 前两位表示容量基数，后一位表示基数后面加上几个0，单位为皮法拉
\item 实际容量为该三位数字乘以1000，单位为皮法拉
\item 实际容量为该三位数字除以1000，单位为法拉
\item 前一位表示容量基数，第二位表示基数后面加上几个0，第三位表示误差级别，单位为皮法拉
\end{enumerate}
\noindent\textbf{解说：}\textbf{}。\\\noindent\textbf{答案：}

\bigskip




\noindent\textbf{试题：}在电路图中或原件上，电阻的阻值标注为3位数字时，读阻值的方法为：
\begin{enumerate}[leftmargin=3em]
\item 前两位表示阻值基数，后一位表示基数后面加上几个0，单位为欧姆
\item 实际阻值为该三位数字乘以100，单位欧姆欧
\item 实际阻值为该三位数字除以100，单位欧姆
\item 前一位表示阻值基数，第二位表示基数后面加上几个0，第三位表示误差级别，单位为千欧姆
\end{enumerate}
\noindent\textbf{解说：}\textbf{}。\\\noindent\textbf{答案：}

\bigskip




\noindent\textbf{试题：}在电路图中或元件上，电阻的阻值标注为2或3位数字和一个字母（例如8R2、27K、1M5）时，读阻值的方法为：
\begin{enumerate}[leftmargin=3em]
\item 把字母当做小数点，和数字连在一起表示阻值基数，字母R、K或M，分别表示单位为欧姆、千欧姆或兆欧姆
\item 把所有数字挑出来连在一起乘以100，单位为欧姆
\item 把所有数字挑出来连在一起，再根据字母为R、K或M，分别取欧姆、千欧姆或兆欧姆为单位
\item 字母前的数字表示阻值基数，字母后的数字表示误差级别，字母R、K或M，分别表示单位为欧姆、千欧姆或兆欧姆
\end{enumerate}
\noindent\textbf{解说：}\textbf{}。\\\noindent\textbf{答案：}

\bigskip




\noindent\textbf{试题：}电路图中，非电解电容器所标注的小于1的小数电容量值的缺省单位是：
\begin{enumerate}[leftmargin=3em]
\item 微法拉
\item 皮法拉
\item 法拉
\item 毫法拉
\end{enumerate}
\noindent\textbf{解说：}\textbf{}。\\\noindent\textbf{答案：}

\bigskip




\noindent\textbf{试题：}无线电设备中的变压器铁芯一般都用矽钢片堆迭或采用不导电的磁性材料，而不像继电器那样采用整块铁芯。这是因为：
\begin{enumerate}[leftmargin=3em]
\item 防止交流电流在导电的铁芯内感应出涡流、造成发热损耗
\item 整块铁芯容易漏电，不利用安全
\item 整块铁芯容易产生机械松动
\item 矽钢片和磁性材料比整块铁芯容易加工
\end{enumerate}
\noindent\textbf{解说：}\textbf{}。\\\noindent\textbf{答案：}

\bigskip




\noindent\textbf{试题：}在电感线圈中加入适当磁性材料做成的磁芯可以减少电流通过时的损耗，是因为：
\begin{enumerate}[leftmargin=3em]
\item 加入磁芯后可用较短导线维持所需的电感量，减少导线电阻引起的发热损耗
\item 磁芯比空气更容易散热，可减少线圈的发热损耗
\item 磁芯可将原来损耗的能量反射回导线
\item 加入磁性可改善导线趋肤效应引起局部电流密度增大而带来的发热损耗
\end{enumerate}
\noindent\textbf{解说：}\textbf{}。\\\noindent\textbf{答案：}

\bigskip




\noindent\textbf{试题：}在很多无线电和电子电路的旁路电容滤波电路中，经常将一大容量电容器与一小容量电容器 并联，这样的作用是：
\begin{enumerate}[leftmargin=3em]
\item 低频率由大电容负责，频率升高时大电容的引线电感不可忽略，则由小电容主要负责，以达到宽频带滤波效果
\item 为防止大电容干涸失效，采用小电容增大保险系数
\item 利用两个电容器容量相加的办法得到精确的容量
\item 增加电容器的总耐压指标
\end{enumerate}
\noindent\textbf{解说：}\textbf{}。\\\noindent\textbf{答案：}

\bigskip




\noindent\textbf{试题：}在直流电路中，用来阻碍电流流动的元件是？
\begin{enumerate}[leftmargin=3em]
\item 电阻
\item 电感
\item 电压表
\item 变压器
\end{enumerate}
\noindent\textbf{解说：}\textbf{}。\\\noindent\textbf{答案：}

\bigskip




\noindent\textbf{试题：}下列哪一个元件经常用来实现音量调节的功能？
\begin{enumerate}[leftmargin=3em]
\item 电位器
\item 功率电阻
\item 定值电阻
\item 变压器
\end{enumerate}
\noindent\textbf{解说：}\textbf{}。\\\noindent\textbf{答案：}

\bigskip




\noindent\textbf{试题：}电位器控制什么电学物理量？
\begin{enumerate}[leftmargin=3em]
\item 电阻
\item 电感
\item 电容
\item 场强
\end{enumerate}
\noindent\textbf{解说：}\textbf{}。\\\noindent\textbf{答案：}

\bigskip




\noindent\textbf{试题：}哪一种电子元件由两个或多个使用绝缘材料分离开的片状导体组成？
\begin{enumerate}[leftmargin=3em]
\item 电容
\item 电位器
\item 振荡器
\item 电阻
\end{enumerate}
\noindent\textbf{解说：}\textbf{}。\\\noindent\textbf{答案：}

\bigskip




\noindent\textbf{试题：}哪一种电子元件用来接通或切断电路？
\begin{enumerate}[leftmargin=3em]
\item 开关
\item 齐纳二极管
\item 电感
\item 可变电阻
\end{enumerate}
\noindent\textbf{解说：}\textbf{}。\\\noindent\textbf{答案：}

\bigskip




\noindent\textbf{试题：}一个充满电的镍镉电池的标称电压是多少？
\begin{enumerate}[leftmargin=3em]
\item 1.2伏
\item 1.0伏
\item 1.5伏
\item 2.2伏
\end{enumerate}
\noindent\textbf{解说：}\textbf{}。\\\noindent\textbf{答案：}

\bigskip




\noindent\textbf{试题：}继电器的功能可以描述为：
\begin{enumerate}[leftmargin=3em]
\item 由电磁铁控制的开关
\item 由电流控制的放大器
\item 一个光学传感器
\item 无线电转发设备
\end{enumerate}
\noindent\textbf{解说：}\textbf{}。\\\noindent\textbf{答案：}

\bigskip




\noindent\textbf{试题：}在业余收发信机电路中，经常用于隔直流或者给交流信号提供旁路的元件是：
\begin{enumerate}[leftmargin=3em]
\item 电容
\item 电阻
\item 电感
\item 半导体开关器件
\end{enumerate}
\noindent\textbf{解说：}\textbf{}。\\\noindent\textbf{答案：}

\bigskip




\noindent\textbf{试题：}在业余收发信机电路中，经常用谐振回路来筛选一定频率的信号。组成谐振回路的主要元器件是：
\begin{enumerate}[leftmargin=3em]
\item 电容和电感的组合
\item 电阻和电容的组合
\item 电感和电阻的组合
\item 半导体三极管和电阻的组合
\end{enumerate}
\noindent\textbf{解说：}\textbf{}。\\\noindent\textbf{答案：}

\bigskip




\noindent\textbf{试题：}一个电容器，在某一频率下测得容抗为若干欧姆。如果频率提高N倍，其容抗将：
\begin{enumerate}[leftmargin=3em]
\item 减少到原来的1/N
\item 增大到原来的N倍
\item 减少到原来的1/（2πN）
\item 增大到原来的2πN倍
\end{enumerate}
\noindent\textbf{解说：}\textbf{}。\\\noindent\textbf{答案：}

\bigskip




\noindent\textbf{试题：}一个电感线圈，在某一频率下测得感抗为若干欧姆。如果频率提高N倍，其感抗将：
\begin{enumerate}[leftmargin=3em]
\item 增大到原来的N倍
\item 减少到原来的1/N
\item 增大到原来的2πN倍
\item 减少到原来的1/（2πN）
\end{enumerate}
\noindent\textbf{解说：}\textbf{}。\\\noindent\textbf{答案：}

\bigskip




\noindent\textbf{试题：}一个电容器，在50Hz频率下实测得容抗为100欧姆。按容抗公式计算，在5MHz时容抗应为0.001欧姆，但实测到的阻抗明显偏大。这是因为：
\begin{enumerate}[leftmargin=3em]
\item 电容器引线在频率变高后带来不可忽略的感抗
\item 仪器的频率特性欠佳
\item 容抗计算公式在高频率下不适用
\item 电容器介质损耗造成非线性失真
\end{enumerate}
\noindent\textbf{解说：}\textbf{}。\\\noindent\textbf{答案：}

\bigskip




\noindent\textbf{试题：}一个电感线圈，在50Hz频率下实测得感抗为10欧姆。按感抗公式计算，在5MHz时感抗应为1兆欧，但实测到的阻抗明显偏小。这是因为：
\begin{enumerate}[leftmargin=3em]
\item 电感线圈的分布电容在频率变高后带来不可忽略的容抗
\item 仪器的频率特性欠佳
\item 感抗计算公式在高频率下不适用
\item 电感线圈导线损耗造成非线性失真
\end{enumerate}
\noindent\textbf{解说：}\textbf{}。\\\noindent\textbf{答案：}

\bigskip




\noindent\textbf{试题：}在电容器两端施加一个正弦交流电压，则流过电容器的电流是正弦交流电流，该电流相位：
\begin{enumerate}[leftmargin=3em]
\item 超前于电压相位90度
\item 落后于电压相位90度
\item 与电压相位相同
\item 与电压相差180度
\end{enumerate}
\noindent\textbf{解说：}\textbf{}。\\\noindent\textbf{答案：}

\bigskip




\noindent\textbf{试题：}使一个正弦交流电流流过电容器，则电容器两端的电压是正弦交流电压，该电压相位：
\begin{enumerate}[leftmargin=3em]
\item 落后于电流相位90度
\item 超前于电流相位90度
\item 与电压相位相同
\item 与电流相差180度
\end{enumerate}
\noindent\textbf{解说：}\textbf{}。\\\noindent\textbf{答案：}

\bigskip




\noindent\textbf{试题：}在电容器两端施加一定幅度的正弦交流电压。流过电容器的电流幅度：
\begin{enumerate}[leftmargin=3em]
\item 与电压和电容量都成正比
\item 与电压和电容量都成反比
\item 与电压成正比，与电容量成反比
\item 与电容量成正比，与电压成反比
\end{enumerate}
\noindent\textbf{解说：}\textbf{}。\\\noindent\textbf{答案：}

\bigskip




\noindent\textbf{试题：}在线圈两端施加一个正弦交流电压，则流过线圈的电流是正弦交流电流，该电流相位：
\begin{enumerate}[leftmargin=3em]
\item 落后于电压相位90度
\item 超前于电压相位90度
\item 与电压相位相同
\item 与电压相差180度
\end{enumerate}
\noindent\textbf{解说：}\textbf{}。\\\noindent\textbf{答案：}

\bigskip




\noindent\textbf{试题：}使一个正弦交流电流流过线圈，则线圈两端的电压是正弦交流电压，该电压相位：
\begin{enumerate}[leftmargin=3em]
\item 超前于电流相位90度
\item 落后于电流相位90度
\item 与电流相位相同
\item 与电流相差180度
\end{enumerate}
\noindent\textbf{解说：}\textbf{}。\\\noindent\textbf{答案：}

\bigskip




\noindent\textbf{试题：}在线圈两端施加一定幅度的正弦交流电压。流过线圈的电流幅度：
\begin{enumerate}[leftmargin=3em]
\item 与电压成正比，与电感量成反比
\item 与电压和电感量都成正比
\item 与电压和电感量都成反比
\item 与电感量成正比，与电压成反比
\end{enumerate}
\noindent\textbf{解说：}\textbf{}。\\\noindent\textbf{答案：}

\bigskip




\noindent\textbf{试题：}在电容、电感串联电路中，如果电容的容抗与电感的感抗大小相等，描述这一状态最适当的术语是：
\begin{enumerate}[leftmargin=3em]
\item 谐振
\item 匹配
\item 叠加
\item 干涉
\end{enumerate}
\noindent\textbf{解说：}\textbf{}。\\\noindent\textbf{答案：}

\bigskip




\noindent\textbf{试题：}由电感为L的线圈、容量为C的电容组成LC谐振电路，其谐振频率f为：（L、C和f的单位分别为微亨、微法拉和兆赫, “x＾m”表示“x的m次方”）
\begin{enumerate}[leftmargin=3em]
\item f = 1/（2π(LC)＾(1/2)）
\item f = 2π(LC)＾(1/2)
\item f = 1/((2πLC)＾2)
\item f = 2π(LC)＾(1/2)
\end{enumerate}
\noindent\textbf{解说：}\textbf{}。\\\noindent\textbf{答案：}

\bigskip




\noindent\textbf{试题：}L和C串联后，其总阻抗为：
\begin{enumerate}[leftmargin=3em]
\item 在谐振频率达到最小值
\item 在谐振频率达到最大值
\item 随频率增高而不断增大
\item 随频率增高而不断减小
\end{enumerate}
\noindent\textbf{解说：}\textbf{}。\\\noindent\textbf{答案：}

\bigskip




\noindent\textbf{试题：}L和C并联后，其总阻抗为：
\begin{enumerate}[leftmargin=3em]
\item 在谐振频率达到最大值
\item 在谐振频率达到最小值
\item 随频率增高而不断增大
\item 随频率增高而不断减小
\end{enumerate}
\noindent\textbf{解说：}\textbf{}。\\\noindent\textbf{答案：}

\bigskip




\noindent\textbf{试题：}构成振荡器的必备元素是：
\begin{enumerate}[leftmargin=3em]
\item 放大倍数大于1的放大器、正反馈电路
\item 放大倍数大于1的放大器、负反馈电路
\item LC或晶体谐振电路、正反馈电路
\item 任意放大器、LC或晶体谐振电路
\end{enumerate}
\noindent\textbf{解说：}\textbf{}。\\\noindent\textbf{答案：}

\bigskip




\noindent\textbf{试题：}能够以电场形式储存能量的元件是什么？
\begin{enumerate}[leftmargin=3em]
\item 电容
\item 电阻
\item 压敏元件
\item 电感
\end{enumerate}
\noindent\textbf{解说：}\textbf{}。\\\noindent\textbf{答案：}

\bigskip




\noindent\textbf{试题：}能够以磁场形式储存能量的元件是什么？
\begin{enumerate}[leftmargin=3em]
\item 电感
\item 电阻
\item 电磁铁
\item 电容
\end{enumerate}
\noindent\textbf{解说：}\textbf{}。\\\noindent\textbf{答案：}

\bigskip




\noindent\textbf{试题：}哪一种电子元件一般由线圈组成？
\begin{enumerate}[leftmargin=3em]
\item 电感
\item 电容
\item 二极管
\item 开关
\end{enumerate}
\noindent\textbf{解说：}\textbf{}。\\\noindent\textbf{答案：}

\bigskip




\noindent\textbf{试题：}下列哪一项和电感一起使用，可以制作一个谐振电路？
\begin{enumerate}[leftmargin=3em]
\item 电容
\item 齐纳二极管
\item 电位器
\item 电阻
\end{enumerate}
\noindent\textbf{解说：}\textbf{}。\\\noindent\textbf{答案：}

\bigskip




\noindent\textbf{试题：}附图中的电路元器件符号代表的是：
\begin{enumerate}[leftmargin=3em]
\item 接地
\item 天线
\item 电阻
\item 二极管
\end{enumerate}
\noindent\textbf{解说：}\textbf{}。\\\noindent\textbf{答案：}

\bigskip

LK0497.jpg


\noindent\textbf{试题：}附图中的电路元器件符号代表的是：
\begin{enumerate}[leftmargin=3em]
\item 天线
\item 接地
\item 电阻
\item 二极管
\end{enumerate}
\noindent\textbf{解说：}\textbf{}。\\\noindent\textbf{答案：}

\bigskip

LK0498.jpg


\noindent\textbf{试题：}附图中的电路元器件符号代表的是：
\begin{enumerate}[leftmargin=3em]
\item 熔断器
\item 电容
\item 电阻
\item 二极管
\end{enumerate}
\noindent\textbf{解说：}\textbf{}。\\\noindent\textbf{答案：}

\bigskip

LK0499.jpg


\noindent\textbf{试题：}附图中的电路元器件符号代表的是：
\begin{enumerate}[leftmargin=3em]
\item 电容器
\item 熔断器
\item 电阻
\item 二极管
\end{enumerate}
\noindent\textbf{解说：}\textbf{}。\\\noindent\textbf{答案：}

\bigskip

LK0500.jpg


\noindent\textbf{试题：}附图中的电路元器件符号代表的是：
\begin{enumerate}[leftmargin=3em]
\item 电阻
\item 电容器
\item 熔断器
\item 压电晶体
\end{enumerate}
\noindent\textbf{解说：}\textbf{}。\\\noindent\textbf{答案：}

\bigskip

LK0501.jpg


\noindent\textbf{试题：}附图中的电路元器件符号代表的是：
\begin{enumerate}[leftmargin=3em]
\item 二极管
\item 电容器
\item 线圈
\item 电阻
\end{enumerate}
\noindent\textbf{解说：}\textbf{}。\\\noindent\textbf{答案：}

\bigskip

LK0502.jpg


\noindent\textbf{试题：}附图中的电路元器件符号代表的是：
\begin{enumerate}[leftmargin=3em]
\item 线圈
\item 二极管
\item 电容器
\item 电阻
\end{enumerate}
\noindent\textbf{解说：}\textbf{}。\\\noindent\textbf{答案：}

\bigskip

LK0503.jpg


\noindent\textbf{试题：}附图中的电路元器件符号代表的是：
\begin{enumerate}[leftmargin=3em]
\item 电池
\item 二极管
\item 线圈
\item 电阻
\end{enumerate}
\noindent\textbf{解说：}\textbf{}。\\\noindent\textbf{答案：}

\bigskip

LK0504.jpg


\noindent\textbf{试题：}附图中的电路元器件符号代表的是：
\begin{enumerate}[leftmargin=3em]
\item 压电晶体
\item 电池
\item 二极管
\item 电阻
\end{enumerate}
\noindent\textbf{解说：}\textbf{}。\\\noindent\textbf{答案：}

\bigskip

LK0505.jpg


\noindent\textbf{试题：}附图中的电路元器件符号代表的是：
\begin{enumerate}[leftmargin=3em]
\item 稳压二极管
\item 压电晶体
\item 发光二极管
\item 电阻
\end{enumerate}
\noindent\textbf{解说：}\textbf{}。\\\noindent\textbf{答案：}

\bigskip

LK0506.jpg


\noindent\textbf{试题：}附图中的电路元器件符号代表的是：
\begin{enumerate}[leftmargin=3em]
\item 发光二极管
\item 稳压二极管
\item 压电晶体
\item 电阻
\end{enumerate}
\noindent\textbf{解说：}\textbf{}。\\\noindent\textbf{答案：}

\bigskip

LK0507.jpg


\noindent\textbf{试题：}附图中的电路元器件符号代表的是：
\begin{enumerate}[leftmargin=3em]
\item PNP双极型半导体三极管
\item NPN双极型半导体三极管
\item 结型场效应半导体三极管
\item 绝缘栅场效应半导体三极管
\end{enumerate}
\noindent\textbf{解说：}\textbf{}。\\\noindent\textbf{答案：}

\bigskip

LK0508.jpg


\noindent\textbf{试题：}附图中的电路元器件符号代表的是：
\begin{enumerate}[leftmargin=3em]
\item NPN双极型半导体三极管
\item PNP双极型半导体三极管
\item 结型场效应半导体三极管
\item 绝缘栅场效应半导体三极管
\end{enumerate}
\noindent\textbf{解说：}\textbf{}。\\\noindent\textbf{答案：}

\bigskip

LK0509.jpg


\noindent\textbf{试题：}附图中的电路元器件符号代表的是：
\begin{enumerate}[leftmargin=3em]
\item 结型场效应半导体三极管
\item PNP双极型半导体三极管
\item NPN双极型半导体三极管
\item 绝缘栅场效应半导体三极管
\end{enumerate}
\noindent\textbf{解说：}\textbf{}。\\\noindent\textbf{答案：}

\bigskip

LK0510.jpg


\noindent\textbf{试题：}附图中的电路元器件符号代表的是：
\begin{enumerate}[leftmargin=3em]
\item 绝缘栅场效应半导体三极管
\item 结型场效应半导体三极管
\item PNP双极型半导体三极管
\item NPN双极型半导体三极管
\end{enumerate}
\noindent\textbf{解说：}\textbf{}。\\\noindent\textbf{答案：}

\bigskip

LK0511.jpg


\noindent\textbf{试题：}下列电路具备的功能是：
\begin{enumerate}[leftmargin=3em]
\item 低通滤波
\item 带通滤波
\item 高通滤波
\item 过流保护
\end{enumerate}
\noindent\textbf{解说：}\textbf{}。\\\noindent\textbf{答案：}

\bigskip

LK0594.jpg


\noindent\textbf{试题：}下列电路具备的功能是：
\begin{enumerate}[leftmargin=3em]
\item 低通滤波
\item 带通滤波
\item 高通滤波
\item 过流保护
\end{enumerate}
\noindent\textbf{解说：}\textbf{}。\\\noindent\textbf{答案：}

\bigskip

LK0595.jpg


\noindent\textbf{试题：}下列电路具备的功能是：
\begin{enumerate}[leftmargin=3em]
\item 低通滤波
\item 带通滤波
\item 高通滤波
\item 过流保护
\end{enumerate}
\noindent\textbf{解说：}\textbf{}。\\\noindent\textbf{答案：}

\bigskip

LK0596.jpg


\noindent\textbf{试题：}下列电路具备的功能是：
\begin{enumerate}[leftmargin=3em]
\item 低通滤波
\item 带通滤波
\item 高通滤波
\item 过流保护
\end{enumerate}
\noindent\textbf{解说：}\textbf{}。\\\noindent\textbf{答案：}

\bigskip

LK0597.jpg


\noindent\textbf{试题：}下列电路具备的功能是：
\begin{enumerate}[leftmargin=3em]
\item 低通滤波
\item 带通滤波
\item 高通滤波
\item 过流保护
\end{enumerate}
\noindent\textbf{解说：}\textbf{}。\\\noindent\textbf{答案：}

\bigskip

LK0598.jpg


\noindent\textbf{试题：}下列电路具备的功能是：
\begin{enumerate}[leftmargin=3em]
\item 高通滤波
\item 低通滤波
\item 带通滤波
\item 过流保护
\end{enumerate}
\noindent\textbf{解说：}\textbf{}。\\\noindent\textbf{答案：}

\bigskip

LK0599.jpg


\noindent\textbf{试题：}下列电路具备的功能是：
\begin{enumerate}[leftmargin=3em]
\item 高通滤波
\item 低通滤波
\item 带通滤波
\item 过流保护
\end{enumerate}
\noindent\textbf{解说：}\textbf{}。\\\noindent\textbf{答案：}

\bigskip

LK0600.jpg


\noindent\textbf{试题：}下列电路具备的功能是：
\begin{enumerate}[leftmargin=3em]
\item 高通滤波
\item 低通滤波
\item 带通滤波
\item 过流保护
\end{enumerate}
\noindent\textbf{解说：}\textbf{}。\\\noindent\textbf{答案：}

\bigskip

LK0601.jpg


\noindent\textbf{试题：}下列电路具备的功能是：
\begin{enumerate}[leftmargin=3em]
\item 高通滤波
\item 带通滤波
\item 低通滤波
\item 过流保护
\end{enumerate}
\noindent\textbf{解说：}\textbf{}。\\\noindent\textbf{答案：}

\bigskip

LK0602.jpg


\noindent\textbf{试题：}下列电路具备的功能是：
\begin{enumerate}[leftmargin=3em]
\item 高通滤波
\item 带通滤波
\item 低通滤波
\item 过流保护
\end{enumerate}
\noindent\textbf{解说：}\textbf{}。\\\noindent\textbf{答案：}

\bigskip

LK0603.jpg


\noindent\textbf{试题：}下列电路具备的功能是：
\begin{enumerate}[leftmargin=3em]
\item 带通滤波
\item 高通滤波
\item 低通滤波
\item 过流保护
\end{enumerate}
\noindent\textbf{解说：}\textbf{}。\\\noindent\textbf{答案：}

\bigskip

LK0604.jpg


\noindent\textbf{试题：}下列电路是一个：
\begin{enumerate}[leftmargin=3em]
\item 电感耦合式振荡器
\item 电容三点式振荡器
\item 电感三点式振荡器
\item RC反馈式振荡器
\end{enumerate}
\noindent\textbf{解说：}\textbf{}。\\\noindent\textbf{答案：}

\bigskip

LK0605.jpg


\noindent\textbf{试题：}下列电路是一个：
\begin{enumerate}[leftmargin=3em]
\item 电感三点式振荡器
\item 电容三点式振荡器
\item 电感反馈式振荡器
\item RC反馈式振荡器
\end{enumerate}
\noindent\textbf{解说：}\textbf{}。\\\noindent\textbf{答案：}

\bigskip

LK0606.jpg


\noindent\textbf{试题：}下列电路是一个：
\begin{enumerate}[leftmargin=3em]
\item 电容三点式振荡器
\item RC反馈式振荡器
\item 电感三点式振荡器
\item 电感反馈式振荡器
\end{enumerate}
\noindent\textbf{解说：}\textbf{}。\\\noindent\textbf{答案：}

\bigskip

LK0607.jpg


\noindent\textbf{试题：}下列电路是一个：
\begin{enumerate}[leftmargin=3em]
\item RC反馈式振荡器
\item 电容三点式振荡器
\item 电感三点式振荡器
\item 电感反馈式振荡器
\end{enumerate}
\noindent\textbf{解说：}\textbf{}。\\\noindent\textbf{答案：}

\bigskip

LK0608.jpg


\noindent\textbf{试题：}下列电路是一个：
\begin{enumerate}[leftmargin=3em]
\item 电容三点式振荡器
\item 电感反馈式振荡器
\item 电感三点式振荡器
\item RC反馈式振荡器
\end{enumerate}
\noindent\textbf{解说：}\textbf{}。\\\noindent\textbf{答案：}

\bigskip

LK0609.jpg


\noindent\textbf{试题：}下列电路是一个：
\begin{enumerate}[leftmargin=3em]
\item 二极管半波整流器
\item 二极管全波整流器
\item 二极管桥式整流器
\item 二极管环形调制器
\end{enumerate}
\noindent\textbf{解说：}\textbf{}。\\\noindent\textbf{答案：}

\bigskip

LK0610.jpg


\noindent\textbf{试题：}下列电路是一个：
\begin{enumerate}[leftmargin=3em]
\item 二极管桥式整流器
\item 二极管全波整流器
\item 二极管半波整流器
\item 二极管环形调制器
\end{enumerate}
\noindent\textbf{解说：}\textbf{}。\\\noindent\textbf{答案：}

\bigskip

LK0611.jpg


\noindent\textbf{试题：}下列电路是一个：
\begin{enumerate}[leftmargin=3em]
\item 二极管全波整流器
\item 二极管半波整流器
\item 二极管桥式整流器
\item 二极管环形调制器
\end{enumerate}
\noindent\textbf{解说：}\textbf{}。\\\noindent\textbf{答案：}

\bigskip

LK0612.jpg


\noindent\textbf{试题：}下列电路是一个：
\begin{enumerate}[leftmargin=3em]
\item 二极管环形调制器
\item 二极管全波整流器
\item 二极管桥式整流器
\item 二极管半波整流器
\end{enumerate}
\noindent\textbf{解说：}\textbf{}。\\\noindent\textbf{答案：}

\bigskip

LK0613.jpg


\noindent\textbf{试题：}图示电路为业余无线电爱好者常用的简易场强表。各部件的作用分别为：
\begin{enumerate}[leftmargin=3em]
\item A-接收电波，C1、C2、L-谐振选频，D-检波，C3-旁路滤波，M-指示
\item A-接收电波，C1、C2旁路滤波，L-升压，D-放大，C3-隔直流，M-指示
\item A-整流，C1、C2-隔直流，L-放大，D-滤波，C3-谐振，M-指示
\item A-接收电波，C1、C2高频旁路，L-放大，D-开关，C3-耦合，M-指示
\end{enumerate}
\noindent\textbf{解说：}\textbf{}。\\\noindent\textbf{答案：}

\bigskip

LK0614.jpg


\noindent\textbf{试题：}图中的电路为：
\begin{enumerate}[leftmargin=3em]
\item 共基极放大器
\item 共发射极放大器
\item 共集电极放大器
\item 共栅极放大器
\end{enumerate}
\noindent\textbf{解说：}\textbf{}。\\\noindent\textbf{答案：}

\bigskip

LK0734.jpg


\noindent\textbf{试题：}图中的电路为：
\begin{enumerate}[leftmargin=3em]
\item 共发射极放大器
\item 共集电极放大器
\item 共基极放大器
\item 共源极放大器
\end{enumerate}
\noindent\textbf{解说：}\textbf{}。\\\noindent\textbf{答案：}

\bigskip

LK0735.jpg


\noindent\textbf{试题：}图中的电路为：
\begin{enumerate}[leftmargin=3em]
\item 共集电极放大器
\item 共发射极放大器
\item 共基极放大器
\item 共漏极放大器
\end{enumerate}
\noindent\textbf{解说：}\textbf{}。\\\noindent\textbf{答案：}

\bigskip

LK0736.jpg


\noindent\textbf{试题：}将电阻R和电容C串联后突然接到直流电压U上，电容C两端的电压会：
\begin{enumerate}[leftmargin=3em]
\item 从0按指数规律逐渐增加到U
\item 从U按指数规律逐渐减小到0
\item 从0突然跳到U，然后再按指数规律逐渐减小到0
\item 从U突然跳到0，然后再按指数规律逐渐增大到U
\end{enumerate}
\noindent\textbf{解说：}\textbf{}。\\\noindent\textbf{答案：}

\bigskip




\noindent\textbf{试题：}将电阻R和电容C串联后突然接到直流电压U上，电容C两端的电压在经过时间常数T后大约会达到U的：
\begin{enumerate}[leftmargin=3em]
\item 63%
\item 99%
\item 37%
\item 6.28%
\end{enumerate}
\noindent\textbf{解说：}\textbf{}。\\\noindent\textbf{答案：}

\bigskip




\noindent\textbf{试题：}将电阻R和电容C串联后突然接到直流电压U上，电阻R两端的电压会：
\begin{enumerate}[leftmargin=3em]
\item 从0突然跳到U，然后再按指数规律逐渐减小到0
\item 从0按指数规律逐渐增加到U
\item 从U按直线规律逐渐减小到0
\item 从U突然跳到0，然后再按直线规律逐渐减小到U
\end{enumerate}
\noindent\textbf{解说：}\textbf{}。\\\noindent\textbf{答案：}

\bigskip




\noindent\textbf{试题：}将电阻R和电容C串联后突然接到直流电压U上，电阻R两端的电压在经过时间常数T后大约会达到U的：
\begin{enumerate}[leftmargin=3em]
\item 37%
\item 99%
\item 63%
\item 6.28%
\end{enumerate}
\noindent\textbf{解说：}\textbf{}。\\\noindent\textbf{答案：}

\bigskip




\noindent\textbf{试题：}将电阻R和电容C串联后突然接到直流电压U上，流过电阻R的电流会：
\begin{enumerate}[leftmargin=3em]
\item 从0突然跳到U/R，然后再按指数规律逐渐减小到0
\item 从0突然跳到U/R并保持
\item 从0按指数规律逐渐增加到U/R
\item 从U/R突然跳到0并保持
\end{enumerate}
\noindent\textbf{解说：}\textbf{}。\\\noindent\textbf{答案：}

\bigskip




\noindent\textbf{试题：}将电阻R和电容C串联后突然接到直流电压U上，流过电容C的电流会：
\begin{enumerate}[leftmargin=3em]
\item 从0突然跳到U/R，然后再按指数规律逐渐减小到0
\item 从0突然跳到U/R并保持
\item 从0按指数规律逐渐增加到U/R
\item 从U/R突然跳到0并保持
\end{enumerate}
\noindent\textbf{解说：}\textbf{}。\\\noindent\textbf{答案：}

\bigskip




\noindent\textbf{试题：}电阻R和电容C并联后接在电压为U的直流电源上。突然断开电源，电容C两端的电压会：
\begin{enumerate}[leftmargin=3em]
\item 从U按指数规律逐渐减小到0
\item 从0按指数规律逐渐增加到U
\item 从0突然跳到U并保持
\item 从U突然跳到0并保持
\end{enumerate}
\noindent\textbf{解说：}\textbf{}。\\\noindent\textbf{答案：}

\bigskip




\noindent\textbf{试题：}电阻R和电容C并联后接在电压为U的直流电源上。突然断开电源，电阻R两端的电压会：
\begin{enumerate}[leftmargin=3em]
\item 从U按指数规律逐渐减小到0
\item 从0按指数规律逐渐增加到U
\item 从0突然跳到U并保持
\item 从U突然跳到0并保持
\end{enumerate}
\noindent\textbf{解说：}\textbf{}。\\\noindent\textbf{答案：}

\bigskip




\noindent\textbf{试题：}电阻R和电容C并联后接在电压为U的直流电源上。突然断开电源，流过电阻R的电流会：
\begin{enumerate}[leftmargin=3em]
\item 从U/R按指数规律逐渐减小到0
\item 从0按指数规律逐渐增加到U/R
\item 从0突然跳到U/R，然后再按指数规律逐渐减小到0
\item 从U突然跳到0，然后再按指数规律逐渐增大到U/R
\end{enumerate}
\noindent\textbf{解说：}\textbf{}。\\\noindent\textbf{答案：}

\bigskip




\noindent\textbf{试题：}电阻R和电容C并联后接在电压为U的直流电源上。突然断开电源，流过电容C的电流会：
\begin{enumerate}[leftmargin=3em]
\item 从0突然跳到U/R，然后再按指数规律逐渐减小到0
\item 从0突然跳到U/R并保持
\item 从0按指数规律逐渐增加到U/R
\item 从U突然跳到0，然后再按指数规律逐渐增大到U/R
\end{enumerate}
\noindent\textbf{解说：}\textbf{}。\\\noindent\textbf{答案：}

\bigskip




\noindent\textbf{试题：}将电阻R和电感L串联后突然接到直流电压U上，电感L两端的电压会：
\begin{enumerate}[leftmargin=3em]
\item 从0突然跳到U，然后再按指数规律逐渐减小到0
\item 从0按指数规律逐渐增加到U
\item 从U按指数规律逐渐减小到0
\item 从U突然跳到0并保持
\end{enumerate}
\noindent\textbf{解说：}\textbf{}。\\\noindent\textbf{答案：}

\bigskip




\noindent\textbf{试题：}将电阻R和电感L串联后突然接到直流电压U上，电感L两端的电压在经过时间常数T后大约会达到U的：
\begin{enumerate}[leftmargin=3em]
\item 37%
\item 99%
\item 63%
\item 6.28%
\end{enumerate}
\noindent\textbf{解说：}\textbf{}。\\\noindent\textbf{答案：}

\bigskip




\noindent\textbf{试题：}将电阻R和电感L串联后突然接到直流电压U上，电阻R两端的电压在经过时间常数T后大约会达到U的：
\begin{enumerate}[leftmargin=3em]
\item 63%
\item 99%
\item 37%
\item 6.28%
\end{enumerate}
\noindent\textbf{解说：}\textbf{}。\\\noindent\textbf{答案：}

\bigskip




\noindent\textbf{试题：}将电阻R和电感L串联后突然接到直流电压U上，电阻R两端的电压会：
\begin{enumerate}[leftmargin=3em]
\item 从0按指数规律逐渐增加到U
\item 从U按指数规律逐渐减小到0
\item 从0突然跳到U，然后再按指数规律逐渐减小到0
\item 从U突然跳到0并保持
\end{enumerate}
\noindent\textbf{解说：}\textbf{}。\\\noindent\textbf{答案：}

\bigskip




\noindent\textbf{试题：}将电阻R和电感L串联后突然接到直流电压U上，流过电阻R的电流会：
\begin{enumerate}[leftmargin=3em]
\item 从0按指数规律逐渐增加到U/R
\item 从0突然跳到U/R，然后再按指数规律逐渐减小到0
\item 从U/R按指数规律逐渐减小到0
\item 从0突然跳到U/R并保持U/R
\end{enumerate}
\noindent\textbf{解说：}\textbf{}。\\\noindent\textbf{答案：}

\bigskip




\noindent\textbf{试题：}将电阻R和电感L串联后突然接到直流电压U上，流过电感L的电流会：
\begin{enumerate}[leftmargin=3em]
\item 从0按指数规律逐渐增加到U/R
\item 从0突然跳到U/R，然后再按指数规律逐渐减小到0
\item 从U/R按指数规律逐渐减小到0
\item 从0突然跳到U/R并保持
\end{enumerate}
\noindent\textbf{解说：}\textbf{}。\\\noindent\textbf{答案：}

\bigskip




\noindent\textbf{试题：}电阻R和电感L并联后接在电流为I的直流电路中。突然断开电路，电感L两端的电压会：
\begin{enumerate}[leftmargin=3em]
\item 从I*R按指数规律逐渐减小到0
\item 从0按指数规律逐渐增加到I*R
\item 保持在I*R
\item 保持在0
\end{enumerate}
\noindent\textbf{解说：}\textbf{}。\\\noindent\textbf{答案：}

\bigskip




\noindent\textbf{试题：}将继电器的线圈接在开关三极管的集电极电路中，用三极管控制继电器的通断。可能发生的现象是：
\begin{enumerate}[leftmargin=3em]
\item 三极管关断的瞬间集电极可能承受远高于电源的电压冲击
\item 三极管导通的瞬间集电极可能承受远高于电源的电压冲击
\item 三极管关断的瞬间集电极可能承受远高于工作电流的电流冲击
\item 三极管导通的瞬间集电极可能承受远高于工作电流的电流冲击
\end{enumerate}
\noindent\textbf{解说：}\textbf{}。\\\noindent\textbf{答案：}

\bigskip




\noindent\textbf{试题：}用开关器件控制继电器线圈的通断时，经常要在继电器线圈两端并联一个二极管或者电容，其作用是：
\begin{enumerate}[leftmargin=3em]
\item 吸收继电器线圈电感在电路关断瞬间感应出的反向电动势以保护开关器件
\item 消除继电器接点通断时产生的火花干扰
\item 减少通过开关器件的工作电流
\item 减少继电器的磨损
\end{enumerate}
\noindent\textbf{解说：}\textbf{}。\\\noindent\textbf{答案：}

\bigskip




\noindent\textbf{试题：}用一个电压为4.2伏的低电压电池和一堆无源电子元件做电路实验，但连接电路时感觉手不慎被高电压击了一下。可能产生这个高电压的元件是：
\begin{enumerate}[leftmargin=3em]
\item 电源变压器的绕组
\item 电解电容器
\item 额定功率为50瓦的大电阻
\item 大电流高反压二极管
\end{enumerate}
\noindent\textbf{解说：}\textbf{}。\\\noindent\textbf{答案：}

\bigskip




\noindent\textbf{试题：}用继电器控制直流负载时，有时在继电器接点线圈两端并联一个电容，其作用是：
\begin{enumerate}[leftmargin=3em]
\item 消除继电器接点通断时产生的火花干扰
\item 吸收继电器线圈电感在电路关断瞬间感应出的反向电动势以保护开关器件
\item 减少通过开关器件的工作电流
\item 减少继电器的磨损
\end{enumerate}
\noindent\textbf{解说：}\textbf{}。\\\noindent\textbf{答案：}

\bigskip




\noindent\textbf{试题：}“Q值”是经常被用来描述谐振回路或者谐振天线电路的一个电气参数。它表示的是：
\begin{enumerate}[leftmargin=3em]
\item 谐振回路中信号的能量损耗情况。损耗越小，Q值越高
\item 谐振回路或者天线所耐受的最大功率。耐受功率越大，Q值越高
\item 谐振回路体积或者天线尺寸与波长的比例。体积或尺寸越紧凑，Q值越高
\item 谐振回路或者天线所能正常工作的最高频率。工作频率上限越高， Q值越高
\end{enumerate}
\noindent\textbf{解说：}\textbf{}。\\\noindent\textbf{答案：}

\bigskip




\noindent\textbf{试题：}根据谐振回路或者谐振天线电路Q值的高低可以做出如下判断：
\begin{enumerate}[leftmargin=3em]
\item Q值越高，谐振曲线越尖锐，选择性越好，通带宽度越窄
\item Q值越高，谐振曲线越平坦，选择性越不明显，通带宽度越宽
\item Q值越高，回路或天线的工作状态受频率变化的影响越小
\item Q值越高，回路或天线的工作频率越高
\end{enumerate}
\noindent\textbf{解说：}\textbf{}。\\\noindent\textbf{答案：}

\bigskip




\noindent\textbf{试题：}谐振回路的通带宽度BW是指：
\begin{enumerate}[leftmargin=3em]
\item 回路两端电压信号幅度从中心频率衰减3dB时上、下限频率的间距
\item 回路两端电压信号幅度从中心频率衰减30%时上、下限频率的间距
\item 回路两端电压信号幅度从中心频率衰减80%时上、下限频率的间距
\item 回路两端电压信号幅度从中心频率衰减95%时上、下限频率的间距
\end{enumerate}
\noindent\textbf{解说：}\textbf{}。\\\noindent\textbf{答案：}

\bigskip




\noindent\textbf{试题：}滤波器的“截止频率”是指：
\begin{enumerate}[leftmargin=3em]
\item 输出频率特性曲线从通带的0dB变化到-3dB的频率
\item 高于该频率的信号将被滤波器完全切除
\item 低于该频率的信号将被滤波器完全切除
\item 高于该频率的信号将会在滤波器中发生非线性失真
\end{enumerate}
\noindent\textbf{解说：}\textbf{}。\\\noindent\textbf{答案：}

\bigskip




\noindent\textbf{试题：}滤波器的“3dB带宽”是指：
\begin{enumerate}[leftmargin=3em]
\item 输出频率特性曲线从通带的0dB变化到-3dB的频率之间的宽度
\item 输出信号相对于输入信号衰减3dB以上（含3dB）的频率范围宽度
\item 输出信号相对于输入信号衰减3dB以下（含3dB）的频率范围宽度
\item 滤波器维持3dB增益的频率比范围宽度
\end{enumerate}
\noindent\textbf{解说：}\textbf{}。\\\noindent\textbf{答案：}

\bigskip




\noindent\textbf{试题：}滤波器有一个称为“通带波动”的参数，是指：
\begin{enumerate}[leftmargin=3em]
\item 通带范围内衰减因频率不同的变化幅度
\item 滤波器频率特性在整个频谱范围内的变化趋势
\item 等幅信号通过滤波器时幅度发生抖动的现象
\item 滤波器因元器件不稳定等随机原因使输出信号幅度发生上下波动的现象
\end{enumerate}
\noindent\textbf{解说：}\textbf{}。\\\noindent\textbf{答案：}

\bigskip




\noindent\textbf{试题：}在下列种类的谐振回路中，通频带最窄的是：
\begin{enumerate}[leftmargin=3em]
\item 松耦合的高Q单调谐或同频双调谐回路
\item 紧耦合高Q同频双调谐回路
\item 松耦合低Q同频双调谐回路
\item 松耦合高Q参差调谐回路
\end{enumerate}
\noindent\textbf{解说：}\textbf{}。\\\noindent\textbf{答案：}

\bigskip




\noindent\textbf{试题：}用SSB接收机的天线引线靠近一个晶体管LC振荡器电路板，接收其信号。振荡器电路接通电源后，发现收到的信号音调会从低到高或者从高到低变化。这主要因为：
\begin{enumerate}[leftmargin=3em]
\item 元器件通电发热，引起相关LC参数变化，造成谐振频率漂移
\item 半导体晶体管处于老化过程
\item 随着射频能量泄露，电路的输出功率下降
\item 接收机的声音能量反馈到电路板引起
\end{enumerate}
\noindent\textbf{解说：}\textbf{}。\\\noindent\textbf{答案：}

\bigskip




\noindent\textbf{试题：}在无线电电路中常用于产生基准频率的元件中，按频率稳定度由低到高的排列为：
\begin{enumerate}[leftmargin=3em]
\item RC定时电路，LC回路，陶瓷谐振器，石英声表面波元件，石英晶体谐振器
\item LC回路，RC定时电路，陶瓷谐振器，石英声表面波元件，石英晶体谐振器
\item RC定时电路，LC回路，石英声表面波元件，石英晶体谐振器，陶瓷谐振器
\item RC定时电路，陶瓷谐振器，LC回路，石英晶体谐振器，石英声表面波元件
\end{enumerate}
\noindent\textbf{解说：}\textbf{}。\\\noindent\textbf{答案：}

\bigskip




\noindent\textbf{试题：}放大器的负反馈是指这样的电路：
\begin{enumerate}[leftmargin=3em]
\item 将放大器输出信号的一部分回输到放大器的输入端，起到抵消输入信号的作用
\item 将放大器输出信号的一部分回输到放大器的输入端，起到加强输入信号的作用
\item 将放大器输入信号的一部分直通到放大器的输出端，起到抵消输出信号的作用
\item 将放大器输入信号的一部分直通到放大器的输出端，起到加强输出信号的作用
\end{enumerate}
\noindent\textbf{解说：}\textbf{}。\\\noindent\textbf{答案：}

\bigskip




\noindent\textbf{试题：}无线电设备的放大器中大量使用了负反馈技术。与没有反馈的相比，负反馈放大器的特点是：
\begin{enumerate}[leftmargin=3em]
\item 失真小，工作稳定，但增益会有所减小
\item 增益增大，但失真会有所增加
\item 输出信号幅度加大，而需要的输入信号幅度减小
\item 放大倍数增大，输出信号含有的新频率分量增多
\end{enumerate}
\noindent\textbf{解说：}\textbf{}。\\\noindent\textbf{答案：}

\bigskip




\noindent\textbf{试题：}采用锁相环频率合成器的收发信机会在信号中带有附加的相位噪声，主要原因是：
\begin{enumerate}[leftmargin=3em]
\item 从检测到本振频率漂移并反馈纠正需要捕捉时间，该延迟造成本振信号相位抖动
\item 频率合成器的数字电路中在信号跳变时产生谐波，造成噪声
\item 频率合成器的振荡器产生非理想正弦波，带有噪声频率分量
\item 锁相环频率合成器电路复杂，大量元器件的热噪声叠加成相位噪声
\end{enumerate}
\noindent\textbf{解说：}\textbf{}。\\\noindent\textbf{答案：}

\bigskip




\noindent\textbf{试题：}很多现代业余无线电收发信机的本机振荡电路采用了直接数字频率合成（DDS）方式。它的主要特点是：
\begin{enumerate}[leftmargin=3em]
\item 电路结构简洁，无锁相捕捉范围限制，不产生相位噪声，跳换频率快
\item 采用同样的频率源振荡器时频率稳定度优于锁相环频率合成方式
\item 与锁相环频率合成方式相比，可以使用速度较低的数字元器件
\item 直接产生纯净的正弦波信号，不需要采用任何滤波器
\end{enumerate}
\noindent\textbf{解说：}\textbf{}。\\\noindent\textbf{答案：}

\bigskip




\noindent\textbf{试题：}电路分析经常用到理想运算放大器的概念。理想运算放大器是指：
\begin{enumerate}[leftmargin=3em]
\item 开环放大倍数无穷大的放大器
\item 内部噪声等于零的放大器
\item 输出电压幅度没有限制的放大器
\item 没有电源消耗的放大器
\end{enumerate}
\noindent\textbf{解说：}\textbf{}。\\\noindent\textbf{答案：}

\bigskip




\noindent\textbf{试题：}小信号放大器是指：
\begin{enumerate}[leftmargin=3em]
\item 工作点始终处于线性工作区的放大器
\item 输出信号电压幅度始终小于10伏的放大器
\item 输入信号电压幅度始终小于10毫伏的放大器
\item 增益不超过30dB的放大器
\end{enumerate}
\noindent\textbf{解说：}\textbf{}。\\\noindent\textbf{答案：}

\bigskip




\noindent\textbf{试题：}大信号放大器是指：
\begin{enumerate}[leftmargin=3em]
\item 工作点可以超出线性工作区的放大器
\item 输出信号电压最大幅度大于10伏的放大器
\item 输入信号电压最大幅度大于10毫伏的放大器
\item 增益超过30dB的放大器
\end{enumerate}
\noindent\textbf{解说：}\textbf{}。\\\noindent\textbf{答案：}

\bigskip




\noindent\textbf{试题：}半导体三极管放大器的输入阻抗由低到高的排列为：
\begin{enumerate}[leftmargin=3em]
\item 共基极放大器、共发射极放大器、共集电极放大器
\item 共发射极放大器、共集电极放大器、共基极放大器
\item 共集电极放大器、共发射极放大器、共基极放大器
\item 共栅极放大器、共源极放大器、共漏极放大器
\end{enumerate}
\noindent\textbf{解说：}\textbf{}。\\\noindent\textbf{答案：}

\bigskip




\noindent\textbf{试题：}半导体三极管放大器的电流放大倍数由低到高的排列为：
\begin{enumerate}[leftmargin=3em]
\item 共基极放大器、共发射极放大器、共集电极放大器
\item 共集电极放大器、共发射极放大器、共基极放大器
\item 共发射极放大器、共集电极放大器、共基极放大器
\item 共源极放大器、共漏极放大器、共栅极放大器
\end{enumerate}
\noindent\textbf{解说：}\textbf{}。\\\noindent\textbf{答案：}

\bigskip




\noindent\textbf{试题：}工作在HF频段的电子管射频线性功率放大器多采用共栅极放大电路，这是因为共栅极放大电路： 
\begin{enumerate}[leftmargin=3em]
\item 输入阻抗低，栅极接地进一步阻断输入输出的极间电容耦合，因此工作稳定，不易自激 
\item 共栅极放大电路的电压放大倍数最大
\item 共栅极放大电路的非线性失真最小
\item 共栅极放大电路的输入阻抗最高，同样的信号电流可得到较大的输入电压
\end{enumerate}
\noindent\textbf{解说：}\textbf{}。\\\noindent\textbf{答案：}

\bigskip




\noindent\textbf{试题：}在下面的三极管放大电路中，Rc的常用名称和作用是：
\begin{enumerate}[leftmargin=3em]
\item 负载电阻，将集电极信号电流转换为信号电压
\item 发射极负反馈电阻，稳定直流工作点
\item 偏流电阻，为晶体管提供偏置电流以工作在适当工作点
\item 集电极负反馈电阻，提供负反馈，减少信号失真
\end{enumerate}
\noindent\textbf{解说：}\textbf{}。\\\noindent\textbf{答案：}

\bigskip

LK0740.jpg


\noindent\textbf{试题：}在下面的三极管放大电路中，Re的常用名称和作用是：
\begin{enumerate}[leftmargin=3em]
\item 发射极负反馈电阻，稳定直流工作点
\item 负载电阻，将集电极信号电流转换为信号电压
\item 负载电阻，将发射极信号电流转换为信号电压
\item 偏流电阻，为晶体管提供偏置电流以工作在适当工作点
\end{enumerate}
\noindent\textbf{解说：}\textbf{}。\\\noindent\textbf{答案：}

\bigskip

LK0741.jpg


\noindent\textbf{试题：}在下面的三极管放大电路中，Rb1的常用名称和作用是：
\begin{enumerate}[leftmargin=3em]
\item 偏流电阻，为晶体管提供偏置电流以工作在适当工作点
\item 负载电阻，将输入信号电流转换为信号电压
\item 发射极负反馈电阻，稳定直流工作点
\item 基极负反馈电阻，提供负反馈，减少信号失真
\end{enumerate}
\noindent\textbf{解说：}\textbf{}。\\\noindent\textbf{答案：}

\bigskip

LK0742.jpg


\noindent\textbf{试题：}在下面的三极管放大电路中，Cb的常用名称和作用是：
\begin{enumerate}[leftmargin=3em]
\item 输入耦合（或隔直流）电容，使输入交流信号电流进入放大器但阻隔直流电流分量出入
\item 输出耦合（或隔直流）电容，使输出交流信号顺利输出但阻隔直流分量电流出入
\item 旁路电容，使发射极交流信号电流分量顺利通过，不致在Re造成压降形成负反馈
\item 负反馈电容，提供负反馈，减少信号失真
\end{enumerate}
\noindent\textbf{解说：}\textbf{}。\\\noindent\textbf{答案：}

\bigskip

LK0743.jpg


\noindent\textbf{试题：}在下面的三极管放大电路中，Ce的常用名称和作用是：
\begin{enumerate}[leftmargin=3em]
\item 旁路电容，使发射极交流信号电流分量顺利通过，不致在Re造成压降形成负反馈
\item 输入耦合（或隔直流）电容，使输入交流信号电流进入放大器但阻隔直流电流分量出入
\item 输出耦合（或隔直流）电容，使输出交流信号顺利输出但阻隔直流分量电流出入
\item 负反馈电容，提供负反馈，减少信号失真
\end{enumerate}
\noindent\textbf{解说：}\textbf{}。\\\noindent\textbf{答案：}

\bigskip

LK0744.jpg


\noindent\textbf{试题：}在下面的三极管放大电路中，Cc的常用名称和作用是：
\begin{enumerate}[leftmargin=3em]
\item 输出耦合（或隔直流）电容，使输出交流信号顺利输出但阻隔直流分量电流出入
\item 输入耦合（或隔直流）电容，使输入交流信号电流进入放大器但阻隔直流电流分量出入
\item 旁路电容，使发射极交流信号电流分量顺利通过，不致在Re造成压降形成负反馈
\item 负反馈电容，提供负反馈，减少信号失真
\end{enumerate}
\noindent\textbf{解说：}\textbf{}。\\\noindent\textbf{答案：}

\bigskip

LK0745.jpg


\noindent\textbf{试题：}在无线电技术中，通常把放大器分为A、B、C、D等类别，这种分类是依据：
\begin{enumerate}[leftmargin=3em]
\item 放大器件的工作点所处的范围
\item 放大器件的质量等级
\item 放大器件的最高工作频率
\item 放大器件的最大输出功率
\end{enumerate}
\noindent\textbf{解说：}\textbf{}。\\\noindent\textbf{答案：}

\bigskip




\noindent\textbf{试题：}根据放大器的工作状态，通常把放大器分为A、B、C、D等类别。A类放大器是指：
\begin{enumerate}[leftmargin=3em]
\item 放大器件在整个信号周期内始终工作在线性区的放大器
\item 放大器件在半个信号周期内工作点处于线性区、另半个信号周期内处于截止区的放大器
\item 放大器件在半个信号周期内处于截止区，另半个周期的部分时间候处于饱和区的放大器
\item 放大器件在半个信号周期内处于截止区，另半个周期处于饱和区的放大器
\end{enumerate}
\noindent\textbf{解说：}\textbf{}。\\\noindent\textbf{答案：}

\bigskip




\noindent\textbf{试题：}根据放大器的工作状态，通常把放大器分为A、B、C、D等类别。B类放大器是指：
\begin{enumerate}[leftmargin=3em]
\item 放大器件在半个信号周期内工作点处于线性区、另半个信号周期内处于截止区的放大器
\item 放大器件在整个信号周期内始终工作在线性区的放大器
\item 放大器件在半个信号周期内处于截止区，另半个周期的部分时间候处于饱和区的放大器
\item 放大器件在半个信号周期内处于截止区，另半个周期处于饱和区的放大器
\end{enumerate}
\noindent\textbf{解说：}\textbf{}。\\\noindent\textbf{答案：}

\bigskip




\noindent\textbf{试题：}根据放大器的工作状态，通常把放大器分为A、B、C、D等类别。C类放大器是指：
\begin{enumerate}[leftmargin=3em]
\item 放大器件在多于半个信号周期的时间内处于截止区，另半个周期的部分时间候处于线性区的放大器
\item 放大器件在整个信号周期内始终工作在线性区的放大器
\item 放大器件在半个信号周期内工作点处于线性区、另半个信号周期内处于截止区的放大器
\item 放大器件在半个信号周期内处于截止区，另半个周期处于饱和区的放大器
\end{enumerate}
\noindent\textbf{解说：}\textbf{}。\\\noindent\textbf{答案：}

\bigskip




\noindent\textbf{试题：}根据放大器的工作状态，通常把放大器分为A、B、C、D等类别。D类放大器是指：
\begin{enumerate}[leftmargin=3em]
\item 放大器件在半个信号周期内处于截止区，另半个周期处于饱和区的放大器
\item 放大器件在整个信号周期内始终工作在线性区的放大器
\item 放大器件在半个信号周期内工作点处于线性区、另半个信号周期内处于截止区的放大器
\item 放大器件在半个信号周期内处于截止区，另半个周期的部分时间候处于饱和区的放大器
\end{enumerate}
\noindent\textbf{解说：}\textbf{}。\\\noindent\textbf{答案：}

\bigskip




\noindent\textbf{试题：}A、B、C、D四类放大器按输出波形失真由小到大的排列顺序是：
\begin{enumerate}[leftmargin=3em]
\item A、B、C、D
\item D、C、B、A
\item A、C、B、D
\item B、A、D、C
\end{enumerate}
\noindent\textbf{解说：}\textbf{}。\\\noindent\textbf{答案：}

\bigskip




\noindent\textbf{试题：}A、B、C、D四类放大器用作射频功率放大时，按电源效率由高到低的排列顺序是：
\begin{enumerate}[leftmargin=3em]
\item D、C、B、A
\item A、B、C、D
\item D、C、A、B
\item B、A、D、C
\end{enumerate}
\noindent\textbf{解说：}\textbf{}。\\\noindent\textbf{答案：}

\bigskip




\noindent\textbf{试题：}A、B、C、D四类放大器中，可以单独用作射频线性功率放大器、但不能单独用作音频线性功率放大器的是：
\begin{enumerate}[leftmargin=3em]
\item B
\item A
\item D
\item C
\end{enumerate}
\noindent\textbf{解说：}\textbf{}。\\\noindent\textbf{答案：}

\bigskip




\noindent\textbf{试题：}A、B、C、D四类放大器中，既可以单独用作射频线性功率放大器、又能单独用作音频线性功率放大器的是： （不考虑电源效率）
\begin{enumerate}[leftmargin=3em]
\item A
\item B
\item C
\item D
\end{enumerate}
\noindent\textbf{解说：}\textbf{}。\\\noindent\textbf{答案：}

\bigskip




\noindent\textbf{试题：}A、B、C、D四类放大器中，可以用作FM射频功率放大器的全部类别有： （不考虑电源效率）
\begin{enumerate}[leftmargin=3em]
\item A、B、C、D
\item A、B、C
\item A、B、D
\item A
\end{enumerate}
\noindent\textbf{解说：}\textbf{}。\\\noindent\textbf{答案：}

\bigskip




\noindent\textbf{试题：}A、B、C、D四类放大器中，可以用作CW射频功率放大器的全部类别有： （不考虑电源效率）
\begin{enumerate}[leftmargin=3em]
\item A、B、C、D
\item A、B、C
\item A、B、D
\item A
\end{enumerate}
\noindent\textbf{解说：}\textbf{}。\\\noindent\textbf{答案：}

\bigskip




\noindent\textbf{试题：}A、B、C、D四类放大器中，适宜于做小信号放大器的是：
\begin{enumerate}[leftmargin=3em]
\item A
\item B
\item C
\item D
\end{enumerate}
\noindent\textbf{解说：}\textbf{}。\\\noindent\textbf{答案：}

\bigskip




\noindent\textbf{试题：}A、B、C、D四类放大器中，属于大信号放大器的全部类别有：
\begin{enumerate}[leftmargin=3em]
\item B、C、D
\item A、B、C、D
\item A、C、D
\item C、D
\end{enumerate}
\noindent\textbf{解说：}\textbf{}。\\\noindent\textbf{答案：}

\bigskip




\noindent\textbf{试题：}很多业余电台的末级和末前级射频输出放大器中采用两个并联的输出半导体功率管，这是为了：
\begin{enumerate}[leftmargin=3em]
\item 双管并联，得到双倍的输出电流和输出功率
\item 构成推挽电路，减小输出波形的失真
\item 双管并联，得到双倍的器件耐压，减少损坏几率
\item 双管并联，使每个功率管的失真互相补偿，减少失真，降低杂散发射
\end{enumerate}
\noindent\textbf{解说：}\textbf{}。\\\noindent\textbf{答案：}

\bigskip




\noindent\textbf{试题：}很多现代无线电设备的音频功率放大电路采用两个串联的输出半导体功率管，分别负责信号正、负半周的放大。这种电路的通用名称和作用是：
\begin{enumerate}[leftmargin=3em]
\item 推挽放大电路，实现极小静态工作点下的高电源效率的线性功率放大
\item 双管串联电路，得到双倍的输出电流和输出功率
\item 双管串联电路，得到较高的输出阻抗以改善与负载的阻抗匹配
\item 双管串联电路，得到较高的输入阻抗以改善与推动级之间的阻抗匹配
\end{enumerate}
\noindent\textbf{解说：}\textbf{}。\\\noindent\textbf{答案：}

\bigskip




\noindent\textbf{试题：}在半导体三极管的基本放大电路中，电压增益最大的是：
\begin{enumerate}[leftmargin=3em]
\item 共发射极放大电路
\item 共基极放大电路
\item 共集电极放大电路
\item 带负反馈的共发射极放大电路
\end{enumerate}
\noindent\textbf{解说：}\textbf{}。\\\noindent\textbf{答案：}

\bigskip




\noindent\textbf{试题：}在半导体三极管的基本放大电路中，输入阻抗最高的是：
\begin{enumerate}[leftmargin=3em]
\item 共集电极放大电路
\item 共发射极放大电路
\item 共基极放大电路
\item 带正反馈的共发射极放大电路
\end{enumerate}
\noindent\textbf{解说：}\textbf{}。\\\noindent\textbf{答案：}

\bigskip




\noindent\textbf{试题：}在半导体三极管的基本放大电路中，输出阻抗最低的是：
\begin{enumerate}[leftmargin=3em]
\item 共集电极放大电路
\item 共发射极放大电路
\item 共基极放大电路
\item 带正反馈的共发射极放大电路
\end{enumerate}
\noindent\textbf{解说：}\textbf{}。\\\noindent\textbf{答案：}

\bigskip




\noindent\textbf{试题：}下列放大器电路中，电源效率最低的是：
\begin{enumerate}[leftmargin=3em]
\item 甲类放大器
\item 乙类放大器
\item 丙类放大器
\item 甲乙类放大器
\end{enumerate}
\noindent\textbf{解说：}\textbf{}。\\\noindent\textbf{答案：}

\bigskip




\noindent\textbf{试题：}下列放大器电路中，电源效率最高的是：
\begin{enumerate}[leftmargin=3em]
\item 丙类放大器
\item 甲类放大器
\item 乙类放大器
\item 甲乙类放大器
\end{enumerate}
\noindent\textbf{解说：}\textbf{}。\\\noindent\textbf{答案：}

\bigskip




\noindent\textbf{试题：}下列放大器电路中，波形失真最小的是：
\begin{enumerate}[leftmargin=3em]
\item 甲类放大器
\item 乙类放大器
\item 丙类放大器
\item 甲乙类放大器
\end{enumerate}
\noindent\textbf{解说：}\textbf{}。\\\noindent\textbf{答案：}

\bigskip




\noindent\textbf{试题：}下列放大器电路中，波形失真最大的是：
\begin{enumerate}[leftmargin=3em]
\item 丙类放大器
\item 甲类放大器
\item 乙类放大器
\item 甲乙类放大器
\end{enumerate}
\noindent\textbf{解说：}\textbf{}。\\\noindent\textbf{答案：}

\bigskip




\noindent\textbf{试题：}下列放大器电路中，最适宜用作音频大功率放大器的是：
\begin{enumerate}[leftmargin=3em]
\item 甲乙类放大器
\item 丙类放大器
\item 甲类放大器
\item 乙类放大器
\end{enumerate}
\noindent\textbf{解说：}\textbf{}。\\\noindent\textbf{答案：}

\bigskip




\noindent\textbf{试题：}下列放大器电路中，最适宜用作音频小信号放大器的是：
\begin{enumerate}[leftmargin=3em]
\item 甲类放大器
\item 甲乙类放大器
\item 丙类放大器
\item 乙类放大器
\end{enumerate}
\noindent\textbf{解说：}\textbf{}。\\\noindent\textbf{答案：}

\bigskip




\noindent\textbf{试题：}下列放大器电路中，最适宜用作CW射频功率放大器的是：
\begin{enumerate}[leftmargin=3em]
\item 丙类放大器
\item 甲类放大器
\item 甲乙类放大器
\item 乙类放大器
\end{enumerate}
\noindent\textbf{解说：}\textbf{}。\\\noindent\textbf{答案：}

\bigskip




\noindent\textbf{试题：}下列放大器电路中，最适宜用作FM射频功率放大器的是：
\begin{enumerate}[leftmargin=3em]
\item 丙类放大器
\item 甲类放大器
\item 甲乙类放大器
\item 乙类放大器
\end{enumerate}
\noindent\textbf{解说：}\textbf{}。\\\noindent\textbf{答案：}

\bigskip




\noindent\textbf{试题：}下列放大器电路中，最适宜用作SSB射频功率放大器的是：
\begin{enumerate}[leftmargin=3em]
\item 接近乙类的甲乙类放大器
\item 接近甲类的甲乙类放大器
\item 丙类放大器
\item 甲类放大器
\end{enumerate}
\noindent\textbf{解说：}\textbf{}。\\\noindent\textbf{答案：}

\bigskip




\noindent\textbf{试题：}频移电报技术（frequency-shift telegraphy）是指：电报信号控制载波频率在预定的范围之内变化的调频电报技术。下述业余通信使用的是移频电报技术：
\begin{enumerate}[leftmargin=3em]
\item RTTY
\item CW
\item PSK31
\item SSTV
\end{enumerate}
\noindent\textbf{解说：}\textbf{}。\\\noindent\textbf{答案：}

\bigskip




\noindent\textbf{试题：}无线电通信选择不同调制方式的主要考虑因素是：
\begin{enumerate}[leftmargin=3em]
\item 信息在传递过程中的保真度、信号的抗干扰能力、尽量节省无线电频谱资源
\item 改善天线阻抗匹配、尽量提高发射频率稳定度、尽量减少杂散发射
\item 有利于提高接收机选择性、尽量提高话筒灵敏度、防止产生谐波干扰
\item 防止与附近发射机产生三阶互调、尽量降低本振相位噪声、采用高中频方案
\end{enumerate}
\noindent\textbf{解说：}\textbf{}。\\\noindent\textbf{答案：}

\bigskip




\noindent\textbf{试题：}接收机解调部件的作用是：
\begin{enumerate}[leftmargin=3em]
\item 从接收到的已调制射频信号中分离出原始信号
\item 对接收到的射频信号进行宽带线性放大
\item 对接收到的射频信号进行选频放大
\item 从接收到的已调制射频信号中提取出载频分量
\end{enumerate}
\noindent\textbf{解说：}\textbf{}。\\\noindent\textbf{答案：}

\bigskip




\noindent\textbf{试题：}选择解调部件的主要应考因素是：
\begin{enumerate}[leftmargin=3em]
\item 尽量忠实地还原原始信号
\item 尽量对已调制信号加以放大
\item 尽量提升已调制信号中的载频分量
\item 尽量补偿接收到的射频信号的频率偏移
\end{enumerate}
\noindent\textbf{解说：}\textbf{}。\\\noindent\textbf{答案：}

\bigskip




\noindent\textbf{试题：}业余电台发射单边带语音信号时，附近连接在某电路上的耳机传出了模糊不清的语音。可以断定这个电路具有下列功能：
\begin{enumerate}[leftmargin=3em]
\item 检波
\item 差拍
\item 变频
\item 振荡
\end{enumerate}
\noindent\textbf{解说：}\textbf{}。\\\noindent\textbf{答案：}

\bigskip




\noindent\textbf{试题：}业余电台发射等幅波电码信号时，附近连接在某电路上的耳机传出了“嘟嘟”的音频电码声。可以断定这个电路具有下列功能：
\begin{enumerate}[leftmargin=3em]
\item 差拍
\item 检波
\item 变频
\item 振荡
\end{enumerate}
\noindent\textbf{解说：}\textbf{}。\\\noindent\textbf{答案：}

\bigskip




\noindent\textbf{试题：}试验某电路，接通电源时，旁边的无线电收信机可以收到它发出的类似鸟叫的声音，断开电源，鸟叫声消失。可以断定这个电路具有下列功能：
\begin{enumerate}[leftmargin=3em]
\item 振荡
\item 差拍
\item 检波
\item 变频
\end{enumerate}
\noindent\textbf{解说：}\textbf{}。\\\noindent\textbf{答案：}

\bigskip




\noindent\textbf{试题：}用收信机的AM挡接收中等强度的SSB语音信号，调谐到最佳状态，听到的将是：
\begin{enumerate}[leftmargin=3em]
\item 模糊不清的类似语音的声音
\item 正常的语音
\item 听不到任何语音
\item 差拍叫声
\end{enumerate}
\noindent\textbf{解说：}\textbf{}。\\\noindent\textbf{答案：}

\bigskip




\noindent\textbf{试题：}用收信机的AM挡接收中等强度的FM语音信号，调谐到最佳状态，听到的将是：
\begin{enumerate}[leftmargin=3em]
\item 听不到语音，信号越强越寂静
\item 正常的语音
\item 差拍叫声
\item 模糊不清的类似语音的声音
\end{enumerate}
\noindent\textbf{解说：}\textbf{}。\\\noindent\textbf{答案：}

\bigskip




\noindent\textbf{试题：}用收信机的SSB挡接收中等强度的FM语音信号，调谐到最佳状态，听到的将是：
\begin{enumerate}[leftmargin=3em]
\item 差拍叫声
\item 模糊不清的类似语音的声音
\item 正常的语音
\item 听不到任何语音
\end{enumerate}
\noindent\textbf{解说：}\textbf{}。\\\noindent\textbf{答案：}

\bigskip




\noindent\textbf{试题：}用收信机的SSB挡接收中等强度的AM语音信号，调谐到最佳状态，听到的将是：
\begin{enumerate}[leftmargin=3em]
\item 正常的语音
\item 差拍叫声
\item 模糊不清的类似语音的声音
\item 听不到任何语音
\end{enumerate}
\noindent\textbf{解说：}\textbf{}。\\\noindent\textbf{答案：}

\bigskip




\noindent\textbf{试题：}用收信机的FM挡接收中等强度的AM语音信号，调谐到最佳状态，听到的将是：
\begin{enumerate}[leftmargin=3em]
\item 听不到语音，有信号时噪声消失，但听不到语音
\item 差拍叫声
\item 模糊不清的类似语音的声音
\item 听不到任何语音
\end{enumerate}
\noindent\textbf{解说：}\textbf{}。\\\noindent\textbf{答案：}

\bigskip




\noindent\textbf{试题：}用收信机的FM挡接收中等强度的SSB语音信号，调谐到最佳状态，听到的将是：
\begin{enumerate}[leftmargin=3em]
\item 听不到语音，但噪声可能会有些变化
\item 差拍叫声
\item 模糊不清的类似语音的声音
\item 听不到任何语音
\end{enumerate}
\noindent\textbf{解说：}\textbf{}。\\\noindent\textbf{答案：}

\bigskip




\noindent\textbf{试题：}在HF业余频段的数据通信段中，用收信机的SSB挡听到一个由两种音调交替组成的信号，这个信号的调制方式最可能属于下述种类：
\begin{enumerate}[leftmargin=3em]
\item FSK
\item PSK
\item ASK
\item SSTV或FAX
\end{enumerate}
\noindent\textbf{解说：}\textbf{}。\\\noindent\textbf{答案：}

\bigskip




\noindent\textbf{试题：}在HF业余频段的数据通信段中，用收信机的SSB挡听到一个音调不变但又似乎不断颤动的信号，这个信号的调制方式最可能属于下述种类：
\begin{enumerate}[leftmargin=3em]
\item PSK
\item FSK
\item ASK
\item SSTV或FAX
\end{enumerate}
\noindent\textbf{解说：}\textbf{}。\\\noindent\textbf{答案：}

\bigskip




\noindent\textbf{试题：}用收信机的SSB挡在业余频段中，听到一个音调大致以约为几分之一秒的重复周期连续变化、并夹有一种规律的“笃、笃”声的信号，。这个信号的调制方式最可能属于下述种类：
\begin{enumerate}[leftmargin=3em]
\item SSTV或FAX
\item PSK
\item FSK
\item ASK
\end{enumerate}
\noindent\textbf{解说：}\textbf{}。\\\noindent\textbf{答案：}

\bigskip




\noindent\textbf{试题：}一个语音信号，经过USB发射机调制成为一个射频信号。从频率域的观点看，上边带调制的过程是发生了：
\begin{enumerate}[leftmargin=3em]
\item 频谱的简单平移
\item 频谱各分量相对幅度的有规律变化
\item 频谱的反转（倒置）
\item 频谱分量完全打乱后重新组合
\end{enumerate}
\noindent\textbf{解说：}\textbf{}。\\\noindent\textbf{答案：}

\bigskip




\noindent\textbf{试题：}对于给定的SSB发射设备，决定其输出信号实际占用带宽的因素是：
\begin{enumerate}[leftmargin=3em]
\item 所传输信号的带宽越宽，射频输出占用带宽越宽，但与其幅度和最高频率无关
\item 所传输信号的最高频率越高，射频输出占用带宽越宽，但与其幅度和带宽无关
\item 所传输信号的幅度越大，射频输出占用带宽越宽，但与其频率和带宽无关
\item 射频输出实际占用带宽为由电路决定的固定值，通信常用的是2.7kHz
\end{enumerate}
\noindent\textbf{解说：}\textbf{}。\\\noindent\textbf{答案：}

\bigskip




\noindent\textbf{试题：}什么叫做“检波”？
\begin{enumerate}[leftmargin=3em]
\item 对调幅信号进行解调的过程称为检波
\item 检查信号波形是否超过允许的幅度范围的过程称为检波
\item 检查信号频率是否发生了不应有的偏离或者漂移过程称为检波
\item 对调频信号进行解调的过程称为检波
\end{enumerate}
\noindent\textbf{解说：}\textbf{}。\\\noindent\textbf{答案：}

\bigskip




\noindent\textbf{试题：}利用个人电脑和用于传送语音的普通收发信机进行业余短波SSTV通信所用到的调制方式为：
\begin{enumerate}[leftmargin=3em]
\item FM（图像信息的音频基带信号调制）和SSB（频率搬移）
\item FSK（图像信息的音频基带信号调制）和SSB（频率搬移）
\item PSK（图像信息的音频基带信号调制）和SSB（频率搬移）
\item 仅使用到SSB调制
\end{enumerate}
\noindent\textbf{解说：}\textbf{}。\\\noindent\textbf{答案：}

\bigskip




\noindent\textbf{试题：}利用个人电脑和用于传送语音的普通收发信机进行业余短波RTTY通信所用到的调制方式有：
\begin{enumerate}[leftmargin=3em]
\item FSK（数字信息的音频基带信号调制）和SSB（频率搬移）
\item AM（数字信息的音频基带信号调制）和SSB（频率搬移）
\item PSK（数字信息的音频基带信号调制）和SSB（频率搬移）
\item 仅使用到SSB调制
\end{enumerate}
\noindent\textbf{解说：}\textbf{}。\\\noindent\textbf{答案：}

\bigskip




\noindent\textbf{试题：}利用个人电脑和用于传送语音的普通收发信机进行业余短波PSK31通信所用到的调制方式有：
\begin{enumerate}[leftmargin=3em]
\item PSK（数字信息的音频基带信号调制）和SSB（频率搬移）
\item FSK（数字信息的音频基带信号调制）和SSB（频率搬移）
\item ASK（数字信息的音频基带信号调制）和SSB（频率搬移）
\item 仅使用到SSB调制
\end{enumerate}
\noindent\textbf{解说：}\textbf{}。\\\noindent\textbf{答案：}

\bigskip




\noindent\textbf{试题：}利用PC电脑和用于传送语音的普通收发信机进行业余短波PACKET通信所用到的调制方式有：
\begin{enumerate}[leftmargin=3em]
\item FSK（数字信息的音频基带信号调制）和SSB（频率搬移）
\item AM（数字信息的音频基带信号调制）和SSB（频率搬移）
\item PSK（数字信息的音频基带信号调制）和SSB（频率搬移）
\item 仅使用到SSB调制
\end{enumerate}
\noindent\textbf{解说：}\textbf{}。\\\noindent\textbf{答案：}

\bigskip




\noindent\textbf{试题：}业余SSTV通信和有些模拟ATV采用调频方式而不是广播电视图像的调幅方式，主要原因是： 
\begin{enumerate}[leftmargin=3em]
\item 业余电台信号较弱，调频解调可以更好地抗拒叠加在信号上的外界噪声所引起的幅度变化
\item 调频方式占用的频带比调幅方式窄
\item 便于与调频通话方式所使用的设备兼容
\item 调频收信设备的灵敏度比调幅的高一个数量级
\end{enumerate}
\noindent\textbf{解说：}\textbf{}。\\\noindent\textbf{答案：}

\bigskip




\noindent\textbf{试题：}模拟ATV和模拟广播电视通信中，图像信号调制方式的差别是：
\begin{enumerate}[leftmargin=3em]
\item ATV为FM或AM，广播电视为AM
\item ATV 为AM，广播电视为FM
\item ATV 为PSK，广播电视为ASK
\item ATV为标清方式，广播电视为高清方式
\end{enumerate}
\noindent\textbf{解说：}\textbf{}。\\\noindent\textbf{答案：}

\bigskip




\noindent\textbf{试题：}如果发射机在不同工作模式时最大射频输出功率相同，无语音调制时，实际射频输出由大到小的排序为：
\begin{enumerate}[leftmargin=3em]
\item FM，AM，SSB
\item AM，SSB，FM
\item SSB，FM，AM
\item SSB，AM，FM
\end{enumerate}
\noindent\textbf{解说：}\textbf{}。\\\noindent\textbf{答案：}

\bigskip




\noindent\textbf{试题：}下列调制得到的信号幅度恒定不变：
\begin{enumerate}[leftmargin=3em]
\item 移频键控FSK
\item 幅度调制AM
\item 单边带幅度调制SSB
\item 幅度键控调制ASK
\end{enumerate}
\noindent\textbf{解说：}\textbf{}。\\\noindent\textbf{答案：}

\bigskip




\noindent\textbf{试题：}下列调制得到的信号幅度恒定不变：
\begin{enumerate}[leftmargin=3em]
\item 频率调制FM
\item 幅度调制AM
\item 单边带幅度调制SSB
\item 幅度键控调制ASK
\end{enumerate}
\noindent\textbf{解说：}\textbf{}。\\\noindent\textbf{答案：}

\bigskip




\noindent\textbf{试题：}下列调制得到的信号幅度恒定不变：
\begin{enumerate}[leftmargin=3em]
\item 相位调制PM
\item 幅度调制AM
\item 单边带幅度调制SSB
\item 幅度键控调制ASK
\end{enumerate}
\noindent\textbf{解说：}\textbf{}。\\\noindent\textbf{答案：}

\bigskip




\noindent\textbf{试题：}下列调制得到的信号幅度恒定不变：
\begin{enumerate}[leftmargin=3em]
\item 移相键控调制PSK
\item 幅度调制AM
\item 单边带幅度调制SSB
\item 幅度键控调制ASK
\end{enumerate}
\noindent\textbf{解说：}\textbf{}。\\\noindent\textbf{答案：}

\bigskip




\noindent\textbf{试题：}下列调制得到的信号周期恒定不变：
\begin{enumerate}[leftmargin=3em]
\item 相位调制PM
\item 频率调制FM
\item 单边带幅度调制SSB
\item 频率键控调制FSK
\end{enumerate}
\noindent\textbf{解说：}\textbf{}。\\\noindent\textbf{答案：}

\bigskip




\noindent\textbf{试题：}下列调制得到的信号中载频分量幅度恒定不变：
\begin{enumerate}[leftmargin=3em]
\item 幅度调制AM
\item 频率调制FM
\item 载波抑制单边带幅度调制SSB
\item 频率键控调制FSK
\end{enumerate}
\noindent\textbf{解说：}\textbf{}。\\\noindent\textbf{答案：}

\bigskip




\noindent\textbf{试题：}信号经过下列调制后，频带宽度可能会大于原有值：
\begin{enumerate}[leftmargin=3em]
\item 频率调制FM
\item 幅度调制AM
\item 单边带幅度调制SSB
\item 幅度键控调制ASK
\end{enumerate}
\noindent\textbf{解说：}\textbf{}。\\\noindent\textbf{答案：}

\bigskip




\noindent\textbf{试题：}乘积检波器的作用是什么？
\begin{enumerate}[leftmargin=3em]
\item 用来给CW或SSB信号检波
\item 解调FM信号
\item 用来给调相信号检波
\item 将语音信号从背景噪声中分离开来
\end{enumerate}
\noindent\textbf{解说：}\textbf{}。\\\noindent\textbf{答案：}

\bigskip




\noindent\textbf{试题：}下列哪一项可以解调FM信号？
\begin{enumerate}[leftmargin=3em]
\item 鉴频器
\item 限幅器
\item 乘积检波器
\item 反相器
\end{enumerate}
\noindent\textbf{解说：}\textbf{}。\\\noindent\textbf{答案：}

\bigskip




\noindent\textbf{试题：}下列哪一种可以归入调幅通信？
\begin{enumerate}[leftmargin=3em]
\item 单边带通信
\item 封包通信
\item 扩频通信
\item 相移键控
\end{enumerate}
\noindent\textbf{解说：}\textbf{}。\\\noindent\textbf{答案：}

\bigskip




\noindent\textbf{试题：}在VHF封包通信中，哪一种调制模式被广泛使用？
\begin{enumerate}[leftmargin=3em]
\item FM
\item SSB
\item AM
\item 扩频通信
\end{enumerate}
\noindent\textbf{解说：}\textbf{}。\\\noindent\textbf{答案：}

\bigskip




\noindent\textbf{试题：}什么是PSK31？
\begin{enumerate}[leftmargin=3em]
\item 一种低速率的数据通信模式
\item 一种可以减少噪音对FM信号干扰的方法
\item 一种压缩数字电视信号的方法
\item 一种高速率的数据通信模式
\end{enumerate}
\noindent\textbf{解说：}\textbf{}。\\\noindent\textbf{答案：}

\bigskip




\noindent\textbf{试题：}220V.AC/13.8V.DC通信开关电源的一般工作过程是：
\begin{enumerate}[leftmargin=3em]
\item 将交流输入整流滤波为高压直流，由半导体开关电路变成高压脉冲电流，由变压器变成低压脉冲，整流滤波为低压直流
\item 将交流输入整流滤波为高压直流，由变压器变成低压脉冲，由半导体开关电路变成低压直流，滤波后输出
\item 由变压器将交流输入变为低压交流，由半导体开关电路变成超音频脉冲电流，经整流滤波为低压直流
\item 由大功率半导体三极管将交流输入变为高压直流，由专用集成电路变成超音频脉冲电流，经整流滤波为低压直流
\end{enumerate}
\noindent\textbf{解说：}\textbf{}。\\\noindent\textbf{答案：}

\bigskip




\noindent\textbf{试题：}现代无线电通信技术大量使用数字处理技术，处理模拟信号时首先要进行模数转换。模数转换器所要求的转换速率取决于：
\begin{enumerate}[leftmargin=3em]
\item 待处理信号的最高频率分量
\item 待处理信号的最低频率分量
\item 待处理信号的带宽占最高频率分量的百分比
\item 待处理信号波形的复杂程度
\end{enumerate}
\noindent\textbf{解说：}\textbf{}。\\\noindent\textbf{答案：}

\bigskip




\noindent\textbf{试题：}现代无线电通信技术大量使用数字处理技术，处理模拟信号时首先要进行模数转换。模数转换器所要求的采样精度位数取决于：
\begin{enumerate}[leftmargin=3em]
\item 待处理信号波形的复杂程度
\item 待处理信号的最高频率分量
\item 待处理信号的最低频率分量
\item 待处理信号的带宽占最高频率分量的百分比
\end{enumerate}
\noindent\textbf{解说：}\textbf{}。\\\noindent\textbf{答案：}

\bigskip




\noindent\textbf{试题：}在模数转换时，要使数字信号能够正确还原出原来的模拟信号，模数转换器的采样速率必须是：
\begin{enumerate}[leftmargin=3em]
\item 大于信号最高频率分量的2倍
\item 等于信号最高频率分量
\item 大于信号最高频率分量的1.2倍
\item 等于信号频率的2倍
\end{enumerate}
\noindent\textbf{解说：}\textbf{}。\\\noindent\textbf{答案：}

\bigskip




\noindent\textbf{试题：}将一个频率为f、电压最大值为U的模拟信号进行转换为数字信号。如要使数字信号能够反映出信号的相当于0.1%U的幅度细节和相当于f的8次谐波的时间细节，则模数转换器应达到如下指标：
\begin{enumerate}[leftmargin=3em]
\item 采样速率大于16f，采样精度不低于10位
\item 采样速率大于8f，采样精度不低于10位
\item 采样速率大于4f，采样精度不低于8位
\item 采样速率大于10f，采样精度不低于9位
\end{enumerate}
\noindent\textbf{解说：}\textbf{}。\\\noindent\textbf{答案：}

\bigskip




\noindent\textbf{试题：}很多现代收发信机都采用DSP技术对信号进行滤波、调制或解调等处理。其基本处理步骤是：
\begin{enumerate}[leftmargin=3em]
\item 模数转换、存储、数字运算、数模转换
\item 数字运算、模数转换、存储、数模转换
\item 数模转换、数字运算、存储、模数转换
\item 存储、数模转换、数字运算、模数转换
\end{enumerate}
\noindent\textbf{解说：}\textbf{}。\\\noindent\textbf{答案：}

\bigskip




\noindent\textbf{试题：}具有两个输入的与门（AND）是最简单的数字逻辑电路之一。如果两个输入信号的组合分别为00、01、10、11，对应的输出信号应为：
\begin{enumerate}[leftmargin=3em]
\item 0、0、0、1
\item 0、1、1、1
\item 0、1、0、1
\item 1、0、1、0
\end{enumerate}
\noindent\textbf{解说：}\textbf{}。\\\noindent\textbf{答案：}

\bigskip




\noindent\textbf{试题：}具有两个输入的或门（OR）是最简单的数字逻辑电路之一。如果两个输入信号组合分别为00、01、10、11，对应的输出信号应为：
\begin{enumerate}[leftmargin=3em]
\item 0、1、1、1
\item 0、1、1、0
\item 0、1、0、1
\item 1、0、1、0
\end{enumerate}
\noindent\textbf{解说：}\textbf{}。\\\noindent\textbf{答案：}

\bigskip




\noindent\textbf{试题：}具有两个输入的异或门（XOR）是最简单的数字逻辑电路之一。如果两个输入信号组合分别为00、01、10、11，对应的输出信号应为：
\begin{enumerate}[leftmargin=3em]
\item 0、1、1、0
\item 0、1、0、1
\item 0、1、1、1
\item 1、0、1、0
\end{enumerate}
\noindent\textbf{解说：}\textbf{}。\\\noindent\textbf{答案：}

\bigskip




\noindent\textbf{试题：}具有两个输入的与非门（NAND）是最简单的数字逻辑电路之一。如果两个输入信号组合分别为00、01、10、11，对应的输出信号应为：
\begin{enumerate}[leftmargin=3em]
\item 1、1、1、0
\item 0、1、1、1
\item 0、1、0、1
\item 1、0、1、0
\end{enumerate}
\noindent\textbf{解说：}\textbf{}。\\\noindent\textbf{答案：}

\bigskip




\noindent\textbf{试题：}具有两个输入的或非门（NOR）是最简单的数字逻辑电路之一。如果两个输入信号组合分别为00、01、10、11，对应的输出信号应为：
\begin{enumerate}[leftmargin=3em]
\item 1、0、0、0
\item 0、1、1、1
\item 0、1、0、1
\item 1、0、1、0
\end{enumerate}
\noindent\textbf{解说：}\textbf{}。\\\noindent\textbf{答案：}

\bigskip




\noindent\textbf{试题：}具有两个输入的异或非门（NXOR）是最简单的数字逻辑电路之一。如果两个输入信号组合分别为00、01、10、11，对应的输出信号应为：
\begin{enumerate}[leftmargin=3em]
\item 1、0、0、1
\item 0、1、1、1
\item 1、1、0、0
\item 0、1、1、0
\end{enumerate}
\noindent\textbf{解说：}\textbf{}。\\\noindent\textbf{答案：}

\bigskip




\noindent\textbf{试题：}数字通信中常用到描述信号流量的“信号速率”和描述信息流量的“数据速率”两个参数。业余RTTY通信方式分别用f0和f1两个不同频率信号来传输数字0和1，最常见的每个码元的时间长度为0.02秒。其信号速率和数据速率分别为：
\begin{enumerate}[leftmargin=3em]
\item 50波特，50位/秒
\item 100波特，100位/秒
\item 50波特，100位/秒
\item 100波特，50位/秒
\end{enumerate}
\noindent\textbf{解说：}\textbf{}。\\\noindent\textbf{答案：}

\bigskip




\noindent\textbf{试题：}数字通信中常用到描述信号流量的“信号速率”和描述信息流量的“数据速率”两个参数。业余BPSK通信方式中，用一个副载频f的两个相反的相位分别传输数字0和1，每个码元的时间长度约为32毫秒。其信号速率和数据速率分别为：
\begin{enumerate}[leftmargin=3em]
\item 31.25波特，31.25位/秒
\item 100波特，100位/秒
\item 50波特，100位/秒
\item 100波特，50位/秒
\end{enumerate}
\noindent\textbf{解说：}\textbf{}。\\\noindent\textbf{答案：}

\bigskip




\noindent\textbf{试题：}数字通信中常用到描述信号流量的“信号速率”和描述信息流量的“数据速率”两个参数。MFK16通信方式中，用总带宽316Hz范围内的16个副载频之一来传输16种状态，信号重复频率为64毫秒。其信号速率和数据速率分别为：
\begin{enumerate}[leftmargin=3em]
\item 15.625波特，62.5位/秒
\item 15.625波特，250位/秒
\item 15.625波特，15.625位/秒
\item 62.5波特，250位/秒
\end{enumerate}
\noindent\textbf{解说：}\textbf{}。\\\noindent\textbf{答案：}

\bigskip




\noindent\textbf{试题：}数字通信中常用到描述信号流量的“信号速率”和描述信息流量的“数据速率”两个参数。QPSK通信方式中，用一个副载频f的四个相差90度的相位分别传输数字00、01、10和11，每个码元的时间长度约为32毫秒。其信号速率和数据速率分别为：
\begin{enumerate}[leftmargin=3em]
\item 31.25波特，62.5位/秒
\item 31.25波特，31.25位/秒
\item 50波特，100位/秒
\item 31.25波特，125位/秒
\end{enumerate}
\noindent\textbf{解说：}\textbf{}。\\\noindent\textbf{答案：}

\bigskip




\noindent\textbf{试题：}在数字通信中，同时传送多位二进制信息的方式属于并行通信。下列业余无线电数字通信方式中，属于并行通信的有：
\begin{enumerate}[leftmargin=3em]
\item MFSK16
\item BPSK
\item RTTY
\item PACKET
\end{enumerate}
\noindent\textbf{解说：}\textbf{}。\\\noindent\textbf{答案：}

\bigskip




\noindent\textbf{试题：}假设以同样的信息传输速率传输报文，MFSK16通信方式的波特率低于信息传输速率，RTTY通信方式的波特率等于信息传输速率。根据这一描述，两者相比可以做以下推论：
\begin{enumerate}[leftmargin=3em]
\item 前者每个信号的宽度比较宽，因此抗御突发干扰的性能好，有利于短波DX通信
\item 后者的实际通信能力可能更强些，任何情况下波特率总是越高越好
\item 既然两者信息传输速率相同，通信效果是完全一样的
\item 这一描述本身有问题，波特率应该总是等于信息传输速率的
\end{enumerate}
\noindent\textbf{解说：}\textbf{}。\\\noindent\textbf{答案：}

\bigskip




\noindent\textbf{试题：}根据数据串行通信收发两端的时钟只需要在传送一个字节的时间内保持同步还是需要在传送一整块数据的时间内保持同步，可以分为“异步”和“同步”两种方式。下列业余无线电数字通信方式中属于异步方式的例子是：
\begin{enumerate}[leftmargin=3em]
\item RTTY
\item QPSK31
\item PACKET
\item PACTOR-II
\end{enumerate}
\noindent\textbf{解说：}\textbf{}。\\\noindent\textbf{答案：}

\bigskip




\noindent\textbf{试题：}根据数据串行通信收发两端的时钟只需要在传送一个字节的时间内保持同步还是需要在传送一整块数据的时间内保持同步，可以分为“异步”和“同步”两种方式。下列业余无线电数字通信方式中不属于同步方式的例子是：
\begin{enumerate}[leftmargin=3em]
\item RTTY
\item PACKET
\item BPSK31
\item AMTOR-FEC
\end{enumerate}
\noindent\textbf{解说：}\textbf{}。\\\noindent\textbf{答案：}

\bigskip




\noindent\textbf{试题：}在进行串行异步数字通信时，双方需要设置相同的波特率，数据位数，校验位数，停止位数。RTTY最常用的设置为：
\begin{enumerate}[leftmargin=3em]
\item 50（或45.45），5，N，1
\item 50（或45.45），8，2，3
\item 31.25，7，170，0.3
\item 2295，2125，170，5
\end{enumerate}
\noindent\textbf{解说：}\textbf{}。\\\noindent\textbf{答案：}

\bigskip




\noindent\textbf{试题：}我国GB2312-80《信息交换用汉字编码字符集-基本集》标准共收录6763个汉字。不考虑误码的情况，如果要以单个汉字为单位，传输由该字符集的汉字汉字组成的总长度为100字的报文，用RTTY和PSK31通信方式至少分别需要发送多少个字符？
\begin{enumerate}[leftmargin=3em]
\item 300,200
\item 100,100
\item 200,200
\item 400,200
\end{enumerate}
\noindent\textbf{解说：}\textbf{}。\\\noindent\textbf{答案：}

\bigskip




\noindent\textbf{试题：}PACKET是业余无线电爱好者利用X.25数据分组通信协议开发的业余无线电通信方式，用于HF频段、VHF频段和卫星通信时通常采用的信号速率分别为：
\begin{enumerate}[leftmargin=3em]
\item 300波特、1200波特、9600波特
\item 1200波特、9600波特、19200波特
\item 300波特、2400波特、19200波特
\item 600波特、1200波特、2400波特
\end{enumerate}
\noindent\textbf{解说：}\textbf{}。\\\noindent\textbf{答案：}

\bigskip




\noindent\textbf{试题：}有些常用的业余无线电数字通信协议只支持5位或者7位数据字节的传送，因而不适合用来传送8位字节的2进制数据文件。常用方式中能够支持2进制数据文件传输的是：
\begin{enumerate}[leftmargin=3em]
\item PACTOR
\item RTTY
\item AMTOR
\item PSK31
\end{enumerate}
\noindent\textbf{解说：}\textbf{}。\\\noindent\textbf{答案：}

\bigskip




\noindent\textbf{试题：}业余无线电数据传输控制差错的技术有“前向纠错”和“自动重传请求”两种。前向纠错（Forward Error Correction，FEC）是指：
\begin{enumerate}[leftmargin=3em]
\item 发送端在发送数据报文前先按一定的算法进行编码处理，加入特定的冗余代码，在接收端按照相应算法进行解码，从而发现传输产生的错误并加以纠正
\item 发送端在发送一定数据后，如收到接收端的重传请求、或者收不到接收端的收妥确认，即重传出错的数据报文，对出错报文进行纠正
\item 接收端收到数据后，发回发送端，由发送端比较核对，如发现出错，将重传出错的数据报文，从而恢复出错的报文
\item 发送端将每组数据报文重复发送足够多的次数，接收端收到后加以互相比较，选取比较一致的报文，以消除传输中的随机错误
\end{enumerate}
\noindent\textbf{解说：}\textbf{}。\\\noindent\textbf{答案：}

\bigskip




\noindent\textbf{试题：}业余无线电数据传输控制差错的技术有“前向纠错”和“自动重传请求”两种。自动重传请求（Automatic Repeat Request，ARQ）是指：
\begin{enumerate}[leftmargin=3em]
\item 发送端在发送一定数据后，如收到接收端的重传请求、或者收不到接收端的收妥确认，即重传出错的数据报文，对出错报文进行纠正
\item 发送端在发送数据报文前先按一定的算法进行编码处理，加入特定的冗余代码，在接收端按照相应算法进行解码，从而发现传输产生的错误并加以纠正
\item 接收端收到数据后，发回发送端，由发送端比较核对，如发现出错，将重传出错的数据报文，从而恢复出错的报文
\item 发送端将每组数据报文重复发送足够多的次数，接收端收到后加以互相比较，选取比较一致的报文，以消除传输中的随机错误
\end{enumerate}
\noindent\textbf{解说：}\textbf{}。\\\noindent\textbf{答案：}

\bigskip




\noindent\textbf{试题：}国际2号电报码（ITA2）的俗称、在业余无线电通信中应用场合及其与计算机常用的数据交换代码相比的主要特点是：
\begin{enumerate}[leftmargin=3em]
\item 博多码（Baudot code），用于RTTY通信，每字节仅包含5位二进制数据
\item ASCII码，用于PSK31通信，每字节仅包含7位二进制数据
\item GB2312编码，用于传输汉字，由两个字节连用代表一个汉字
\item 莫尔斯电码，用于CW通信，每个字符包含的信号数量不等（非对称码）
\end{enumerate}
\noindent\textbf{解说：}\textbf{}。\\\noindent\textbf{答案：}

\bigskip




\noindent\textbf{试题：}以CW方式进行速度为25WPM的摩尔斯电码通信，如果考虑传播衰落的影响，必要带宽约为：
\begin{enumerate}[leftmargin=3em]
\item 200Hz 
\item 2700Hz
\item 6.25kHz
\item 12.5kHz
\end{enumerate}
\noindent\textbf{解说：}\textbf{}。\\\noindent\textbf{答案：}

\bigskip




\noindent\textbf{试题：}用单边带方式进行RTTY通信，速度为50波特，频偏为170Hz，必要带宽约为：
\begin{enumerate}[leftmargin=3em]
\item 450Hz 
\item 200Hz
\item 2700Hz
\item 12.5kHz
\end{enumerate}
\noindent\textbf{解说：}\textbf{}。\\\noindent\textbf{答案：}

\bigskip




\noindent\textbf{试题：}用单边带方式进行话音通信，速度为50波特，频偏为70Hz，必要带宽约为：
\begin{enumerate}[leftmargin=3em]
\item 2700Hz
\item 200Hz
\item 6.25kHz
\item 12.5kHz
\end{enumerate}
\noindent\textbf{解说：}\textbf{}。\\\noindent\textbf{答案：}

\bigskip




\noindent\textbf{试题：}国际业余无线电联盟（IARU）业余无线电测向锦标赛规则要求144MHz测向发射机采用调幅度为70%-80%的A2A调制方式发射8-15WPM的音频摩尔斯电码发射电台标识。该类发射占用带宽不得大于下列数值：
\begin{enumerate}[leftmargin=3em]
\item 6kHz
\item 12.5kHz
\item 25kHz
\item 100kHz
\end{enumerate}
\noindent\textbf{解说：}\textbf{}。\\\noindent\textbf{答案：}

\bigskip




\noindent\textbf{试题：}无线电通信一般不用调频话方式传输摩尔斯电码，这是因为：
\begin{enumerate}[leftmargin=3em]
\item 调频话信道带宽超过传输摩尔斯电码的必要带宽，不符合我国《无线电频率划分规定》“电台的技术特性”章关于 “把带宽保持在技术状态和该项业务的性质所允许的最低值上”的要求
\item 调频话方式不能传输摩尔斯电码
\item 调频话方式传输摩尔斯电码可靠性差
\item 调频话方式传输摩尔斯电码时速度受限制
\end{enumerate}
\noindent\textbf{解说：}\textbf{}。\\\noindent\textbf{答案：}

\bigskip




\noindent\textbf{试题：}为了满足我国《无线电频率划分规定》“电台的技术特性”关于无线电通信“把带宽保持在技术状态和该项业务的性质所允许的最低值上”的要求，业余电台操作者应了解各种通信方式的必要带宽。决定必要带宽的因素是：
\begin{enumerate}[leftmargin=3em]
\item  传输的信息速率越高、系统的噪声干扰越大，必要带宽越宽
\item  发射设备功率越大，必要带宽越宽
\item  接收设备灵敏度越高，必要带宽越宽
\item  通信距离越近，必要带宽越宽
\end{enumerate}
\noindent\textbf{解说：}\textbf{}。\\\noindent\textbf{答案：}

\bigskip




\noindent\textbf{试题：}下列发射模式中拥有最窄带宽的是：
\begin{enumerate}[leftmargin=3em]
\item CW
\item 单边带话音
\item 调频话音
\item 慢扫描电视
\end{enumerate}
\noindent\textbf{解说：}\textbf{}。\\\noindent\textbf{答案：}

\bigskip




\noindent\textbf{试题：}“频率失真”是指电路的输出信号波形与输入信号相比，发生了下列变化：
\begin{enumerate}[leftmargin=3em]
\item 各频率分量的比例发生了改变
\item 信号的幅度发生了改变
\item 产生了新的频率分量
\item 不同频率分量的相位延迟差发生了改变
\end{enumerate}
\noindent\textbf{解说：}\textbf{}。\\\noindent\textbf{答案：}

\bigskip




\noindent\textbf{试题：}“非线性失真”是指电路的输出信号波形与输入信号相比，发生了下列变化：
\begin{enumerate}[leftmargin=3em]
\item 产生了新的频率分量
\item 各频率分量的比例发生了改变
\item 信号的幅度发生了改变
\item 不同频率分量的相位延迟差发生了改变
\end{enumerate}
\noindent\textbf{解说：}\textbf{}。\\\noindent\textbf{答案：}

\bigskip




\noindent\textbf{试题：}“相位失真”是指电路的输出信号波形与输入信号相比，发生了下列变化：
\begin{enumerate}[leftmargin=3em]
\item 不同频率分量的相位延迟差发生了改变
\item 产生了新的频率分量
\item 各频率分量的比例发生了改变
\item 信号的幅度发生了改变
\end{enumerate}
\noindent\textbf{解说：}\textbf{}。\\\noindent\textbf{答案：}

\bigskip




\noindent\textbf{试题：}在射频电路分析中，能产生信号频率以外的新频率分量的元器件属于有源元器件，可能成为形成干扰的重要环节。下列元器件中属于有源元器件的有：
\begin{enumerate}[leftmargin=3em]
\item 半导体二极管
\item 电解电容器
\item 宽带变压器
\item 碱性干电池
\end{enumerate}
\noindent\textbf{解说：}\textbf{}。\\\noindent\textbf{答案：}

\bigskip




\noindent\textbf{试题：}在射频电路分析中，能产生信号频率以外的新频率分量的元器件属于有源元器件，可能成为形成干扰的重要环节。下列元器件中属于有源元器件的有：
\begin{enumerate}[leftmargin=3em]
\item 半导体三极管
\item 可调电感器
\item 电阻假负载
\item 晶体滤波器
\end{enumerate}
\noindent\textbf{解说：}\textbf{}。\\\noindent\textbf{答案：}

\bigskip




\noindent\textbf{试题：}频率为f1、f2的两个正弦交流信号流过一个非线性元件，会发生“混频”。混频产物中属于三阶互调的干扰信号的频率是：
\begin{enumerate}[leftmargin=3em]
\item 2f1±f2、2f2±f1
\item 2f1、3f1、2f2、3f2
\item f1±f2、2f1±2f2、3f1±3f2
\item 4f1±f2、5f1±2f2、6f1±3f2……
\end{enumerate}
\noindent\textbf{解说：}\textbf{}。\\\noindent\textbf{答案：}

\bigskip




\noindent\textbf{试题：}“差拍”现象是指：
\begin{enumerate}[leftmargin=3em]
\item 两个不同频率信号经过非线性电路得到频率为两者之差的新频率信号
\item 凡是接收机收到的连续音频叫声都叫做“差拍”
\item 发射机用以调制载波的连续音频信号
\item 一个单频率信号经过非线性电路得到一系列谐波信号，相邻信号之间的频率差等于单频信号的频率。这样一族谐波信号的集合总称“差拍”
\end{enumerate}
\noindent\textbf{解说：}\textbf{}。\\\noindent\textbf{答案：}

\bigskip




\noindent\textbf{试题：}可以作为倍频器的电路是：
\begin{enumerate}[leftmargin=3em]
\item 非线性电路
\item 甲类放大电路
\item 调谐在倍频上的LC谐振电路
\item 谐振频率等于倍频的石英晶体谐振器
\end{enumerate}
\noindent\textbf{解说：}\textbf{}。\\\noindent\textbf{答案：}

\bigskip




\noindent\textbf{试题：}正弦交流信号通过下列电路时会产生高次谐波：
\begin{enumerate}[leftmargin=3em]
\item 二极管整流器或三极管开关放大器
\item 复杂电阻网络
\item 复杂电容电感网络
\item 调谐在信号的整倍频上的谐振电路
\end{enumerate}
\noindent\textbf{解说：}\textbf{}。\\\noindent\textbf{答案：}

\bigskip




\noindent\textbf{试题：}一个重复频率为F的非正弦周期信号的频谱包含有：
\begin{enumerate}[leftmargin=3em]
\item 频率为F的整数倍的无穷多个频率分量
\item 频率为F的一个频率分量
\item 频率为F以外的无穷多个频率分量
\item 无穷多个连续的频率分量
\end{enumerate}
\noindent\textbf{解说：}\textbf{}。\\\noindent\textbf{答案：}

\bigskip




\noindent\textbf{试题：}一个PSK31信号在频谱仪上显示为：
\begin{enumerate}[leftmargin=3em]
\item 一条固定的垂直线
\item 一条固定的直线，左右伴随一组对称的随语音变化的直线
\item 多条固定的直线
\item 一条复杂的周期性曲线
\end{enumerate}
\noindent\textbf{解说：}\textbf{}。\\\noindent\textbf{答案：}

\bigskip




\noindent\textbf{试题：}下列几种图表中，最容易用来表达和解释PSK调制原理的是：
\begin{enumerate}[leftmargin=3em]
\item 相位矢量图
\item 波形图
\item 频谱图
\item 频谱瀑布图
\end{enumerate}
\noindent\textbf{解说：}\textbf{}。\\\noindent\textbf{答案：}

\bigskip




\noindent\textbf{试题：}下列几种图表中，最容易用来表达和解释SSB调制原理的是：
\begin{enumerate}[leftmargin=3em]
\item 频谱图
\item 波形图
\item 相位矢量图
\item 频谱瀑布图
\end{enumerate}
\noindent\textbf{解说：}\textbf{}。\\\noindent\textbf{答案：}

\bigskip




\noindent\textbf{试题：}下列几种图表中，最容易用来表达和解释RTTY调制原理的是：
\begin{enumerate}[leftmargin=3em]
\item 频率-时间曲线
\item 功率-时间曲线
\item 波形图
\item 相位矢量图
\end{enumerate}
\noindent\textbf{解说：}\textbf{}。\\\noindent\textbf{答案：}

\bigskip




\noindent\textbf{试题：}下列几种图表中，最容易用来表达和解释SSTV调制原理的是：
\begin{enumerate}[leftmargin=3em]
\item 频率-时间曲线
\item 功率-时间曲线
\item 波形图
\item 相位矢量图
\end{enumerate}
\noindent\textbf{解说：}\textbf{}。\\\noindent\textbf{答案：}

\bigskip




\noindent\textbf{试题：}下列几种图表中，最容易用来表达和解释A、B、C类放大器工作点的差别的是：
\begin{enumerate}[leftmargin=3em]
\item 波形图
\item 频谱图
\item 相位矢量图
\item 频谱瀑布图
\end{enumerate}
\noindent\textbf{解说：}\textbf{}。\\\noindent\textbf{答案：}

\bigskip




\noindent\textbf{试题：}下列几种图表中，最容易用来表达和解释A、B、C类放大器工作点的差别的是：
\begin{enumerate}[leftmargin=3em]
\item 输入-输出特性曲线
\item 功率-时间曲线
\item 频率-时间曲线
\item 输入阻抗特性曲线
\end{enumerate}
\noindent\textbf{解说：}\textbf{}。\\\noindent\textbf{答案：}

\bigskip




\noindent\textbf{试题：}一个CW电报信号在频谱仪上显示为：
\begin{enumerate}[leftmargin=3em]
\item 一条闪动的垂直线
\item 一条固定的直线
\item 两条水平直线
\item 一条正弦波曲线
\end{enumerate}
\noindent\textbf{解说：}\textbf{}。\\\noindent\textbf{答案：}

\bigskip




\noindent\textbf{试题：}一个RTTY信号在频谱仪上显示为：
\begin{enumerate}[leftmargin=3em]
\item 两条闪动的垂直线
\item 一条垂直线
\item 一条正弦波曲线
\item 一条复杂的周期性曲线
\end{enumerate}
\noindent\textbf{解说：}\textbf{}。\\\noindent\textbf{答案：}

\bigskip




\noindent\textbf{试题：}一个SSB话音信号在频谱仪上显示为：
\begin{enumerate}[leftmargin=3em]
\item 一组随语音出现和变化的非对称垂直线
\item 一条随语音闪烁的直线
\item 一条固定的直线
\item 一条复杂的周期性曲线
\end{enumerate}
\noindent\textbf{解说：}\textbf{}。\\\noindent\textbf{答案：}

\bigskip




\noindent\textbf{试题：}一个AM话音信号在频谱仪上显示为：
\begin{enumerate}[leftmargin=3em]
\item 一条固定的垂直线，左右伴随一组对称的随语音出现和变化的垂直线
\item 一条随语音闪烁的直线
\item 多条固定的直线
\item 一条复杂的周期性曲线
\end{enumerate}
\noindent\textbf{解说：}\textbf{}。\\\noindent\textbf{答案：}

\bigskip




\noindent\textbf{试题：}下列几种图表中，最容易用来表达和解释LC振荡器温度漂移程度的是：
\begin{enumerate}[leftmargin=3em]
\item 频谱瀑布图
\item 频谱图
\item 相位矢量图
\item 波形图
\end{enumerate}
\noindent\textbf{解说：}\textbf{}。\\\noindent\textbf{答案：}

\bigskip




\noindent\textbf{试题：}将接收机的输出与实时频谱仪或运行频谱显示软件的电脑相连，观察HF频段的RTTY或PACKET模式的DX信号，通常看到两个幅度相同的谱峰。但有时两个谱峰的高度会变得不等并且不断变化。造成这种现象的原因是：
\begin{enumerate}[leftmargin=3em]
\item 电离层的选择性衰落
\item 仪器或电脑工作点飘移
\item 发射机的发射功率电平不稳定
\item 接收机的自动增益控制不稳定
\end{enumerate}
\noindent\textbf{解说：}\textbf{}。\\\noindent\textbf{答案：}

\bigskip




\noindent\textbf{试题：}比较考究的接收机说明书中常列出一项“通带矩形系数”或“通带形状系数”的指标，是用来描述：
\begin{enumerate}[leftmargin=3em]
\item 通带滤波器频率特性曲线斜坡的陡峭程度
\item 矩形波信号通过接收机所产生的延迟时间
\item 通带滤波器对矩形波信号的通过能力
\item 通带晶体滤波器中石英晶片的形状和切割方式
\end{enumerate}
\noindent\textbf{解说：}\textbf{}。\\\noindent\textbf{答案：}

\bigskip




\noindent\textbf{试题：}根据说明书列出的技术指标，两台业余无线电接收机在USB方式下选择带宽为2.7kHz的滤波器时具有不同的“通带矩形系数”。接收机A的“60dB带宽对6dB带宽的矩形系数”为3.8，接收机B的 为5。由此可作出推论：
\begin{enumerate}[leftmargin=3em]
\item A机对邻近频道干扰的抑制能力比B机强
\item A机对邻近频道干扰的抑制能力比B机差
\item A机对镜像频道干扰的抑制能力比B机强
\item A机对偏离工作频率±10kHz以外的干扰信号的抑制能力比B机强
\end{enumerate}
\noindent\textbf{解说：}\textbf{}。\\\noindent\textbf{答案：}

\bigskip




\noindent\textbf{试题：}根据说明书列出的技术指标，两台业余无线电接收机的“通带波动”分别为1dB和3dB。由此可作出推论：
\begin{enumerate}[leftmargin=3em]
\item A机对通道内信号的频率响应比B机平坦
\item A机对通道内信号的频率响应不如B机平坦
\item A机接收到的信号的幅度颤动比B机小
\item A机的通带增益的可调范围比B机小
\end{enumerate}
\noindent\textbf{解说：}\textbf{}。\\\noindent\textbf{答案：}

\bigskip




\noindent\textbf{试题：}无线电发射设备参数和业余无线电原理书籍中经常出现缩写为ppm的度量单位。其中文含义和最经常的用处是：
\begin{enumerate}[leftmargin=3em]
\item “百万分比”，常用于描述频率的相对稳定度
\item “每分钟点数”，常用于描述随机干扰的出现次数
\item “每分钟点数”，常用于描述太阳黑子的发生频度
\item “每分钟平均脉冲数”，常用于描述单个脉冲的出现次数
\end{enumerate}
\noindent\textbf{解说：}\textbf{}。\\\noindent\textbf{答案：}

\bigskip




\noindent\textbf{试题：}用功率电平表示的接收机灵敏度的意义以及常用单位是：
\begin{enumerate}[leftmargin=3em]
\item 输出信号达到一定质量标准时输入信号的最小功率电平，单位dBm或dBμW
\item 输出信号达到一定质量标准时输出信号与输入信号的功率电平比，单位dB
\item 输出信号达到一定质量标准时输入信号与输出信号的功率电平比，单位dB
\item 输出信号维持一定质量标准时输入信号的最小和最大功率电平之比，单位dB
\end{enumerate}
\noindent\textbf{解说：}\textbf{}。\\\noindent\textbf{答案：}

\bigskip




\noindent\textbf{试题：}用电压电平表示的接收机灵敏度在无线电理论中的意义以及常用单位是：
\begin{enumerate}[leftmargin=3em]
\item 输出信号达到一定质量标准时输入信号的最小电动势，单位μV，或换算成dBμV、dBmV
\item 输出信号达到一定质量标准时输入信号的最小电压，单位μV，或换算成dBμV、dBmV
\item 输出信号达到一定质量标准时输入信号的最大电动势，单位μV或mV
\item 输出信号达到一定质量标准时输出信号与输入信号的电压比，单位dB/V或dB/mV
\end{enumerate}
\noindent\textbf{解说：}\textbf{}。\\\noindent\textbf{答案：}

\bigskip




\noindent\textbf{试题：}甲乙两种业余无线电台设备资料列出接收机灵敏度指标分别为0.1μV和0.15μV。正确的推论为：
\begin{enumerate}[leftmargin=3em]
\item 凭此指标还无法比较两者接收微弱信号的能力，因没有给出测量灵敏度时的输出信号质量条件
\item 可以推断甲机接收微弱信号的能力比乙机的低，因为灵敏度数值比较小
\item 可以推断甲机接收微弱信号的能力比乙机的高，因为可以接收的信号更微弱
\item 可以推断甲机承受强信号的能力比乙机的低，因为灵敏度数值比较小
\end{enumerate}
\noindent\textbf{解说：}\textbf{}。\\\noindent\textbf{答案：}

\bigskip




\noindent\textbf{试题：}输出信号达到一定质量标准时，在某输入阻抗为50欧的接收机天线输入端测得输入信号的最小电压为1μV，该接收机敏度用电压μV表示的数值为：
\begin{enumerate}[leftmargin=3em]
\item 2μV
\item 1μV
\item 50μV
\item 0.5μV
\end{enumerate}
\noindent\textbf{解说：}\textbf{}。\\\noindent\textbf{答案：}

\bigskip




\noindent\textbf{试题：}输出信号达到一定质量标准时，在某输入阻抗为50欧的接收机天线输入端测得输入信号的最小电压为1μV，该接收机灵敏度的功率电平dBm数值为：
\begin{enumerate}[leftmargin=3em]
\item -107dBm
\item -103dBm
\item -73dBm
\item -113dBm
\end{enumerate}
\noindent\textbf{解说：}\textbf{}。\\\noindent\textbf{答案：}

\bigskip




\noindent\textbf{试题：}输出信号达到一定质量标准时，在某输入阻抗为50欧的接收机天线输入端测得输入信号的最小电压为1μV，该接收机敏度用dBμV表示的数值为：
\begin{enumerate}[leftmargin=3em]
\item 6dBμV
\item -2dBμV
\item 3dBμV
\item 0dBμV
\end{enumerate}
\noindent\textbf{解说：}\textbf{}。\\\noindent\textbf{答案：}

\bigskip




\noindent\textbf{试题：}输出信号达到一定质量标准时，在某输入阻抗为50欧的接收机天线输入端测得输入信号的最小电压为0.5μV，该接收机敏度用电压μV表示的数值为：
\begin{enumerate}[leftmargin=3em]
\item 1μV
\item 0.5μV
\item 50μV
\item 2.5μV
\end{enumerate}
\noindent\textbf{解说：}\textbf{}。\\\noindent\textbf{答案：}

\bigskip




\noindent\textbf{试题：}输出信号达到一定质量标准时，在某输入阻抗为50欧的接收机天线输入端测得输入信号的最小电压为0.5 μV，该接收机敏度的功率电平dBm数值为：
\begin{enumerate}[leftmargin=3em]
\item -113 dBm
\item -107dBm
\item -103 dBm
\item -73 dBm
\end{enumerate}
\noindent\textbf{解说：}\textbf{}。\\\noindent\textbf{答案：}

\bigskip




\noindent\textbf{试题：}输出信号达到一定质量标准时，在某输入阻抗为50欧的接收机天线输入端测得输入信号的最小电压为0.5μV，该接收机敏度用dBμV表示的数值为：
\begin{enumerate}[leftmargin=3em]
\item 0dBμV
\item -2dBμV
\item 3dBμV
\item 6dBμV
\end{enumerate}
\noindent\textbf{解说：}\textbf{}。\\\noindent\textbf{答案：}

\bigskip




\noindent\textbf{试题：}已知接收机天线输入阻抗为50Ω，电压灵敏度为2μV，其功率灵敏度为：
\begin{enumerate}[leftmargin=3em]
\item -107dBm
\item 0.02dBm
\item -103dBm
\item -113dBm
\end{enumerate}
\noindent\textbf{解说：}\textbf{}。\\\noindent\textbf{答案：}

\bigskip




\noindent\textbf{试题：}已知接收机天线输入阻抗为50Ω，电压灵敏度为1μV，其功率灵敏度为：
\begin{enumerate}[leftmargin=3em]
\item -113dBm
\item -107dBm
\item 0.01dBm
\item -103dBm
\end{enumerate}
\noindent\textbf{解说：}\textbf{}。\\\noindent\textbf{答案：}

\bigskip




\noindent\textbf{试题：}决定接收机抗拒工作频率附近QRM的主要选择性指标是：
\begin{enumerate}[leftmargin=3em]
\item 信道带宽、信道选择性和信道滤波器特性矩形系数
\item 带内波动和信道带宽
\item 镜像抑制比
\item 前端带宽
\end{enumerate}
\noindent\textbf{解说：}\textbf{}。\\\noindent\textbf{答案：}

\bigskip




\noindent\textbf{试题：}决定接收机接收SSB信号时语音频响均匀度的主要选择性指标是：
\begin{enumerate}[leftmargin=3em]
\item 带内波动和信道带宽
\item 信道带宽、信道选择性和信道滤波器特性矩形系数
\item 镜像抑制比
\item 前端带宽
\end{enumerate}
\noindent\textbf{解说：}\textbf{}。\\\noindent\textbf{答案：}

\bigskip




\noindent\textbf{试题：}决定接收机抗拒与工作频率相距较远的强信号干扰的主要选择性指标是：
\begin{enumerate}[leftmargin=3em]
\item 前端带宽
\item 带内波动和信道带宽
\item 信道带宽、信道选择性和信道滤波器特性矩形系数
\item 镜像抑制比
\end{enumerate}
\noindent\textbf{解说：}\textbf{}。\\\noindent\textbf{答案：}

\bigskip




\noindent\textbf{试题：}决定接收机抗拒与工作频率相距两倍于中频的频率上强信号干扰的主要选择性指标是：
\begin{enumerate}[leftmargin=3em]
\item 镜像抑制比
\item 前端带宽
\item 带内波动和信道带宽
\item 信道带宽、信道选择性和信道滤波器特性矩形系数
\end{enumerate}
\noindent\textbf{解说：}\textbf{}。\\\noindent\textbf{答案：}

\bigskip




\noindent\textbf{试题：}“信噪比（SNR）”是无线电技术大量使用的名词，它是指：
\begin{enumerate}[leftmargin=3em]
\item 有用信号功率对噪声功率的比值
\item 有用信号电压幅度对噪声电压幅度的比值
\item 有用信号功率对有用信号功率及噪声功率之和的比值
\item 有用信号电压幅度对有用信号电压幅度及噪声电压幅度之和的比值
\end{enumerate}
\noindent\textbf{解说：}\textbf{}。\\\noindent\textbf{答案：}

\bigskip




\noindent\textbf{试题：}对于需要接收微弱信号的业余通信，接收机的噪声系数Fn （以比值的数值形式表达）是一项重要技术指标。它是指：
\begin{enumerate}[leftmargin=3em]
\item 接收机输入端信噪比Si/Ni对输出端信噪比So/No的比值
\item 接收机输入端无信号时，输出端的噪声功率电平
\item 接收机输出端噪声功率电平与输入端噪声功率电平的比值
\item 接收机输入端噪声功率电平与输出端噪声功率电平的比值
\end{enumerate}
\noindent\textbf{解说：}\textbf{}。\\\noindent\textbf{答案：}

\bigskip




\noindent\textbf{试题：}对于需要接收微弱信号的业余通信，接收机的噪声系数Fn （以比值的数值形式表达）是一项重要技术指标。关于Fn的一些基本常识是：
\begin{enumerate}[leftmargin=3em]
\item Fn一定大于1；在同样的灵敏度下Fn越接近于1越好
\item Fn一定小于1；在同样的灵敏度下Fn越接近于0越好
\item Fn一定大于1；在同样的灵敏度下Fn越大越好
\item Fn一定小于1；在同样的灵敏度下Fn越接近于1越好
\end{enumerate}
\noindent\textbf{解说：}\textbf{}。\\\noindent\textbf{答案：}

\bigskip




\noindent\textbf{试题：}对于需要接收微弱信号的业余通信，接收机的噪声指数NF（以对数形式表达）是一项重要技术指标。它是指：
\begin{enumerate}[leftmargin=3em]
\item 接收机输入端信噪比Si/Ni对输出端信噪比So/No的比值的对数表达形式
\item 接收机输入端无信号时，输出端的噪声功率电平的对数表达形式
\item 接收机输出端噪声功率电平与输入端噪声功率电平的比值的对数表达形式
\item 接收机输入端噪声功率电平与输出端噪声功率电平的比值的对数表达形式
\end{enumerate}
\noindent\textbf{解说：}\textbf{}。\\\noindent\textbf{答案：}

\bigskip




\noindent\textbf{试题：}对于需要接收微弱信号的业余通信，接收机的噪声指数NF（以对数形式表达）是一项重要技术指标。关于NF的一些基本常识是：
\begin{enumerate}[leftmargin=3em]
\item NF一定大于0；在同样的灵敏度下NF越接近于0越好
\item NF一定小于0；在同样的灵敏度下NF越接近于0越好
\item NF一定大于1；在同样的灵敏度下NF越大越好
\item NF一定处于0和1之间；在同样的灵敏度下NF越接近于1越好
\end{enumerate}
\noindent\textbf{解说：}\textbf{}。\\\noindent\textbf{答案：}

\bigskip




\noindent\textbf{试题：}制约现代无线电接收机灵敏度的主要因素是：
\begin{enumerate}[leftmargin=3em]
\item 机内噪声
\item 放大电路的增益
\item 放大电路的稳定性
\item 电源噪声
\end{enumerate}
\noindent\textbf{解说：}\textbf{}。\\\noindent\textbf{答案：}

\bigskip




\noindent\textbf{试题：}业余通信接收机大多带有接收信号强度指示。VHF/UHF频段的最小刻度S1对应于输入信号功率电平-141dBm（50Ω输入电压0.02μV）标为S1，而HF频段的S1则对应于输入信号功率电平-121dBm（0.2μV）。这是因为：
\begin{enumerate}[leftmargin=3em]
\item V/UHF频段较寂静而HF频段外界背景噪声电平较高，前者可感知的最小信号电平比后者低约20dB
\item 由于电路技术的原因，HF频段接收机的灵敏度只能做到比VHF/UHF频段低大约20dB
\item HF业余电台功率一般比较大，VHF/UHF电台功率比较小，因此需要不同的刻度标准
\item HF业余电台主要用于DX通信，VHF/UHF电台主要做本地通信，因此需要不同的刻度标准
\end{enumerate}
\noindent\textbf{解说：}\textbf{}。\\\noindent\textbf{答案：}

\bigskip




\noindent\textbf{试题：}无线电通信领域中，在描述信号源、放大器、接收机等系统部件的内部噪声大小时，经常用“噪声温度”来度量。以接收机为例，等效噪声温度Te的意义是：
\begin{enumerate}[leftmargin=3em]
\item 接收机的内部噪声功率等于一个接在天线输入端的优质匹配电阻产生相同的热噪声功率时电阻所具有的绝对温度
\item 接收机内部噪声在输出端的功率可以使一个接在输出端上的匹配电阻加热升高的相对温度
\item 接收机内部噪声在输出端的功率可以使一个接在输出端上的匹配电阻加热到的绝对温度
\item 接收机信噪比符合技术指标时所要求的设备环境温度
\end{enumerate}
\noindent\textbf{解说：}\textbf{}。\\\noindent\textbf{答案：}

\bigskip




\noindent\textbf{试题：}不产生任何内部噪声的理想放大器或接收机的噪声系数Fn、噪声指数NF和噪声温度Te分别为：
\begin{enumerate}[leftmargin=3em]
\item 1，0dB，0°K
\item 0，0dB，-273°K
\item 0，1dB，17°K
\item 0，0dB，-275°K
\end{enumerate}
\noindent\textbf{解说：}\textbf{}。\\\noindent\textbf{答案：}

\bigskip




\noindent\textbf{试题：}如果一个接收业余卫星用的天线放大器工作在标准温度（17℃）下，其输入端接一副良好匹配的天线。如果放大器所产生的内部噪声与输入的热噪声等效，该放大器的噪声系数Fn、噪声指数NF和噪声温度Te分别为：
\begin{enumerate}[leftmargin=3em]
\item 2，3dB，290°K
\item 1，1dB，0°K
\item 2，0dB，17°K
\item 1，0dB，-273°K
\end{enumerate}
\noindent\textbf{解说：}\textbf{}。\\\noindent\textbf{答案：}

\bigskip




\noindent\textbf{试题：}无线电接收机的灵敏度是指：
\begin{enumerate}[leftmargin=3em]
\item 接收机正常工作所需的最小输入信号强度
\item 接收机正常工作所需的最大输入信号强度
\item 接收机正常工作所需的最大电源电压
\item 接收机正常工作所需的最小电源功率
\end{enumerate}
\noindent\textbf{解说：}\textbf{}。\\\noindent\textbf{答案：}

\bigskip




\noindent\textbf{试题：}单工操作（simplex operation）是指：在一条电信通路的两个方向上交替进行传输的一种工作方式。下列通信采用单工操作方式：
\begin{enumerate}[leftmargin=3em]
\item 两台由PTT键转换收发状态的业余对讲机之间的通信
\item 两台有线电话座机之间的通信
\item 两台公众移动电话终端（手机）之间的通信
\item 一台业余中继台与一台由PTT键转换收发状态的业余对讲机之间的通信
\end{enumerate}
\noindent\textbf{解说：}\textbf{}。\\\noindent\textbf{答案：}

\bigskip




\noindent\textbf{试题：}双工操作（duplex operation）是指：一条电信通路的两个方向能同时进行传输的工作方式。下列通信采用双工操作方式：
\begin{enumerate}[leftmargin=3em]
\item 由一台连续工作的430MHz接收机和一台144MHz发信机组成的地面业余电台与业余卫星之间的通信
\item 两台具有QSK功能的业余CW收发信之间的通信
\item 一台业余中继台与一台由PTT键转换收发状态的业余对讲机之间的通信
\item 两台由PTT键转换收发状态的业余对讲机之间的通信
\end{enumerate}
\noindent\textbf{解说：}\textbf{}。\\\noindent\textbf{答案：}

\bigskip




\noindent\textbf{试题：}半双工操作（semi-duplex operation）是指：电路的一端用单工操作，另一端用双工操作的一种工作方式。下列通信采用半双工操作方式：
\begin{enumerate}[leftmargin=3em]
\item 一台业余中继台与一台由PTT键转换收发状态的业余对讲机之间的通信
\item 两台具有QSK功能的业余CW收发信之间的通信
\item 两台由PTT键转换收发状态的业余对讲机之间的通信
\item 两台公众移动电话终端（手机）之间的通信
\end{enumerate}
\noindent\textbf{解说：}\textbf{}。\\\noindent\textbf{答案：}

\bigskip




\noindent\textbf{试题：}业余无线电技术描述收信机预选滤波电路和发信机输出滤波电路时，经常会提到LC滤波器的“阶数”或“极点数”。确定LC滤波器阶数的简单办法是：
\begin{enumerate}[leftmargin=3em]
\item 先合并简化串、并联的电容元件或串、并联的电感元件，得到的电容元件和电感元件总个数即为阶数
\item 滤波电路中分叉节的数量即为阶数
\item 滤波电路中分叉节的数量减1即为阶数
\item 滤波电路中电感元件的总个数减去电容元件的总个数，即得阶数
\end{enumerate}
\noindent\textbf{解说：}\textbf{}。\\\noindent\textbf{答案：}

\bigskip




\noindent\textbf{试题：}接收月面反射通信等VHF/UHF微弱信号时常需要在八木天线端加接加接低噪声前置放大器，而HF频段通信很少在八木天线端加装前置放大器。这是因为：
\begin{enumerate}[leftmargin=3em]
\item VHF/UHF频段信号微弱，在天线端放大可减少电缆固有噪声影响；HF频段环境背景噪声电平高，电缆固有噪声可以忽略
\item 电缆固有噪声只产生在VHF/UHF频段
\item HF频段的放大器体积太大，不适合于安装在天线端
\item 在HF频段的八木天线体积较大，把放大器安装得太近容易引起自激
\end{enumerate}
\noindent\textbf{解说：}\textbf{}。\\\noindent\textbf{答案：}

\bigskip




\noindent\textbf{试题：}接收机前置放大器的主要作用是：
\begin{enumerate}[leftmargin=3em]
\item 提高信噪比
\item 提高接收机最终的音频输出功率电平
\item 改善接收机音频输出的保真度
\item 改善接收机的动态范围
\end{enumerate}
\noindent\textbf{解说：}\textbf{}。\\\noindent\textbf{答案：}

\bigskip




\noindent\textbf{试题：}在EME通信等需要接收极微弱信号的场合，安装前置放大器时应做到：
\begin{enumerate}[leftmargin=3em]
\item 尽量安装在靠近天线的地方，信号特别微弱时对放大器冷却以降低热噪声
\item 尽量安装在靠近接收机的地方，放大器最好加热以防止元器件温度过低
\item 尽量安装在靠近接收机的地方，放大器尽量采用较高的电源电压以求更高的增益
\item 尽量安装在远离天线的地方，放大器采用射频正反馈电路以提高增益
\end{enumerate}
\noindent\textbf{解说：}\textbf{}。\\\noindent\textbf{答案：}

\bigskip




\noindent\textbf{试题：}业余调频中继台发射机只要被上行信号正常启动，就会一直继续发射载波，上行信号消失不能使其停止。可能的原因是：
\begin{enumerate}[leftmargin=3em]
\item 中继台上下行隔离不良，中继台发射的载波窜入中继台接收机造成自锁
\item 肯定受到人为恶意干扰
\item 中继台接收机电源电压不稳
\item 中继台发射机电源电压不稳
\end{enumerate}
\noindent\textbf{解说：}\textbf{}。\\\noindent\textbf{答案：}

\bigskip




\noindent\textbf{试题：}业余调频中继台发射机被上行信号正常启动，但上行信号消失后经常会继续发射一段或长或短的时间并夹杂有一些不清楚的语音。可能的原因是：
\begin{enumerate}[leftmargin=3em]
\item 中继台下行信号与附近的其他通信发射机形成对中继台上行频率的互调干扰
\item 肯定受到人为恶意干扰
\item 中继台接收机电源电压不稳
\item 中继台发射机电源电压不稳
\end{enumerate}
\noindent\textbf{解说：}\textbf{}。\\\noindent\textbf{答案：}

\bigskip




\noindent\textbf{试题：}业余中继台上下行共用一副天线时，需要在接收机、发信机和天线之间插入一个：
\begin{enumerate}[leftmargin=3em]
\item 双工器（duplexer）
\item 收发信转换开关
\item 功率分配器（power divider）
\item 环形器（circulator）
\end{enumerate}
\noindent\textbf{解说：}\textbf{}。\\\noindent\textbf{答案：}

\bigskip




\noindent\textbf{试题：}业余中继台上下行共用一副天线时，需要在接收机、发信机和天线之间插入一个双工器，其基本构造和作用为：
\begin{enumerate}[leftmargin=3em]
\item 一组滤波器，阻止中继台发射信号反馈进入中继台接收机
\item 一组半导体开关，在中继台发射时关断中继台接收机
\item 一个匹配网络，使天线、中继发射机、中继接收机三者之间都满足阻抗匹配条件
\item 一个环形器，使信号只能沿中继发射机-天线-中继接收机的方向前进
\end{enumerate}
\noindent\textbf{解说：}\textbf{}。\\\noindent\textbf{答案：}

\bigskip




\noindent\textbf{试题：}天线的接收单元直接连到适当的放大器就构成有源天线。对于同样结构的天线，使用有源接收天线的目的是：
\begin{enumerate}[leftmargin=3em]
\item 提高天线输出的信号电平，以克服传输线路和接收机前级固有噪声的影响
\item 提高天线接收到的有用信号电平，以克服天线周围环境的外部噪声影响
\item 提高天线接收到的有用信号电平，以克服临近频道电台信号的干扰
\item 展宽天线的工作带宽
\end{enumerate}
\noindent\textbf{解说：}\textbf{}。\\\noindent\textbf{答案：}

\bigskip




\noindent\textbf{试题：}射频前置放大器通常安装在什么地方？
\begin{enumerate}[leftmargin=3em]
\item 在天线和接收机之间
\item 在发射机功放的后面
\item 在发射机和天调之间
\item 在接收机音频输出处
\end{enumerate}
\noindent\textbf{解说：}\textbf{}。\\\noindent\textbf{答案：}

\bigskip




\noindent\textbf{试题：}亚音调静噪(CTCSS)是指附加在发射端信号中的一个亚音频控制音调。这个信号的频率范围大致是：
\begin{enumerate}[leftmargin=3em]
\item 67Hz - 250.3Hz
\item 16Hz - 20kHz
\item 16kHz - 20kHz
\item 220Hz - 2503Hz
\end{enumerate}
\noindent\textbf{解说：}\textbf{}。\\\noindent\textbf{答案：}

\bigskip




\noindent\textbf{试题：}学习现代无线电技术或从事业余无线电制作时，经常遇到一种与超外差式收信机具有不同原理的“DC式收信机”。其中缩写DC是指：
\begin{enumerate}[leftmargin=3em]
\item 直接变换（Direct-Covertion），即接收到的射频信号在解调处理之前不经过频率变换
\item 直流（Direct Current），指用直流电源供电的收信机
\item 介质电容（Dielectric Capacitor），指采用了特殊介质电容器的收信机
\item 数模变换（Digital-Analog Conversion），指包含有数模变换器的收信机
\end{enumerate}
\noindent\textbf{解说：}\textbf{}。\\\noindent\textbf{答案：}

\bigskip




\noindent\textbf{试题：}在超外差式收信机电路中，信号通道的有用信号频率比本振频率低（或者高）一个中频频率。但比本振频率高（或者低）一个中频频率的信号也可能窜入信号通道，称为“镜像频率干扰”或“镜频干扰”。某VHF对讲机使用说明书的技术指标部分给出了第一中频（IF）为45.05MHz，但没有更多的资料。由此可推测当接收145.00MHz信号时下述频率之一的强信号可能造成镜频干扰：
\begin{enumerate}[leftmargin=3em]
\item 235.10MHz或54.90MHz
\item 190.05MHz或99.95MHz
\item 45.05MHz或90.10MHz
\item 90.10MHz或.180.20MHz
\end{enumerate}
\noindent\textbf{解说：}\textbf{}。\\\noindent\textbf{答案：}

\bigskip




\noindent\textbf{试题：}在超外差式收信机电路中，信号通道的有用信号频率比本振频率低（或者高）一个中频频率。但比本振频率高（或者低）一个中频频率的信号也可能窜入信号通道，称为“镜像频率干扰”或“镜频干扰”。某对讲机使用说明书的技术指标部分给出了在NFM方式时第一中频（IF）为47.25MHz，但没有更多的资料。由此可推测当接收145.00MHz信号时下述频率之一的强信号可能造成镜频干扰：：
\begin{enumerate}[leftmargin=3em]
\item 239.50MHz或50.50MHz
\item 192.25MHz或97.75MHz
\item 50.50MHz或101.00MHz
\item 151.50MHz或.202.00MHz
\end{enumerate}
\noindent\textbf{解说：}\textbf{}。\\\noindent\textbf{答案：}

\bigskip




\noindent\textbf{试题：}在超外差式收信机电路中，信号通道的有用信号频率比本振频率低（或者高）一个中频频率。但比本振频率高（或者低）一个中频频率的信号也可能窜入信号通道，称为“镜像频率干扰”或“镜频干扰”。某对讲机使用说明书的技术指标部分给出了接收NFM信号时第一中频（IF）为47.25MHz，但没有更多的资料。由此可推测当接收435.00MHz信号时下述频率之一的强信号可能造成镜频干扰：
\begin{enumerate}[leftmargin=3em]
\item 340.50MHz或529.50MHz
\item 387.75MHz或482.25MHz
\item 47.25MHz或94.50MHz
\item 141.70MHz或.236.25MHz
\end{enumerate}
\noindent\textbf{解说：}\textbf{}。\\\noindent\textbf{答案：}

\bigskip




\noindent\textbf{试题：}在超外差式收信机电路中，信号通道的有用信号频率比本振频率低（或者高）一个中频频率。但比本振频率高（或者低）一个中频频率的信号也可能窜入信号通道，称为“镜像频率干扰”或“镜频干扰”。某UHF对讲机的使用说明书技术指标部分给出了第一中频（IF）为58.525MHz，但没有更多的资料。由此可推测当接收435.00MHz信号时下述频率之一的强信号可能造成镜频干扰：
\begin{enumerate}[leftmargin=3em]
\item 317.95MHz或552.05MHz
\item 376.475MHz或493.525MHz
\item 58.525MHz或117.05MHz
\item 234.10.05MHz或.468.20MHz
\end{enumerate}
\noindent\textbf{解说：}\textbf{}。\\\noindent\textbf{答案：}

\bigskip




\noindent\textbf{试题：}现代超外差式业余收发信机面板上常设有选择中频滤波器带宽的控制部件。这些中频滤波器负责抑制的干扰种类为：
\begin{enumerate}[leftmargin=3em]
\item 邻近频率干扰
\item 镜像频率干扰
\item 中频频率干扰
\item 突发脉冲干扰
\end{enumerate}
\noindent\textbf{解说：}\textbf{}。\\\noindent\textbf{答案：}

\bigskip




\noindent\textbf{试题：}超外差式业余收发信机中负责抑制镜像频率干扰的部件是：
\begin{enumerate}[leftmargin=3em]
\item 变频级之前的预选滤波器
\item 变频级之后的中频滤波器
\item 中频放大级中的限幅电路
\item 带有音调控制的音频滤波器
\end{enumerate}
\noindent\textbf{解说：}\textbf{}。\\\noindent\textbf{答案：}

\bigskip




\noindent\textbf{试题：}在超外差式变频装置中，所谓“上变频”和“下变频”分别是指：
\begin{enumerate}[leftmargin=3em]
\item 中频频率高于输入频率为上变频方式，中频频率低于输入频率为下变频方式
\item 输入频率高于中频频率为上变频方式，输入频率低于中频频率为下变频方式
\item 本振频率高于输入频率为上变频方式，本振频率低于输入频率为下变频方式
\item 输入频率高于本振频率为上变频方式，输入频率低于本振频率为下变频方式
\end{enumerate}
\noindent\textbf{解说：}\textbf{}。\\\noindent\textbf{答案：}

\bigskip




\noindent\textbf{试题：}很多具有静噪功能的FM通信接收机在对方松开话筒PTT键后，会听到一声很明显的“嘶啦”或“喀拉”噪声拖尾，其原因是：
\begin{enumerate}[leftmargin=3em]
\item 此类电路根据鉴频输出中的强高音频噪声分量判断电台信号是否消失，从而关断音频输出。该项检测需占用一定时间，造成静噪的延迟，短时间漏出鉴频噪声
\item 该噪声由发射台话筒PTT接点跳动造成，发射到接收端
\item 该噪声是发射设备有目的地发射，作为结束发射的一种标志
\item 接收机自动增益电路的时间常数造成
\end{enumerate}
\noindent\textbf{解说：}\textbf{}。\\\noindent\textbf{答案：}

\bigskip




\noindent\textbf{试题：}超外差式接收机中混频器的作用是？
\begin{enumerate}[leftmargin=3em]
\item 将接收到的信号转换到中频
\item 将几个电台的信号混合在一起
\item 阻止通带以外不必要的信号
\item 将电台连接到一些辅助设备上，如TNC
\end{enumerate}
\noindent\textbf{解说：}\textbf{}。\\\noindent\textbf{答案：}

\bigskip




\noindent\textbf{试题：}一部业余无线电台，工作电压直流13.8伏，FM发射方式的射频输出载波功率为N瓦，电源效率约80%。发射时的工作电流约为：
\begin{enumerate}[leftmargin=3em]
\item 0.091×Ｎ（安）
\item 13.8×Ｎ（安）
\item 13.8/Ｎ×80%（安）
\item 0.058×Ｎ（安）
\end{enumerate}
\noindent\textbf{解说：}\textbf{}。\\\noindent\textbf{答案：}

\bigskip




\noindent\textbf{试题：}一部业余无线电台，工作电压交流220伏，FM发射方式的射频输出载波功率为N瓦，电源效率约80%。发射时的工作电流约为：
\begin{enumerate}[leftmargin=3em]
\item 0.0057×Ｎ（安）
\item 220×Ｎ（安）
\item 200/Ｎ×80%（安）
\item 0.0036×Ｎ（安）
\end{enumerate}
\noindent\textbf{解说：}\textbf{}。\\\noindent\textbf{答案：}

\bigskip




\noindent\textbf{试题：}一部业余无线电台，FM发射方式的射频输出载波功率为N瓦，电源效率约80%。通话时每发射10秒钟的电源消耗约为：
\begin{enumerate}[leftmargin=3em]
\item 0.0000035×Ｎ（千瓦小时）
\item 0.0768 /Ｎ（千瓦小时）
\item 0.0022×Ｎ（千瓦小时）
\item 220 / Ｎ（千瓦小时）
\end{enumerate}
\noindent\textbf{解说：}\textbf{}。\\\noindent\textbf{答案：}

\bigskip




\noindent\textbf{试题：}一部业余无线电台，FM发射方式的射频输出载波功率为10瓦，电源效率约80%。连续发话10秒钟，在此期间发射到空间的平均功率：
\begin{enumerate}[leftmargin=3em]
\item 约为10瓦
\item 约为12.5瓦
\item 约为8瓦
\item 肯定高于10瓦
\end{enumerate}
\noindent\textbf{解说：}\textbf{}。\\\noindent\textbf{答案：}

\bigskip




\noindent\textbf{试题：}一部业余无线电台，CW发射方式的射频输出载波功率为10瓦，电源效率约80%。连续发报10秒钟，在此期间发射到空间的平均功率：
\begin{enumerate}[leftmargin=3em]
\item 显著低于10瓦
\item 约为12.5瓦
\item 约为10瓦
\item 约为8瓦
\end{enumerate}
\noindent\textbf{解说：}\textbf{}。\\\noindent\textbf{答案：}

\bigskip




\noindent\textbf{试题：}一部业余无线电台，SSB发射方式的射频输出峰包功率为10瓦，电源效率约80%。连续发话10秒钟，在此期间发射到空间的平均功率：
\begin{enumerate}[leftmargin=3em]
\item 显著低于10瓦
\item 约为12.5瓦
\item 约为10瓦
\item 约为8瓦
\end{enumerate}
\noindent\textbf{解说：}\textbf{}。\\\noindent\textbf{答案：}

\bigskip




\noindent\textbf{试题：}一部业余无线电台，SSB发射方式的射频输出峰包功率为10瓦，电源效率约80%。在此方式下连续发送RTTY信号10秒钟，在此期间发射到空间的平均功率：
\begin{enumerate}[leftmargin=3em]
\item 约为10瓦
\item 显著低于10瓦
\item 约为12.5瓦
\item 约为8瓦
\end{enumerate}
\noindent\textbf{解说：}\textbf{}。\\\noindent\textbf{答案：}

\bigskip




\noindent\textbf{试题：}一部业余无线电台，SSB发射方式的射频输出峰包功率为10瓦，电源效率约80%。在此方式下连续发送PSK31信号10秒钟，在此期间发射到空间的平均功率：
\begin{enumerate}[leftmargin=3em]
\item 约为10瓦
\item 显著低于10瓦
\item 约为12.5瓦
\item 约为8瓦
\end{enumerate}
\noindent\textbf{解说：}\textbf{}。\\\noindent\textbf{答案：}

\bigskip




\noindent\textbf{试题：}一部业余无线电台，SSB发射方式的射频输出峰包功率为10瓦，电源效率约80%。在此方式下连续发送SSTV信号10秒钟，在此期间发射到空间的平均功率：
\begin{enumerate}[leftmargin=3em]
\item 约为10瓦
\item 显著低于10瓦
\item 约为12.5瓦
\item 约为8瓦
\end{enumerate}
\noindent\textbf{解说：}\textbf{}。\\\noindent\textbf{答案：}

\bigskip




\noindent\textbf{试题：}从能量转换的观点，“匹配”是指：
\begin{enumerate}[leftmargin=3em]
\item 选择电路参数，使负载能够得到最高实际输出功率的状态
\item 选择电路参数，使电源内阻能够达到最小功率损耗的状态
\item 选择电路参数，使负载能够得到最高实际输出电压的状态
\item 选择电路参数，使负载能够得到最高实际输出电流的状态
\end{enumerate}
\noindent\textbf{解说：}\textbf{}。\\\noindent\textbf{答案：}

\bigskip




\noindent\textbf{试题：}信号源和负载达到阻抗“匹配”时，信号源内阻损耗的功率为：
\begin{enumerate}[leftmargin=3em]
\item 与负载得到的输出功率相等
\item 是负载得到的输出功率的一半
\item 与负载得到的输出功率相比可以忽略不计
\item 损耗功率为零
\end{enumerate}
\noindent\textbf{解说：}\textbf{}。\\\noindent\textbf{答案：}

\bigskip




\noindent\textbf{试题：}一个放大器具有20dB的信号增益，其意义是：
\begin{enumerate}[leftmargin=3em]
\item 放大器把相当于输入信号的100倍的能量从电源转移到了输出负载
\item 放大器产生了当于输入信号的100倍的能量，供给了输出负载
\item 放大器把相输入信号的能量放大了100倍
\item 放大器把相输入信号的能量放大了99倍
\end{enumerate}
\noindent\textbf{解说：}\textbf{}。\\\noindent\textbf{答案：}

\bigskip




\noindent\textbf{试题：}射频信号通过某电路时产生了20dB的损耗。这部分被损耗的能量：
\begin{enumerate}[leftmargin=3em]
\item 在电路中被转化为热能等其他形式，或者通过电磁辐射等转移到了其他地方
\item 在电路中消失了
\item 返回了信号源
\item 一部分在电路中消失了，一部分返回了信号源
\end{enumerate}
\noindent\textbf{解说：}\textbf{}。\\\noindent\textbf{答案：}

\bigskip




\noindent\textbf{试题：}某电路输出信号功率是输入信号功率的100倍，该电路的增益为：
\begin{enumerate}[leftmargin=3em]
\item 20dB
\item 10dB
\item 100dB
\item 1dB
\end{enumerate}
\noindent\textbf{解说：}\textbf{}。\\\noindent\textbf{答案：}

\bigskip




\noindent\textbf{试题：}某电路输出信号功率是输入信号功率的100万倍，该电路的增益为：
\begin{enumerate}[leftmargin=3em]
\item 60dB
\item 100dB
\item 99万dB
\item 100万dB
\end{enumerate}
\noindent\textbf{解说：}\textbf{}。\\\noindent\textbf{答案：}

\bigskip




\noindent\textbf{试题：}某电路输出信号功率是输入信号功率的5倍，该电路的增益约为：
\begin{enumerate}[leftmargin=3em]
\item 7dB
\item 3.5dB
\item 5dB
\item 14dB
\end{enumerate}
\noindent\textbf{解说：}\textbf{}。\\\noindent\textbf{答案：}

\bigskip




\noindent\textbf{试题：}某电路输出信号功率是输入信号功率的2倍，该电路的增益约为：
\begin{enumerate}[leftmargin=3em]
\item 3dB
\item 1dB
\item 2dB
\item 0.5dB
\end{enumerate}
\noindent\textbf{解说：}\textbf{}。\\\noindent\textbf{答案：}

\bigskip




\noindent\textbf{试题：}某电路输出信号电压是输入信号电压的100倍，该电路的增益为：
\begin{enumerate}[leftmargin=3em]
\item 40dB
\item 10dB
\item 100dB
\item 20dB
\end{enumerate}
\noindent\textbf{解说：}\textbf{}。\\\noindent\textbf{答案：}

\bigskip




\noindent\textbf{试题：}某电路输出信号电压是输入信号电压的1万倍，该电路的增益为：（”x＾m”表示“x的m次方”）
\begin{enumerate}[leftmargin=3em]
\item 80dB
\item 10,000 dB
\item 9,999 dB
\item 10＾4dB
\end{enumerate}
\noindent\textbf{解说：}\textbf{}。\\\noindent\textbf{答案：}

\bigskip




\noindent\textbf{试题：}某电路输出信号电压是输入信号电压的10倍，该电路的增益约为：
\begin{enumerate}[leftmargin=3em]
\item 20dB
\item 7dB
\item 14dB
\item 15dB
\end{enumerate}
\noindent\textbf{解说：}\textbf{}。\\\noindent\textbf{答案：}

\bigskip




\noindent\textbf{试题：}某电路输出信号电压是输入信号电压的2倍，该电路的增益约为：
\begin{enumerate}[leftmargin=3em]
\item 6dB
\item 2dB
\item 4dB
\item 0.5dB
\end{enumerate}
\noindent\textbf{解说：}\textbf{}。\\\noindent\textbf{答案：}

\bigskip




\noindent\textbf{试题：}某电路输出信号功率是输入信号功率的1/100，该电路的增益为：
\begin{enumerate}[leftmargin=3em]
\item -20dB
\item -10dB
\item -100dB
\item 100dB
\end{enumerate}
\noindent\textbf{解说：}\textbf{}。\\\noindent\textbf{答案：}

\bigskip




\noindent\textbf{试题：}某电路输出信号功率是输入信号功率的百万分之一，该电路的增益为：
\begin{enumerate}[leftmargin=3em]
\item -60dB
\item -100dB
\item 990000dB
\item -1000000dB
\end{enumerate}
\noindent\textbf{解说：}\textbf{}。\\\noindent\textbf{答案：}

\bigskip




\noindent\textbf{试题：}某电路输出信号功率是输入信号功率的1/5，该电路的增益约为：
\begin{enumerate}[leftmargin=3em]
\item -7dB
\item 3.5dB
\item -5dB
\item -14dB
\end{enumerate}
\noindent\textbf{解说：}\textbf{}。\\\noindent\textbf{答案：}

\bigskip




\noindent\textbf{试题：}某电路输出信号功率是输入信号功率的1/2，该电路的增益约为：
\begin{enumerate}[leftmargin=3em]
\item -3dB
\item -1dB
\item -2dB
\item 0.5dB
\end{enumerate}
\noindent\textbf{解说：}\textbf{}。\\\noindent\textbf{答案：}

\bigskip




\noindent\textbf{试题：}某电路输出信号电压是输入信号电压的1/100，该电路的增益为：
\begin{enumerate}[leftmargin=3em]
\item -40dB
\item -10dB
\item -100dB
\item -20dB
\end{enumerate}
\noindent\textbf{解说：}\textbf{}。\\\noindent\textbf{答案：}

\bigskip




\noindent\textbf{试题：}某电路输出信号电压是输入信号电压的万分之一，该电路的增益为：（”x＾m”表示“x的m次方”）
\begin{enumerate}[leftmargin=3em]
\item -80dB
\item -10,000 dB
\item 1/10,000 dB
\item 10＾-4dB
\end{enumerate}
\noindent\textbf{解说：}\textbf{}。\\\noindent\textbf{答案：}

\bigskip




\noindent\textbf{试题：}某电路输出信号电压是输入信号电压的1/10倍，该电路的增益约为：
\begin{enumerate}[leftmargin=3em]
\item -20dB
\item -7dB
\item -14dB
\item 0.143dB
\end{enumerate}
\noindent\textbf{解说：}\textbf{}。\\\noindent\textbf{答案：}

\bigskip




\noindent\textbf{试题：}某电路输出信号电压是输入信号电压的1/2，该电路的增益约为：
\begin{enumerate}[leftmargin=3em]
\item -6dB
\item -2dB
\item 0.5dB
\item -0.5dB
\end{enumerate}
\noindent\textbf{解说：}\textbf{}。\\\noindent\textbf{答案：}

\bigskip




\noindent\textbf{试题：}信号依次通过增益分别为 x dB、y dB和 z dB的三个电路，总增益为：
\begin{enumerate}[leftmargin=3em]
\item （x + y + z）dB
\item （x × y × z）dB
\item （x + y + z）倍
\item （x × y × z）倍
\end{enumerate}
\noindent\textbf{解说：}\textbf{}。\\\noindent\textbf{答案：}

\bigskip




\noindent\textbf{试题：}信号依次通过增益分别为 x dB、y dB和 z dB的三个电路，总增益为：（”x＾m”表示“x的m次方”）
\begin{enumerate}[leftmargin=3em]
\item 10＾(（x + y + z）/10) 倍
\item 10＾（x × y × z） 倍
\item （x + y + z）倍
\item （x × y × z）倍
\end{enumerate}
\noindent\textbf{解说：}\textbf{}。\\\noindent\textbf{答案：}

\bigskip




\noindent\textbf{试题：}接收机的接收信号强度表每两档的信号强度相差6dB。接收某电台信号，发射功率为20dBW时读数为S9。该台减小发射功率后，接收机读数变为S4。此时该台的发射功率约为（以W为单位）：
\begin{enumerate}[leftmargin=3em]
\item 0.098W
\item 10.24W
\item 0.156W
\item 1.73W
\end{enumerate}
\noindent\textbf{解说：}\textbf{}。\\\noindent\textbf{答案：}

\bigskip




\noindent\textbf{试题：}接收机的接收信号强度表每两档的信号强度相差6dB。接收某电台信号，发射功率为10dBW时读数为S8。该台减小发射功率后，接收机读数变为S5。此时该台的发射功率约为（以W为单位）：
\begin{enumerate}[leftmargin=3em]
\item 0.156W
\item 0.098W
\item 10.24W
\item 1.73W
\end{enumerate}
\noindent\textbf{解说：}\textbf{}。\\\noindent\textbf{答案：}

\bigskip




\noindent\textbf{试题：}功率为0 dBm的射频信号通过增益为 23 dB的电路后，输出功率为：
\begin{enumerate}[leftmargin=3em]
\item 0.2W
\item 23mW
\item 23W
\item 0.23W
\end{enumerate}
\noindent\textbf{解说：}\textbf{}。\\\noindent\textbf{答案：}

\bigskip




\noindent\textbf{试题：}功率为0dBμ的射频信号通过增益为 36 dB的电路后，输出功率为：
\begin{enumerate}[leftmargin=3em]
\item 4mW
\item 3.6W
\item 36mW
\item 360μW
\end{enumerate}
\noindent\textbf{解说：}\textbf{}。\\\noindent\textbf{答案：}

\bigskip




\noindent\textbf{试题：}功率为0 dBW的射频信号通过增益为 -36 dB的电路后，输出功率为：
\begin{enumerate}[leftmargin=3em]
\item 0.25 mW
\item 3.6mW
\item 25mW
\item 360μW
\end{enumerate}
\noindent\textbf{解说：}\textbf{}。\\\noindent\textbf{答案：}

\bigskip




\noindent\textbf{试题：}功率为0 dBW的射频信号通过衰减量为 40 dB的衰减器后，输出功率为：
\begin{enumerate}[leftmargin=3em]
\item 100μW
\item 40mW
\item 140μW
\item 0.40 W
\end{enumerate}
\noindent\textbf{解说：}\textbf{}。\\\noindent\textbf{答案：}

\bigskip




\noindent\textbf{试题：}某业余电台以100瓦功率发射时，对方接收机的信号强度指示为S8。现双方天线不变，将发射功率降到25瓦，对方接收的信号强度指示将变为：【提示：收信机信号强度指示S1至S9每档相差6dB】
\begin{enumerate}[leftmargin=3em]
\item S7
\item S6
\item S5
\item S4
\end{enumerate}
\noindent\textbf{解说：}\textbf{}。\\\noindent\textbf{答案：}

\bigskip




\noindent\textbf{试题：}某业余电台以80瓦功率发射时，对方接收机的信号强度指示为S8。现双方天线不变，将发射功率降为5瓦QRP，对方接收的信号强度指示将变为：【提示：收信机信号强度指示S1至S9每档相差6dB】
\begin{enumerate}[leftmargin=3em]
\item S6
\item S7
\item S4
\item S2
\end{enumerate}
\noindent\textbf{解说：}\textbf{}。\\\noindent\textbf{答案：}

\bigskip




\noindent\textbf{试题：}在特定方向上具有主辐射瓣的水平偶极天线，其振子的总长度应为：
\begin{enumerate}[leftmargin=3em]
\item 1/2波长的奇数倍
\item 1/2波长的任意整数倍
\item 1/2波长的偶数倍
\item 1/4波长的奇数倍
\end{enumerate}
\noindent\textbf{解说：}\textbf{}。\\\noindent\textbf{答案：}

\bigskip




\noindent\textbf{试题：}偶极天线与工作频率发生谐振的充分和必要条件是：
\begin{enumerate}[leftmargin=3em]
\item 两臂总电气长度为1/2工作波长的整数倍
\item 两臂总电气长度为1/4工作波长的整数倍
\item 两臂总电气长度为1/2工作波长的奇数倍
\item 两臂总电气长度为工作波长的整数倍
\end{enumerate}
\noindent\textbf{解说：}\textbf{}。\\\noindent\textbf{答案：}

\bigskip




\noindent\textbf{试题：}偶极天线两臂总长度选取下列电气长度时，在垂直于天线轴线方向的增益达到峰值：
\begin{enumerate}[leftmargin=3em]
\item 1/2工作波长的奇数倍
\item 1/2工作波长的偶数倍
\item 1/2工作波长的整数倍
\item 1/4工作波长的奇数倍
\end{enumerate}
\noindent\textbf{解说：}\textbf{}。\\\noindent\textbf{答案：}

\bigskip




\noindent\textbf{试题：}一副两臂电气长度各为1/4波长的偶极天线，断开中点，通过平衡-不平衡变换器进行馈电，测得谐振频率为f，输入阻抗为Z。如果总长度不变，但将断开的馈电点向一侧偏移1/8波长，最明显的影响将是：
\begin{enumerate}[leftmargin=3em]
\item 阻抗显著变大，相比之下谐振频率变化不大
\item 阻抗显著变小，相比之下谐振频率变化不大
\item 谐振频率显著变高，相比之下阻抗变化不大
\item 谐振频率显著变低，相比之下阻抗变化不大
\end{enumerate}
\noindent\textbf{解说：}\textbf{}。\\\noindent\textbf{答案：}

\bigskip




\noindent\textbf{试题：}制作工作频率为f（单位：兆赫兹）的某相控天线阵列需要长度为1/4波长的同轴电缆。其大致长度（单位：米）为： 
\begin{enumerate}[leftmargin=3em]
\item 48.8 / f
\item 149.8 / f
\item 75 / f
\item 71.3 / f
\end{enumerate}
\noindent\textbf{解说：}\textbf{}。\\\noindent\textbf{答案：}

\bigskip




\noindent\textbf{试题：}制作工作频率为f（单位：兆赫兹）的半波长偶极天线。每个振子的大致长度（单位：米）为：
\begin{enumerate}[leftmargin=3em]
\item 71.3 / f
\item 48.8 / f
\item 142.6 / f
\item 150 / f
\end{enumerate}
\noindent\textbf{解说：}\textbf{}。\\\noindent\textbf{答案：}

\bigskip




\noindent\textbf{试题：}南北走向的水平极化偶极天线，中点馈电，通过特性阻抗为50欧的电缆连接到输入/输出阻抗为50欧的收发信机，通信对象在东西方向。选择天线长度的原则是：
\begin{enumerate}[leftmargin=3em]
\item 天线臂长为四分之一波长的奇数倍时，通信效果肯定最好
\item 天线臂长为四分之一波长的偶数倍时，通信效果肯定最好
\item 只要天线与工作频率谐振，通信效果一定好
\item 只要电压驻波比达到最低，通信效果一定好
\end{enumerate}
\noindent\textbf{解说：}\textbf{}。\\\noindent\textbf{答案：}

\bigskip




\noindent\textbf{试题：}对称半波振子每一臂的长度为波长的：
\begin{enumerate}[leftmargin=3em]
\item 1/4倍
\item 1/2倍
\item 2倍
\item 4倍
\end{enumerate}
\noindent\textbf{解说：}\textbf{}。\\\noindent\textbf{答案：}

\bigskip




\noindent\textbf{试题：}对一个偶极子天线怎么做，才能让它的谐振频率升高一些？
\begin{enumerate}[leftmargin=3em]
\item 将振子截短一些
\item 在振子某部位串联一个线圈
\item 将振子加长一些
\item 在振子的两端加上电容帽
\end{enumerate}
\noindent\textbf{解说：}\textbf{}。\\\noindent\textbf{答案：}

\bigskip




\noindent\textbf{试题：}甲偶极天线的增益为6.15dBi，乙偶极天线的增益为1dBd。当两副天线按同样条件架设、用同样功率驱动时、在它们最大发射方向的同一远方地点接收时，两天线产生的信号功率的关系为：
\begin{enumerate}[leftmargin=3em]
\item 甲天线的信号功率为乙天线的两倍
\item 甲天线的信号功率为乙天线的1/2
\item 甲天线的信号功率为乙天线的5.15倍
\item 甲天线的信号功率为乙天线的6.15倍
\end{enumerate}
\noindent\textbf{解说：}\textbf{}。\\\noindent\textbf{答案：}

\bigskip




\noindent\textbf{试题：}甲天线的增益为0dBd，乙天线的增益为2dBi。当两副天线按同样条件架设、用同样功率驱动时、在它们最大发射方向的同一远方地点接收并比较收到的信号功率强度，正确的说法为：
\begin{enumerate}[leftmargin=3em]
\item 甲天线的效果与半波长偶极天线相当，乙天线比甲天线略差。
\item 甲天线效果为零，不能工作，乙天线效果比甲天线好2倍。
\item 甲天线的效果与半波长偶极天线相当，乙天线发射的信号强度比甲天线好2dB。
\item 甲、乙天线的效果实际相同
\end{enumerate}
\noindent\textbf{解说：}\textbf{}。\\\noindent\textbf{答案：}

\bigskip




\noindent\textbf{试题：}业余条件测试天线增益的典型方法如图。用场强表或接收机接收设置在远处同一地点、最大辐射方向朝向自己的半波偶极天线（上）和待测天线（下）。调整送到两副天线的射频功率Po和P，使接收到的场强相同。待测天线的增益dBd值为：
\begin{enumerate}[leftmargin=3em]
\item 10 lg(Po/P)
\item 10 lg(P/Po)
\item P - Po
\item 10 lg(P-Po)
\end{enumerate}
\noindent\textbf{解说：}\textbf{}。\\\noindent\textbf{答案：}

\bigskip

LK0927.jpg


\noindent\textbf{试题：}业余条件测试天线增益的典型方法如图。用场强表或接收机接收设置在远处同一地点、最大辐射方向朝向自己的半波偶极天线（上）和待测天线（下）。调整送到两副天线的射频功率Po和P，使接收到的场强相同。待测天线的增益dBi值为：
\begin{enumerate}[leftmargin=3em]
\item 10 lg(Po/P) + 2.15
\item 10 lg(P/Po) + 2.15
\item 20 lg(P/Po) + 2.15
\item 10 lg(P-Po) + 2.15
\end{enumerate}
\noindent\textbf{解说：}\textbf{}。\\\noindent\textbf{答案：}

\bigskip

LK0928.jpg


\noindent\textbf{试题：}某业余电台使用半波长偶极天线发射时，对方接收机的信号强度指示为S4。现发射功率不变，发信端改用增益为 8.15 dBi的八木天线（最大辐射方向不变），对方接收的信号强度指示将变为：【提示：收信机信号强度指示S1至S9每档相差6dB】
\begin{enumerate}[leftmargin=3em]
\item S5
\item S6
\item S7
\item S8
\end{enumerate}
\noindent\textbf{解说：}\textbf{}。\\\noindent\textbf{答案：}

\bigskip




\noindent\textbf{试题：}某业余电台使用半波长偶极天线发射时，对方接收机的信号强度指示为S4。现发射功率不变，发信端改用增益为 12 dBd的八木天线（最大辐射方向不变），对方接收的信号强度指示将变为：【提示：收信机信号强度指示S1至S9每档相差6dB】
\begin{enumerate}[leftmargin=3em]
\item S6
\item S5
\item S7
\item S8
\end{enumerate}
\noindent\textbf{解说：}\textbf{}。\\\noindent\textbf{答案：}

\bigskip




\noindent\textbf{试题：}某业余电台使用半波长偶极天线发射时，对方亦使用半波长偶极天线接收，接收机的信号强度指示为S4。现发射功率不变，收发双方都改用增益为 8.15 dBi的八木天线（最大辐射方向不变），对方接收的信号强度指示将变为：【提示：收信机信号强度指示S1至S9每档相差6dB】
\begin{enumerate}[leftmargin=3em]
\item S6
\item S4
\item S5
\item S7
\end{enumerate}
\noindent\textbf{解说：}\textbf{}。\\\noindent\textbf{答案：}

\bigskip




\noindent\textbf{试题：}甲、乙业余电台相距2000公里，均使用1/2波长水平偶极天线，正在HF频段进行稳定的通信。现其中一方改用1/2波长垂直偶极天线，改变前后的通信效果的比较将是：
\begin{enumerate}[leftmargin=3em]
\item 通信效果的变化不确定，取决于当时天波反射途中极化方向的旋转情况
\item 通信效果变差
\item 通信效果变好
\item 通信效果没有改变
\end{enumerate}
\noindent\textbf{解说：}\textbf{}。\\\noindent\textbf{答案：}

\bigskip




\noindent\textbf{试题：}甲、乙业余电台相距10公里，均使用1/2波长水平偶极天线，正在UHF频段进行稳定的通信。现其中一方改用1/2波长垂直偶极天线，改变前后的通信效果的比较将是：
\begin{enumerate}[leftmargin=3em]
\item 通信效果变差
\item 通信效果没有改变
\item 通信效果变好
\item 通信效果的变化不确定
\end{enumerate}
\noindent\textbf{解说：}\textbf{}。\\\noindent\textbf{答案：}

\bigskip




\noindent\textbf{试题：}甲、乙业余电台相距10公里，分别使用1/2波长水平和垂直偶极天线，正在UHF频段进行稳定的通信。现改为双方都用1/2波长垂直偶极天线，改变前后的通信效果的比较将是：
\begin{enumerate}[leftmargin=3em]
\item 通信效果变好
\item 通信效果没有改变
\item 通信效果变差
\item 通信效果的变化不确定
\end{enumerate}
\noindent\textbf{解说：}\textbf{}。\\\noindent\textbf{答案：}

\bigskip




\noindent\textbf{试题：}甲、乙业余电台相距10公里，分别使用1/2波长水平和垂直偶极天线，正在UHF频段进行稳定的通信。现改为双方都用1/2波长水平偶极天线，改变前后的通信效果的比较将是：
\begin{enumerate}[leftmargin=3em]
\item 通信效果变好
\item 通信效果没有改变
\item 通信效果变差
\item 通信效果的变化不确定
\end{enumerate}
\noindent\textbf{解说：}\textbf{}。\\\noindent\textbf{答案：}

\bigskip




\noindent\textbf{试题：}甲、乙业余电台相距10公里，分别使用左旋圆极化和右旋圆极化天线，互相对指，正在UHF频段进行通信。现改为双方都用左旋圆极化天线，比较改变前后的通信效果是： 
\begin{enumerate}[leftmargin=3em]
\item 通信效果变好
\item 通信效果没有改变
\item 通信效果变差
\item 通信效果的变化不确定
\end{enumerate}
\noindent\textbf{解说：}\textbf{}。\\\noindent\textbf{答案：}

\bigskip




\noindent\textbf{试题：}甲、乙业余电台相距10公里，分别使用左旋圆极化和右旋圆极化天线，互相对指，正在UHF频段进行通信。现改为双方都用右旋圆极化天线，改变前后的通信效果的比较将是： 
\begin{enumerate}[leftmargin=3em]
\item 通信效果变好
\item 通信效果没有改变
\item 通信效果变差
\item 通信效果的变化不确定
\end{enumerate}
\noindent\textbf{解说：}\textbf{}。\\\noindent\textbf{答案：}

\bigskip




\noindent\textbf{试题：}甲、乙业余电台相距10公里，分别使用左旋圆极化和半波长水平极化天线，互相对指，正在UHF频段进行通信。现乙台改用半波长垂直极化天线，改变前后的通信效果的比较将是： 
\begin{enumerate}[leftmargin=3em]
\item 通信效果没有改变
\item 通信效果变差
\item 通信效果变好
\item 通信效果的变化不确定
\end{enumerate}
\noindent\textbf{解说：}\textbf{}。\\\noindent\textbf{答案：}

\bigskip




\noindent\textbf{试题：}甲、乙业余电台相距10公里，分别使用左旋圆极化和半波长水平极化天线，正在UHF频段进行通信。现将甲台天线改为右旋圆极化天线，改变前后的通信效果的比较将是： 
\begin{enumerate}[leftmargin=3em]
\item 通信效果没有改变
\item 通信效果变差
\item 通信效果变好
\item 通信效果的变化不确定
\end{enumerate}
\noindent\textbf{解说：}\textbf{}。\\\noindent\textbf{答案：}

\bigskip




\noindent\textbf{试题：}甲、乙业余电台相距10公里，甲台使用朝向乙台的左旋圆极化天线，乙台使用半波长水平极化天线，正在UHF频段进行通信。现将乙台天线改为朝向甲台的右旋圆极化天线，改变前后的通信效果的比较将是： 
\begin{enumerate}[leftmargin=3em]
\item 通信效果变差
\item 通信效果变好
\item 通信效果没有改变
\item 通信效果的变化不确定
\end{enumerate}
\noindent\textbf{解说：}\textbf{}。\\\noindent\textbf{答案：}

\bigskip




\noindent\textbf{试题：}甲、乙业余电台相距10公里，甲台使用朝向乙台的左旋圆极化天线，乙台使用半波长水平极化天线，正在UHF频段进行通信。现将乙台天线改为朝向甲台的左旋圆极化天线，改变前后的通信效果的比较将是： 
\begin{enumerate}[leftmargin=3em]
\item 通信效果变好
\item 通信效果变差
\item 通信效果没有改变
\item 通信效果的变化不确定
\end{enumerate}
\noindent\textbf{解说：}\textbf{}。\\\noindent\textbf{答案：}

\bigskip




\noindent\textbf{试题：}甲、乙业余电台相距10公里，分别使用左旋圆极化和右旋圆极化天线，互相对指，正在UHF频段进行通信。现将双方天线均改为半波长水平极化天线，改变前后的通信效果的比较将是： 
\begin{enumerate}[leftmargin=3em]
\item 通信效果变好
\item 通信效果没有改变
\item 通信效果变差
\item 通信效果的变化不确定
\end{enumerate}
\noindent\textbf{解说：}\textbf{}。\\\noindent\textbf{答案：}

\bigskip




\noindent\textbf{试题：}垂直偶极天线所发射的无线电波的极化方式为：
\begin{enumerate}[leftmargin=3em]
\item 垂直极化波
\item 水平极化波
\item 右旋圆极化波
\item 左旋圆极化波
\end{enumerate}
\noindent\textbf{解说：}\textbf{}。\\\noindent\textbf{答案：}

\bigskip




\noindent\textbf{试题：}水平偶极天线所发射的无线电波的极化方式为：
\begin{enumerate}[leftmargin=3em]
\item 水平极化波
\item 垂直极化波
\item 右旋圆极化波
\item 左旋圆极化波
\end{enumerate}
\noindent\textbf{解说：}\textbf{}。\\\noindent\textbf{答案：}

\bigskip




\noindent\textbf{试题：}假设收发天线均采用半波长偶极天线。在地面电台之间的近距离通信中，发射天线和接收天线的最佳极化方式为：
\begin{enumerate}[leftmargin=3em]
\item 收发天线都处于垂直于两台连线的平面内并且极化方向互相一致
\item 收发天线都处于垂直于两台连线的平面内收发天线极化方向互相垂直
\item 收发天线的极化方向都平行于两台之间的连线
\item 发信天线垂直极化，收信天线的极化方向平行于两台之间的连线
\end{enumerate}
\noindent\textbf{解说：}\textbf{}。\\\noindent\textbf{答案：}

\bigskip




\noindent\textbf{试题：}假设收发天线均采用半波长偶极天线。在依靠电离层反射的远距离通信中，发射天线和接收天线的最佳极化方式为：
\begin{enumerate}[leftmargin=3em]
\item 不确定，根据具体传播情况而经常变化
\item 收发天线都处于垂直于两台连线的平面内并且极化方向互相一致
\item 收发天线都处于垂直于两台连线的平面内收发天线极化方向互相垂直
\item 收发天线的极化方向都平行于两台之间的连线
\end{enumerate}
\noindent\textbf{解说：}\textbf{}。\\\noindent\textbf{答案：}

\bigskip




\noindent\textbf{试题：}右旋（或顺时针）极化波[right-hand（clockwise）polarized wave]是指：在任何一个垂直于传播方向的固定平面上，顺着传播方向看去，其电场向量随时间向右（顺时针方向）旋转的椭圆极化波或圆极化波。从地面某业余电台看一颗业余卫星，发现从卫星到达该台的无线电波的电场是按顺时针方向旋转的，其极化方式为：
\begin{enumerate}[leftmargin=3em]
\item 左旋（或逆时针）椭圆极化波或圆极化波
\item 右旋（或顺时针）椭圆极化波或圆极化波
\item 垂直极化波
\item 水平极化波
\end{enumerate}
\noindent\textbf{解说：}\textbf{}。\\\noindent\textbf{答案：}

\bigskip




\noindent\textbf{试题：}在视距通信中，已知发射天线为指向接收点的左旋圆极化天线，接收天线的最佳极化方式为： 
\begin{enumerate}[leftmargin=3em]
\item 指向发射点的左旋圆极化
\item 指向发射点的右旋圆极化
\item 垂直极化
\item 水平极化
\end{enumerate}
\noindent\textbf{解说：}\textbf{}。\\\noindent\textbf{答案：}

\bigskip




\noindent\textbf{试题：}卫星采用右旋圆极化天线，从北向南飞行，天线一直指向地球的南极。地面业余电台采用自动跟踪的圆极化天线指向卫星。地面业余电台天线的最佳极化旋转方向为： 
\begin{enumerate}[leftmargin=3em]
\item 卫星过顶前为右旋圆极化，过顶后为左旋圆极化
\item 卫星过顶后为右旋圆极化，过顶前为左旋圆极化
\item 最佳方向始终为右旋圆极化
\item 最佳方向始终为左旋圆极化
\end{enumerate}
\noindent\textbf{解说：}\textbf{}。\\\noindent\textbf{答案：}

\bigskip




\noindent\textbf{试题：}在视距通信中，已知发射天线是处于垂直于两台连线的平面内的偶极天线，但极化方向不断变化。能够保证在任何时刻信号都不会因极化问题而中断通信的接收天线极化方式为：
\begin{enumerate}[leftmargin=3em]
\item 指向发射点、右旋或左旋圆极化
\item 垂直极化
\item 水平极化
\item 极化方向平行于两台之间的连线
\end{enumerate}
\noindent\textbf{解说：}\textbf{}。\\\noindent\textbf{答案：}

\bigskip




\noindent\textbf{试题：}下列哪一项描述了一个导线与地球表面平行架设的偶极子天线的性质？
\begin{enumerate}[leftmargin=3em]
\item 是水平极化的天线
\item 是地波天线
\item 是菱形天线
\item 是垂直极化的天线
\end{enumerate}
\noindent\textbf{解说：}\textbf{}。\\\noindent\textbf{答案：}

\bigskip




\noindent\textbf{试题：}短波水平偶极类天线（如偶极天线和八木天线等）的发射仰角主要由下列因素决定：
\begin{enumerate}[leftmargin=3em]
\item 由天线的辐射和大地的反射叠加造成，仰角高低与天线离地高度与波长的比值有关
\item 由天线振子导体所指的方向决定
\item 由八木天线主梁所指的方向决定
\item 由天线振子的长度所决定
\end{enumerate}
\noindent\textbf{解说：}\textbf{}。\\\noindent\textbf{答案：}

\bigskip




\noindent\textbf{试题：}架设短波天线时，天线发射仰角的大致选择原则是：
\begin{enumerate}[leftmargin=3em]
\item 远距离通信选择低发射仰角，近距离通信选择高发射仰角
\item 近距离通信选择低发射仰角，远距离通信选择高发射仰角
\item 近处开阔时选择低发射仰角，近处有建筑物时选择高发射仰角
\item 较低频率通信选择低发射仰角，较高频率通信选择高发射仰角
\end{enumerate}
\noindent\textbf{解说：}\textbf{}。\\\noindent\textbf{答案：}

\bigskip




\noindent\textbf{试题：}架设短波天线时，天线高度的大致选择原则是：
\begin{enumerate}[leftmargin=3em]
\item 远距离通信选择较高的高度，近距离通信选择较低的高度
\item 远距离通信选择较低的高度，近距离通信选择较高的高度
\item 近处有建筑物时选择较低的高度，近处开阔时选择较高的高度
\item 较低频率通信选择较高的高度，较高频率通信选择较低的高度
\end{enumerate}
\noindent\textbf{解说：}\textbf{}。\\\noindent\textbf{答案：}

\bigskip




\noindent\textbf{试题：}在针对特定对象的DX通信中，计算天线最佳发射仰角的基本方法是：
\begin{enumerate}[leftmargin=3em]
\item 根据所使用电离层的大致高度、通信对象的大致距离、电波在传播途经中经电离层反射的次数，用简单几何方法计算
\item 根据通信对象所在的方位、地球半径、对方天线高度、实际工作频率、太阳平均黑子数，查表计算
\item 根据通信对象所在的方位、通信方向上障碍物所遮挡的仰角、本台周围的大地导电率、实际工作频率，找公式计算
\item 根据通信双方的发射功率、天线极化方向、通信方向上障碍物所遮挡的仰角、太阳10.7 cm 射电流量，找公式计算
\end{enumerate}
\noindent\textbf{解说：}\textbf{}。\\\noindent\textbf{答案：}

\bigskip




\noindent\textbf{试题：}通过目视判断全尺寸八木天线发射方向的办法是：
\begin{enumerate}[leftmargin=3em]
\item 比主振子短者为引向振子，比主振子长者为反射振子，引向振子朝向最大辐射方向
\item 比主振子长者为引向振子，比主振子短者为反射振子，引向振子朝向最大辐射方向
\item 比主振子短者为引向振子，比主振子长者为反射振子，反射振子朝向最大辐射方向
\item 主振子两端所指方向为最大辐射方向
\end{enumerate}
\noindent\textbf{解说：}\textbf{}。\\\noindent\textbf{答案：}

\bigskip




\noindent\textbf{试题：}北京的水平极化半波长偶极天线，通信对象为纽约的业余电台。按电波的最短传输途径考虑，天线的最佳走向应大致为：
\begin{enumerate}[leftmargin=3em]
\item 东-西
\item 南-北
\item 东偏南30度-西偏北30度
\item 西偏南30度-东偏北30度
\end{enumerate}
\noindent\textbf{解说：}\textbf{}。\\\noindent\textbf{答案：}

\bigskip




\noindent\textbf{试题：}在导电良好的地面上，决定短波天线辐射仰角的主要参数是：
\begin{enumerate}[leftmargin=3em]
\item 天线离地面的相对于波长的高度，即离地高度除以波长
\item 天线的绝对高度，与波长无关
\item 天线离海平面的绝对高度
\item 天线导线或者八木天线主梁与地面之间的夹角
\end{enumerate}
\noindent\textbf{解说：}\textbf{}。\\\noindent\textbf{答案：}

\bigskip




\noindent\textbf{试题：}什么是八木天线？
\begin{enumerate}[leftmargin=3em]
\item 一种可以集中聚集某一方向信号的天线
\item 一种由日本八木秀次发明的全向天线
\item 用八根木头制作的天线
\item 一种可以将接收到的信号反向的天线
\end{enumerate}
\noindent\textbf{解说：}\textbf{}。\\\noindent\textbf{答案：}

\bigskip




\noindent\textbf{试题：}在自由空间中的半波偶极子天线，哪个方向的辐射强度最大？
\begin{enumerate}[leftmargin=3em]
\item 垂直于导体的方向
\item 沿着导体的方向
\item 各个方向强度相同
\item 沿着馈线的方向
\end{enumerate}
\noindent\textbf{解说：}\textbf{}。\\\noindent\textbf{答案：}

\bigskip




\noindent\textbf{试题：}已知某天线的增益为3.85 dBd，馈入的功率为10W，其等效全向辐射功率（e.i.r.p）为：
\begin{enumerate}[leftmargin=3em]
\item 40W
\item 13.85W
\item 38.5W
\item 100W
\end{enumerate}
\noindent\textbf{解说：}\textbf{}。\\\noindent\textbf{答案：}

\bigskip




\noindent\textbf{试题：}《中华人民共和国无线电频率划分规定》脚注5.67A 说“使用135.7-137.8kHz频段内频率的业余业务台站，其最大有效辐射功率不得超过1瓦（e.i.r.p.）”。其意思为，该频率的业余电台：
\begin{enumerate}[leftmargin=3em]
\item 实际馈送到发射天线的功率不得超过1瓦除以天线的绝对增益比值Gi后得到的功率值
\item 发射机标称输出功率不得超过1瓦
\item 扣除馈线损耗后实际馈送到发射天线的功率不得超过1瓦
\item 扣除馈线和天线的总损耗后天线实际辐射的功率最大不得超过1瓦
\end{enumerate}
\noindent\textbf{解说：}\textbf{}。\\\noindent\textbf{答案：}

\bigskip




\noindent\textbf{试题：}《中华人民共和国无线电频率划分规定》脚注5.67A 说“使用135.7-137.8kHz频段内频率的业余业务台站，其最大辐射功率不得超过1瓦（e.i.r.p.）”。其意思为，该频率的业余电台：
\begin{enumerate}[leftmargin=3em]
\item 实际馈送到发射天线的功率不得超过1瓦除以天线的绝对增益比值Gi后得到的功率值
\item 发射机标称输出功率不得超过1瓦
\item 扣除馈线损耗后实际馈送到发射天线的功率不得超过1瓦
\item 扣除馈线和天线的总损耗后天线实际辐射的功率最大不得超过1瓦
\end{enumerate}
\noindent\textbf{解说：}\textbf{}。\\\noindent\textbf{答案：}

\bigskip




\noindent\textbf{试题：}已知某天线的增益为5.15 dBi，馈入的功率为10W，其有效辐射功率（e.r.p）为：
\begin{enumerate}[leftmargin=3em]
\item 20W
\item 51.5W
\item 10W
\item 15.15W
\end{enumerate}
\noindent\textbf{解说：}\textbf{}。\\\noindent\textbf{答案：}

\bigskip




\noindent\textbf{试题：}在HF远距离通信中，收发信天线周围的大地导电率对辐射效果的影响大致是：
\begin{enumerate}[leftmargin=3em]
\item 大地导电率越高，大地反射比较接近于理想情况，效果越好
\item 大地导电率越低，大地中的感应电流越弱，效果越好
\item HF远距离通信依靠天波反射，与大地导电率无关
\item 大地导电率太高、太低都不好，最好是处于中间值
\end{enumerate}
\noindent\textbf{解说：}\textbf{}。\\\noindent\textbf{答案：}

\bigskip




\noindent\textbf{试题：}在因空间限制、HF偶极天线达不到自然谐振所需长度而必须缩短时，可以在天线电路的某一点串入一个线圈，或者将整个振子做成一个线圈，加以补偿，构成“加感天线”。在充分优化设计的条件下，它们与全尺寸半波长偶极天线相比较的结果是：
\begin{enumerate}[leftmargin=3em]
\item 可以和全尺寸天线一样满足谐振和阻抗匹配条件，但辐射效率肯定不如全尺寸天线
\item 无论是谐振和阻抗匹配条件还是辐射效率，都不能得到全尺寸天线的效果
\item 只要设计得当、制作精心、妥善调试，可以得到和全尺寸天线完全一样的相同效果
\item 如果设计得当、制作精心、妥善调试，有可能得到比全尺寸偶极天线更好的效果
\end{enumerate}
\noindent\textbf{解说：}\textbf{}。\\\noindent\textbf{答案：}

\bigskip




\noindent\textbf{试题：}行波天线是一种依靠射频电流在长导线组成的连续匹配回路中流动产生辐射的天线，业余通信常用的有T2FD天线和贝伐列奇天线等。与驻波天线相比，行波天线的特点有：
\begin{enumerate}[leftmargin=3em]
\item 是没有自然谐振点的宽带天线，对长度要求不严格但一般需接近或大于波长，
\item 对长度要求比较严格，是没有自然谐振点的宽带天线
\item 在较宽的频带内具有一些列自然谐振频率点，但对长度要求比较严格
\item 可以依靠天线调谐器实现在较宽的频带内谐振，对长度没有要求
\end{enumerate}
\noindent\textbf{解说：}\textbf{}。\\\noindent\textbf{答案：}

\bigskip




\noindent\textbf{试题：}业余通信常用的T2FD行波天线和贝伐列奇等行波天线都在天线的终端接了一个无感电阻。关于它的作用，正确说法是：
\begin{enumerate}[leftmargin=3em]
\item 长度足够时大部分功率边行进边辐射到空间，只剩小部分到达电阻并消耗而不反射回去
\item 电阻相当于一个假负载，虽实现了宽带工作，但绝大部分能量消耗在这个电阻上
\item 在匹配状态下，输入功率的一半消耗在电阻上，所以天线的辐射效率大约是50%
\item 串联电阻的损耗降低了天线谐振时的Q值，使谐振响应显得平坦，但辐射效率比较低下
\end{enumerate}
\noindent\textbf{解说：}\textbf{}。\\\noindent\textbf{答案：}

\bigskip




\noindent\textbf{试题：}图示为我国业余电台野外通信常见的14MHz/21MHz/29MHz三频段偶极天线原理图。其中各段天线的电气长度应分别为：
\begin{enumerate}[leftmargin=3em]
\item A为28MHz的1/4波长，A+B为21MHz的1/4波长，A+B+C为14MHz的1/4波长
\item A为28MHz的半波长，A+B为21MHz的半波长，A+B+C为14MHz的半波长
\item A为28MHz的1/4波长，B为21MHz的1/4波长，C为14MHz的1/4波长
\item A为28MHz的半波长，B为21MHz的半波长，C为14MHz的半波长
\end{enumerate}
\noindent\textbf{解说：}\textbf{}。\\\noindent\textbf{答案：}

\bigskip

LK0941.jpg


\noindent\textbf{试题：}图示为我国业余电台野外通信常见的14MHz/21MHz/29MHz三频段偶极天线原理图。其中部件A、B、C、D的作用分别为：
\begin{enumerate}[leftmargin=3em]
\item A、D为21MHz陷波器，B、C为28MHz陷波器，即分别谐振于21MHz和28MHz的并联谐振回路
\item A、D为21MHz陷波器，B、C为28MHz陷波器，即分别谐振于21MHz和28MHz的串联谐振回路
\item A、D为14MHz陷波器，B、C为21MHz陷波器，即分别谐振于14MHz和21MHz的并联谐振回路
\item A、D为14MHz陷波器，B、C为21MHz陷波器，即分别谐振于14MHz和21MHz的串联谐振回路
\end{enumerate}
\noindent\textbf{解说：}\textbf{}。\\\noindent\textbf{答案：}

\bigskip

LK0942.jpg


\noindent\textbf{试题：}水平偶极天线振子的最短理想长度是工作频率的半波长。因条件限制而必须缩短振子时，不得不在天线的某个地方串入额外的电感将振子的自然谐振频率补偿到工作频率，但天线的效率比半波长天线有所降低。图为常用的方案，如果振子（红色部分）总长度相等，各方案的名称以及按发射效率由高到低的排序为： 
\begin{enumerate}[leftmargin=3em]
\item C-顶部加感,B-中间加感,A-底部加感
\item A-底部加感,B-中间加感, C-顶部加感
\item A-底部加感,C-顶部加感, B-中间加感
\item B-中间加感, A-底部加感, C-顶部加感
\end{enumerate}
\noindent\textbf{解说：}\textbf{}。\\\noindent\textbf{答案：}

\bigskip

LK0943.jpg


\noindent\textbf{试题：}垂直GP天线振子的最短理想长度是工作频率的1/4波长。因条件限制而必须缩短振子时，不得不在天线的某个地方串入额外的电感将振子的自然谐振频率补偿到工作频率，但天线的效率比1/4波长天线有所降低。图为常用的方案，如果振子（红色部分）总长度相等，各方案的名称以及按发射效率由高到低的排序为： 
\begin{enumerate}[leftmargin=3em]
\item C-顶部加感, B-中间加感, A-底部加感
\item A-底部加感, B-中间加感, C-顶部加感
\item A-底部加感, C-顶部加感, B-中间加感
\item B-中间加感, A-底部加感, C-顶部加感
\end{enumerate}
\noindent\textbf{解说：}\textbf{}。\\\noindent\textbf{答案：}

\bigskip

LK0944.jpg


\noindent\textbf{试题：}业余和专业无线电通信常用工作于不同频率的多副偶极天线的馈电点并联在一起构成多频段天线。对图示的由7MHz、14MHz和28MHz半波长振子组成的多频段天线的分析为：
\begin{enumerate}[leftmargin=3em]
\item 7MHz振子可以以3/2波长方式工作于21MHz，可在7/14/21/28MHz四个频段工作；不工作的振子呈现高阻抗，并联后总阻抗为50欧左右
\item 只能工作于7/14/28MHz三个频段。不工作的振子呈现高阻抗；并联后的总阻抗为50欧左右
\item 7MHz振子可以3/2波长方式工作于21MHz，可在7/14/21/28MHz四个频段工作；每个偶极天线的阻抗都是50欧，并联后总阻抗为17欧左右
\item 只能工作于7/14/28MHz三个频段；每个偶极天线的阻抗都是50欧，并联后总阻抗为17欧左右
\end{enumerate}
\noindent\textbf{解说：}\textbf{}。\\\noindent\textbf{答案：}

\bigskip

LK0945.jpg


\noindent\textbf{试题：}无线电发射设备和市电网的发电机都属于交流电源。接到无线电发射设备的平行导线上可能会出现电压和电流的驻波，造成各处电压不同，而接到市电网上的电线却看不到这种现象。原因是：
\begin{enumerate}[leftmargin=3em]
\item 市电频率很低，导线长度的差别与波长相比微不足道，驻波现象很不明显而已
\item 无线电设备所接的馈线和市电网用的电线性能不同，所以现象不同
\item 市电网的发电能力远大于无线电发射设备，致使电路各处电压恒定
\item 供电工程和无线电工程所适用的电学理论基础不相同
\end{enumerate}
\noindent\textbf{解说：}\textbf{}。\\\noindent\textbf{答案：}

\bigskip




\noindent\textbf{试题：}同轴电缆的绝缘介质相同时，影响特性阻抗的因素是：
\begin{enumerate}[leftmargin=3em]
\item 外导体内径和内导体外径的比越大，特性阻抗越高
\item 外导体内径越大，特性阻抗越高
\item 内导体外径越大，特性阻抗越高
\item 电缆的长度越长，特性阻抗越高
\end{enumerate}
\noindent\textbf{解说：}\textbf{}。\\\noindent\textbf{答案：}

\bigskip




\noindent\textbf{试题：}一根特性阻抗为50欧的同轴电缆，一端接一个50欧电阻负载，用天线阻抗分析仪测量另一端在某工作频率下的实际阻抗。
\begin{enumerate}[leftmargin=3em]
\item 读数为50欧，与电缆长度无关
\item 读数取决于电缆长度，长1/4波长的奇数倍时接近无穷大，1/4波长偶数倍时接近零
\item 读数取决于电缆长度，长1/4波长的偶数倍时接近无穷大，1/4波长奇数倍时接近零
\item 读数取决于负载电阻的额定功率，额定功率越小，读数越大
\end{enumerate}
\noindent\textbf{解说：}\textbf{}。\\\noindent\textbf{答案：}

\bigskip




\noindent\textbf{试题：}一根特性阻抗为50欧的同轴电缆，一端开路，用天线阻抗分析仪测量另一端在某工作频率下的实际阻抗。
\begin{enumerate}[leftmargin=3em]
\item 读数取决于电缆长度，长1/4波长的偶数倍时接近无穷大，1/4波长奇数倍时接近零
\item 读数取决于电缆长度，长1/4波长的奇数倍时接近无穷大，1/4波长偶数倍时接近零
\item 读数取决于负载电阻的额定功率，额定功率越小，读数越大
\item 读数为50欧，与电缆长度无关
\end{enumerate}
\noindent\textbf{解说：}\textbf{}。\\\noindent\textbf{答案：}

\bigskip




\noindent\textbf{试题：}一根特性阻抗为50欧的同轴电缆，一端芯线和外屏蔽层短路，用天线阻抗分析仪测量另一端在某工作频率下的实际阻抗。
\begin{enumerate}[leftmargin=3em]
\item 读数取决于电缆长度，长1/4波长的奇数倍时接近无穷大，1/4波长偶数倍时接近零
\item 读数取决于电缆长度，长1/4波长的偶数倍时接近无穷大，1/4波长奇数倍时接近零
\item 读数取决于负载电阻的额定功率，额定功率越小，读数越大
\item 读数为50欧，与电缆长度无关
\end{enumerate}
\noindent\textbf{解说：}\textbf{}。\\\noindent\textbf{答案：}

\bigskip




\noindent\textbf{试题：}用天线阻抗分析仪测量一副失配天线在某工作频率下的实际阻抗，发现配用不同的电缆连接仪器和天线时，得到的读数相差很大。为保证读数准确，可采取以下措施：
\begin{enumerate}[leftmargin=3em]
\item 使用电气长度正好等于波长的连接电缆
\item 换用接触件镀金的高档电缆接头
\item 用高档纯银音响线材代替普通铜质同轴电缆
\item 将电缆外皮妥善接地，并将仪器放入屏蔽室
\end{enumerate}
\noindent\textbf{解说：}\textbf{}。\\\noindent\textbf{答案：}

\bigskip




\noindent\textbf{试题：}如果在驻波表上读到了4:1，这意味着？
\begin{enumerate}[leftmargin=3em]
\item 阻抗匹配得不好
\item 良好的阻抗匹配
\item 天线的增益是4dB
\item 天线有4dB的损失
\end{enumerate}
\noindent\textbf{解说：}\textbf{}。\\\noindent\textbf{答案：}

\bigskip




\noindent\textbf{试题：}同样材料、同样直径、同样长度的实心铜线和空心铜管，在交流电路中的发热损耗情况为：
\begin{enumerate}[leftmargin=3em]
\item 在低频率下实心铜线损耗较小，在高频率下两者损耗一样
\item 在低频率下实心铜线损耗较大，在高频率下两者损耗一样
\item 在不同频率两者的发热损耗大小不好说，取决于具体散热条件
\item 在各种频率下两者耗差都一样
\end{enumerate}
\noindent\textbf{解说：}\textbf{}。\\\noindent\textbf{答案：}

\bigskip




\noindent\textbf{试题：}把实心导线接到频率为数十兆赫兹的高频率射频电路中，则会有下列现象：
\begin{enumerate}[leftmargin=3em]
\item 电流集中在导线表层，导线内部没有电流
\item 导线截面各处的电流密度均匀分布
\item 导线外层和内层都有电流，但两者方向相反
\item 表层的电流沿导线方向流动，内层电流形成螺旋状涡流
\end{enumerate}
\noindent\textbf{解说：}\textbf{}。\\\noindent\textbf{答案：}

\bigskip




\noindent\textbf{试题：}工作在高频率下的射频部件积灰或受潮后，即使没有漏电，也可能因绝缘物体的物理变化而带来额外的：
\begin{enumerate}[leftmargin=3em]
\item 介质损耗
\item 磁滞损耗
\item 涡流损耗
\item 磁阻损耗
\end{enumerate}
\noindent\textbf{解说：}\textbf{}。\\\noindent\textbf{答案：}

\bigskip




\noindent\textbf{试题：}天线和馈线之间经常接一个俗称“巴伦（BALUN）”的部件。“巴伦”的由来是：
\begin{enumerate}[leftmargin=3em]
\item 平衡和不平衡两个英文字头的组合
\item 发明平衡不平衡转换器的人的名字
\item 著名天线阻抗匹配理论家的名字
\item 宽带阻抗变压器的英文缩写
\end{enumerate}
\noindent\textbf{解说：}\textbf{}。\\\noindent\textbf{答案：}

\bigskip




\noindent\textbf{试题：}天线和馈线之间经常接一个俗称“巴伦（BALUN）”的部件。它的主要功能是：
\begin{enumerate}[leftmargin=3em]
\item 在平衡电路和不平衡电路之间传递射频能量，并阻断两者之间的任何寄生耦合
\item 实现天线和馈线之间的自动阻抗匹配
\item 展宽天线的工作频带
\item 降低天线的驻波比
\end{enumerate}
\noindent\textbf{解说：}\textbf{}。\\\noindent\textbf{答案：}

\bigskip




\noindent\textbf{试题：}天线调谐器（俗称“天调”）的作用是：
\begin{enumerate}[leftmargin=3em]
\item 补偿不匹配系统，向收发信机提供谐振的、阻抗匹配的负载，但不能改善天线本身的辐射效率
\item 对不匹配的天馈系统进行补偿，使整个天馈系统的传输和辐射效率达到匹配天线的水平
\item 对天馈系统进行优化调谐，使整个系统的辐射功率获得一个附加的增益
\item 对不匹配的天馈系统进行补偿，虽然不能改善馈线的损耗，但能使天线本身的辐射效率达到匹配天线的水平
\end{enumerate}
\noindent\textbf{解说：}\textbf{}。\\\noindent\textbf{答案：}

\bigskip




\noindent\textbf{试题：}天线通过50欧同轴馈线与输出阻抗为50欧的收发信机相连接，并打算在天线电路中串入天线调谐器和通过式驻波功率计来监测和补偿天线的失配。理论上最理想的连接顺序为：
\begin{enumerate}[leftmargin=3em]
\item 天线-天线调谐器-驻波功率计-馈线-收发信机
\item 天线-馈线-天线调谐器-驻波功率计-收发信机
\item 天线-天线调谐器-馈线-驻波功率计-收发信机
\item 天线-驻波功率计-天线调谐器-馈线-收发信机
\end{enumerate}
\noindent\textbf{解说：}\textbf{}。\\\noindent\textbf{答案：}

\bigskip




\noindent\textbf{试题：}塔上的天线通过50欧同轴馈线与输出阻抗为50欧的收发信机相连接，在天线电路中串入天线调谐器ATU和通过式驻波功率计M来监测和补偿天线的失配。有四种方案：1、ATU和M均在塔顶，2、ATU和M均在塔底，3、ATU在塔底、M在机房，4、ATU和M均在机房。当ATU调到最佳状态时，各方案按天线系统发射效率由高到低的排序为：
\begin{enumerate}[leftmargin=3em]
\item 方案1最好，方案2、3其次，方案4最差
\item 方案4最好，方案2、3其次，方案1最差
\item 方案3最好，方案4其次，方案2再其次，方案1最差
\item 方案2最好，方案1其次，方案4再其次，方案3最差
\end{enumerate}
\noindent\textbf{解说：}\textbf{}。\\\noindent\textbf{答案：}

\bigskip

LK0938.jpg


\noindent\textbf{试题：}下列哪一种导体最适合射频接地使用？
\begin{enumerate}[leftmargin=3em]
\item 镀银软铜丝编织扁带
\item 圆形铜包钢单股线
\item 双绞线
\item 圆形多股线
\end{enumerate}
\noindent\textbf{解说：}\textbf{}。\\\noindent\textbf{答案：}

\bigskip




\noindent\textbf{试题：}收发信机天线调谐器（天调）的作用是什么？
\begin{enumerate}[leftmargin=3em]
\item 它将发射机的输出阻抗和天线的输入阻抗进行良好的匹配
\item 它帮助接收机在遇到弱信号时能自动调准频率
\item 它可以使一部天线既能作为发射天线，又能作为接收天线
\item 它能够依照发射机当前工作的波段自动连接合适的天线
\end{enumerate}
\noindent\textbf{解说：}\textbf{}。\\\noindent\textbf{答案：}

\bigskip




\noindent\textbf{试题：}电磁波在电能和磁能相互交换和贮存的过程中向四周传播。影响电磁波辐射场能量的因素有：
\begin{enumerate}[leftmargin=3em]
\item 工作频率越高、产生磁场的电流的大小越大，向四周传播的辐射场能量越多
\item 工作频率越低、产生磁场的电流的大小越大，向四周传播的辐射场能量越多
\item 产生磁场的电流的大小越大，向四周传播的辐射场能量越多，与工作频率无关
\item 工作频率与产生磁场的电流的比值越大，向四周传播的辐射场能量越多
\end{enumerate}
\noindent\textbf{解说：}\textbf{}。\\\noindent\textbf{答案：}

\bigskip




\noindent\textbf{试题：}自由空间中甲乙两台相距100公里，均采用增益为0dBi的天线，发射端功率1W，工作频率为145MHz，接收端天线得到的信号功率约为-115.6dBm。如将距离增加为500km，接收到的信号功率将为：【提示：无线电波的自由空间路径损耗为：L=32.4+20logR+20logF，其中R为发射端到接收端的距离（km），F为工作频率（MHz）】
\begin{enumerate}[leftmargin=3em]
\item -129.6 dBm
\item -101.6 dBm
\item -135.6 dBm
\item -95.6 dBm
\end{enumerate}
\noindent\textbf{解说：}\textbf{}。\\\noindent\textbf{答案：}

\bigskip




\noindent\textbf{试题：}自由空间中甲乙两台相距100公里，均采用增益为0dBi的天线，发射端功率1W，工作频率为145MHz，接收端天线得到的信号功率约为-115.6dBm。如将距离增加为1000km，接收到的信号功率将为：【提示：无线电波的自由空间路径损耗为：L=32.4+20logR+20logF，其中R为发射端到接收端的距离（km），F为工作频率（MHz）】
\begin{enumerate}[leftmargin=3em]
\item -135.6 dBm
\item -95.6 dBm
\item -129.6 dBm
\item -101.6 dBm
\end{enumerate}
\noindent\textbf{解说：}\textbf{}。\\\noindent\textbf{答案：}

\bigskip




\noindent\textbf{试题：}自由空间中甲乙两台相距100公里，均采用增益为0dBi的天线，发射端功率1W，工作频率为145MHz，接收端天线得到的信号功率约为-115.6dBm。如将频率改为435MHz，接收到的信号功率将约为：【提示：无线电波的自由空间路径损耗为：L=32.4+20logR+20logF，其中R为发射端到接收端的距离（km），F为工作频率（MHz）】
\begin{enumerate}[leftmargin=3em]
\item -125.1 dBm
\item -109.6 dBm
\item -129.6 dBm
\item -135.6 dBm
\end{enumerate}
\noindent\textbf{解说：}\textbf{}。\\\noindent\textbf{答案：}

\bigskip




\noindent\textbf{试题：}自由空间中甲乙两台相距100公里，均采用增益为0dBi的天线，发射端功率1W，工作频率为145MHz，接收端天线得到的信号功率约为-115.6dBm。如将发射功率减少为0.25W，接收到的信号功率将为：【提示：无线电波的自由空间路径损耗为：L=32.4+20logR+20logF，其中R为发射端到接收端的距离（km），F为工作频率（MHz）】
\begin{enumerate}[leftmargin=3em]
\item -121.6 dBm
\item -109.6 dBm
\item -129.6 dBm
\item -135.6 dBm
\end{enumerate}
\noindent\textbf{解说：}\textbf{}。\\\noindent\textbf{答案：}

\bigskip




\noindent\textbf{试题：}自由空间中相距100km的甲乙两台通信，工作频率145MHz，接收端的信号强度指示为S8。如将距离增加为500km，接收信号强度指示将变为：【提示：无线电波的自由空间路径损耗为：L=32.4+20logR+20logF，其中R为发射端到接收端的距离（km），F为工作频率（MHz）；收信机信号强度指示S1至S9每挡相差6dB】
\begin{enumerate}[leftmargin=3em]
\item 略低于S6
\item 略高于S6
\item 略低于S5
\item 略高于S4
\end{enumerate}
\noindent\textbf{解说：}\textbf{}。\\\noindent\textbf{答案：}

\bigskip




\noindent\textbf{试题：}自由空间中相距100km的甲乙两台通信，工作频率145MHz，接收端的信号强度指示为S8。如将距离增加为1000km，接收信号强度指示将变为：【提示：无线电波的自由空间路径损耗为：L=32.4+20logR+20logF，其中R为发射端到接收端的距离（km），F为工作频率（MHz）；收信机信号强度指示S1至S9每挡相差6dB】
\begin{enumerate}[leftmargin=3em]
\item 略低于S5
\item 略高于S5
\item 略高于S6
\item S7
\end{enumerate}
\noindent\textbf{解说：}\textbf{}。\\\noindent\textbf{答案：}

\bigskip




\noindent\textbf{试题：}自由空间中相距100km的甲乙两台通信，工作频率145MHz，接收端的信号强度指示为S8。如将频率改为435MHz，接收信号强度指示将变为：【提示：无线电波的自由空间路径损耗为：L=32.4+20logR+20logF，其中R为发射端到接收端的距离（km），F为工作频率（MHz）；收信机信号强度指示S1至S9每挡相差6dB】
\begin{enumerate}[leftmargin=3em]
\item 略高于S6
\item 略低于S5
\item 略低于S6
\item S7
\end{enumerate}
\noindent\textbf{解说：}\textbf{}。\\\noindent\textbf{答案：}

\bigskip




\noindent\textbf{试题：}自由空间中相距100km的甲乙两台通信，工作频率为145MHz，接收端的信号强度指示为S8。如将发射功率减少为原来的1/4，接收信号强度指示将变为：【提示：无线电波的自由空间路径损耗为：L=32.4+20logR+20logF，其中R为发射端到接收端的距离（km），F为工作频率（MHz）；收信机信号强度指示S1至S9每档相差6dB】
\begin{enumerate}[leftmargin=3em]
\item S7
\item 略高于S5
\item 略高于S6
\item 略低于S5
\end{enumerate}
\noindent\textbf{解说：}\textbf{}。\\\noindent\textbf{答案：}

\bigskip




\noindent\textbf{试题：}什么是波阻抗？
\begin{enumerate}[leftmargin=3em]
\item 电磁波传到远方时，电场E和磁场H在空间的比值为常量，称为波阻抗。
\item 射频电流在同轴电缆中通过时遇到的阻抗就是波阻抗
\item 射频电流在进入发射天线馈电点时遇到的阻抗就是波阻抗
\item 射频电流在发射设备输出端遇到的外界阻抗就是波阻抗
\end{enumerate}
\noindent\textbf{解说：}\textbf{}。\\\noindent\textbf{答案：}

\bigskip




\noindent\textbf{试题：}什么叫“理想点源天线”？对业余无线电爱好者有什么意义？	
\begin{enumerate}[leftmargin=3em]
\item 一种小到一个点的、能把发射机的射频能量全部转换为电磁波并加以各向同性均匀辐射的理论假想天线。经常作为衡量实际天线辐射性能的比较基准
\item 一种专业通信使用的增益极高的专用天线。在业余无线电中没有应用价值
\item 一种仅用于无线电测试的标准接收天线，但发射效果不好，对业余无线电无用处
\item 一种带宽几乎无穷宽的高级天线专利的名称。业余无线电不需要这样的宽带天线
\end{enumerate}
\noindent\textbf{解说：}\textbf{}。\\\noindent\textbf{答案：}

\bigskip




\noindent\textbf{试题：}真空中理想点源天线发出的电磁波是一个均匀球面波，任何一个与传播途径垂直的空间截面上的功率通量密度的分布规律为：
\begin{enumerate}[leftmargin=3em]
\item 功率通量密度与离辐射源的距离的平方成反比
\item 功率通量密度与离辐射源的距离成反比
\item 功率通量密度与离辐射源的距离成正比
\item 功率通量密度与离辐射源的距离成指数关系
\end{enumerate}
\noindent\textbf{解说：}\textbf{}。\\\noindent\textbf{答案：}

\bigskip




\noindent\textbf{试题：}自由空间中电磁波随距离而发散减弱的路径损耗为L＝20 lg d + 20 lg f + 32.44（其中损耗L的单位为dB，距离的单位为 km，频率f的单位为MHz）。由此可知，当频率一定时：（“X＾M”表示“X的二次方”）
\begin{enumerate}[leftmargin=3em]
\item 距离增加到N倍，衰减增大到N＾2倍
\item 距离增加到N倍，衰减增大到N倍
\item 距离增加到N倍，衰减增大到2×N倍
\item 距离增加到N倍，衰减增大到20×N倍
\end{enumerate}
\noindent\textbf{解说：}\textbf{}。\\\noindent\textbf{答案：}

\bigskip




\noindent\textbf{试题：}自由空间中电磁波随距离而发散减弱的路径损耗为L＝20 lg d + 20 lg f + 32.44（其中损耗L的单位为dB，距离的单位为 km，频率f的单位为MHz）。由此可知，当距离一定时：（“X＾M”表示“X的二次方”）
\begin{enumerate}[leftmargin=3em]
\item 频率增加到N倍，衰减增大到N＾2倍
\item 频率增加到N倍，衰减增大到N倍
\item 频率增加到N倍，衰减增大到2×N倍
\item 频率增加到N倍，衰减减少到1/N，所以430MHz信号比144MHz强
\end{enumerate}
\noindent\textbf{解说：}\textbf{}。\\\noindent\textbf{答案：}

\bigskip




\noindent\textbf{试题：}在电磁场理论中，关于“强度随离辐射源距离的平方成反比”的规律适用条件是：
\begin{enumerate}[leftmargin=3em]
\item 仅适用于离辐射源大约10倍波长外的“远区场”
\item 仅适用于离辐射源一定距离范围内的“远区场”
\item 适用于辐射源周围空间内的任何一点
\item 仅适用于从LF到UHF的频率范围
\end{enumerate}
\noindent\textbf{解说：}\textbf{}。\\\noindent\textbf{答案：}

\bigskip




\noindent\textbf{试题：}在空旷场地用场强计比较不同型号的辐射仰角为0的全向垂直天线在无线电通信中的实际辐射能力，需要注意：
\begin{enumerate}[leftmargin=3em]
\item 测试点应选择在离天线至少10倍波长远的地方
\item 场强计外壳应妥善接地
\item 场强计应与地面绝缘
\item 待测天线的馈电电缆外面应套上磁环
\end{enumerate}
\noindent\textbf{解说：}\textbf{}。\\\noindent\textbf{答案：}

\bigskip




\noindent\textbf{试题：}80米波段业余无线电测向机多采用环形天线和鞭状天线分别接收信号的磁场和电场分量，组合得到心形图单向特性。但远处调好的方向性，往往在接近电台后变差而需要重新调整鞭状天线的长度。其原因是：
\begin{enumerate}[leftmargin=3em]
\item 远场区内电场和磁场的比值不变，而近场区内恒比关系不再存在
\item 接近电台时信号场强导致接收机前级半导体器件过载
\item 接近电台时信号场强导致鞭状天线进入非线性工作区
\item 接近电台时信号场强导致环形天线发热变形而改变参数
\end{enumerate}
\noindent\textbf{解说：}\textbf{}。\\\noindent\textbf{答案：}

\bigskip




\noindent\textbf{试题：}影响短波电离层传播的主要因素有：
\begin{enumerate}[leftmargin=3em]
\item 太阳黑子活动、太阳耀斑活动和地磁活动
\item 当地天气
\item 地表温度
\item 发射机功率
\end{enumerate}
\noindent\textbf{解说：}\textbf{}。\\\noindent\textbf{答案：}

\bigskip




\noindent\textbf{试题：}影响短波电离层传播的主要因素有：
\begin{enumerate}[leftmargin=3em]
\item 季节和昼夜
\item 当地地面气压
\item 发射机功率
\item 接收天线增益
\end{enumerate}
\noindent\textbf{解说：}\textbf{}。\\\noindent\textbf{答案：}

\bigskip




\noindent\textbf{试题：}影响短波电离层传播的主要因素有：
\begin{enumerate}[leftmargin=3em]
\item 工作频率和通信距离
\item 发射机功率
\item 高空云量
\item 对流层气压
\end{enumerate}
\noindent\textbf{解说：}\textbf{}。\\\noindent\textbf{答案：}

\bigskip




\noindent\textbf{试题：}进行短波电离层传播预测所必需的参数为：
\begin{enumerate}[leftmargin=3em]
\item 太阳黑子平均数、地磁活动指数、通信双方的位置、通信时间
\item 太阳黑子活动平均周期长度、通信双方的天线高度、通信双方位置的磁偏角、通信时间
\item 上一太阳黑子活动周期的黑子平均数极大值、地磁活动指数、通信双方的相对直线距离、通信双方的高空温度
\item 太阳耀斑数量、地磁活动指数、通信双方的相对直线距离、通信时间
\end{enumerate}
\noindent\textbf{解说：}\textbf{}。\\\noindent\textbf{答案：}

\bigskip




\noindent\textbf{试题：}对短波电离层传播发生主要影响的各电离层按高度自低到高分别称为：
\begin{enumerate}[leftmargin=3em]
\item D、E、F1、F2
\item A、B、C、D
\item F、E2、E1、D
\item F2、F1、E2、E1
\end{enumerate}
\noindent\textbf{解说：}\textbf{}。\\\noindent\textbf{答案：}

\bigskip




\noindent\textbf{试题：}对短波电离层传播发生主要影响的各电离层按离子密度自高到低分别为：
\begin{enumerate}[leftmargin=3em]
\item F2、F1、E、D
\item A、B、C、D
\item F、E2、E1、D
\item E1、E2、F1、F2
\end{enumerate}
\noindent\textbf{解说：}\textbf{}。\\\noindent\textbf{答案：}

\bigskip




\noindent\textbf{试题：}各电离层对短波电离层传播所起的主要影响为：
\begin{enumerate}[leftmargin=3em]
\item F2、F1、E层可反射电波，D层不能反射但衰减电波
\item D、E层可反射电波，F1、F2层不能反射但衰减电波
\item F2、D层可反射电波，E、F1层不能反射但衰减电波
\item F1、E层可反射电波，D、F2层不能反射但衰减电波
\end{enumerate}
\noindent\textbf{解说：}\textbf{}。\\\noindent\textbf{答案：}

\bigskip




\noindent\textbf{试题：}太阳黑子活动的平均周期约为：
\begin{enumerate}[leftmargin=3em]
\item 11.2年
\item 38年
\item 6.7年
\item 5.7年
\end{enumerate}
\noindent\textbf{解说：}\textbf{}。\\\noindent\textbf{答案：}

\bigskip




\noindent\textbf{试题：}太阳黑子活动的强弱用“太阳黑子平均数（SSN）”来描述。一般说来：
\begin{enumerate}[leftmargin=3em]
\item SSN大,有利于短波远程通信
\item SSN小,有利于短波远程通信
\item SSN与短波远程通信效果无直接关系
\item 只有在发生太阳耀斑的情况下SSN 才会影响短波远程通信效果
\end{enumerate}
\noindent\textbf{解说：}\textbf{}。\\\noindent\textbf{答案：}

\bigskip




\noindent\textbf{试题：}在业余无线电HF传播预测中，经常遇到缩写为F107的参数，其意义是：
\begin{enumerate}[leftmargin=3em]
\item 太阳10.7cm波长射电辐射通量指数
\item FM广播频段107MHz的典型传播距离
\item 电离层对10.7MHz电波的衰减指数
\item 最高可用频率与10.7MHz的比值
\end{enumerate}
\noindent\textbf{解说：}\textbf{}。\\\noindent\textbf{答案：}

\bigskip




\noindent\textbf{试题：}在业余无线电HF传播预测中，经常遇到缩写为F107的参数，其数值范围大致为60-300。数值越大，其意义为：
\begin{enumerate}[leftmargin=3em]
\item 宁静太阳的辐射强度越强，电离层密度越大，有利于F层对DX电波的反射
\item 电离层对10.7MHz电波的衰减越大，说明DX传播的条件越差
\item FM广播频段107MHz的典型传播距离越远，可用作300公里以内的VHF传播距离预测
\item 宁静太阳的辐射强度越强，电离层密度越大，不利于F层对DX电波的反射
\end{enumerate}
\noindent\textbf{解说：}\textbf{}。\\\noindent\textbf{答案：}

\bigskip




\noindent\textbf{试题：}当最高可用频率（MUF）为20MHz时，具有较大DX通联机会的业余频段是：
\begin{enumerate}[leftmargin=3em]
\item 18MHz
\item 14MHz
\item 21MHz
\item 24MHz
\end{enumerate}
\noindent\textbf{解说：}\textbf{}。\\\noindent\textbf{答案：}

\bigskip




\noindent\textbf{试题：}太阳耀斑引起的电离层扰动（SID）对短波通信的影响是：
\begin{enumerate}[leftmargin=3em]
\item 低频率受到的影响超过高频率
\item 高纬度的传播路径影响会超过低纬度
\item 卫星通信受到的影响会超过地面直射波通信
\item 地球黑夜一面受到的影响会超过白昼面
\end{enumerate}
\noindent\textbf{解说：}\textbf{}。\\\noindent\textbf{答案：}

\bigskip




\noindent\textbf{试题：}HF频段远距离通信主要依靠下列传播方式：
\begin{enumerate}[leftmargin=3em]
\item 电离层反射
\item 中继台网转发
\item 业余卫星转发
\item 对流层散射传播
\end{enumerate}
\noindent\textbf{解说：}\textbf{}。\\\noindent\textbf{答案：}

\bigskip




\noindent\textbf{试题：}“静寂区”或者“越距”是指：
\begin{enumerate}[leftmargin=3em]
\item HF频段天波和地波都传播不到的中间区域
\item VHF/UHF频段超过视距电波传播不到的区域
\item VHF/UHF频段视距范围内但受障碍物阻挡电波传播不到的区域
\item 卫星通信中覆盖区以外电波传播不到的区域
\end{enumerate}
\noindent\textbf{解说：}\textbf{}。\\\noindent\textbf{答案：}

\bigskip




\noindent\textbf{试题：}进行传播预测时，太阳通量的值最低可以取为50左右，最高可以取为：
\begin{enumerate}[leftmargin=3em]
\item 300
\item 280
\item 250
\item 200
\end{enumerate}
\noindent\textbf{解说：}\textbf{}。\\\noindent\textbf{答案：}

\bigskip



%\noindent\textbf{试题：}反映地磁活动的程度有A指数和K指数，K指数是一个0到___之间变化的数字：
\noindent\textbf{试题：}反映地磁活动的程度有A指数和K指数，K指数是一个0到$\_\_\_$之间变化的数字：
\begin{enumerate}[leftmargin=3em]
\item 9
\item 8
\item 7
\item 6
\end{enumerate}
\noindent\textbf{解说：}\textbf{}。\\\noindent\textbf{答案：}

\bigskip




\noindent\textbf{试题：}电磁波在地球与电离层间反射后，通常会发生什么？
\begin{enumerate}[leftmargin=3em]
\item 信号的极化方向发生随机改变
\item 每一次反射都会使边带倒转
\item 电磁波的频率被随机增加或减少一定的值
\item 频率在30MHz以上的电磁波信号会在每一次反射后变得更强
\end{enumerate}
\noindent\textbf{解说：}\textbf{}。\\\noindent\textbf{答案：}

\bigskip




\noindent\textbf{试题：}大气层中的哪一部分使得无线电信号可以在全世界范围内传播？
\begin{enumerate}[leftmargin=3em]
\item 电离层
\item 对流层
\item 平流层
\item 磁层
\end{enumerate}
\noindent\textbf{解说：}\textbf{}。\\\noindent\textbf{答案：}

\bigskip




\noindent\textbf{试题：}业余电台和业余卫星之间通信，VHF/UHF频段影响信号强度的主要因素是：
\begin{enumerate}[leftmargin=3em]
\item 无线电波在自由空间的衰减
\item 电离层对无线电波的吸收
\item 电离层对无线电波的反射
\item 大地对无线电波的吸收
\end{enumerate}
\noindent\textbf{解说：}\textbf{}。\\\noindent\textbf{答案：}

\bigskip




\noindent\textbf{试题：}业余卫星通信和业余无线电测向都会出现“多普勒效应”的概念。多普勒效应是指：
\begin{enumerate}[leftmargin=3em]
\item 收发信机之间相对距离的变化使接收信号频率产生偏移的现象
\item 传播条件随时间的扰动使接收信号幅度产生波动的现象
\item 传播条件随空间位置的扰动使接收信号幅度产生波动的现象
\item 电波极化方向在传播途中变化使接收信号幅度产生波动的现象
\end{enumerate}
\noindent\textbf{解说：}\textbf{}。\\\noindent\textbf{答案：}

\bigskip




\noindent\textbf{试题：}地面业余电台进行业余卫星通信时，接收到的卫星信号频率的多普勒频移规律是：
\begin{enumerate}[leftmargin=3em]
\item 卫星飞来时频率偏高，飞离时频率偏低，越远频偏越大，过顶时频偏最小
\item 卫星飞来时频率偏低，飞离时频率偏高，越远频偏越大，过顶时频偏最小
\item 卫星飞来时频率偏高，飞离时频率偏低，越远频偏越小，过顶时频偏最大
\item 卫星飞来时频率偏低，飞离时频率偏高，越远频偏越小，过顶时频偏最大
\end{enumerate}
\noindent\textbf{解说：}\textbf{}。\\\noindent\textbf{答案：}

\bigskip




\noindent\textbf{试题：}按照频率划分表，122.25—123GHz业余频段不能用于卫星业余业务，但这对业余卫星通信活动影响不大。这是因为：
\begin{enumerate}[leftmargin=3em]
\item 60GHz、120GHz、183GHz附近存在衰减很大的大气吸收频带，并不适合用作业余卫星通信
\item 120GHz附近频段存在特别强烈的宇宙射线干扰，不适合用作卫星通信
\item 业余无线电爱好者不具备制造工作在120GHz附近频段的设备的能力
\item 120GHz附近频段的天线尺寸太大，不适合用作业余卫星通信
\end{enumerate}
\noindent\textbf{解说：}\textbf{}。\\\noindent\textbf{答案：}

\bigskip




\noindent\textbf{试题：}无线电波传播受降雨影响最严重的频段是：
\begin{enumerate}[leftmargin=3em]
\item 极高频EHF（毫米波）
\item 高频HF（短波）
\item 低频LF（长波）
\item 特高频UHF（分米波）
\end{enumerate}
\noindent\textbf{解说：}\textbf{}。\\\noindent\textbf{答案：}

\bigskip




\noindent\textbf{试题：}无线电波传播受降雨影响最严重的频段是：
\begin{enumerate}[leftmargin=3em]
\item 极高频EHF（毫米波）
\item 中频MF（中波）
\item 甚低频VLF（甚长波）
\item 甚高频VHF（米波）
\end{enumerate}
\noindent\textbf{解说：}\textbf{}。\\\noindent\textbf{答案：}

\bigskip




\noindent\textbf{试题：}做EME通信试验最好的时机是：
\begin{enumerate}[leftmargin=3em]
\item 月亮处于近地点
\item 满月
\item 通信两端天气都好
\item 太阳黑子数高
\end{enumerate}
\noindent\textbf{解说：}\textbf{}。\\\noindent\textbf{答案：}

\bigskip




\noindent\textbf{试题：}在10米、6米、2米波段经常能够接收到的“超视距传播”信号与下列哪种传播现象密切相关？
\begin{enumerate}[leftmargin=3em]
\item 突发E电离层
\item 流星余迹反向散射
\item D层的吸收
\item 灰线传播
\end{enumerate}
\noindent\textbf{解说：}\textbf{}。\\\noindent\textbf{答案：}

\bigskip




\noindent\textbf{试题：}在通常情况下，下列哪种传播模式可能使500公里左右范围的VHF和UHF超视距通信变得可行？
\begin{enumerate}[leftmargin=3em]
\item 对流层散射
\item D层折射
\item F2层折射
\item 法拉第旋转
\end{enumerate}
\noindent\textbf{解说：}\textbf{}。\\\noindent\textbf{答案：}

\bigskip




\noindent\textbf{试题：}哪一个业余无线电波段最适合流星余迹散射通信？
\begin{enumerate}[leftmargin=3em]
\item 6米
\item 10米
\item 2米
\item 70厘米
\end{enumerate}
\noindent\textbf{解说：}\textbf{}。\\\noindent\textbf{答案：}

\bigskip




\noindent\textbf{试题：}是什么导致了对流层的大气波导传播？
\begin{enumerate}[leftmargin=3em]
\item 大气高空逆温
\item 太阳黑子和太阳耀斑
\item 飓风或龙卷风导致的上升气流
\item 雷暴时闪电的放电
\end{enumerate}
\noindent\textbf{解说：}\textbf{}。\\\noindent\textbf{答案：}

\bigskip




\noindent\textbf{试题：}用甲乙两块电压表检查一节新干电池两端电压，均测得1.5伏。检查一节旧干电池，读数分别为1.2伏和1.3伏。正确的结论是：
\begin{enumerate}[leftmargin=3em]
\item 甲电压表的内阻小于乙电压表
\item 甲电压表的内阻大于乙电压表
\item 甲电压表的精度高于乙电压表
\item 甲电压表的精度低于乙电压表
\end{enumerate}
\noindent\textbf{解说：}\textbf{}。\\\noindent\textbf{答案：}

\bigskip




\noindent\textbf{试题：}用万用电表的交流电压档测量简单正弦交流电压，得到的读数是该电压的：
\begin{enumerate}[leftmargin=3em]
\item 有效值
\item 最大值
\item 峰-峰值
\item 平均值
\end{enumerate}
\noindent\textbf{解说：}\textbf{}。\\\noindent\textbf{答案：}

\bigskip




\noindent\textbf{试题：}用万用电表的直流电压档测量简单正弦交流电压，得到的读数是该电压的：
\begin{enumerate}[leftmargin=3em]
\item 零
\item 最大值
\item 峰-峰值
\item 有效值
\end{enumerate}
\noindent\textbf{解说：}\textbf{}。\\\noindent\textbf{答案：}

\bigskip




\noindent\textbf{试题：}用万用电表的直流电压档测量正负对称的方波电压（占空比为1），得到的读数是该电压的：
\begin{enumerate}[leftmargin=3em]
\item 零
\item 最大值
\item 最大值的1/2
\item 最大值的2倍
\end{enumerate}
\noindent\textbf{解说：}\textbf{}。\\\noindent\textbf{答案：}

\bigskip




\noindent\textbf{试题：}用磁电式万用电表的交流电压档测量正负对称的方波电压，其指针指示的电压读数：
\begin{enumerate}[leftmargin=3em]
\item 比最大值大
\item 为最大值
\item 为峰-峰值
\item 为平均值
\end{enumerate}
\noindent\textbf{解说：}\textbf{}。\\\noindent\textbf{答案：}

\bigskip




\noindent\textbf{试题：}用万用电表的直流电压档测量最小值为0的方波电压（占空比为1），得到的读数是该电压的：
\begin{enumerate}[leftmargin=3em]
\item 最大值的1/2
\item 最大值
\item 最大值的2倍
\item 最小值
\end{enumerate}
\noindent\textbf{解说：}\textbf{}。\\\noindent\textbf{答案：}

\bigskip




\noindent\textbf{试题：}测量一个电解电容器是否完好，应该使用：
\begin{enumerate}[leftmargin=3em]
\item 万用电表的电阻挡
\item 万用电表的电压挡
\item 万用电表的电流挡
\item 蓄电池充电器
\end{enumerate}
\noindent\textbf{解说：}\textbf{}。\\\noindent\textbf{答案：}

\bigskip




\noindent\textbf{试题：}用万用电表的电阻档测量一支阻抗为50欧姆的四分之一波长接地天线，读数为0欧。则可以做出如下结论：
\begin{enumerate}[leftmargin=3em]
\item 该天线与地之间可能存在着由电感线圈构成的直流通路
\item 该天线肯定已经短路损坏
\item 该天线肯定无法与特性阻抗为50欧的馈线相匹配
\item 该天线肯定无法与输出阻抗为50欧姆的收发信机相匹配
\end{enumerate}
\noindent\textbf{解说：}\textbf{}。\\\noindent\textbf{答案：}

\bigskip




\noindent\textbf{试题：}用万用电表的电阻挡测量一支阻抗为50欧姆的四分之一波长接地天线，读数为无穷大。则可以做出如下结论：
\begin{enumerate}[leftmargin=3em]
\item 该天线与地之间不存在由电感线圈构成的直流通路
\item 该天线肯定已经开路损坏
\item 该天线肯定无法与特性阻抗为50欧的馈线相匹配
\item 该天线肯定无法与输出阻抗为50欧姆的收发信机相匹配
\end{enumerate}
\noindent\textbf{解说：}\textbf{}。\\\noindent\textbf{答案：}

\bigskip




\noindent\textbf{试题：}万用电表设在电阻挡，将表笔分别接到尚未连接BALUN和电缆等其它部件的每臂长度均为四分之一波长的偶极天线的中心馈电点两端，读数应为：
\begin{enumerate}[leftmargin=3em]
\item 无穷大
\item 50欧姆
\item 读数受表笔连线造成的驻波比的影响
\item 读数受振子周围物体的影响
\end{enumerate}
\noindent\textbf{解说：}\textbf{}。\\\noindent\textbf{答案：}

\bigskip




\noindent\textbf{试题：}万用电表设在电阻挡，将表笔分别接到一条终端开路的任意长度的50欧同轴电缆的中心导体和屏蔽层，读数应为：
\begin{enumerate}[leftmargin=3em]
\item 无穷大
\item 50欧姆
\item 0欧姆
\item 读数与电缆长度有关
\end{enumerate}
\noindent\textbf{解说：}\textbf{}。\\\noindent\textbf{答案：}

\bigskip




\noindent\textbf{试题：}万用电表设在电阻挡，将表笔分别接到一条终端短路的任意长度的理想50欧同轴电缆的中心导体和屏蔽层，读数应为：
\begin{enumerate}[leftmargin=3em]
\item 0欧姆
\item 无穷大
\item 50欧姆
\item 读数与电缆长度有关
\end{enumerate}
\noindent\textbf{解说：}\textbf{}。\\\noindent\textbf{答案：}

\bigskip




\noindent\textbf{试题：}用万用电表的电阻挡测量一个常用元件的电阻，在R×1挡、R×10挡和R×100挡测的的读数有很大差别，分别为40欧、240欧和2000欧。正确的推论是：
\begin{enumerate}[leftmargin=3em]
\item 该元件是一个非线性元件（例如二极管）
\item 该元件肯定是一只损坏了的电阻
\item 万用表的准确度肯定有问题
\item 万用表的电池肯定电量不足了
\end{enumerate}
\noindent\textbf{解说：}\textbf{}。\\\noindent\textbf{答案：}

\bigskip




\noindent\textbf{试题：}用指针式万用电表的电阻挡测量一个常用元件，表笔第一次连接时表针挥动了一下，但很快回到零点，断开表笔重新连接，表针却不再有任何指示。正确的推论是：
\begin{enumerate}[leftmargin=3em]
\item 该元件可能是一个电解电容器
\item 该元件肯定是一只损坏了的电阻
\item 万用表的准确度肯定有问题
\item 万用表的电池肯定电量不足了
\end{enumerate}
\noindent\textbf{解说：}\textbf{}。\\\noindent\textbf{答案：}

\bigskip




\noindent\textbf{试题：}万用电表的电压挡有一个“欧姆/伏（每伏欧姆）”的指标。每伏欧姆数越大，说明：
\begin{enumerate}[leftmargin=3em]
\item 测量电压时电表对电路工作状态的影响越小
\item 测量电压时电表对电路工作状态的影响越大
\item 可以耐受的电压过载能力越强
\item 可以耐受的电压过载能力越弱
\end{enumerate}
\noindent\textbf{解说：}\textbf{}。\\\noindent\textbf{答案：}

\bigskip




\noindent\textbf{试题：}用万用电表的交流电压档测量电路中某两点间的等幅度音频信号电压，发现交换两支表笔时虽然都有指示，但得到的度数不相同。正确的推论是：
\begin{enumerate}[leftmargin=3em]
\item 音频信号上叠加有直流成分，且电表电路采用了半波整流
\item 音频信号上叠加有直流成分，且电表电路采用了全波整流
\item 电表精度太差
\item 电表导线之间的分布电容造成了读数的误差
\end{enumerate}
\noindent\textbf{解说：}\textbf{}。\\\noindent\textbf{答案：}

\bigskip




\noindent\textbf{试题：}用万用电表的交流电压档测量电路中某两点间的等幅度音频信号电压，得到一个读数，但在交换两支表笔后表针就根本没有偏转了，反复试验都是如此。正确的推论是：
\begin{enumerate}[leftmargin=3em]
\item 音频信号上叠加有直流成分，且电表电路采用了半波整流
\item 音频信号上叠加有直流成分，且电表电路采用了全波整流
\item 电表精度太差
\item 电表导线之间的分布电容造成了读数的误差
\end{enumerate}
\noindent\textbf{解说：}\textbf{}。\\\noindent\textbf{答案：}

\bigskip




\noindent\textbf{试题：}无线电通信和测试设备、电视设备和音频设备常用的传输接口标准阻抗分别为：
\begin{enumerate}[leftmargin=3em]
\item 50欧、75欧和600欧
\item 50欧、600欧和75欧
\item 50欧、50欧和75欧
\item 75欧、50欧和16欧
\end{enumerate}
\noindent\textbf{解说：}\textbf{}。\\\noindent\textbf{答案：}

\bigskip




\noindent\textbf{试题：}哪一种焊锡材料比较适用于无线电制作和修理？
\begin{enumerate}[leftmargin=3em]
\item 松香芯焊锡丝
\item 银焊锡丝
\item 酸心焊锡丝
\item 铝焊锡丝
\end{enumerate}
\noindent\textbf{解说：}\textbf{}。\\\noindent\textbf{答案：}

\bigskip




\noindent\textbf{试题：}如果用欧姆表测量一个电路，最开始欧姆表显示阻抗很小，接着显示电阻变得越来越大，这是因为：
\begin{enumerate}[leftmargin=3em]
\item 这个电路两端可能存在一个很大的电容
\item 欧姆表有了故障
\item 这个电路两端可能存在一个很大的电感
\item 该电路是一个张弛振荡电路
\end{enumerate}
\noindent\textbf{解说：}\textbf{}。\\\noindent\textbf{答案：}

\bigskip




\noindent\textbf{试题：}在使用欧姆表测量一个电路的某两个端点间电阻值之前，要先注意什么？
\begin{enumerate}[leftmargin=3em]
\item 确保电路没有连接至电源
\item 确保电路已接通了正常工作时需要的电源
\item 确保电路已经正常接地
\item 确保电路在正常的频率下工作
\end{enumerate}
\noindent\textbf{解说：}\textbf{}。\\\noindent\textbf{答案：}

\bigskip




\noindent\textbf{试题：}防雷装置的作用是防止雷电危害。传统防雷装置的主要组成部分是：
\begin{enumerate}[leftmargin=3em]
\item 接闪器（避雷针）、引下线、接地体
\item 天线、限流器、地线
\item 避雷针、高压指示灯、过流保护器
\item 避雷针、过压保护器、熔丝
\end{enumerate}
\noindent\textbf{解说：}\textbf{}。\\\noindent\textbf{答案：}

\bigskip




\noindent\textbf{试题：}防雷接地的作用是：
\begin{enumerate}[leftmargin=3em]
\item 把接闪器引入的雷击电流有效地泄入大地
\item 用接闪器感应到的雷电高压启动过压保护电路
\item 有效地阻断接闪器引入的雷击电流使其不致流入大地
\item 当接闪器引入雷击电流时迅速烧断熔丝，阻断其流动
\end{enumerate}
\noindent\textbf{解说：}\textbf{}。\\\noindent\textbf{答案：}

\bigskip




\noindent\textbf{试题：}关于对防雷接地基本要求的正确说法：
\begin{enumerate}[leftmargin=3em]
\item 要有单独的接地体，接地电阻要小，接闪器到接地体之间的引下线应尽量短而粗
\item 可以利用与埋地金属管线相连的自来水管作为接地体
\item 交流电网的“零线”在配电系统中已经接地，因此可代替防雷接地体及其引下线
\item 接闪器到接地体之间的引下线平时没有电流流过，采用直径0.5毫米的导线为好
\end{enumerate}
\noindent\textbf{解说：}\textbf{}。\\\noindent\textbf{答案：}

\bigskip




\noindent\textbf{试题：}单支避雷针的保护范围大致有多大：
\begin{enumerate}[leftmargin=3em]
\item 以避雷针为顶点的45度圆锥体以内空间
\item 避雷针周围水平方圆30米内的任何物体
\item 避雷针周围所有比避雷针低的空间
\item 以避雷针为顶点、避雷针高度为半径的半球体以内空间
\end{enumerate}
\noindent\textbf{解说：}\textbf{}。\\\noindent\textbf{答案：}

\bigskip




\noindent\textbf{试题：}安全电压是指不致使人直接致死或致残的电压。标准国家标准GB3805－83《安全电压》，一般环境条件下允许持续接触的“安全特低电压”为：
\begin{enumerate}[leftmargin=3em]
\item 24V
\item 48V
\item 72V
\item 6V
\end{enumerate}
\noindent\textbf{解说：}\textbf{}。\\\noindent\textbf{答案：}

\bigskip




\noindent\textbf{试题：}触及裸露的射频导线时，与触及相同电压的直流或50Hz交流导线相比，对人身安全危险的大致差别为：
\begin{enumerate}[leftmargin=3em]
\item 致死危险性下降，但皮肤容易灼伤
\item 更容易出现呼吸或心脏麻痹、更易致死
\item 射频电流对人体没有危险
\item 触及不同频率、相同电压的导线时，对人体的危害没有差别
\end{enumerate}
\noindent\textbf{解说：}\textbf{}。\\\noindent\textbf{答案：}

\bigskip




\noindent\textbf{试题：}如遇设备、电线或者电源引起失火，正确的处置为：
\begin{enumerate}[leftmargin=3em]
\item 立即切断电源，使用二氧化碳灭火器灭火
\item 立即切断电源，使用化学泡沫灭火器灭火
\item 立即切断电源，用水灭火
\item 迅速起身逃离火场
\end{enumerate}
\noindent\textbf{解说：}\textbf{}。\\\noindent\textbf{答案：}

\bigskip




\noindent\textbf{试题：}必须带电检修由市电供电的无线电设备时，应做到：
\begin{enumerate}[leftmargin=3em]
\item 双脚与地绝缘，单手操作，另一只手不触摸机壳等任何与电路设备有关的金属物品
\item 双脚与地绝缘，单手操作，另一只手通过触摸机壳或水管等途径良好接地
\item 只要设备外壳良好接地，双脚是否与地绝缘、双手是否同时操作都没有关系
\item 只要确保设备外壳与地绝缘，双脚是否与地绝缘、双手是否同时操作都没有关系
\end{enumerate}
\noindent\textbf{解说：}\textbf{}。\\\noindent\textbf{答案：}

\bigskip




\noindent\textbf{试题：}两手分别接触电压有效值相同但频率不同的电路两端时，对人体生命安全威胁由大到小的排序为：
\begin{enumerate}[leftmargin=3em]
\item 工频交流电、HF射频交流电、UHF射频交流电
\item UHF射频交流电、HF射频交流电、工频交流电
\item HF射频交流电、工频交流电、UHF射频交流电
\item HF射频交流电、UHF射频交流电、工频交流电
\end{enumerate}
\noindent\textbf{解说：}\textbf{}。\\\noindent\textbf{答案：}

\bigskip




\noindent\textbf{试题：}电路中的保险丝起到什么作用？
\begin{enumerate}[leftmargin=3em]
\item 过载时切断电路
\item 防止电源纹波损害电路
\item 限制电流，防止触电
\item 以上三项全部正确
\end{enumerate}
\noindent\textbf{解说：}\textbf{}。\\\noindent\textbf{答案：}

\bigskip




\noindent\textbf{试题：}为什么在需要安装5安培保险丝的地方安装一个20安培的保险丝是不可取的？
\begin{enumerate}[leftmargin=3em]
\item 过大的电流可能导致火灾
\item 电源的纹波会显著增大
\item 因为20安培的保险丝电流更大，所以它更容易熔断
\item 其他三项全部正确
\end{enumerate}
\noindent\textbf{解说：}\textbf{}。\\\noindent\textbf{答案：}

\bigskip




\noindent\textbf{试题：}防止设备外壳带电危险的措施包括：
\begin{enumerate}[leftmargin=3em]
\item 其他三项全部正确
\item 将所有的交流供电设备全部连接至一个安全地线
\item 安装漏电保护断路器
\item 所有使用交流供电的设备的电源线都使用带有单独保护地线端的三线插头
\end{enumerate}
\noindent\textbf{解说：}\textbf{}。\\\noindent\textbf{答案：}

\bigskip




\noindent\textbf{试题：}在为同轴电缆馈线安装避雷器时应当注意什么？
\begin{enumerate}[leftmargin=3em]
\item 将所有避雷器的地线接到同一个金属板上，然后将这个金属板接到室外的接地极
\item 在每一个避雷器的接地线上安装开关，以防止射频过载损伤避雷器
\item 将每一个避雷器单独引出接地线，并且将它们都连接至电台的地线
\item 要在每一个避雷器处两端并联一个开关，以便在使用大功率输出的时候可以将避雷器旁路掉
\end{enumerate}
\noindent\textbf{解说：}\textbf{}。\\\noindent\textbf{答案：}

\bigskip




\noindent\textbf{试题：}常规的12伏酸铅蓄电池通常有什么潜在的危险？
\begin{enumerate}[leftmargin=3em]
\item 如果通风不良，有爆炸风险的气体会聚集
\item 有高电压，存在触电的风险
\item 它会释放臭氧，进而污染大气层
\item 长时间不使用可能会引起自燃
\end{enumerate}
\noindent\textbf{解说：}\textbf{}。\\\noindent\textbf{答案：}

\bigskip




\noindent\textbf{试题：}以下哪一项是天线铁塔安装防雷接地的良好措施？
\begin{enumerate}[leftmargin=3em]
\item 接地连接要尽可能地短而直接
\item 确保接地线的所有需要弯折的地方都是干净的直角弯
\item 为接地系统做好防水处理
\item 打入地面的接地极要尽量远离天线塔
\end{enumerate}
\noindent\textbf{解说：}\textbf{}。\\\noindent\textbf{答案：}

\bigskip




\noindent\textbf{试题：}设备电源拔掉电源线以后，检修时还有什么安全风险？
\begin{enumerate}[leftmargin=3em]
\item 充满高电压的电容器可能造成电击
\item 地磁场可能在变压器中激起感应电流导致电源损坏
\item 打开电源外壳可能引起保险丝烧断
\item 静电可能损坏接地系统
\end{enumerate}
\noindent\textbf{解说：}\textbf{}。\\\noindent\textbf{答案：}

\bigskip




\noindent\textbf{试题：}自制一台由220伏交流供电的设备，推荐采用的安全措施是：
\begin{enumerate}[leftmargin=3em]
\item 交流电源入口火线端串联安装保险丝
\item 在交流供电入口处并联安装一个交流电压表
\item 在交流供电入口处串联安装一个电感
\item 在交流供电入口处并联安装一个电容
\end{enumerate}
\noindent\textbf{解说：}\textbf{}。\\\noindent\textbf{答案：}

\bigskip




\noindent\textbf{试题：}我国业余电台应该遵守的关于电磁辐射污染的具体管理规定文件为： 
\begin{enumerate}[leftmargin=3em]
\item 国家标准GB/T 51391《通信工程建设环境保护技术标准》
\item 《业余电台管理办法》
\item 国际非电离辐射防护委员会《限制时变电场和磁场暴露的导则》
\item 美国FCC《射频电磁场人员暴露准则的测评方法规》
\end{enumerate}
\noindent\textbf{解说：}\textbf{}。\\\noindent\textbf{答案：}

\bigskip




\noindent\textbf{试题：}按照我国国家标准GB/T 51391《通信工程建设环境保护技术标准》，向没有屏蔽空间发射电磁场可以免于管理的有： 
\begin{enumerate}[leftmargin=3em]
\item 等效辐射功率在0.1-3MHz不大于300W、在3MHz-300GHz不大于100瓦的无线通信设施（设备） 
\item 所有业余电台
\item 发射频率30MHz以下的所有业余电台
\item 发射频率30MHz以上的所有业余电台
\end{enumerate}
\noindent\textbf{解说：}\textbf{}。\\\noindent\textbf{答案：}

\bigskip




\noindent\textbf{试题：}我国国家标准《电磁环境控制限值》规定公众暴露控制限值中，对环境电磁辐射场强任意连续6分钟内的方均根值要求最严格的频率范围为： 
\begin{enumerate}[leftmargin=3em]
\item 30MHz- 3GHz
\item 3MHz- 30MHz
\item 300kHz-30MHz
\item 1200MHz-30GHz
\end{enumerate}
\noindent\textbf{解说：}\textbf{}。\\\noindent\textbf{答案：}

\bigskip




\noindent\textbf{试题：}为什么电磁辐射防护规定国家标准中的照射限值随着频率的变化而不同？
\begin{enumerate}[leftmargin=3em]
\item 人体会对某些特定频率的电磁波吸收量更大
\item 较低频率的无线电波不会穿透人体
\item 在自然中高频电磁波不常见
\item 较低频率的无线电波相对高频率的无线电波拥有更高的能量
\end{enumerate}
\noindent\textbf{解说：}\textbf{}。\\\noindent\textbf{答案：}

\bigskip




\noindent\textbf{试题：}要防止HF发射机的杂散发射干扰天线附近的VHF电视机，应该发射机和天线之间串联：
\begin{enumerate}[leftmargin=3em]
\item 截止频率为30MHz左右的低通滤波器
\item 截止频率为30MHz左右的高通滤波器
\item 截止频率为300MHz左右的低通滤波器
\item 中心频率为30MHz左右的带通滤波器
\end{enumerate}
\noindent\textbf{解说：}\textbf{}。\\\noindent\textbf{答案：}

\bigskip




\noindent\textbf{试题：}假设中继台的收、发信机共用天线，上下行频率分别为F1和F2。要防止中继台发射机对接收机产生干扰，应该对中继台设备采取下列措施：
\begin{enumerate}[leftmargin=3em]
\item 在发信机与天线间串联中心频率为F1的带阻滤波器，在收信机与天线间串接中心频率为F2的带阻滤波器
\item 在发信机与天线间串联中心频率为F1的带通滤波器，在收信机与天线间串接中心频率为F2的带通滤波器
\item 在发信机与天线间串联中心频率为F2的带阻滤波器，在收信机与天线间串接中心频率为F2的带阻滤波器
\item 在发信机与天线间串联中心频率为F1的带阻滤波器，在收信机与天线间串接中心频率为F1的带阻滤波器
\end{enumerate}
\noindent\textbf{解说：}\textbf{}。\\\noindent\textbf{答案：}

\bigskip




\noindent\textbf{试题：}要防止业余HF发射机的杂散发射干扰天线附近的电话机，应该在电话机和电话线之间之间串联：
\begin{enumerate}[leftmargin=3em]
\item 截止频率不高于1MHz的低通滤波器
\item 截止频率约为3kHz的高通滤波器
\item 截止频率约为3kHz的带阻滤波器
\item 中心频率约为3kHz的带通滤波器
\end{enumerate}
\noindent\textbf{解说：}\textbf{}。\\\noindent\textbf{答案：}

\bigskip




\noindent\textbf{试题：}A、B两部HF业余电台相距很近，分别工作在A、B两个频段。为减少B电台受到来自A电台的干扰，可以在B电台与天线之间串联：
\begin{enumerate}[leftmargin=3em]
\item 中心频率为A的带阻滤波器
\item 中心频率为A的带通滤波器
\item 截止频率为A的高通滤波器
\item 截止频率为B的高通滤波器
\end{enumerate}
\noindent\textbf{解说：}\textbf{}。\\\noindent\textbf{答案：}

\bigskip




\noindent\textbf{试题：}为了减少发射设备的谐波干扰近在咫尺的接收机，可以在发射设备和天线之间串联一个LC低通滤波器。正确的说法是：
\begin{enumerate}[leftmargin=3em]
\item 滤波器的阶数越高，抑制倍频干扰的效果越好
\item 滤波器的阶数越低，抑制倍频干扰的效果越好
\item 滤波器的阶数越高，损耗的功率越小
\item 滤波器的阶数越高，耐受的功率越大
\end{enumerate}
\noindent\textbf{解说：}\textbf{}。\\\noindent\textbf{答案：}

\bigskip




\noindent\textbf{试题：}架设业余中继台前应确定台址附近没有能与中继台下行频率形成三阶互调的发射台。如果中继台的上、下行频率分别为fR和fT，可能造成这种三阶互调的干扰频率fX是：
\begin{enumerate}[leftmargin=3em]
\item 2fT - fR 或 (fT + fR ) / 2
\item fT - fR 或 fT + fR
\item 2（fT - fR） 或 2(fT + fR )
\item 2fT 或 2fR
\end{enumerate}
\noindent\textbf{解说：}\textbf{}。\\\noindent\textbf{答案：}

\bigskip




\noindent\textbf{试题：}如果别人报告说你的发射干扰了相邻频率的通信，此时你应当做的是：
\begin{enumerate}[leftmargin=3em]
\item 检查发射机的频率指示是否准确、发射机的杂散发射指标是否合格
\item 换用另一种调制模式工作
\item 将这种情况通知你的设备制造商
\item 加大发射功率
\end{enumerate}
\noindent\textbf{解说：}\textbf{}。\\\noindent\textbf{答案：}

\bigskip




\noindent\textbf{试题：}滤去杂散发射的滤波器应该安装在什么地方？
\begin{enumerate}[leftmargin=3em]
\item 发信机和天线之间
\item 收信机和发信机之间
\item 电台的电源处
\item 话筒上
\end{enumerate}
\noindent\textbf{解说：}\textbf{}。\\\noindent\textbf{答案：}

\bigskip




\noindent\textbf{试题：}在解决电视接收机被附近的144MHz业余电台的过载干扰问题的时候，应当先尝试什么措施？
\begin{enumerate}[leftmargin=3em]
\item 在电视接收机的天线端子前安装144MHz带阻滤波器
\item 在电视接收机的天线端子前安装144MHz带通波器
\item 在业余电台的射频输出端安装144MHz带通滤波器
\item 在业余电台的射频输出端安装144MHz带阻滤波器
\end{enumerate}
\noindent\textbf{解说：}\textbf{}。\\\noindent\textbf{答案：}

\bigskip




\noindent\textbf{试题：}在汽车上安装的移动电台中能听到的随着引擎转速变化的高频啸叫声的来源是？
\begin{enumerate}[leftmargin=3em]
\item 发电机
\item 火花塞系统
\item 电动油泵
\item 防抱死刹车装置的控制器
\end{enumerate}
\noindent\textbf{解说：}\textbf{}。\\\noindent\textbf{答案：}

\bigskip




\noindent\textbf{试题：}移动车载台的直流电源负极应当接在哪里？
\begin{enumerate}[leftmargin=3em]
\item 连接在电池的负极或发动机的接地带
\item 连接在天线座上
\item 可以连接在汽车的任意的金属部分
\item 连接在固定住电台的挂置架上
\end{enumerate}
\noindent\textbf{解说：}\textbf{}。\\\noindent\textbf{答案：}

\bigskip




\noindent\textbf{试题：}下列哪一项可以有效减小火花塞干扰？
\begin{enumerate}[leftmargin=3em]
\item 打开电台的抑噪（NB）功能
\item 降低静噪（SQL）阀值
\item 将频率稍稍偏离一些
\item 调节电台的RIT旋钮
\end{enumerate}
\noindent\textbf{解说：}\textbf{}。\\\noindent\textbf{答案：}

\bigskip




\noindent\textbf{试题：}一般来说，如果要解决发射机对附近有线电话的干扰，最先做的应当是：
\begin{enumerate}[leftmargin=3em]
\item 在有线电话进线处安装射频滤波器
\item 在发射机射频输出端安装高通滤波器
\item 在发射机射频输出端安装低通滤波器
\item 改善发射机的接地情况
\end{enumerate}
\noindent\textbf{解说：}\textbf{}。\\\noindent\textbf{答案：}

\bigskip




\noindent\textbf{试题：}下列哪一种方法可以用来定位无线电噪音源或者恶意干扰源？
\begin{enumerate}[leftmargin=3em]
\item 无线电测向
\item 多普勒雷达
\item 回波定位
\item 相位锁定
\end{enumerate}
\noindent\textbf{解说：}\textbf{}。\\\noindent\textbf{答案：}


%所有附录
\chapter{附录}

\section{业余无线电台分区表}

\begin{center}
  \begin{tabular}[t]{|c|l|}
    \hline
    \multicolumn{1}{|c|}{\textbf{分区号}} & \multicolumn{1}{|c|}{\textbf{地区}} \\
    \hline
    第1区                                & 北京                                \\
    \hline
    第2区                                & 黑龙江、吉林、辽宁                         \\
    \hline
    第3区                                & 天津、内蒙古、河北、山西                      \\
    \hline
    第4区                                & 上海、山东、江苏                          \\
    \hline
    第5区                                & 浙江、江西、福建                          \\
    \hline
    第6区                                & 安徽、河南、湖北                          \\
    \hline
    第7区                                & 湖南、广东、广西、海南                       \\
    \hline
    第8区                                & 四川、重庆、贵州、云南                       \\
    \hline
    第9区                                & 陕西、甘肃、宁夏、青海                       \\
    \hline
    第0区                                & 新疆、西藏                             \\
    \hline
  \end{tabular}
\end{center}

\newpage

\section{字母解释法（节录）}

\begin{center}
  \begin{tabular}{cc}%
    \begin{tabular}[t]{|c|l|}
      \hline
      \textbf{字母} & \textbf{单词} \\
      \hline
      A           & Alfa        \\
      \hline
      B           & Bravo       \\
      \hline
      C           & Charlie     \\
      \hline
      D           & Delta       \\
      \hline
      E           & Echo        \\
      \hline
      F           & Foxtrot     \\
      \hline
      G           & Golf        \\
      \hline
      H           & Hotel       \\
      \hline
      I           & India       \\
      \hline
      J           & Juliett     \\
      \hline
      K           & Kilo        \\
      \hline
      L           & Lima        \\
      \hline
      M           & Mike        \\
      \hline
      N           & November    \\
      \hline
      O           & Oscar       \\
      \hline
      P           & Papa        \\
      \hline
      Q           & Quebec      \\
      \hline
      R           & Romeo       \\
      \hline
      S           & Sierra      \\
      \hline
      T           & Tango       \\
      \hline
      U           & Uniform     \\
      \hline
      V           & Victor      \\
      \hline
      W           & Whiskey     \\
      \hline
      X           & X-ray       \\
      \hline
      Y           & Yankee      \\
      \hline
      Z           & Zulu        \\
      \hline
    \end{tabular} &
    \begin{tabular}[t]{|c|l|}
      \hline
      \textbf{字母} & \textbf{单词} \\
      \hline
      0           & Zero        \\
      \hline
      1           & One         \\
      \hline
      2           & Two         \\
      \hline
      3           & Three       \\
      \hline
      4           & Four        \\
      \hline
      5           & Five        \\
      \hline
      6           & Six         \\
      \hline
      7           & Seven       \\
      \hline
      8           & Eight       \\
      \hline
      9           & Nine        \\
      \hline
    \end{tabular} \tabularnewline
  \end{tabular}
\end{center}

\newpage

\section{常用Q简语（节录）}

\begin{longtable}{|l|l|l|}
  \hline
                         & \textbf{提问}                       & \textbf{回答或建议}                    \\
  \hline
  QRL                    & 你正忙着吗                             & 我正忙着                              \\
  \hline
  QRM                    & 你遇到他台干扰吗                          & 我遇到他台干扰                           \\
  \hline
  QRN                    & 你遇到天电干扰吗                          & 我遇到天电干扰                           \\
  \hline
  QRO                    & 要我增加功率吗                           &                                   \\
  \hline
  QRP                    & 要我减小功率吗                           &                                   \\
  \hline
  QRQ                    & 要我加快发送速度吗                         & 请加快发送速度                           \\
  \hline
  QRS                    & 要我减慢发送速度吗                         & 请减慢发送速度                           \\
  \hline
  QRT                    & 要我停止发送吗                           & 请停止发送                             \\
  \hline
  QRU                    & 你和我还有事吗                           & 我和你无事了                            \\
  \hline
  QRV                    & 你是否已准备好                           & 我已准备好                             \\
  \hline
  QRZ                    & 谁在呼叫我                             &                                   \\
  \hline
  QSA                    & 我的信号强度如何                          & 你的信号强度为 × 级（1-5 级）                \\
  \hline
  QSB                    & 我的信号有衰落吗                          & 你的信号有衰落                           \\
  \hline
  QSD                    & 我发报的手法有毛病吗                        & 你发报的手法有毛病                         \\
  \hline
  \multirow{2}{1em}{QSK} & 能在你的信号间隙中接收吗                      & 我在发射的信号间隙中接收                      \\
                         & （即 QSK 插入方式）                      & （即 QSK 插入方式）                      \\
  \hline
  QSL                    & 你能给我收据（或 QSL 卡片）吗                 & 我给你收据（QSL 卡片）、我已收妥                \\
  \hline
  QSO                    & 你能直接和 ××× 电台通信吗                   & 我能直接和 ××× 电台通信                    \\
  \hline
  QSP                    & 你能传信到 ××× 电台吗                     & 我能传信到 ××× 电台                      \\
  \hline
  \multirow{2}{1em}{QSX} & 你将在 nnnn \si{\kHz}（或 \si{\MHz}）频率 & 我将在 nnnn \si{\kHz}（或 \si{\MHz}）频率 \\
                         & 守听 ××× 电台吗                        & 守听 ××× 电台                         \\
  \hline
  QSY                    & 要我将频率改到 nnnn 频率吗                  & 请将频率改到 nnnn 频率                    \\
  \hline
  QTH                    & 你的电台位置在哪里                         & 我的电台位置是 ××××                      \\
  \hline
\end{longtable}

\newpage

\section{业余通信常用缩语表（节录）}

说明：黑体字的缩语为考到的重要缩语。

\begin{longtable}[l]{llll}
    & \textbf{缩语}                       & \textbf{原词}                             & \textbf{含义}         \\
    & \endfirsthead
    & \textbf{缩语}                       & \textbf{原词}                             & \textbf{含义}         \\
    & \endhead
  A & \textbf{ABT}                      & About                                   & 大约                  \\
    & AC                                & Alternating current                     & 交流                  \\
    & \textbf{ADR}                      & Address                                 & 地址                  \\
    & \textbf{ADDR}                     & Address                                 & 地址                  \\
    & \textbf{AGC}                      & Automatic gain control                  & 收信机自动增益控制           \\
    & \textbf{AGN}                      & Again                                   & 再、再来一次              \\
    & \textbf{AHR}                      & Another                                 & 另一个                 \\
    & \textbf{ALC}                      & Automatic level control                 & 发信自动电平控制            \\
    & \textbf{AM}                       & Amplitude modulation                    & 幅度调制                \\
    & \textbf{ANT}                      & Antenna                                 & 天线                  \\
    & \textbf{ARDF}                     & Amateur radio direction finding         & 业余无线电测向             \\
    & $\overline{\mathbf{A}\mathbf{S}}$ &                                         & 请稍等                 \\
    & \textbf{AS}                       &                                         & 亚洲、如同               \\
    & ASK                               & Amplitude-shift keying                  & 移幅键控                \\
    & \textbf{AT}                       & Antenna tuner                           & 自动天线调谐              \\
    & \textbf{ATT}                      & Attenuated / Attenuator                 & 衰减/收信机输入衰减器         \\
    & \textbf{ATU}                      & Antenna tuning unit                     & 天线调谐器               \\
  B & \textbf{BALUN}                    & balanced to unbalanced / balancing unit & 平衡和不平衡/平衡单元         \\
    & \textbf{BEAM}                     & Beam antenna                            & 定向天线                \\
    & \textbf{BEST}                     &                                         & 最好的                 \\
    & \textbf{BJT}                      & Beijing Time                            & 北京时间                \\
    & \textbf{BK}                       & Break                                   & 插入、打断               \\
    & BREAK                             &                                         & 我还没有说完              \\
    & BREAK BREAK                       &                                         & 我还没有说完              \\
    & Break in                          &                                         & 能在你的信号间隙中接收吗        \\
    & \textbf{BURO}                     & QSL Bureau                              & QSL卡片管理局            \\
  C & \textbf{C}                        &                                         & 遇到、见面               \\
    & \textbf{CFM}                      & Confirm                                 & 确认                  \\
    & \textbf{CHEERIO}                  &                                         & 再会、祝贺               \\
    & \textbf{CL}                       & Close, call                             & 关闭（或呼叫）             \\
    & \textbf{CLG}                      & Calling                                 & 呼叫                  \\
    & \textbf{CLS}                      & Call sign                               & 呼号                  \\
    & \textbf{CPI}                      & Copy                                    & 抄收                  \\
    & CQ                                &                                         & 普遍呼叫                \\
    & \textbf{CW}                       & Continuous wave                         & 等幅电报                \\
    & \textbf{CTCSS}                    & Continuous Tone-Coded Squelch System    & 亚音调静噪               \\
  D & \textbf{DATE}                     &                                         & 日期                  \\
    & dB                                & decibel                                 & 增益（单位）              \\
    & DC                                & Direct current                          & 直流                  \\
    & DR                                & Dear                                    & 亲爱的                 \\
    & DTMF                              & Dual-tone multi-frequency signaling     & 双音多频编码              \\
    & \textbf{DP}                       & Dipole antenna                          & 偶级天线                \\
  E & \textbf{EHF}                      & Extremely high frequency                & 极高频                 \\
    & \textbf{EL}                       &                                         & 单元（常用于天线振子）         \\
    & \textbf{ELE}                      &                                         & 单元（常用于天线振子）         \\
    & \textbf{ELS}                      &                                         & 单元（常用于天线振子）         \\
    & \textbf{ES}                       &                                         & 和                   \\
  F & FAX                               & Facsimile                               & 传真                  \\
    & \textbf{FB}                       & Fine business                           & 很好的                 \\
    & \textbf{FER}                      &                                         & 为了，对于               \\
    & \textbf{FINE}                     &                                         & 好的，精细的              \\
    & \textbf{FM}                       & Frequency modulation                    & 频率调制、调频             \\
    & \textbf{FR}                       &                                         & 为了，对于               \\
    & \textbf{FREQ}                     & Frequency                               & 频率                  \\
    & \textbf{FSK}                      & Frequency-shift keying                  & 移频键控                \\
  G & GA                                & Go ahead                                & 继续、请过来              \\
    & \textbf{GA}                       & Good afternoon                          & 下午好                 \\
    & \textbf{GB}                       & Goodbye                                 & 再见                  \\
    & \textbf{GE}                       & Good evening                            & 晚上好                 \\
    & \si{\GHz}                         & gigahertz                               & 吉赫兹（单位）             \\
    & \textbf{GL}                       & Good luck                               & 好运气                 \\
    & \textbf{GLD}                      & Glad                                    & 高兴                  \\
    & \textbf{GM}                       & Good morning                            & 早晨好                 \\
    & \textbf{GMT}                      & Greenwich Mean Time                     & 格林威治时间              \\ %格林尼治标准时间
    & \textbf{GN}                       & Good night                              & 晚安                  \\
    & \textbf{GND}                      & Ground                                  & 地线、地面               \\
    & \textbf{GP}                       & Ground-plane antenna                    & 垂直接地天线              \\
  H & HF                                & High frequency                          & 高频、即SW（Shortwave）   \\
    & \textbf{HPE}                      & Hope                                    & 希望                  \\
    & \textbf{HPI}                      & Happy                                   & 幸福                  \\
    & \textbf{HPY}                      & Happy                                   & 幸福                  \\
    & \textbf{HR}                       & Here                                    & 这里                  \\
    & \textbf{HR}                       & Hear                                    & 听到                  \\
    & HST                               & High-speed telegraphy                   & 快速收发报               \\
    & \textbf{HW}                       & How                                     & 怎样、如何               \\
    & Hz                                & hertz                                   & 赫兹（单位）              \\
  I & \textbf{ITA1}                     & Baudot code                             & 博多码、ITA2的前身         \\
    & \textbf{ITA2}                     &                                         & 国际2号电报码             \\
  K & \si{\kHz}                         & kilohertz                               & 千赫兹（单位）             \\
    & KP                                &                                         & 收听                  \\
  L & \textbf{LCD}                      & Liquid-crystal display                  & 液晶显示器               \\
    & \textbf{LED}                      & Light-emitting diode                    & 发光二极管               \\
    & \textbf{LP}                       & Log-periodic antenna                    & 对数周期天线              \\
    & \textbf{LW}                       & Long whip antenna                       & 长线天线                \\
  M & \si{\MHz}                         & megahertz                               & 兆赫兹（单位）             \\
    & \textbf{MNI}                      & Many                                    & 很多                  \\
    & \textbf{MNY}                      & Many                                    & 很多                  \\
    & \textbf{MODE}                     &                                         & 方式                  \\
    & \textbf{MTRS}                     & Meters                                  & 米                   \\
    & MUF                               & Maximum usable frequency                & 最高可用频率              \\
  N & \textbf{NAME}                     &                                         & 名字                  \\
    & \textbf{NB}                       & Noise blanker                           & 抑噪                  \\
    & NFM                               & Narrowband FM                           & 窄带调频                \\
    & \textbf{NICE}                     &                                         & 良好的                 \\
    & \textbf{NW}                       & Now                                     & 现在                  \\
  O & \textbf{OM}                       & Old man                                 & 老朋友                 \\%朋友
    & \textbf{OP}                       & Operator                                & 操作员                 \\
    & \textbf{OPR}                      & Operator                                & 操作员                 \\
    & OVER                              &                                         & 我的传输已结束，期待回复        \\
  P & \textbf{PM}                       & Phase modulation                        & 相位调制                \\
    & \textbf{P O BOX}                  & Post office box                         & 邮政信箱                \\
    & \textbf{PRE}                      & Preamplifier                            & 收信机前置放大器            \\
    & \textbf{PROC}                     & Processor                               & 发信语音压缩              \\
    & \textbf{PSK}                      & Phase-shift keying                      & 相位偏移调制              \\
    & PSK31                             & Phase Shift Keying, 31 Baud             & 移相键控、31波特           \\
    & \textbf{PTT}                      & Push-to-talk                            & 按键发射                \\
    & \textbf{PWR}                      & Power                                   & 功率                  \\
  Q & QPSK31                            & Phase Shift Keying, 31 Baud             & 移相键控、31波特           \\
  R & \textbf{RCVR}                     & Receiver                                & 收信机                 \\
    & \textbf{RF}                       & Radio frequency                         & 无线电波                \\
    & \textbf{RIG}                      & Equipment                               & 电台设备                \\
    & \textbf{RIT}                      & Receiver incremental tuning             & 接收增量调谐              \\
    & \textbf{RMKS}                     & Remarks                                 & 备注、注释               \\
    & \textbf{ROGER}                    &                                         & 回答起始语，相当于“明白”，      \\
    &                                   &                                         & 仅在已完全抄收对方刚才发送的信息时使用 \\
    & \textbf{RPRT}                     & Report                                  & 报告                  \\
    & RST                               & Readability, Strength, Tone             & 可辨度、信号强度、音调系统       \\
    & \textbf{RTTY}                     & Radioteletype                           &                     \\
    & \textbf{RX}                       & Receiver                                & 收信机                 \\
  S & \textbf{SASE}                     &                                         & 写好收信人地址的信封          \\
    & \textbf{SHF}                      & Super high frequency                    & 超高频                 \\
    & \textbf{SID}                      & Sudden ionospheric disturbance          & 突发电离层扰动             \\
    & $\overline{\mathbf{S}\mathbf{K}}$ &                                         & 结束通信                \\
    & \textbf{SQL}                      & Squelch                                 & 静噪                  \\
    & \textbf{SRI}                      & Sorry                                   & 对不起                 \\
    & \textbf{SRY}                      & Sorry                                   & 对不起                 \\
    & \textbf{SSB}                      & Single-sideband modulation              & 单边带调制               \\
    & \textbf{SSN}                      & Sunspot number                          & 太阳黑子平均数             \\
    & \textbf{SSTV}                     & Slow-scan television                    & 慢扫描电视               \\
    & \textbf{STN}                      & Station                                 & 电台                  \\
    & STANDBY                           &                                         & 等待你的呼叫              \\
    & \textbf{SURE}                     &                                         & 确实                  \\
    & \textbf{SWL}                      & Short waves listener                    & 短波收听者               \\
  T & \textbf{TEMP}                     & Temperature                             & 温度                  \\
    & THIS IS                           &                                         & 此消息发自以下xxxx电台       \\
    & TIME                              &                                         & 时间                  \\
    & \textbf{TKS}                      & Thanks                                  & 谢谢                  \\
    & \textbf{TNX}                      & Thanks                                  & 谢谢                  \\
    & \textbf{TU}                       & Thank you                               & 谢谢你                 \\
    & \textbf{TX}                       & Transmitter                             & 发信机                 \\
  U & \textbf{UHF}                      & Ultra high frequency                    & 特高频                 \\
    & \textbf{UR}                       & Your                                    & 你的，你是               \\
    & \textbf{UTC}                      & Coordinated Universal Time              & 世界协调时               \\%协调世界时
  V & \textbf{VER}                      & Vertical antenna                        & 垂直天线                \\
    & \textbf{VHF}                      & Very high frequency                     & 甚高频                 \\
    & \textbf{VIA}                      &                                         & 经、由                 \\
    & \textbf{VOX}                      & Voice-operated switch                   & 发信机声控               \\
    & \textbf{VY}                       & Very                                    & 很、非常                \\
  W & WFM                               & Wideband FM                             & 宽带调频                \\
    & \textbf{WK}                       & Week                                    & 星期                  \\
    & \textbf{WK}                       & Work                                    & 工作                  \\
    & \textbf{WKD}                      & Worked                                  & 联络过、工作过             \\
    & \textbf{WTS}                      & Watts                                   & 瓦特                  \\
    & \textbf{WX}                       & Weather                                 & 天气                  \\
  X & \textbf{XIT}                      & Transmitter Incremental Tuning          & 发射增量调谐              \\
    & \textbf{XMAS}                     & Christmas                               & 圣诞节                 \\
    & \textbf{XMTR}                     & Transmitter                             & 发信机                 \\
    & \textbf{XCVR}                     & Transceiver                             & 收发信机                \\
    & \textbf{XYL}                      & Wife                                    & 妻子、已婚女子             \\
  Y & \textbf{YAGI}                     & Yagi–Uda antenna                        & 八木天线                \\
    & \textbf{YL}                       & Young lady                              & 小姐、女士               \\
    & \textbf{73}                       &                                         & 向对方的致意、美好的祝愿        \\
    & \textbf{88}                       &                                         & 向对方异性操作员的致意、美好的祝愿   \\
\end{longtable}

\newpage

\section{术语表}

\begin{longtable}[l]{ll}
  \textbf{术语}                & \textbf{含义} \\
  Class of emission          & 发射类别        \\
  Frequency-shift telegraphy & 频移电报技术      \\
  Necessary bandwidth        & 必要带宽        \\
\end{longtable}

\newpage

\section{发射类别的表示方法（节录）}

\begin{tabular}{|c|c|}
  \hline
  \multicolumn{2}{|c|}{\textbf{主载波类型（第一个符号）}} \\
  \hline
  \textbf{符号} & \textbf{描述}                   \\
  \hline
  A           & 双边带                           \\
  \hline
  F           & 调频                            \\
  \hline
  G           & 调相                            \\
  \hline
  J           & 单边带、抑制载波                      \\
  \hline
\end{tabular}

\bigskip

\begin{tabular}{|c|c|}
  \hline
  \multicolumn{2}{|c|}{\textbf{调制主载波的信号的性质（第二个符号）}} \\
  \hline
  \textbf{符号} & \textbf{描述}                         \\
  \hline
  1           & 不用调制副载波但包含量化或数字信息的单个通路              \\
  \hline
  2           & 利用调制副载波且包含量化或数字信息的单个通路              \\
  \hline
  3           & 包含模拟信息的单个通路                         \\
  \hline
\end{tabular}

\bigskip

\begin{tabular}{|c|c|}
  \hline
  \multicolumn{2}{|c|}{\textbf{被发送信息类型（第三个符号）}} \\
  \hline
  \textbf{符号} & \textbf{描述}                     \\
  \hline
  A           & 电报──用于人工收听                      \\
  \hline
  B           & 电报──用于自动接收                      \\
  \hline
  E           & 电话（包括声音广播）                      \\
  \hline
  F           & 电视（视频）                          \\
  \hline
\end{tabular}

\newpage

\section{无线电频段和波段表（节录）}

\begin{tabular}{|c|c|c|c|c|}
  \hline
  \textbf{频段名称} & \textbf{缩写} & \textbf{频率范围}       & \textbf{波长范围} & \textbf{波段名称} \\
  %\hline
  %反正不考
  %极低频 & ELF & 3–30 Hz & 100,000–10,000 km & 极长波 \\
  %\hline
  %超低频 & SLF & 30–300 Hz & 10,000–1,000 km & 超长波 \\
  %\hline
  %特低频 & ULF & 300–3,000 Hz & 1,000–100 km & 特长波 \\
  %\hline
  %甚低频 & VLF & 3–30 \si{\kHz} & 100–10 km & 甚长波 \\
  \hline
  低频            & LF          & 30–300 \si{\kHz}    & 10–1 km       & 长波            \\
  \hline
  中频            & MF          & 300–3,000 \si{\kHz} & 1,000–100 m   & 中波            \\
  \hline
  高频            & HF          & 3–30 \si{\MHz}      & 100–10 m      & 短波            \\
  \hline
  甚高频           & VHF         & 30–300 \si{\MHz}    & 10–1 m        & 米波            \\
  \hline
  特高频           & UHF         & 300–3,000 \si{\MHz} & 1–0.1 m       & 分米波           \\
  \hline
  超高频           & SHF         & 3–30 \si{\GHz}      & 100–10 mm     & 厘米波           \\
  \hline
  极高频           & EHF         & 30–300 \si{\GHz}    & 10–1 mm       & 毫米波           \\
  \hline
  %反正不考
  %至高频 & THF & 300–3,000 \si{\GHz} & 1–0.1 mm & 丝米波 \\
  %\hline
\end{tabular}

\newpage

\section{各业余无线电频段名称（节录）}

\begin{longtable}{|c|c|}
  \hline
  \textbf{波段名称} & \textbf{频率范围}        \\
  \hline
  160米          & 1800–2000 \si{\kHz}  \\
  \hline
  80米           & 3.5–4.0 \si{\MHz}    \\
  \hline
  40米           & 7.0–7.3 \si{\MHz}    \\
  \hline
  20米           & 14.0–14.35 \si{\MHz} \\
  \hline
  15米           & 21–21.45 \si{\MHz}   \\
  \hline
  10米           & 28–29.7 \si{\MHz}    \\
  \hline
  6米            & 50–54 \si{\MHz}      \\
  \hline
  2米            & 144–148 \si{\MHz}    \\
  \hline
  0.7米（70厘米）    & 420–450 \si{\MHz}    \\
  \hline
\end{longtable}

\newpage

\section{ITU字头（节录）}

\begin{longtable}{|l|l|}
  \hline
  \textbf{字头范围}       & \textbf{分配给国家/地区}                         \\
                      & \endfirsthead
  \textbf{字头范围}       & \textbf{分配给国家/地区}                         \\
                      & \endhead
  \hline
  AA--AL              & United States                             \\
  \hline
  AM--AO              & Spain                                     \\
  \hline
  AP--AS              & Pakistan                                  \\
  \hline
  AT--AW              & India                                     \\
  \hline
  AX                  & Australia                                 \\
  \hline
  AY--AZ              & Argentina                                 \\
  \hline
  A2                  & Botswana                                  \\
  \hline
  A3                  & Tonga                                     \\
  \hline
  A4                  & Oman                                      \\
  \hline
  A5                  & Bhutan                                    \\
  \hline
  A6                  & United Arab Emirates                      \\
  \hline
  A7                  & Qatar                                     \\
  \hline
  A8                  & Liberia                                   \\
  \hline
  A9                  & Bahrain                                   \\
  \hline
  B                   & China                                     \\
  \hline
  B (BM-BQ, BU-BX)    & Taiwan                                    \\
  \hline
  CA--CE              & Chile                                     \\
  \hline
  CF--CK              & Canada                                    \\
  \hline
  CN                  & Morocco                                   \\
  \hline
  CO                  & Cuba                                      \\
  \hline
  CP                  & Bolivia                                   \\
  \hline
  CQ--CU              & Portugal                                  \\
  \hline
  CV--CX              & Uruguay                                   \\
  \hline
  CY--CZ              & Canada                                    \\
  \hline
  C2                  & Nauru                                     \\
  \hline
  C3                  & Andorra                                   \\
  \hline
  C4                  & Cyprus                                    \\
  \hline
  C5                  & The Gambia                                \\
  \hline
  C6                  & The Bahamas                               \\
  \hline
  C7                  & World Meteorological Organization         \\
  \hline
  C8--C9              & Mozambique                                \\
  \hline
  DA--DR              & Germany                                   \\
  \hline
  DS--DT              & South Korea                               \\
  \hline
  DU--DZ              & Philippines                               \\
  \hline
  D2--D3              & Angola                                    \\
  \hline
  D4                  & Cape Verde                                \\
  \hline
  D5                  & Liberia                                   \\
  \hline
  D6                  & Comoros                                   \\
  \hline
  D7--D9              & South Korea                               \\
  \hline
  EA--EH              & Spain                                     \\
  \hline
  EI--EJ              & Ireland                                   \\
  \hline
  EK                  & Armenia                                   \\
  \hline
  EL                  & Liberia                                   \\
  \hline
  EM--EO              & Ukraine                                   \\
  \hline
  EP--EQ              & Iran                                      \\
  \hline
  ER                  & Moldova                                   \\
  \hline
  ES                  & Estonia                                   \\
  \hline
  ET                  & Ethiopia                                  \\
  \hline
  EU--EW              & Belarus                                   \\
  \hline
  EX                  & Kyrgyzstan                                \\
  \hline
  EY                  & Tajikistan                                \\
  \hline
  EZ                  & Turkmenistan                              \\
  \hline
  E2                  & Thailand                                  \\
  \hline
  E3                  & Eritrea                                   \\
  \hline
  E4                  & Palestinian Authority                     \\
  \hline
  E5                  & Cook Islands                              \\
  \hline
  E6                  & Niue                                      \\
  \hline
  E7                  & Bosnia and Herzegovina                    \\
  \hline
  F                   & France                                    \\
  \hline
  G                   & United Kingdom                            \\
  \hline
  HA                  & Hungary                                   \\
  \hline
  HB                  & Switzerland                               \\
  \hline
  HB (HB0, HB3Y, HBL) & Liechtenstein                             \\
  \hline
  HC--HD              & Ecuador                                   \\
  \hline
  HE                  & Switzerland                               \\
  \hline
  HF                  & Poland                                    \\
  \hline
  HG                  & Hungary                                   \\
  \hline
  HH                  & Haiti                                     \\
  \hline
  HI                  & Dominican Republic                        \\
  \hline
  HJ--HK              & Colombia                                  \\
  \hline
  HL                  & South Korea                               \\
  \hline
  HM                  & North Korea                               \\
  \hline
  HN                  & Iraq                                      \\
  \hline
  HO--HP              & Panama                                    \\
  \hline
  HQ--HR              & Honduras                                  \\
  \hline
  HS                  & Thailand                                  \\
  \hline
  HT                  & Nicaragua                                 \\
  \hline
  HU                  & El Salvador                               \\
  \hline
  HV                  & Vatican City                              \\
  \hline
  HW--HY              & France                                    \\
  \hline
  HZ                  & Saudi Arabia                              \\
  \hline
  H2                  & Cyprus                                    \\
  \hline
  H3                  & Panama                                    \\
  \hline
  H4                  & Solomon Islands                           \\
  \hline
  H6--H7              & Nicaragua                                 \\
  \hline
  H8--H9              & Panama                                    \\
  \hline
  I                   & Italy                                     \\
  \hline
  JA--JS              & Japan                                     \\
  \hline
  JT--JV              & Mongolia                                  \\
  \hline
  JW--JX              & Norway                                    \\
  \hline
  JY                  & Jordan                                    \\
  \hline
  JZ                  & Indonesia                                 \\
  \hline
  J2                  & Djibouti                                  \\
  \hline
  J3                  & Grenada                                   \\
  \hline
  J4                  & Greece                                    \\
  \hline
  J5                  & Guinea-Bissau                             \\
  \hline
  J6                  & Saint Lucia                               \\
  \hline
  J7                  & Dominica                                  \\
  \hline
  J8                  & Saint Vincent and the Grenadines          \\
  \hline
  K                   & United States                             \\
  \hline
  LA--LN              & Norway                                    \\
  \hline
  LO--LW              & Argentina                                 \\
  \hline
  LX                  & Luxembourg                                \\
  \hline
  LY                  & Lithuania                                 \\
  \hline
  LZ                  & Bulgaria                                  \\
  \hline
  L2--L9              & Argentina                                 \\
  \hline
  M                   & United Kingdom                            \\
  \hline
  N                   & United States                             \\
  \hline
  OA--OC              & Peru                                      \\
  \hline
  OD                  & Lebanon                                   \\
  \hline
  OE                  & Austria                                   \\
  \hline
  OF--OJ              & Finland                                   \\
  \hline
  OK--OL              & Czech Republic                            \\
  \hline
  OM                  & Slovakia                                  \\
  \hline
  ON--OT              & Belgium                                   \\
  \hline
  OU--OZ              & Denmark                                   \\
  \hline
  PA--PI              & Netherlands                               \\
  \hline
  PJ                  & Netherlands Antilles                      \\
  \hline
  PK--PO              & Indonesia                                 \\
  \hline
  PP--PY              & Brazil                                    \\
  \hline
  PZ                  & Suriname                                  \\
  \hline
  P2                  & Papua New Guinea                          \\
  \hline
  P3                  & Cyprus                                    \\
  \hline
  P4                  & Aruba                                     \\
  \hline
  P5--P9              & North Korea                               \\
  \hline
  R                   & Russia                                    \\
  \hline
  SA--SM              & Sweden                                    \\
  \hline
  SN--SR              & Poland                                    \\
  \hline
  SSA--SSM            & Egypt                                     \\
  \hline
  SSN--STZ            & Sudan                                     \\
  \hline
  SU                  & Egypt                                     \\
  \hline
  SV--SZ              & Greece                                    \\
  \hline
  S2--S3              & Bangladesh                                \\
  \hline
  S5                  & Slovenia                                  \\
  \hline
  S6                  & Singapore                                 \\
  \hline
  S7                  & Seychelles                                \\
  \hline
  S8                  & South Africa                              \\
  \hline
  S9                  & São Tomé and Príncipe                     \\
  \hline
  TA--TC              & Turkey                                    \\
  \hline
  TD                  & Guatemala                                 \\
  \hline
  TE                  & Costa Rica                                \\
  \hline
  TF                  & Iceland                                   \\
  \hline
  TG                  & Guatemala                                 \\
  \hline
  TH                  & France                                    \\
  \hline
  TI                  & Costa Rica                                \\
  \hline
  TJ                  & Cameroon                                  \\
  \hline
  TK                  & France                                    \\
  \hline
  TL                  & Central African Republic                  \\
  \hline
  TM                  & France                                    \\
  \hline
  TN                  & Republic of the Congo                     \\
  \hline
  TO--TQ              & France                                    \\
  \hline
  TR                  & Gabon                                     \\
  \hline
  TS                  & Tunisia                                   \\
  \hline
  TT                  & Chad                                      \\
  \hline
  TU                  & Ivory Coast                               \\
  \hline
  TV--TX              & France                                    \\
  \hline
  TY                  & Benin                                     \\
  \hline
  TZ                  & Mali                                      \\
  \hline
  T2                  & Tuvalu                                    \\
  \hline
  T3                  & Kiribati                                  \\
  \hline
  T4                  & Cuba                                      \\
  \hline
  T5                  & Somalia                                   \\
  \hline
  T6                  & Afghanistan                               \\
  \hline
  T7                  & San Marino                                \\
  \hline
  T8                  & Palau                                     \\
  \hline
  UJ--UM              & Uzbekistan                                \\
  \hline
  UN--UQ              & Kazakhstan                                \\
  \hline
  UR--UZ              & Ukraine                                   \\
  \hline
  VA--VG              & Canada                                    \\
  \hline
  VH--VN              & Australia                                 \\
  \hline
  VO                  & Canada                                    \\
  \hline
  VP--VQ              & United Kingdom                            \\
  \hline
  VR                  & Hong Kong                                 \\
  \hline
  VS                  & United Kingdom                            \\
  \hline
  VT--VW              & India                                     \\
  \hline
  VX--VY              & Canada                                    \\
  \hline
  VZ                  & Australia                                 \\
  \hline
  V2                  & Antigua and Barbuda                       \\
  \hline
  V3                  & Belize                                    \\
  \hline
  V4                  & Saint Kitts and Nevis                     \\
  \hline
  V5                  & Namibia                                   \\
  \hline
  V6                  & Federated States of Micronesia            \\
  \hline
  V7                  & Marshall Islands                          \\
  \hline
  V8                  & Brunei                                    \\
  \hline
  W                   & United States                             \\
  \hline
  XA--XI              & Mexico                                    \\
  \hline
  XJ--XO              & Canada                                    \\
  \hline
  XP                  & Denmark                                   \\
  \hline
  XQ--XR              & Chile                                     \\
  \hline
  XS                  & China                                     \\
  \hline
  XT                  & Burkina Faso                              \\
  \hline
  XU                  & Cambodia                                  \\
  \hline
  XV                  & Vietnam                                   \\
  \hline
  XW                  & Laos                                      \\
  \hline
  XX                  & Macao                                     \\
  \hline
  XY--XZ              & Burma                                     \\
  \hline
  YA                  & Afghanistan                               \\
  \hline
  YB--YH              & Indonesia                                 \\
  \hline
  YI                  & Iraq                                      \\
  \hline
  YJ                  & Vanuatu                                   \\
  \hline
  YK                  & Syria                                     \\
  \hline
  YL                  & Latvia                                    \\
  \hline
  YM                  & Turkey                                    \\
  \hline
  YN                  & Nicaragua                                 \\
  \hline
  YO--YR              & Romania                                   \\
  \hline
  YS                  & El Salvador                               \\
  \hline
  YT--YU              & Serbia                                    \\
  \hline
  YV--YY              & Venezuela                                 \\
  \hline
  Y2--Y9              & Germany                                   \\
  \hline
  ZA                  & Albania                                   \\
  \hline
  ZB--ZJ              & United Kingdom                            \\
  \hline
  ZK--ZM              & New Zealand                               \\
  \hline
  ZN--ZO              & United Kingdom                            \\
  \hline
  ZP                  & Paraguay                                  \\
  \hline
  ZQ                  & United Kingdom                            \\
  \hline
  ZR--ZU              & South Africa                              \\
  \hline
  ZV--ZZ              & Brazil                                    \\
  \hline
  Z2                  & Zimbabwe                                  \\
  \hline
  Z3                  & North Macedonia                           \\
  \hline
  Z8                  & South Sudan                               \\
  \hline
  2                   & United Kingdom                            \\
  \hline
  3A                  & Monaco                                    \\
  \hline
  3B                  & Mauritius                                 \\
  \hline
  3C                  & Equatorial Guinea                         \\
  \hline
  3DA--3DM            & Eswatini                                  \\
  \hline
  3DN--3DZ            & Fiji                                      \\
  \hline
  3E--3F              & Panama                                    \\
  \hline
  3G                  & Chile                                     \\
  \hline
  3H--3U              & China                                     \\
  \hline
  3V                  & Tunisia                                   \\
  \hline
  3W                  & Vietnam                                   \\
  \hline
  3X                  & Guinea                                    \\
  \hline
  3Y                  & Norway                                    \\
  \hline
  3Z                  & Poland                                    \\
  \hline
  4A--4C              & Mexico                                    \\
  \hline
  4D--4I              & Philippines                               \\
  \hline
  4J--4K              & Azerbaijan                                \\
  \hline
  4L                  & Georgia                                   \\
  \hline
  4M                  & Venezuela                                 \\
  \hline
  4O                  & Montenegro                                \\
  \hline
  4P--4S              & Sri Lanka                                 \\
  \hline
  4T                  & Peru                                      \\
  \hline
  4U                  & United Nations                            \\
  \hline
  4V                  & Haiti                                     \\
  \hline
  4W                  & Timor-Leste                               \\
  \hline
  4X                  & Israel                                    \\
  \hline
  4Y                  & International Civil Aviation Organization \\
  \hline
  4Z                  & Israel                                    \\
  \hline
  5A                  & Libya                                     \\
  \hline
  5B                  & Cyprus                                    \\
  \hline
  5C--5G              & Morocco                                   \\
  \hline
  5H--5I              & Tanzania                                  \\
  \hline
  5J--5K              & Colombia                                  \\
  \hline
  5L--5M              & Liberia                                   \\
  \hline
  5N--5O              & Nigeria                                   \\
  \hline
  5P--5Q              & Denmark                                   \\
  \hline
  5R--5S              & Madagascar                                \\
  \hline
  5T                  & Mauritania                                \\
  \hline
  5U                  & Niger                                     \\
  \hline
  5V                  & Togo                                      \\
  \hline
  5W                  & Western Samoa                             \\
  \hline
  5X                  & Uganda                                    \\
  \hline
  5Y--5Z              & Kenya                                     \\
  \hline
  6A--6B              & Egypt                                     \\
  \hline
  6C                  & Syria                                     \\
  \hline
  6D--6J              & Mexico                                    \\
  \hline
  6K--6N              & South Korea                               \\
  \hline
  6O                  & Somalia                                   \\
  \hline
  6P--6S              & Pakistan                                  \\
  \hline
  6T--6U              & Sudan                                     \\
  \hline
  6V--6W              & Senegal                                   \\
  \hline
  6X                  & Madagascar                                \\
  \hline
  6Y                  & Jamaica                                   \\
  \hline
  6Z                  & Liberia                                   \\
  \hline
  7A--7I              & Indonesia                                 \\
  \hline
  7J--7N              & Japan                                     \\
  \hline
  7O                  & Yemen                                     \\
  \hline
  7P                  & Lesotho                                   \\
  \hline
  7Q                  & Malawi                                    \\
  \hline
  7R                  & Algeria                                   \\
  \hline
  7S                  & Sweden                                    \\
  \hline
  7T--7Y              & Algeria                                   \\
  \hline
  7Z                  & Saudi Arabia                              \\
  \hline
  8A--8I              & Indonesia                                 \\
  \hline
  8J--8N              & Japan                                     \\
  \hline
  8O                  & Botswana                                  \\
  \hline
  8P                  & Barbados                                  \\
  \hline
  8Q                  & Maldives                                  \\
  \hline
  8R                  & Guyana                                    \\
  \hline
  8S                  & Sweden                                    \\
  \hline
  8T--8Y              & India                                     \\
  \hline
  8Z                  & Saudi Arabia                              \\
  \hline
  9A                  & Croatia                                   \\
  \hline
  9B--9D              & Iran                                      \\
  \hline
  9E--9F              & Ethiopia                                  \\
  \hline
  9G                  & Ghana                                     \\
  \hline
  9H                  & Malta                                     \\
  \hline
  9I--9J              & Zambia                                    \\
  \hline
  9K                  & Kuwait                                    \\
  \hline
  9L                  & Sierra Leone                              \\
  \hline
  9M                  & Malaysia                                  \\
  \hline
  9N                  & Nepal                                     \\
  \hline
  9O--9T              & Democratic Republic of the Congo          \\
  \hline
  9U                  & Burundi                                   \\
  \hline
  9V                  & Singapore                                 \\
  \hline
  9W                  & Malaysia                                  \\
  \hline
  9X                  & Rwanda                                    \\
  \hline
  9Y--9Z              & Trinidad and Tobago                       \\
  \hline
\end{longtable}

\newpage

\section{计算公式}

直流电路欧姆定律：
\[\mbox{电流}I(\si{\ampere})=\frac{\mbox{电压}U(\si{\volt})}{\mbox{电阻值}R(\si{\ohm})}\]

峰值电压的计算公式：
\[\mbox{峰值电压}(\si{\volt})=\frac{\mbox{峰--峰值}}{2}\]

电功率的计算公式：
\begin{equation*}
  \begin{aligned}
    \mbox{电功率}P(\si{\watt}) & = U\times I     \\
                            & = I^2 \times R  \\
                            & = \frac{U^2}{R}
  \end{aligned}
\end{equation*}

正弦交流电压的有效值计算公式：
\begin{equation*}
  \begin{aligned}
    \mbox{交流电压的有效值}(\si{\volt}) & =\frac{\mbox{峰值}}{\sqrt{2}}   \\
                                & \approx \mbox{峰值}\times 0.707
  \end{aligned}
\end{equation*}

功率增益的计算公式：
\[\mbox{功率增益}=10 \log_{10} \left( {\frac{P_{ \mbox{输出} }}{P_{ \mbox{输入} }}}\right) \si{\dB}\]

电压增益的计算公式：
\[\mbox{电压增益}=20 \log \left( {\frac{V_{ \mbox{输出} }}{V_{ \mbox{输入} }}} \right) \si{\dB}\]

%电流增益的计算公式：
%$$\mbox{电流增益}=20 \log \left( {\frac{I_{ \mbox{输出} }}{I_{ \mbox{输入} }}} \right)\ \mathrm{dB}$$

绝对增益与相对增益的关系：
\[\mbox{绝对增益} = \mbox{相对增益} + \SI{2.15}{\dB}\]

两数的乘积的对数等于两数的对数之和：
\[\log a\: b=\log a+\log b\]

两数的商的对数等于两数的对数之差：
\[\log\frac{a}{b}=\log a-\log b\]

某数的n次幂的对数等于该数的对数的n倍：
\[\log a^n =n\log a\]

\newpage

\section{常用对数表（节录）}

说明：强烈建议背诵$\log_{10} 1.0$、$\log 2.0$、……$\log 10.0$等第0列的10个数值到小数点后三位，考试闭卷，无法使用计算器，如果不背诵，将无法计算。

\begin{longtable}[c]{|c|c|}
  \hline
  \textbf{N} & \textbf{0}  \\
  \hline
  %	\endfirsthead
  %	\hline
  %	\textbf{N} & \textbf{1.0} \\
  \endhead
  \num{1.0}  & \num{.000}  \\ \hline
  \num{2.0}  & \num{.301}  \\ \hline
  \num{3.0}  & \num{.477}  \\ \hline
  \num{4.0}  & \num{.602}  \\ \hline
  \num{5.0}  & \num{.699}  \\ \hline
  \num{6.0}  & \num{.778}  \\ \hline
  \num{7.0}  & \num{.845}  \\ \hline
  \num{8.0}  & \num{.903}  \\ \hline
  \num{9.0}  & \num{.954}  \\ \hline
  \num{10.0} & \num{1.000} \\ \hline
\end{longtable}

\newpage

\section{逻辑门的真值表}

\begin{tabular}{cc}%
  \begin{tabular}{|c|c|c|}
    \multicolumn{3}{c}{\textbf{与门}}                           \\
    \hline
    \multicolumn{2}{|c|}{\textbf{输入}} & \textbf{输出}           \\
    \hline
    A                                 & B           & A AND B \\
    \hline
    \num{0}                           & \num{0}     & \num{0} \\
    \hline
    \num{0}                           & \num{1}     & \num{0} \\
    \hline
    \num{1}                           & \num{0}     & \num{0} \\
    \hline
    \num{1}                           & \num{1}     & \num{1} \\
    \hline
  \end{tabular} &
  \begin{tabular}{|c|c|c|}
    \multicolumn{3}{c}{\textbf{与非门}}                           \\
    \hline
    \multicolumn{2}{|c|}{\textbf{输入}} & \textbf{输出}            \\
    \hline
    A                                 & B           & A NAND B \\
    \hline
    \num{0}                           & \num{0}     & \num{1}  \\
    \hline
    \num{0}                           & \num{1}     & \num{1}  \\
    \hline
    \num{1}                           & \num{0}     & \num{1}  \\
    \hline
    \num{1}                           & \num{1}     & \num{0}  \\
    \hline
  \end{tabular} \tabularnewline
\end{tabular}

\bigskip

\begin{tabular}{cc}%
  \begin{tabular}{|c|c|c|}
    \multicolumn{3}{c}{\textbf{或门}}                           \\
    \hline
    \multicolumn{2}{|c|}{\textbf{输入}} & \textbf{输出}           \\
    \hline
    A                                 & B           & A OR B  \\
    \hline
    \num{0}                           & \num{0}     & \num{0} \\
    \hline
    \num{0}                           & \num{1}     & \num{1} \\
    \hline
    \num{1}                           & \num{0}     & \num{1} \\
    \hline
    \num{1}                           & \num{1}     & \num{1} \\
    \hline
  \end{tabular} &
  \begin{tabular}{|c|c|c|}
    \multicolumn{3}{c}{\textbf{或非门}}                          \\
    \hline
    \multicolumn{2}{|c|}{\textbf{输入}} & \textbf{输出}           \\
    \hline
    A                                 & B           & A NOR B \\
    \hline
    \num{0}                           & \num{0}     & \num{1} \\
    \hline
    \num{0}                           & \num{1}     & \num{0} \\
    \hline
    \num{1}                           & \num{0}     & \num{0} \\
    \hline
    \num{1}                           & \num{1}     & \num{0} \\
    \hline
  \end{tabular}  \tabularnewline
\end{tabular}

\bigskip

\begin{tabular}{cc}%
  \begin{tabular}{|c|c|c|}
    \multicolumn{3}{c}{\textbf{异或门}}                          \\
    \hline
    \multicolumn{2}{|c|}{\textbf{输入}} & \textbf{输出}           \\
    \hline
    A                                 & B           & A XOR B \\
    \hline
    \num{0}                           & \num{0}     & \num{0} \\
    \hline
    \num{0}                           & \num{1}     & \num{1} \\
    \hline
    \num{1}                           & \num{0}     & \num{1} \\
    \hline
    \num{1}                           & \num{1}     & \num{0} \\
    \hline
  \end{tabular} &
  \begin{tabular}{|c|c|c|}
    \multicolumn{3}{c}{\textbf{异或非门}}                          \\
    \hline
    \multicolumn{2}{|c|}{\textbf{输入}} & \textbf{输出}            \\
    \hline
    A                                 & B           & A XNOR B \\
    \hline
    \num{0}                           & \num{0}     & \num{1}  \\
    \hline
    \num{0}                           & \num{1}     & \num{0}  \\
    \hline
    \num{1}                           & \num{0}     & \num{0}  \\
    \hline
    \num{1}                           & \num{1}     & \num{1}  \\
    \hline
  \end{tabular} \tabularnewline
\end{tabular}

\newpage

\section{计算逻辑门电路输出信号题型答案的C程序}

\lstinputlisting[language=C]{jisuan.c}

\newpage

\section{将呼号转换成字母解释法的单词组合的C程序}

\lstinputlisting[language=C]{icao.c}

\newpage

\section{有用网址}

\begin{longtable}{|p{8cm}|p{8cm}|}
  \hline
  \textbf{名称}          & \textbf{网址}                                                     \\
  \hline
  业余无线电台操作技术能力验证考核报名系统 & \url{http://82.157.138.16:8091/CRAC/crac/login_student.html}    \\
  \hline
  复习试题和模拟考试系统下载        & \url{http://82.157.138.16:8091/CRAC/crac/pages/list_files.html} \\
  \hline
  中国无线电协会业余无线电分会资料下载   & \url{http://www.crac.org.cn/News/List?type=6}                   \\
  \hline
\end{longtable}


% 后记
\chapter*{后记}

本书在撰写过程中得到了家人的大力支持，在此感谢家人的付出，如果没有他们无微不至的照顾，是不会有这本书的。

作为非无线电相关专业人员出身，笔者能力有限，此书仍有许多不足，望读者指正。

本书的最新版PDF文件可在项目网站\url{https://github.com/mike2718/ham}免费下载，也可扫描下方二维码访问。

\begin{figure}[h]
  \includegraphics{qr}
  \centering
\end{figure}

\newpage


% 参考文献
%% \chapter*{参考文献}
\bibliographystyle{gbt7714-numerical}  % 参考文献用
\bibliography{cankaowenxian}


\end{document}